\documentclass[11pt,a4paper]{article}
\usepackage[T1]{fontenc}
\usepackage[utf8]{inputenc}
\usepackage{lmodern,microtype}
\usepackage[margin=28mm]{geometry}
\usepackage{amsmath,amssymb,amsthm,mathtools}
\usepackage{booktabs,longtable,array,enumitem}
\usepackage[numbers,sort&compress]{natbib}
\usepackage[hyphens]{url}
\usepackage{xurl}
\usepackage[hidelinks]{hyperref}
\usepackage{bookmark}
\setlength{\emergencystretch}{2em}
\setcounter{tocdepth}{1}
\newcolumntype{P}[1]{>{\raggedright\arraybackslash}p{#1}}
\newtheorem{theorem}{Theorem}[section]
\newtheorem{lemma}[theorem]{Lemma}
\newtheorem{proposition}[theorem]{Proposition}
\newtheorem{corollary}[theorem]{Corollary}
\theoremstyle{definition}
\newtheorem{definition}[theorem]{Definition}
\newtheorem{condition}[theorem]{Condition}
\theoremstyle{remark}
\newtheorem{remark}[theorem]{Remark}
\newcommand{\E}{\mathbb E}
\newcommand{\Pp}{\mathbb P}
\newcommand{\R}{\mathbb R}
\newcommand{\N}{\mathbb N}
\newcommand{\CE}{C_{\mathrm E}}
\newcommand{\supp}{\operatorname{supp}}
\newcommand{\Var}{\operatorname{Var}}
\newcommand{\pos}[1]{(#1)_+}
\newcommand{\ind}{\mathbf1}
\newcommand{\iid}{\mathrm{iid}}
\title{The sharp iid Berry--Esseen inequality:\\
the mathematical argument and its certificates\thanks{AI systems were used
extensively in mathematical exploration, checking programs and manuscript
revision. Their output and agreement are not substitutes for mathematical
proof or independent human review. The authors are responsible for the content.}}
\author{Bartosz Ko{\l}odziejek \and Kacper Rodziewicz}
\date{4 November 2026}
\begin{document}
\maketitle
\begin{center}\small
Revision for mathematical review.
\end{center}
\begin{abstract}
We give a computer-assisted proof of the sharp iid Berry--Esseen
inequality with constant $(3+\sqrt{10})/(6\sqrt{2\pi})$.
The argument excludes a first failing index using a finite sharp base
and two local comparisons. Extremal contact and a nonnegative support-gap
defect force a possible maximizer near a two-point law; conditional
binomial counts then measure the effect of its remaining fluctuations.
Fourier bounds and exact certificates establish the quantitative
inequalities on complete parameter regions. We prove their analytic
reductions and explain certificate construction, acceptance and
independent reproduction. Equality in the uniform local comparison
also gives a strict bound for each sample size at least $47$.
\end{abstract}
\begingroup\small\tableofcontents\endgroup
\section{The problem and the proof strategy}
\label{j:introduction}

Let $\mathcal P$ be the probability laws on $\R$ with mean zero,
variance one and finite third absolute moment $\rho(P)=\E_P|X|^3$.
For independent copies with this same law, write $S_n=X_1+\cdots+X_n$
and define
\[
 \Delta_n(P)=\sup_{H\in\R}|\Pp_P(S_n\le H)-\Phi(H/\sqrt n)|,
 \qquad C_n=\sup_{P\in\mathcal P}\frac{\sqrt n\Delta_n(P)}{\rho(P)}.
\]
The threshold $H$ is in the units of the raw sum; $H/\sqrt n$ is its
normal-standardized value. The Berry--Esseen inequality
\cite{Berry1941,Esseen1942,Esseen1945} asserts a universal finite upper
bound on the $C_n$. The sharp-constant problem asks whether this bound
can be taken to be
\[
 \CE=\frac{3+\sqrt{10}}{6\sqrt{2\pi}}=0.409732\ldots.
\]
The lower bound comes from one fixed skew Bernoulli law, rather than
from a sequence of progressively more complicated distributions.
The difficulty is to bound every possible marginal, uniformly in the
sample size.

\begin{theorem}[Sharp iid Berry--Esseen bound]\label{j:main}
For every $P\in\mathcal P$ and every integer $n\ge1$,
\[
 \Delta_n(P)\le\frac{\CE\rho(P)}{\sqrt n}.
\]
The constant is optimal: $\sup_{n\ge1}C_n=\CE$.
Moreover $C_n<\CE$ for every $n\ge47$.
\end{theorem}
The proof combines the analytic arguments below with the finite
inequalities specified by their certificates. Section~\ref{j:composition}
assembles these results and proves strictness using the equality case
of the uniform local exclusion.

\subsection{The three parts of the upper-bound argument}

The proof first establishes $C_j\le\CE$ for $1\le j\le42$.
Suppose a later failure exists and let $n$ be its least index. Then both
deleted predecessors satisfy $C_{n-1},C_{n-2}\le\CE$. A Gaussian
one-step estimate, proved in Section~\ref{j:smoothing:section}, gives
\[
 0<n(C_n-\CE)<\frac{3003}{10250}<0.3,\qquad
 C_n<\frac{73}{175},\qquad n(C_{n-1}-C_n)_+=0.
\]
Two local comparisons exclude precisely such an extremizer: one covers
the finite range $43\le n\le48$, and one is uniform for $n\ge49$.
Thus the finite base and these two comparisons prove the upper bound
at every index. The role of the recurrence is to supply the hypotheses
of the local comparisons, rather than to propagate the sharp constant
by itself.

The mathematical object shared by the comparisons is a nonnegative
support-gap defect. At a fixed index, the optimal ratio is attained.
Its first variation gives a piecewise cubic function which bounds the
conditional probability $y\mapsto\Pp(S_{n-1}\le H-y)$ and agrees
with it on the support. Integrating this contact relation against the
zero-bias density gives a conservation identity. Its nonnegative terms
measure the score above $\CE$, the departure from the extremal
two-point geometry, and the probability entering support gaps. Its
remainder is bounded by Fourier estimates.

Small remainder therefore forces the same extremizing law near two
clusters. Conditioning on the signs of its summands gives an exact
binomial count plus centered conditional noise. The local comparison
keeps this noise and shows that it cannot improve the score enough to
violate the sharp bound. At zero noise, the actual law is Bernoulli and
a published binomial theorem applies. At positive noise, the errors
must remain proportional to the noise, however small it is.

The smallest indices require different quantitative estimates, but
use the same logic: first acquire a valid support and moment domain,
then majorize a probability on that domain. No step assumes that an
unrestricted extremizer has two atoms, or replaces the iid marginal by
an optimizer of a moment relaxation.

\subsection{Where the computer enters}

The reductions leave explicit inequalities on bounded parameter
regions. These include polynomial majorants, enclosures of Fourier
integrals, and inequalities in the probability, noise and phase of a
binomial count. A certificate supplies finite data allowing these
inequalities to be checked on the whole region. Typical data are
rational polynomial coefficients and a finite tree of closed boxes.
The checker proves a valid inequality or an exclusion on each leaf
and verifies that no part of the root domain was lost.

Two mathematical questions precede this computation: why every
possible witness belongs to the specified domain, and why the checked
inequality gives the desired probability bound. Both are proved in
the text. The certificates supply the remaining finite arithmetic.
They are not samples from a class of probability laws. Nor does
numerical precision alone make them rigorous: rational arithmetic and
proved outward enclosures control every rounding error.

Finding a certificate and verifying it are different tasks. For some
majorants, floating-point linear programming proposes coefficients;
for enclosure and count bounds, adaptive subdivision proposes a
finite proof tree. These procedures may fail and need not be trusted.
The acceptance program reconstructs the precise inequality, checks
all of its cells using exact arithmetic, and rejects unresolved cases.
Section~\ref{j:cert:section} explains these constructions and the
common soundness argument. The reproduction package contains the
large coefficient arrays and complete checking programs. Their size
does not enlarge the mathematical assumptions of the proof.

\subsection{Published inputs and sharpness}

For clarity, the external probability estimates used below are stated
here with their full law class.
\begin{proposition}[Affine Berry--Esseen estimates]\label{j:external:affine}
For every $P\in\mathcal P$ and integer $n\ge1$,
\[
 \sqrt n\Delta_n(P)\le
 \min\left\{0.4690\rho,\ 0.3322(\rho+0.429),\
                         0.3031(\rho+0.646)\right\}.
\]
\end{proposition}
These are the iid estimates of \citet{Shevtsova2013}. Standardizing
the summands converts the total-variance normalization there into the
displayed one. Only a finite third absolute moment is assumed. A
putative large normalized score consequently has an explicit upper
bound on $\rho$, which starts several support arguments. All
terminating decimals in this paper denote the corresponding rational
numbers.

\begin{proposition}[Finite binomial comparison and the lower limit]
\label{j:external:binomial}\label{j:external:limit}
For $0<p<1$, let $X=(B-p)/\sqrt{p(1-p)}$ with
$B\sim\operatorname{Bernoulli}(p)$. At every finite $n\ge1$ its
normalized discrepancy is strictly smaller than $\CE$.
For $p_*=(4-\sqrt{10})/2$ and this fixed marginal,
\[
 \lim_{n\to\infty}\frac{\sqrt n\Delta_n(P_*)}{\rho(P_*)}=\CE.
\]
\end{proposition}
The finite assertion is the binomial result of
\citet[Theorem~1]{Schulz2016}; its fixed-law limit is Esseen's lattice
asymptotic \cite{Esseen1956}, in the normalization recorded in
\citet[Remarks~4.1--4.2]{Schulz2016}. Indeed, with $q_*=1-p_*$,
\[
 \rho(P_*)=\frac{p_*^2+q_*^2}{\sqrt{p_*q_*}},\qquad
 \frac{3+q_*-p_*}{6\sqrt{2\pi}(p_*^2+q_*^2)}=\CE.
\]
Thus every coefficient smaller than $\CE$ fails for a sufficiently
large sum of this one law. No interchange of a supremum over laws
with a limit is needed. Together with the upper bound, this proves
that the universal constant is sharp.

The present argument reorganizes the earlier certificate proof
\cite{kolodziejekrodziewicz2026exact}. Its purpose is to give each
analytic reduction once and to expose the precise interfaces between
those reductions and the finite checks. The classical zero-bias
transform \cite{GoldsteinReinert1997} explains the conservation
identity; explicit contact and same-law domain control are needed to
apply it here. The final result concerns the optimal iid constant,
not eventual exact two-point optimality.

\section{Extremizers, two-point geometry and conservation}\label{j:core:section}
Throughout, $X,X_1,X_2,\ldots$ are iid, $\E X=0$, $\E X^2=1$,
$\rho=\E|X|^3<\infty$, and $S_n=\sum_{i=1}^nX_i$. Write
$\beta=\E[X|X|]$, $\gamma=\E X^3$, and $\phi=\Phi'$.
A threshold $H$ always refers to the raw sum; $z=H/\sqrt n$ is its
standardized value. Define the ratio and an auxiliary excess by
\[
 r(P,H)=\frac{\sqrt n}{\rho}\{\Pp(S_n\le H)-\Phi(z)\},\qquad
 a_c(P,H)=\sqrt n\{\Pp(S_n\le H)-\Phi(z)\}-c\rho.
\]
We maximize the ratio. The auxiliary quantity $a_c$ is used only to
state an identity for arbitrary coefficients and off-contact laws.

\subsection{Ratio attainment and contact, including atoms}
\begin{theorem}\label{j:core:attainment}
For every $n\ge2$, $C_n$ is attained as a positive closed-CDF ratio by
a compactly supported standardized law and a real raw threshold $H$.
At such a pair put $c=C_n$, $F=\Pp(S_n\le H)$,
$h(y)=\Pp(S_{n-1}\le H-y)$ and $\phi_z=\phi(z)$. Then
\begin{align}
 I_c(y)={}&n^{3/2}(h(y)-F)+ny\phi_z
 +\frac{H\phi_z+3c\rho}{2}(y^2-1)
 -c(|y|^3-\rho-3y\beta)\le0,\label{j:core:contact}\\
 Q(y)={}&F+\frac{H\phi_z+c\rho}{2n^{3/2}}
 -\frac{n\phi_z+3c\beta}{n^{3/2}}y
 -\frac{H\phi_z+3c\rho}{2n^{3/2}}y^2
 +\frac{c}{n^{3/2}}|y|^3.\label{j:core:Q}
\end{align}
Thus $\Delta=Q-h\ge0$ everywhere and $\Delta=0$ on the entire closed
support. No sign assumption on $H$ is needed.
\end{theorem}
\begin{proof}
Reflection takes a negative error to a positive closed error because
\[
 \Pp(-S_n\le-H)-\Phi(-z)
 =\Phi(z)-\Pp(S_n<H)\ge\Phi(z)-\Pp(S_n\le H).
\]
Hence $C_n=\sup r(P,H)$. Since $\rho\ge1$, $C_n\le\sqrt n$;
symmetric Bernoulli at $H=n$ shows $C_n>0$.
Choose $r_j\uparrow C_n$ with $r_j\ge C_n/2$. Then
$\rho_j\le2\sqrt n/C_n$, and the positive discrepancies are bounded
away from zero. Chebyshev on the negative tail and the normal upper
tail on the positive side bound $H_j$. Pass to $P_j\Rightarrow P$,
$H_j\to H$, $\rho_j\to b$. A common third-moment ceiling $M$ bounds
the first- and second-moment tails by $M/K^2$ and $M/K$, so the limit
retains mean zero and variance one. Lower semicontinuity gives
$\rho(P)\le b$. For the sum laws, the closed-set inequality, first at
$H+t$ and then as $t\downarrow0$, gives
\[
 \Pp(S_n\le H)-\Phi(H/\sqrt n)
 \ge\limsup_j\{\Pp(S_{n,j}\le H_j)-\Phi(H_j/\sqrt n)\}
 =C_n b/\sqrt n.
\]
Thus $r(P,H)\ge C_n b/\rho(P)\ge C_n$, proving attainment and
$\rho(P)=b$. This does not assume uniform integrability of third moments.

Contaminate the raw law by $(1-\epsilon)P+\epsilon\delta_y$, then
recenter and rescale. Its mean is $\epsilon y$, variance
$v_\epsilon=1+\epsilon(y^2-1)-\epsilon^2y^2$, and centered third
absolute moment has derivative $|y|^3-\rho-3y\beta$.
Dominated differentiation uses $3|y|(|x|+|y|)^2$.
Evaluate the standardized law at the image of the \emph{fixed raw event}
$S_n\le H$. Its probability is a polynomial in $\epsilon$, with
derivative $n(h(y)-F)$ even at an atom collision. The normal argument
$(H-n\epsilon y)/\sqrt{nv_\epsilon}$ has derivative
$-\sqrt n y-H(y^2-1)/(2\sqrt n)$. The ratio derivative is
$I_c(y)/\rho$, and maximality makes it nonpositive.
Furthermore $\int I_c\,dP=0$. Since the closed-event function $h$ is
upper semicontinuous, strict negativity at a support point would hold
on a neighborhood of positive mass. Thus $I_c=0$ on the closed support.
Finally $I_c(y)\le-c|y|^3+O(y^2+|y|+1)$, with constants depending on
this pair, excludes both sufficiently distant rays.
\end{proof}

At a ratio maximizer, put $D_0=F-\Phi(z)$ and $\lambda=D_0/\rho=C_n/\sqrt n$.
The majorant used in the finite calculations is simply
\begin{equation}\label{j:core:scaled-Q}
 nQ(y)=\lambda|y|^3-\frac{z\phi_z+3D_0}{2}y^2
 -(\sqrt n\phi_z+3\lambda\beta)y+nF+\frac{z\phi_z+D_0}{2}.
\end{equation}
Thus its equality on support is $n h=nQ$; it is not a majorant for a
different marginal. If $A\ge C_{n-1}$, this equality implies, at every
support point $t$,
\begin{equation}\label{j:core:predecessor}
 0\le\mathcal H_n(H,t)+C_n\{-|t|^3+\tfrac32\rho t^2+3\beta t+\rho e_n\}
       +\rho s_n n(A-C_n),
\end{equation}
where
\[
 s_n=\sqrt{n/(n-1)},\quad e_n=n(s_n-1)-\tfrac12,
\]
\[
 \mathcal H_n(H,t)=n^{3/2}\!\left\{\Phi\!\left(\frac{H-t}{\sqrt{n-1}}\right)
 -\Phi\!\left(\frac H{\sqrt n}\right)\right\}
 +nt\phi_z+\tfrac H2(t^2-1)\phi_z.
\]
Indeed substitute $h(t)\le\Phi((H-t)/\sqrt{n-1})+A\rho/\sqrt{n-1}$
in \eqref{j:core:contact} and use $\sqrt n(F-\Phi(z))=C_n\rho$.
The signed difference $A-C_n$ may be negative; replacing it by a
nonnegative upper bound gives a weaker permitted inequality.

For later induction, if $C_j\le c$ for $j<N$ and some $n\ge N$ has
$C_n\ge c$, the least such $n$ has an attained ratio pair with
$C_n\ge c$ and $C_j\le c$ for every $j<n$. This includes equality at the
selected index; it requires no supremum over varying sample sizes.

\subsection{One sign normalization and its moment constraints}
Put $p=\Pp(X>0)$, $q=1-p$, $v=pq$, and
$m=\E X_+=\E X_->0$. Assign a zero atom to the nonpositive class.
For $B=\ind_{\{X>0\}}$, let $V=\E[X\mid B]$, $W=X-V$, and
$\theta=\E W^2$. Then $V=-m/q$ or $m/p$, and
\begin{equation}\label{j:core:normalization}
 \ell=\frac m{pq}=\sqrt{\frac{1-\theta}{pq}},\quad
 U=W/\ell,\quad X=\ell(B+U-p)=\frac{B+U-p}{\sqrt{pq+s}},\quad
 s=\E U^2=\frac{pq\theta}{1-\theta}.
\end{equation}
Here $0<p,q<1$, $0\le\theta<1$ and $\E[U\mid B]=0$.
Write $U_i$ for the conditional laws, $w_i=\E U_i^2$,
$\mu_i=\E U_i^3$, and $\lambda=ns$. Thus $s=qw_0+pw_1$; if $s>0$,
$\alpha=pw_1/s\in[0,1]$ gives $w_0=s(1-\alpha)/q$, $w_1=s\alpha/p$.
If $s=0$, both conditional laws are zero and $X$ is standardized Bernoulli.
The exact common-marginal identities are
\begin{align}
 \rho/\ell^3&=v(1-2v)+3v(w_0+w_1)-q\mu_0+p\mu_1,\notag\\
 \gamma/\ell^3&=v(q-p)+3v(w_1-w_0)+q\mu_0+p\mu_1,\label{j:core:conditional-moments}\\
 \beta/\ell^2&=v(q-p)+pw_1-qw_0.\notag
\end{align}
They follow by expanding $\ell^3(p-U_0)^3$ and
$\ell^3(q+U_1)^3$ on the two sign classes and using conditional centering.
In particular the two cubic moments in the first two lines cannot be
selected independently. Higher moments require an established support
bound or another explicit integrability hypothesis.

\begin{lemma}\label{j:core:moments}
Every standardized law with finite third moment satisfies
\[
 \rho\ge\frac{\gamma^2+2}{\sqrt{\gamma^2+4}},\qquad
 |\beta|<1,\qquad \rho\ge\frac{1+\beta^2}{\sqrt{1-\beta^2}}.
\]
\end{lemma}
\begin{proof}
For the first inequality set $s_0=\sqrt{\gamma^2+4}$,
$a_0=(s_0-\gamma)/2$, $b_0=(s_0+\gamma)/2$ and
$c_0=\gamma(\gamma^2+6)/s_0^3$. The polynomial on each half-line
\[
 w(y)=|y|^3-c_0y^3-12y^2/s_0^3+6\gamma y/s_0^3+4/s_0^3
\]
has the factorizations, for $y,t\ge0$,
\[
 w(y)=\frac{2a_0^2}{s_0^3}(y-b_0)^2[2+(a_0+3b_0)y],\quad
 w(-t)=\frac{2b_0^2}{s_0^3}(t-a_0)^2[2+(3a_0+b_0)t].
\]
Taking expectations gives $0\le\E w=\rho-(\gamma^2+2)/s_0$.
For the second inequality let $a_\pm=\E X_\pm^2=(1\pm\beta)/2>0$.
Cauchy gives $m^2/a_++m^2/a_-\le1$, hence
$2m\le\sqrt{1-\beta^2}$, and
$m\rho\ge a_+^2+a_-^2=(1+\beta^2)/2$. These imply the result;
a zero atom only strengthens the first inequality.
\end{proof}

\subsection{A defect measuring departure from two points}
Let $K=\supp P$ and $\tau(t)=\E[X\ind_{\{X\ge t\}}]$.
Centering and signwise Tonelli give
\[
 \tau\ge0,\qquad \int\tau=1,\qquad 2\int|t|\tau(t)\,dt=\rho.
\]
On every bounded gap $(u,v)$ of $K$, $\tau$ has a constant interior value
$\tau_{u,v}$. Set
\begin{equation}\label{j:core:defect-definition}
 J=\sum_{(u,v)}\tau_{u,v}(v-u)^2,\qquad
 \sigma=\sqrt{10},\qquad D=(3+\sigma)\rho-\gamma-3J.
\end{equation}
The gaps are countable; the next bounds justify their sum even for
unbounded support.
\begin{theorem}\label{j:core:defect}
For every such law,
\[
 \gamma^2+4m^2\le\rho^2,\qquad J\le\rho+2m,\qquad
 D=(\sigma\rho-\gamma-6m)+3(\rho+2m-J)\ge0.
\]
Equality $D=0$ occurs exactly for
$X=(B-p_*)/\sqrt{p_*q_*}$, where
$p_*=(4-\sqrt{10})/2$ and $q_*=1-p_*$.
If $0\in K$, then $D\ge(\sigma-1)\rho$. Otherwise write the central
gap as $(-A_-,A_+)$ and define
\[
 T_\pm=\E X_\pm^3-mA_\pm^2,\quad
 R_\pm=\E X_\pm^2-mA_\pm,\quad J_o=J-m(A_-+A_+)^2.
\]
The quantities $T_\pm,R_\pm,J_o$ are nonnegative, and
\begin{align}
 J_o^\pm&\le T_\pm-2A_\pm R_\pm,\label{j:core:outer-gaps}\\
 D&=m\{\sqrt{\sigma+1}A_--\sqrt{\sigma-1}A_+\}^2
 +(\sigma-1)T_++(\sigma+1)T_-+3(T_++T_--J_o).\label{j:core:gap-defect}
\end{align}
All four terms in the last expression are nonnegative.
Moreover $1/(2\rho)\le m\le1/2$; if $D<2$, then
$T_++T_-\le D/2$ and $A_-+A_+\ge\sqrt{2-D}$.
\end{theorem}
\begin{proof}
Put $a_\pm=\E X_\pm^2$, $b_\pm=\E X_\pm^3$.
The signwise Cauchy inequalities give $a_\pm^2\le mb_\pm$ and
$m^2\le a_+a_-$. Thus $b_+b_-\ge m^2$, proving the circle inequality.
A gap lying on one side of zero satisfies
$(v-u)^2\le2\int_u^v|t|\,dt$. A central gap has weight $m$ and excess
$2mA_-A_+$ over this integral; since
$0\le\E[(X+A_-)(X-A_+)]=1-A_-A_+$, summation gives $J\le\rho+2m$.
Without a central gap it gives $J\le\rho$. Cauchy in $\R^2$ now gives
$\gamma+6m\le\sigma\rho$. Equality requires equality in both signwise
Cauchy bounds and in $m^2\le a_+a_-$: each sign is a single nonzero
atom and there is no zero atom. Writing its positive probability as $p$,
the remaining equality is $3(1-2p)=2p(1-p)$, whose unique solution
in $(0,1)$ is $p_*$. Conversely this two-point law gives equality.

For a positive outer gap,
$(v-u)^2\le2\int_u^v(t-A_+)\,dt$; summing yields
$J_o^+\le\E[X_+(X_+-A_+)^2]=T_+-2A_+R_+$.
The negative side is identical after reflection. Expansion using
$\sqrt{\sigma+1}\sqrt{\sigma-1}=3$ proves \eqref{j:core:gap-defect}.
Finally Cauchy gives $1\le2m\rho$ and $2m\le1$.
Since $\sigma-1>2$, $T_++T_-\le D/2$; hence
$m(A_-^2+A_+^2)=\rho-T_+-T_-\ge1-D/2$, which gives the last claim.
All gap summations use nonnegative Tonelli; endpoint atoms change no
Lebesgue integral.
\end{proof}
For the standardized Bernoulli law itself,
\[
 J=1/\sqrt{pq},\quad \rho_B=(1-2pq)/\sqrt{pq},\quad
 \gamma_B=(q-p)/\sqrt{pq},\quad
 D_B=(6+2\sigma)(p-p_*)^2/\sqrt{pq}.
\]
Thus the defect charges a quadratic displacement of $p$ as well as
conditional noise. The following bound charges them simultaneously.

\begin{lemma}[Joint defect budget]\label{j:core:joint-defect}
In a central gap put $\bar A_-=\ell p$, $\bar A_+=\ell q$,
$\Delta_\pm=\bar A_\pm-A_\pm\ge0$,
$a_0=\Delta_-/\ell$, $a_1=\Delta_+/\ell$, $d=a_0+a_1<1$, and
$w=1-\theta=pq\ell^2$. Then $U_0\le a_0$, $U_1\ge-a_1$, and
\begin{align}
 D\ge{}&\ell w(6+2\sigma)(p-p_*)^2
 +6\ell w\{d(1-d)+a_0a_1\}\notag\\
 &+\E[\{(\sigma-1)(X+2\bar A_+)+6A_+\}W^2;X>0]\notag\\
 &+\E[\{(\sigma+1)(-X+2\bar A_-)+6A_-\}W^2;X<0].\label{j:core:joint}
\end{align}
In particular, with $\theta_\pm=\E[W^2;\pm X>0]$,
\begin{align}
 D\ge{}&\frac{(6+2\sigma)(1-\theta)^{3/2}}{\sqrt{pq}}(p-p_*)^2
 +6m\{(A_-+A_+)(\Delta_-+\Delta_+)+\Delta_-\Delta_+\}\notag\\
 &+3(\sigma+1)A_+\theta_++3(\sigma+3)A_-\theta_-.\label{j:core:simultaneous}
\end{align}
\end{lemma}
\begin{proof}
Insert \eqref{j:core:outer-gaps} in the definition of $D$. The result is
\[
 D\ge(\sigma-1)M_{3,+}+(\sigma+1)M_{3,-}
 +6A_+M_{2,+}+6A_-M_{2,-}-6m(A_-^2+A_-A_++A_+^2),
\]
where $M_{j,\pm}=\E X_\pm^j$. Conditional centering gives
$M_{2,+}=m\bar A_++\E[W^2;X>0]$ and
$M_{3,+}=m\bar A_+^2+\E[(X+2\bar A_+)W^2;X>0]$,
and the reflected identities. Substitute $A_-=\ell(p-a_0)$,
$A_+=\ell(q-a_1)$; the deterministic part is precisely the first line
of \eqref{j:core:joint}. All kernels are nonnegative. Bounding them
below by $3(\sigma+1)A_+$ and $3(\sigma+3)A_-$ proves
\eqref{j:core:simultaneous}. Finite third moments suffice throughout.
\end{proof}

\subsection{Conservation: one identity, with and without contact}
For a compactly supported law retain $a=a_c$, $Q,h,\Delta$ from above,
and put $\nu=\mathcal L(H-S_{n-1})$. Over its bounded support gaps define
\[
 J_1=\sum\tau_{u,v}\!\int_u^v(t-u)t\,dt,\quad
 J_2=\sum\tau_{u,v}\!\int_u^v(t-u)t|t|\,dt,\quad
 L=\sum\tau_{u,v}\!\int_{[u,v)}(t-u)\,d\nu(t)\ge0.
\]
The half-open interval is essential: an atom at $v$ belongs to $h(v)$,
whereas an atom at $u$ receives zero weight. If $K\subset[-R,R]$, then
$J\le2R$, $|J_1|\le RJ/2$, $|J_2|\le R^2J/2$ and $L\le2R$.
Let $M_4=\E X^4$, $M_5=\E|X|^5$, and set
\[
 K_q=M_5/4-\rho M_4/2-3\beta\gamma/2+\rho/2,\quad
 K_1=\E[X|X|^3]-3\beta-3\rho\gamma/2,\quad
 \Omega=J_1+(M_4-3)/6.
\]
For $Z=S_n/\sqrt n$, a standard normal $Z_0$, and $q_z(x)=(x-z)|x-z|$, put
\[
 \varepsilon=\E q_z(Z)-\E q_z(Z_0)-2\gamma\phi_z/(3\sqrt n),\qquad
 E_n=-n^{3/2}\varepsilon/2,
\]
\begin{equation}\label{j:core:remainder}
 \mathcal R_c=c(3\beta J/2+3\rho J_1-3J_2-K_q)
 +H(cK_1+\phi_z\Omega)-H^2(c\rho+\gamma\phi_z/2)+E_n.
\end{equation}
Define the residuals
\[
 \xi=\E[X\Delta(X)],\qquad
 R_K=\int_K\tau(t)\Delta(t)\,dt
       +\sum_{(u,v)}\tau_{u,v}(v-u)\Delta(v).
\]
The integral on $K$ is a Lebesgue integral, not a $P$-integral. All terms
converge absolutely since $\Delta$ is bounded on the convex hull and
$\sum\tau_{u,v}(v-u)\le1$.

\begin{theorem}\label{j:core:conservation}
For every compactly supported standardized law and every $c>0$, $H\in\R$,
\begin{equation}\label{j:core:general-identity}
 (n+H^2)a+nc\rho-\frac{n\phi_z}{6}(\gamma+3J)+n^{3/2}L
 =\mathcal R_c+n^{3/2}(R_K-H\xi).
\end{equation}
At a ratio maximizer set $c=C_n$, so $a=0$. With
$e=C_n-\CE$ and $w=H^2/(2n)$, the residuals vanish and
\begin{equation}\label{j:core:positive-identity}
 \mathcal R_{C_n}=ne\rho+n\CE\rho(1-e^{-w})
       +n\phi(0)e^{-w}D/6+n^{3/2}L.
\end{equation}
All four terms are nonnegative when $C_n\ge\CE$, including equality.
\end{theorem}
\begin{proof}
Write $\mathcal A=\int\tau(Q-h)$. On a gap, $h(t)=h(v)+\nu([t,v))$.
Fubini and integration of $Q'$ therefore give
\begin{equation}\label{j:core:gap-integral}
 \mathcal A=R_K+\frac{\phi_zJ}{2\sqrt n}
 +\frac{3c\beta J/2+(H\phi_z+3c\rho)J_1-3cJ_2}{n^{3/2}}-L.
\end{equation}
For absolutely continuous $g$, signwise Fubini gives
$\int\tau g'=\E[Xg(X)]$. Thus
$\int t\tau=\gamma/2$, $\int t^2\tau=M_4/3$,
$\int|t|^3\tau=M_5/4$. Applying the same identity conditionally to
$g_s(x)=\min(s+x,H)$ and using exchangeability gives
\[
 \int\tau h=\frac1n\E[S_n\min(S_n,H)]
 =T,\qquad T=\E[(Z^2-zZ)\ind_{\{Z\le z\}}].
\]
The integrand at the threshold atom is zero. Consequently
\[
 \mathcal A=F-\frac{\gamma\phi_z}{2\sqrt n}
 +\frac{H\phi_z(3-M_4)/6+cK_q}{n^{3/2}}-T.
\]
If $M=\E[Z\ind_{\{Z\le z\}}]$, integrating $X(Q-h)$ against $P$ gives
\[
 M+\phi_z=-\gamma H\phi_z/(2n)+cK_1/n-\sqrt n\xi.
\]
Finally $q_z(x)=(x-z)^2-2(x-z)^2\ind_{\{x\le z\}}$, for both $Z$ and
$Z_0$, implies
\[
 T-\Phi(z)=-\varepsilon/2-\gamma\phi_z/(3\sqrt n)
               +z(M+\phi_z)-z^2(F-\Phi(z)).
\]
Substitute these two identities, use $\sqrt n(F-\Phi(z))=c\rho+a$,
and equate with \eqref{j:core:gap-integral}; collecting terms yields
\eqref{j:core:general-identity}. At contact $\Delta=0$ on the entire
support, so $R_K=\xi=0$. Substituting $a=0$, $c=C_n$ and
$\CE=(3+\sqrt{10})\phi(0)/6$ gives \eqref{j:core:positive-identity}.
\end{proof}
The off-contact identity permits checks on arbitrary bounded laws; it
is not a positivity statement there. In particular the support integral
in $R_K$ cannot be dropped for a continuous law. The proof uses contact
only when setting the two residuals to zero. Section~\ref{j:fourier:section}
proves uniform bounds for this exact $E_n$, including changing laws and
thresholds.

\begin{corollary}[The sign license]\label{j:core:license}
At a ratio maximizer with $C=C_n\ge\CE$, suppose a first pass proves
$\mathcal R_C\le T<nC\rho$. Then $\gamma+3J>0$ and
\begin{equation}\label{j:core:phase-price}
 n(C-\CE)\rho+n\phi(0)D/6+n^{3/2}L\le\mathcal R_C.
\end{equation}
\end{corollary}
\begin{proof}
The contact identity before splitting the defect gives
$\phi_z(\gamma+3J)/6=C\rho+\sqrt nL-\mathcal R_C/n>0$.
Thus $M=\CE\rho-\phi(0)D/6>0$. Regrouping
\eqref{j:core:positive-identity} writes its right side as the left side
of \eqref{j:core:phase-price} plus $n(1-e^{-w})M\ge0$.
\end{proof}
The initial bound is on the raw remainder. Only after it licenses the
positive phase term may that term improve a later resource bound.
No $D<2$ assumption or positive lower bound on $C_n-\CE$ was used.

For context, on a lattice with maximal span $h>0$, its gaps have lengths
$k_gh$, $k_g\ge1$. Since $\tau$ vanishes outside the convex hull and the
support has Lebesgue measure zero,
$\sum_g\tau_g k_gh=1$ and $J=\sum_g\tau_g k_g^2h^2\ge h$.
Reflection preserves $J,\rho,h$ and reverses $\gamma$, whence
$|\gamma|+3h\le(3+\sqrt{10})\rho$. Equality is precisely the zero-defect
Bernoulli law or its reflection. Thus the gap functional extends the
lattice span to arbitrary supports; no lattice hypothesis enters the proof.

\section{What a first failing index must satisfy}
\label{j:smoothing:section}

A simple recurrence supplies the
same small excess at any first failure of the sharp bound. This lets us
use local exclusions directly, without first proving a separate sequence
of positive upper barriers.

\begin{lemma}[Gaussian derivative bounds]\label{j:smoothing:gaussian}
The standard normal density satisfies
\[
 \|\phi''\|_\infty=\phi(0)<\frac25,
 \qquad \|\phi'''\|_\infty\le\frac2\pi,
 \qquad \frac1{4\pi}<\frac2{25},\qquad \CE<\frac{41}{100}.
\]
\end{lemma}
\begin{proof}
Since $\phi''(x)=(x^2-1)\phi(x)$, its absolute maximum is at zero:
the other candidates $x=\pm\sqrt3$ have value
$2e^{-3/2}\phi(0)<\phi(0)$. The absolutely convergent Gaussian
Fourier integral gives
\[
 |\phi'''(x)|\le\frac1{2\pi}
       \int_{\R}|t|^3e^{-t^2/2}\,dt=\frac2\pi.
\]
The bounds for $\phi(0)$ and $1/(4\pi)$ follow from $\pi>25/8$.
For the last comparison one may use
$\sqrt{10}<31623/10000$ and $\pi>314159/100000$, square the
positive sides, and compare rational numbers.
\end{proof}

\begin{theorem}[One-step smoothing]\label{j:smoothing:recurrence}
For every integer $n\ge2$,
\begin{equation}\label{j:smoothing:section-recurrence}
 C_n\le\sqrt{\frac n{n-1}}\,C_{n-1}
       +\frac{\sqrt n}{15(n-1)^{3/2}}
       +\frac{2\sqrt n}{25(n-1)^2}.
\end{equation}
\end{theorem}
\begin{proof}
Fix one standardized marginal $P$, write $v=n-1$, and put
$G_v(x)=\Phi(x/\sqrt v)$. Conditioning on the last summand yields
\[
 |\Pp(S_n\le H)-\E G_v(H-X)|\le C_{n-1}\rho/\sqrt v.
\]
Taylor expansion of the smooth Gaussian CDF, followed by expectation
under this same marginal, gives
\[
 |\E G_v(H-X)-G_v(H)-\tfrac12G_v''(H)|
 \le\frac{\phi(0)\rho}{6v^{3/2}}<\frac\rho{15v^{3/2}}.
\]
Here $\E X=0$ and $\E X^2=1$ cancel the first two Taylor terms.
For the change of variance, $\partial_vG_v=G_v''/2$ and
$\partial_v^2G_v=G_v^{(4)}/4$. Thus
\[
 G_{v+1}(H)-G_v(H)-\tfrac12G_v''(H)
 =\frac14\int_0^1(1-t)G_{v+t}^{(4)}(H)\,dt.
\]
Since $G_w^{(4)}(H)=w^{-2}\phi'''(H/\sqrt w)$, this remainder
has absolute value at most $1/(4\pi v^2)<2/(25v^2)$.
Combine the three estimates, multiply by $\sqrt n/\rho$, and use
$\rho\ge1$. The bound is uniform in $P$ and $H$ and treats both
signs of the discrepancy, so it passes to the two suprema defining $C_n$.
\end{proof}

\begin{corollary}[Entry from a sharp predecessor]
\label{j:smoothing:entry}
If $n\ge40$ and $C_{n-1}\le\CE\le C_n$, then
\[
 0\le n(C_n-\CE)<\frac{1433}{4875}<\frac3{10},
 \qquad C_n<\frac{167}{400}<0.4226,
 \qquad n(C_{n-1}-C_n)_+=0.
\]
For $n\ge42$ the excess bound improves to $3003/10250$ and
the score bound to $73/175$.
For every standardized marginal $P$, the same-law predecessor loss
$n(\sqrt{n-1}\Delta_{n-1}(P)/\rho(P)-C_n)_+$ is zero.
At a first failing index $n\ge43$ above a sharp base, both deleted
constants $C_{n-1}$ and $C_{n-2}$ are below $0.4253$.
\end{corollary}
\begin{proof}
Put $s=\sqrt{n/(n-1)}$. The recurrence and
$n(s-1)=\{n/(n-1)\}/(s+1)$ give
\[
 n(C_n-\CE)\le\frac n{n-1}
 \left\{\frac{\CE}{s+1}+\frac{s}{15}
                  +\frac{2s}{25\sqrt{n-1}}\right\}.
\]
For $n\ge40$, $n/(n-1)\le40/39$, $\sqrt{n-1}>6$, and
$1<s<51/50$, since $39\cdot51^2-40\cdot50^2=1439>0$.
The right side is therefore less than
\[
 \frac{40}{39}\left(\frac{41}{200}
                  +\frac{17}{250}+\frac{17}{1250}\right)
 =\frac{1433}{4875}<\frac3{10}.
\]
This gives $C_n<0.41+0.3/40=167/400$.
For $n\ge42$, $n/(n-1)\le42/41$ and
$1<s<81/80$; the last comparison follows from
$41\cdot81^2-42\cdot80^2=201>0$.
Using $\CE<41/100$ and $s+1>2$, the right side is less than
\[
 \frac{42}{41}\left(\frac{41}{200}
                    +\frac{27}{400}+\frac{27}{2000}\right)
 =\frac{3003}{10250}<\frac3{10}.
\]
Consequently $C_n<0.41+0.3/42=73/175$.
Every law's predecessor ratio is bounded by $C_{n-1}\le\CE\le C_n$,
so both loss assertions follow. At a first failure the two deleted
indices are sharp and $\CE<0.41<0.4253$.
\end{proof}

\subsection{The loose ceilings needed after index twenty-four}
\label{j:smoothing:finite-ceilings}
\begin{lemma}\label{j:smoothing:entry25}
If $n\ge25$ and $C_{n-1}\le\CE\le C_n$, then
\[
 n(C_n-\CE)<\frac{4177}{13824}<\frac{303}{1000},
 \qquad C_n<\frac{145873}{345600}<\frac{423}{1000}.
\]
For each marginal the same-law predecessor allowance is zero.
\end{lemma}
\begin{proof}
Set $s=\sqrt{n/(n-1)}$. Theorem~\ref{j:smoothing:recurrence} gives
\[
 n(C_n-\CE)\le\frac{n}{n-1}
 \left(\frac{\CE}{s+1}+\frac{s}{15}
                       +\frac{2s}{25\sqrt{n-1}}\right).
\]
For $n\ge25$ one has $n/(n-1)\le25/24$,
$1<s<49/48$ and $\sqrt{n-1}>24/5$.
The latter strict comparisons follow from
$24\cdot49^2-25\cdot48^2=24>0$ and $24>576/25$.
Together with $\CE<41/100$, these imply
\[
 n(C_n-\CE)<\frac{25}{24}
  \left(\frac{41}{200}+\frac{49}{720}+\frac{49}{2880}\right)
 =\frac{4177}{13824}.
\]
Dividing by $n\ge25$ proves the score bound. The exact positive
comparison margins are
$303/1000-4177/13824=1459/1728000$ and
$423/1000-145873/345600=1579/1728000$.
Finally each marginal's predecessor score is at most $C_{n-1}\le\CE\le C_n$.
\end{proof}

The entry bounds concern the same marginal as the attained ratio.
They license weaker own-score and predecessor-loss inputs used in the
finite calculations. Transferring a calculation from a floor index to
a later index additionally requires the actual-index support and weighted
resource statements proved below; the scalar recurrence alone does not
transfer a support bound.

\section{What a certificate proves}
\label{j:cert:section}

A certificate in this proof is finite data from which a small collection
of exact inequalities can be checked. It may specify coefficients of a
polynomial majorant, a tree of parameter boxes, or a schedule of bounds
which successively restricts an extremizer. The numerical statements are
about whole regions, not sampled distributions. The mathematical
reduction must first show that every still possible extremizer maps into
one of those regions.

Here is a typical example. Suppose a preceding argument has confined the
support of the same marginal to $K$ and its moment vector to $M$.
For a threshold interval $[a,b]$, a certificate specifies a function $T$
such that
\[
 T(x_1,\ldots,x_n)\ge\ind_{\{\sum_i x_i\le b\}}
 \quad(x_1,\ldots,x_n\in K),\qquad
 \E T(X_1,\ldots,X_n)\le L_a+\frac{c\rho}{\sqrt n},
\]
where the second inequality holds for every moment vector in $M$ and
$L_a\le\Phi(a/\sqrt n)$. For every $H\in[a,b]$, monotonicity gives
\[
 \Pp(S_n\le H)-\Phi(H/\sqrt n)
 \le\E T-L_a\le c\rho/\sqrt n.
\]
Thus verification proves a probability inequality, including threshold
atoms. Choosing $a$ for the normal lower bound and $b$ for the closed
event is essential. This is how the small-summand certificates work;
the count comparisons apply the same principle after conditioning on
the signs of the summands.

\subsection{A common acceptance argument}

\begin{lemma}[Expectation under one iid marginal]\label{j:cert:tensor}
Let $f_0=1,f_1,\ldots,f_d$ be absolutely integrable under $P$, and put
$m_j=\E f_j(X)$. For a finite tensor sum
\[
 T(x_1,\ldots,x_n)=\sum_\alpha t_\alpha
                         \prod_{i=1}^n f_{\alpha_i}(x_i),
 \qquad M_T(m)=\sum_\alpha t_\alpha\prod_i m_{\alpha_i},
\]
one has $\E T(X_1,\ldots,X_n)=M_T(m)$. If $T$ majorizes the
closed sum event on the support product, a uniform upper bound on
$M_T$ gives the corresponding upper bound on its probability.
\end{lemma}
\begin{proof}
Independence factors the absolute expectation of each product into a
product of finite absolute moments. Fubini and linearity therefore
evaluate every term. Integrating the pointwise majorant proves the
last assertion.
\end{proof}
The same $m$ occurs in every factor. In particular, a factor with zero
mean makes its tensor term have zero expectation. Such terms can improve
the pointwise inequality without changing its expected value. Features
of degree above three are used only after support has been bounded, so
their integrability does not strengthen the theorem's moment assumption.

\begin{lemma}[Exact positivity on a box]\label{j:cert:bernstein}
Let $p(t)=\sum_\alpha a_\alpha t^\alpha$ have coordinate degrees
at most $d_i$. Its tensor Bernstein coefficients on $[0,1]^r$ are
\[
 b_\nu=\sum_{\alpha\le\nu}a_\alpha
       \prod_{i=1}^r\frac{\binom{\nu_i}{\alpha_i}}{\binom{d_i}{\alpha_i}}.
\]
If all $b_\nu\ge0$, then $p\ge0$ on the entire closed box.
\end{lemma}
\begin{proof}
The binomial theorem gives
\[
 t^k=\sum_{j=k}^d\frac{\binom jk}{\binom dk}
                       \binom dj t^j(1-t)^{d-j}.
\]
Apply this in each coordinate. The resulting basis functions are
nonnegative and sum to one, so $p(t)$ is a convex combination of
the $b_\nu$.
\end{proof}
An affine rational substitution treats any rational box. On a simplex,
the analogous multinomial basis in barycentric coordinates is
nonnegative and sums to one, giving the same criterion after an exact
change of basis. Subdivision can improve an inconclusive coefficient
test. It never turns that test into sampling.

\begin{proposition}[Coverage and exclusion]\label{j:cert:cover}
Suppose an analytic argument maps every law and threshold under
consideration into a finite union of closed root cells. At each root,
let a finite tree subdivide a parent into closed children whose union
contains the parent. On each leaf suppose either a proved necessary
condition excludes every actual mapped point, or a valid inequality
establishes the desired score bound. Then the bound holds for every
mapped law and threshold. If all leaves exclude a putative violating
extremizer, no such extremizer exists.
\end{proposition}
\begin{proof}
A point in a root belongs to a child at each subdivision, including on
a shared face. Finiteness leads to a leaf, where its conclusion
applies. The root cover includes every actual point.
\end{proof}

There are three distinct obligations in this proposition: the map from
laws to parameters, validity of the leaf inequalities, and coverage.
A relaxed moment domain may contain unrealizable points; this only
makes verification harder. Missing an actual point invalidates the
argument. In particular, we split sign orthants and event planes
explicitly: on $\sum x_i=b$ the required inequality is $T-1\ge0$,
not just $T\ge0$. Non-strict inequalities in the domain retain endpoint
atoms. Resource limits, unresolved leaves and failed divisions cannot
be interpreted as acceptance.

\subsection{Exact evaluation and rigorous enclosures}
\label{j:cert:arithmetic}
The acceptance programs evaluate rational numbers by integer arithmetic.
Irrational constants are enclosed between rationals with proved error
bounds. For example, alternating arctangent series in Machin's identity
\[
 \frac\pi4=4\arctan(1/5)-\arctan(1/239)
\]
bound $\pi$; integer square roots then bound its required square roots.
Integrating the Taylor expansion of $e^{-t^2/2}$ gives rational bounds
for the Gaussian CDF at rational endpoints. Taylor's remainder, rather
than a printed decimal approximation, supplies the error sign and size.
The same principle bounds exponential, trigonometric and binomial
expressions used below. Each application states its finite interval or
separate tail estimate. Products and quotients use outward bounds, with
denominator positivity established on the whole cell. Multiplying a
negative expression by a lower bound requires the corresponding reversal
of inequality; the signed comparisons retain this distinction.

These statements are proved inequalities in ordinary analysis. Their
finite arithmetic evaluation is performed by the programs. Exact
arithmetic removes round-off uncertainty, but does not by itself prove
coverage, correctness of the formulas, or their applicability to a
particular marginal. The preceding lemmas and the applications below
supply those mathematical obligations. This separation is also used in
other computer-assisted inequality proofs; formal interval verification
provides a stronger implementation assurance when available
\cite{SolovyevHales2013}. The certificate format here is specific to this
argument, and is not a standard proof-assistant format.

\subsection{How the certificates are found and checked}
\label{j:cert:discovery}
For the two-summand majorants, the unknown tensor coefficients enter
linearly. The construction program solves floating-point linear
programs imposing the proposed majorant and expectation inequalities
on an initial finite set of points. It searches for violated constraints,
adds them, and solves again. Coefficients are converted to rational
numbers and passed to the exact checker. The three-summand construction
uses the same linear-programming and separation approach, with extra
zero-expectation tensor features. Success of the sampled optimization
is only a way of proposing a polynomial; the whole-domain check is the
proof. If rationalization loses a margin, the candidate must be improved
or rejected.

The resource and count certificates are obtained by adaptive subdivision.
Starting with the analytically licensed root domains, construction
applies valid bounds, discards infeasible cells, and splits cells whose
inequality is still undecided. A successful finite forest records a
complete cover and enough information to re-evaluate every leaf. Some
modules recompute their entire enclosure schedule during acceptance
rather than trusting recorded outputs. The schedules may have been
chosen by experimentation; their acceptance depends only on the stated
inequalities and coverage. The method does not require a theorem that
this search always terminates. The particular finite accepted objects
suffice.

The distributed candidates are treated as untrusted data. Acceptance
reconstructs the defined polynomials or bound expressions, checks their
identities and complete domains, and evaluates the exact inequalities.
It does not accept a stored success flag or an optimizer's declaration
of infeasibility. The public count candidates were verified in this
way; their historical search was not rerun. This has no bearing on the
logical validity of a verified certificate, just as finding a factor
and checking its product are different tasks.

\subsection{Proof data and reproduction}
\label{j:cert:reproduction}
The companion reproduction package identifies each certificate family,
its mathematical conclusion, checking source and complete execution.
The fixed reference is commit
\texttt{c31d0cbe2af867d26871fbae37a3a7f4066f4a55} of the
\href{https://github.com/kolodziejek/berry-esseen-certificates/tree/c31d0cbe2af867d26871fbae37a3a7f4066f4a55/manuscript/extended-v17/reproduction}
{public certificate repository}, in
\texttt{manuscript/extended-v17/reproduction}.
Its \texttt{REPRODUCE.md} gives the restoration and acceptance commands;
\texttt{COMPLETION\_LEDGER.json} identifies the proof objects and
completed result records. The mathematical interfaces used in this
revision are stated in the numerical propositions below.
It preserves the exact accepted candidates and their domain interfaces.
A reader can inspect the analytic arguments independently of running
the programs, and can verify the finite arithmetic independently of
repeating the search. Both parts are necessary for the whole proof.
The completed executions use Python with assertions enabled, with integer and rational
arithmetic; the three-summand kernel uses C++17 and
Boost.Multiprecision integers. The compiler, language runtimes and
faithful execution are part of the trusted computational environment.
The implementation is not a proof-assistant formalization.

The full execution history, command lines and machine identities belong
to that package rather than the mathematical argument. They distinguish
complete native calculations from later checks of their input
interfaces, and preserve failed or interrupted attempts. All selected
numerical obligations in the reference proof have completed; the
separate 114-command recipe has not been rerun as one relocated
campaign. This revision changes the exposition, not those certificate
bytes or acceptance targets. Any new mathematical extension of a
certificate's scope must be justified explicitly and cannot follow from
its previous success alone.

\section{Fourier bounds for one marginal, uniformly in the index}\label{j:fourier:section}
Two transformations are needed. A compensated inversion formula bounds the
smooth-payoff error $E_n$ in conservation; Prawitz's inequality bounds the
CDF itself. Both use characteristic functions of the same marginal. The
formulas below are pointwise inequalities with explicit moment and support
hypotheses, so they also cover triangular arrays of changing laws. The finite
certificates evaluate these formulas on closed cells; they do not supply an
asymptotic or infinite-domain analytic argument.

\begin{lemma}[Gaussian normalizations]\label{j:fourier:constants}
One has $\|\phi''\|_\infty=\phi(0)<2/5$,
$\|\phi'''\|_\infty\le2/\pi$, and
\begin{gather*}
 \frac{333}{106}<\pi<\frac{22}{7},\quad
 \sqrt\pi>\frac{4431}{2500},\quad
 \sqrt{\pi/2}<\frac{251}{200},\\
 \phi(0)>0.39886>\frac{997}{2500},\quad
 \frac{819}{2000}<\CE<\frac{20487}{50000}.
\end{gather*}
\end{lemma}
\begin{proof}
Differentiation and integration by parts in the Gaussian integral give
$\phi(x)=(2\pi)^{-1}\int_\R e^{-t^2/2-itx}dt$; its normalization follows
by squaring the integral and using polar coordinates. Absolute convergence gives
\[
 \|\phi^{(k)}\|_\infty\le\frac1{2\pi}\int_\R |t|^ke^{-t^2/2}dt.
\]
This proves the third derivative bound. For the second derivative,
$(x^2-1)e^{-x^2/2}$ has absolute maximum one (the other candidate is
$2e^{-3/2}<1$).
For the rational constants put
$s_k(x)=\sum_{j=0}^{k-1}(-1)^jx^{2j+1}/(2j+1)$.
Tangent addition, with $0<4\arctan(1/5)-\arctan(1/239)<\pi/2$,
gives Machin's identity $\pi=16\arctan(1/5)-4\arctan(1/239)$.
The geometric remainder in $1/(1+t^2)$, integrated from zero, gives
\[
 \pi_-=16s_6(1/5)-4s_1(1/239)<\pi
 <16s_5(1/5)-4s_2(1/239)=\pi_+.
\]
Exact rational comparison gives $333/106<\pi_-<\pi_+<22/7$ and
$\pi_->(4431/2500)^2$, $\pi_+<2(251/200)^2$.
Also $(0.39886)^2<7/44$.
Finally, for $r_-=3162277660/10^9$, $r_+=3162277661/10^9$,
$r_-^2<10<r_+^2$ and
\[
 (819/2000)^2<\frac{19+6r_-}{72\pi_+}
 <\CE^2<\frac{19+6r_+}{72\pi_-}<(20487/50000)^2.
\]
All denominators are positive; the displayed finite sums specify the
comparisons as integer cross-multiplications.
\end{proof}

For $n\ge2$, write
\[
 f(t)=\E e^{itX},\quad e_n(t)=e^{-nt^2/2}(1-in\gamma t^3/6),\quad
 V_n(t)=f(t)^n-e_n(t),
\]
\[
 I_0(H)=\int_0^{1/2}\frac{|\Im(e^{-itH}V_n(t))|}{t^3}\,dt,\qquad
 I_\infty=\int_{1/2}^\infty\frac{|V_n(t)|}{t^3}\,dt.
\]
\subsection{Compensated inversion, including atoms}

The compensation below makes absolute Fubini valid, including at atoms.

\begin{lemma}\label{j:fourier:inversion}
If a finite signed measure $\mu$ has mass and first moment zero and
$\int x^2\,d|\mu|(x)<\infty$, then, for every $H\in\mathbb R$,
\[
 \int(x-H)|x-H|\,d\mu(x)
 =-\frac4\pi\int_0^\infty
 \frac{\Im(e^{-itH}\widehat\mu(t))}{t^3}\,dt.
\]
The integral on the right is absolutely convergent.
\end{lemma}
\begin{proof}
For real $y$, substitution and integration by parts give
\[
 \int_0^\infty\frac{ty-\sin(ty)}{t^3}\,dt=\frac\pi4 y|y|.
\]
The boundary terms vanish: near zero the numerator is $O(t^3)$, and
at infinity it is $O(t)$. The integrand has the sign of $y$, since
$v-\sin v\ge0$ for $v\ge0$. To calculate the constant, nonnegative
Tonelli and elementary Laplace integration give
\[
 \int_0^\infty\frac{1-\cos t}{t^2}\,dt
 =\int_0^\infty s\left(\frac1s-\frac{s}{s^2+1}\right)ds
 =\frac\pi2.
\]
The absolute double integral of the compensated expression against
$d|\mu|(x)$ equals
\[
 \frac\pi4\int(x-H)^2d|\mu|(x)<\infty.
\]
Absolute Fubini is therefore valid. The linear compensation disappears
because $\int(x-H)d\mu=0$, proving both assertions. No continuity of a
CDF at $H$ is involved.
\end{proof}

Apply this to the sum law minus the scale-$\sqrt n$ pushforward of
the signed density
\[
 \phi(x)\{1+\gamma(x^3-3x)/(6\sqrt n)\}.
\]
This signed density has mass one, first moment zero and finite absolute
second moment. Its transform after scaling is $e_n$. Gaussian integration
by parts, using $q_z(x)=(x-z)|x-z|$ and $q_z'''=4\delta_z$, yields
\[
 \varepsilon_n:=\E q_{H/\sqrt n}(S_n/\sqrt n)
 -\E q_{H/\sqrt n}(Z)-\frac{2\gamma\phi(H/\sqrt n)}{3\sqrt n}
 =-\frac4{\pi n}\int_0^\infty
 \frac{\Im(e^{-itH}V_n(t))}{t^3}\,dt.
\]
Consequently the error $E_n=-n^{3/2}\varepsilon_n/2$ satisfies
\begin{equation}\label{j:fourier:ledger}
 |E_n|\le\frac{2\sqrt n}{\pi}\{I_0(H)+I_\infty\}.
\end{equation}
The identity and bound hold for each $n$ and $H$ separately.

\subsection{Characteristic-function bounds from common moments}
Let $X$ be centered with $\E X^2=1$, $\rho=\E|X|^3<\infty$,
$\gamma=\E X^3$, and $f(u)=\E e^{iuX}$. In this section put
\[
 \begin{gathered}
 h=\E|X|,\quad m=\E X_+=\E X_-=h/2,\quad
 A=\E X_+^2,\quad B=\E X_-^2,\\
 \zeta=2AB/m,\qquad b(r)=\frac{\sqrt{r^2+8}-r}{2}.
 \end{gathered}
\]
Here $X_-=\max\{-X,0\}$, so $A+B=1$ and all of $h,m,A,B$ are
positive. The letter $b$ is used for this root function to avoid confusion
with the signed second moment $\beta=\E[X|X|]$.

\begin{lemma}[Anchored characteristic function]\label{j:fourier:anchor}
For $0\le u\le\pi$,
\[
 \cos u\le\Re f(u)\le\cos u+(\rho-1)u^3/6,
 \qquad
 |\Im f(u)|^2\le\frac{u^6}{36}(\rho-1)(\rho+5/3).
\]
\end{lemma}
\begin{proof}
Write $\operatorname{sinc}t=\int_0^1\cos(st)\,ds$. For real $a,b$,
\[
 |\operatorname{sinc}a-\operatorname{sinc}b|
 \le |a-b|/2,\qquad
 |\operatorname{sinc}a-\operatorname{sinc}b|
 \le |a^2-b^2|/6.
\]
The first follows from the Lipschitz bound for cosine. For the second,
$v\mapsto\cos\sqrt v$ has derivative of absolute value at most $1/2$
on $v\ge0$, including its limit at zero; integrate with respect to $s$.
For $r\ge0$, integration from $1$ to $r$ of the first estimate gives
\[
 \left|\cos(ur)-\cos u+\frac{u\sin u}{2}(r^2-1)\right|
 \le \frac{u^3}{12}(r-1)^2(2r+1).
\]
Since $\E(|X|-1)^2(2|X|+1)=2(\rho-1)$, averaging proves the
upper real bound. For the lower bound, when $0<u\le\pi$ set
$a=\sin u/(2u)$ and $J(v)=\cos v+av^2$. The function
$\operatorname{sinc}v$ decreases on $[0,\pi]$, so $J$ decreases up to
$u$ and increases from $u$ to $\pi$. For $v\ge\pi$,
$J(v)\ge-1+a\pi^2=J(\pi)\ge J(u)$. Evenness gives the same conclusion
for negative $v$. Average $J(uX)\ge J(u)$ and use $\E X^2=1$.

The second sinc estimate, together with $\E X=0$, gives
$|\Im f(u)|\le u^3\E[|X|\,||X|^2-1|]/6$.
Put $W(r)=(r-1)^2(2r+1)$. Cauchy's inequality applies to the factorization
$r|r^2-1|=\sqrt{W(r)}\,r(r+1)/\sqrt{2r+1}$.
The exact identity
\[
 r^3+\frac{29r^2+1}{18}
 -\frac{2r^2(r+1)^2}{2r+1}
 =\frac{(r-1)^2(4r+1)}{18(2r+1)}\ge0
\]
bounds the second squared factor in expectation by $(\rho+5/3)/2$;
the first has expectation $2(\rho-1)$. All expectations are finite under
the stated third-moment hypothesis. The value at $u=0$ is immediate.
\end{proof}

\begin{lemma}[One common moment surrogate]\label{j:fourier:surrogate}
For an independent copy $X'$, put
$\eta=\E|X-X'|^3/2-\rho$. Then
\[
 \eta\le h\le1,\qquad \zeta\ge h,\qquad
 \rho+\zeta\ge2/h,\qquad \zeta\ge b(\rho),
 \qquad \gamma^2+\zeta^2\le\rho^2.
\]
For every real $u$ one also has
$36|\Im f(u)|^2\le u^6(\rho^2-\zeta^2)$.
If $|X|\le R$ almost surely, then additionally
\[
 \eta\le b(\rho)+\frac{(R-\sqrt{\rho^2+8})_+^2}{4R}.
\]
In particular $R\le\sqrt{\rho^2+8}$ implies $\eta\le b(\rho)$.
\end{lemma}
\begin{proof}
The scalar inequality
$|x-y|^3\le(x-y)^2(|x|+|y|)$, independence and centering give
$\E|X-X'|^3\le2\rho+2h$; Cauchy gives $h\le1$.
Signwise Cauchy gives
\[
 A^2\le m\E X_+^3,\quad B^2\le m\E X_-^3,
 \quad \frac{m^2}{A}+\frac{m^2}{B}\le\Pp(X\ne0)\le1.
\]
Thus $AB\ge m^2$, $\zeta\ge2m=h$, and
$\rho+\zeta\ge(A+B)^2/m=2/h$. Multiplication yields
$\zeta(\rho+\zeta)\ge2$, proving $\zeta\ge b(\rho)$.

For completeness, the sine-circle bound uses an exact mixture, not an
extreme-point assertion about convolution. Let $\mu_+$ and $\mu_-$ be the
restrictions of the law of $|X|$ to its positive and negative sign classes.
On $a,b>0$ define
\[
 d\nu(a,b)=\frac{a+b}{m}\,d\mu_+(a)d\mu_-(b).
\]
Integrating the centered law with atoms $a,-b$ and probabilities
$b/(a+b),a/(a+b)$ against $\nu$ recovers the nonzero part of the original
law; the zero atom is added separately. For this two-point law the second
sinc estimate implies
\[
 \frac{6|\Im f_{a,b}(u)|}{|u|^3}
 \le\frac{ab|a^2-b^2|}{a+b}.
\]
Together with $e_{a,b}=2a^2b^2/(a+b)$, the right side is one coordinate
of a vector of length $ab(a^2+b^2)/(a+b)$. Minkowski's inequality in
$\R^2$ therefore bounds the vector
$(\zeta,6\Im f(u)/u^3)$ by $\rho$. The third-moment assumption
justifies all integrals. Letting $u\to0$ proves the assertion for $\gamma$:
Taylor's remainder is dominated by $|X|^3/6$, so dominated convergence
applies after using $\E X=0$.

For the final assertion the exact same-sign deficit is
\[
 h-\eta=\E\!\left[
 \min(|X|,|X'|)(X-X')^2\ind_{\{XX'>0\}}\right].
\]
On a common sign class, $\min(a,b)\ge ab/R$. Expansion of the two
independent class integrals gives
\[
 h-\eta\ge\frac{h\rho-2(A^2+B^2)}{R}
 =\frac{h(\rho+\zeta)-2}{R}.
\]
If $h\le b(\rho)$, the already proved $\eta\le h$ suffices. Otherwise,
write $b=b(\rho)$, use $\zeta\ge h$ and $b(\rho+b)=2$ to obtain
\[
 \eta-b\le\frac{(h-b)(R-\rho-h-b)}{R}\le0,
\]
when $R\le\rho+2b=\sqrt{\rho^2+8}$. For general $R$, the quadratic
$y(R-\sqrt{\rho^2+8}-y)$ on $y\ge0$ is at most
$(R-\sqrt{\rho^2+8})_+^2/4$. This proves the stated remainder too.
In particular radius three suffices for every $\rho\ge1$.
\end{proof}

For a universal calculation put $H_* =\max\{h,b(\rho)\}$.
Then $b(\rho)\le H_*\le1$, $\eta\le H_*$, $\zeta\ge H_*$ and
$\rho+\zeta\ge2/H_*$. On a closed cell
$\rho\in[\ell,r]$, $H_*\in[a,b]$, it is consequently safe to use
\[
 \eta\le b,\qquad
 \zeta\ge\max\{a,2/b-r,b(r)\}.
\]
Every cell uses both bounds from the same marginal law. A lower bound on
$H_*$ is not a lower bound on $h$; no later sign-probability argument may
substitute it for one. If support $|X|\le R\le\sqrt{\ell^2+8}$ has independently been acquired,
the last part of Lemma~\ref{j:fourier:surrogate} instead permits
$\eta\le b(\ell)$ and $\zeta\ge b(r)$. For larger $R$ its displayed
nonnegative remainder must be retained.

\subsection{Keeping the cubic compensation in the high frequencies}
\label{j:fourier:compensated}
Suppose now that $|X|\le R$ and $\rho\le r$. We may use
\begin{align*}
 M_4&=\min\left\{Rr,1+\frac{(R+1)^2}{R+1/2}(r-1)\right\},\\
 M_5&=\min\left\{RM_4,
 1+\frac{R^3+2R^2+3R+3/2}{R+1/2}(r-1)\right\}
\end{align*}
as upper bounds on $\E X^4$ and $\E|X|^5$ of the same law. To verify the
less immediate branches, set $y=|X|$ and expand
\[
 \rho-1=\E(y-1)^2(y+1/2),\quad
 \E X^4-1=\E(y-1)^2(y+1)^2,
\]
\[
 \E|X|^5-1=\E(y-1)^2(y^3+2y^2+3y+3/2).
\]
The ratios of the last two factors to $y+1/2$ increase on $y\ge0$:
their derivatives are respectively
$y(y+1)/(y+1/2)^2$ and $y(4y^2+7y+4)/(2(y+1/2)^2)$.
The other branches follow from $y\le R$. These moment bounds require
no bound on the product-gap defect.

Put $G=\min\{(r-1)(r+5/3),r^2-Z^2\}$, where $Z>0$ is a proved
lower bound for $\zeta$ on the whole cell; $Z=b(r)$ always suffices.
Then $\gamma^2\le G$: the first branch follows by letting $u\to0$ in
Lemma~\ref{j:fourier:anchor}, and the second from
Lemma~\ref{j:fourier:surrogate}.
Assume that a proved nonnegative envelope $m(u)$ bounds $|f(u)|$ on
the interval being integrated. Set
\[
 g=e^{-u^2/2},\quad v=-i\gamma u^3/6,\quad
 b_u=\sqrt G\,u^3/6,\quad W=g\sqrt{1+b_u^2},
\]
and, for $0\le u\le\pi$,
\[
 \begin{gathered}
 A_u=\max\{|\cos u-g|,
 |\cos u+(r-1)u^3/6-g|\},\\
 B_u=\min\{M_4u^4/24, M_5u^5/120\}+b_u(1-g).
 \end{gathered}
\]
Taylor's theorem gives
\[
 |\Im f(u)+\gamma u^3/6|\le\min\{M_4u^4/24,M_5u^5/120\}.
\]
The anchored real estimate therefore gives
$|f-g(1+v)|\le\sqrt{A_u^2+B_u^2}$. Telescoping the $n$th powers and
then applying the integral Taylor remainder to $(1+sv)^n$, $0\le s\le1$,
proves, for every integer $n\ge2$,
\begin{align*}
 |f(u)^n-g^n(1+nv)|
 &\le\sqrt{A_u^2+B_u^2}\sum_{j=0}^{n-1}m(u)^jW^{n-1-j}
 +\binom n2 b_u^2g^n(1+b_u^2)^{(n-2)/2},\\
 |f(u)^n-g^n(1+nv)|
 &\le m(u)^n+g^n\sqrt{1+n^2b_u^2}.
\end{align*}
One may take their pointwise minimum. These formulas use no quotient of
$m(u)-W$, so coincident and zero values cause no exceptional cases.

On a closed frequency cell $[a,b]\subset[1/2,T]$, a rigorous upper
bound for this minimum is multiplied by the exact weight
$\int_a^b u^{-3}\,du=(a^{-2}-b^{-2})/2$. The infinite remainder is
bounded without a periodicity assumption: for $T>0$,
\[
 \int_T^\infty\frac{|f(u)^n-g(u)^n(1+nv(u))|}{u^3}\,du
 \le\frac1{2T^2}
 +\frac{e^{-nT^2/2}}{nT}
 \left(T^{-3}+\frac{n\sqrt G}{6}\right).
\]
Indeed $|f|\le1$, $\sqrt{1+x^2}\le1+|x|$, $u^{-3}\le T^{-3}$
on the normal part, and
$\int_T^\infty e^{-nu^2/2}\,du\le e^{-nT^2/2}/(nT)$.

To preserve cancellation in the real residual on a closed frequency cell,
put $d_0(u)=\cos u-e^{-u^2/2}$ and
$d_1(u)=d_0(u)+(r-1)u^3/6$. Their derivatives are
\[
 d_0'(u)=-\sin u+ue^{-u^2/2},\qquad
 d_1'(u)=d_0'(u)+(r-1)u^2/2.
\]
If $D_j$ encloses $d_j'$ throughout a cell of width $w$ and midpoint
$u_0$, the fundamental theorem of calculus gives
$d_j(u)\in d_j(u_0)+[-w/2,w/2]D_j$.
Intersecting this with a direct interval evaluation remains valid.

\subsection{The compatible modulus profile}
\label{j:fourier:profile}
The scalar certificate uses $\kappa=62/625$ and proves
\begin{equation}\label{j:fourier:scalar-cosine}
 \cos x\le1-x^2/2+\kappa|x|^3\qquad(x\in\R).
\end{equation}
Here is the entire reduction to its finite check. For $x\ge625/124$ the
right side is at least one. On the remaining interval, the degree-sixteen
Taylor polynomial $P_{16}$ bounds cosine above: the sixteenth derivative
of $\cos-P_{16}$ is $\cos-1\le0$, and its lower derivatives vanish at
zero. The degree-thirteen polynomial
\[
 Q(x)=\kappa+\sum_{j=2}^{8}\frac{(-1)^{j+1}x^{2j-3}}{(2j)!}
\]
satisfies $x^3Q(x)=1-x^2/2+\kappa x^3-P_{16}(x)$.
On the eighty cells $[i/16,(i+1)/16]$, $0\le i<80$, and the last
cell $[5,625/124]$, the exact verifier checks the inverse Bernstein
identity and all fourteen coefficients. The smallest coefficient is
$2479189/65585520000>0$. Lemma~\ref{j:cert:bernstein} proves positivity
on every closed cell; evenness completes \eqref{j:fourier:scalar-cosine}.

Put $a=1/2-4\kappa=129/1250$, $C=1-\cos4$ and $A_0=16a/C$.
The profile
\[
 q(z)=\begin{cases}
 1/2-\kappa z,&0\le z\le4,\\
 A_0(1-\cos z)/z^2,&4\le z\le2\pi,\\
 0,&z\ge2\pi
 \end{cases}
\]
is nonnegative, convex and nonincreasing on $[0,\infty)$, and satisfies
$\cos z\le1-z^2q(z)$. Indeed \eqref{j:fourier:scalar-cosine} gives
$0<A_0\le1$. The pieces agree at their endpoints. On the middle interval
the first derivative is nonpositive and the second has the sign of
\[
 B_0(z)=6+(z^2-6)\cos z-4z\sin z.
\]
On $[4,3\pi/2]$, $B_0'(z)=-2z\cos z+(2-z^2)\sin z\ge0$ and
$B_0(4)\ge4$: indeed $7\pi/6\le4\le3\pi/2$ implies
$\sin4\le-1/2$, and $\cos4\ge-1$, so
$6+10\cos4-16\sin4\ge4$. On $[3\pi/2,2\pi]$, $B_0(z)\ge6$.
The derivative at $2\pi$ is zero. The derivative jump at $4$ has the
correct sign because $258\sin4+119(1-\cos4)>0$. Its exact lower bound is
obtained by replacing sine by its degree-fifteen Taylor lower bound and
cosine by $P_{16}$. The same scalar checker verifies this rational
inequality and the needed $25/8<\pi<22/7$ by Machin's formula.

To apply the profile, weight the distribution of $X-X'$ by
$(X-X')^2/2$. This is a probability measure whose first absolute moment
is $\rho+\eta$. Jensen's inequality gives, for $u\ge0$,
\[
 |f(u)|^2=\E\cos(u(X-X'))
 \le1-2u^2q(u(\rho+\eta)).
\]
If $\rho+\eta\le \mathcal D$, monotonicity permits $\mathcal D$ on the right. Thus the
linear part gives $1-u^2+2\kappa \mathcal Du^3$; when $4\le \mathcal Du\le2\pi$ the
curved part gives $1-2A_0(1-\cos(\mathcal Du))/\mathcal D^2$.
Take the minimum with one and with the coordinate-disk estimate from
Lemmas~\ref{j:fourier:anchor}--\ref{j:fourier:surrogate}. The upper
square of a real interval is the larger of the squares of its two
endpoints, even when that interval crosses zero or lies below it.

\subsection{A Fourier majorant of a closed half-line}

\begin{lemma}\label{j:fourier:closed-majorant}
For every integrable real random variable $Y$, every $h\in\R$ and $A>0$,
\begin{equation}\label{j:fourier:closed-majorant-formula}
 \Pp(Y\le h)\le\frac12+
 2\lim_{\epsilon\downarrow0}\Re\int_\epsilon^1
 K(t)e^{-iAth}\E e^{iAtY}\,dt,
\end{equation}
where
\[
 K(t)=\frac{1-t}{2}+\frac i2
 \left((1-t)\cot(\pi t)+\frac1\pi\right).
\]
The limit concerns the real part; the complex integral itself may diverge.
\end{lemma}
\begin{proof}
Put $W(x)=(\sin(\pi x)/(\pi x))^2$, with $W(0)=1$, and
\[
 H_B(x)=\sum_{m\ge1}\{W(x-m)-W(x+m)\}+2xW(x).
\]
The series converges absolutely and locally uniformly. It is odd, vanishes
at zero, and takes value $\operatorname{sgn}(j)$ at every nonzero integer.
The partial-fraction identity
\[
 \sum_{m\in\mathbb Z}(x-m)^{-2}=\pi^2\csc^2(\pi x)
 \quad(x\notin\mathbb Z)
\]
follows by integrating $\pi\cot(\pi z)/(z-x)^2$ around the square
with vertices $\pm(N+1/2)\pm i(N+1/2)$: the boundary integral is
$O(N^{-1})$, and the residues are $(m-x)^{-2}$ and
$-\pi^2\csc^2(\pi x)$. For positive noninteger $x$ this identity gives
\[
 H_B(x)-1=\left(\frac{\sin\pi x}{\pi}\right)^2
 \left(\frac2x-\frac1{x^2}-2\sum_{m\ge1}(x+m)^{-2}\right).
\]
Integral comparison of the decreasing function $(x+s)^{-2}$ yields
\[
 \sum_{m\ge1}(x+m)^{-2}\le\frac1x
 \le\frac1{x^2}+\sum_{m\ge1}(x+m)^{-2}.
\]
Thus $|H_B(x)-\operatorname{sgn}(x)|\le W(x)$; oddness and the integer
values supply the remaining points. Consequently
\[
 M(x):=\frac{1-H_B(x)+W(x)}2\ge\ind_{\{x\le0\}},
 \qquad M(0)=1,\qquad 0\le M(x)-\ind_{\{x\le0\}}\le W(x).
\]
Elementary integration gives
\[
 W(x)=2\int_0^1(1-t)\cos(2\pi tx)\,dt,\qquad
 2xW(x)=\frac2\pi\int_0^1\sin(2\pi tx)\,dt.
\]
Inserting the first formula into the shifted series with Abel weights
$r^m$, $0<r<1$, gives
\[
 \sum_{m\ge1}r^m\{W(x-m)-W(x+m)\}
 =4\int_0^1(1-t)\sin(2\pi tx)
 \frac{r\sin(2\pi t)}{1-2r\cos(2\pi t)+r^2}\,dt.
\]
The fraction tends to $\tfrac12\cot(\pi t)$ and is bounded in absolute
value by $\tfrac12|\cot(\pi t)|$, since its denominator is
$(1-r)^2+4r\sin^2(\pi t)$. The resulting majorant is integrable:
the sine vanishes at zero and $(1-t)\cot(\pi t)$ is bounded at one.
Dominated convergence therefore proves
\[
 H_B(x)=2\int_0^1
 \left((1-t)\cot(\pi t)+\frac1\pi\right)\sin(2\pi tx)\,dt.
\]
Combining the formulas yields
\[
 M(x)=\frac12+2\lim_{\epsilon\downarrow0}\Re\int_\epsilon^1
 K(t)e^{2\pi itx}\,dt
 =1+2\Re\int_0^1K(t)(e^{2\pi itx}-1)\,dt.
\]
The last integral is absolutely convergent. For integrable $U$,
$\E|e^{2\pi itU}-1|\le2\pi t\E|U|$; together with $K(t)=O(t^{-1})$
at zero and boundedness at one, this permits absolute Fubini.
Apply the identity to $U=A(Y-h)/(2\pi)$ and take expectations.
The value $M(0)=1$ proves the claimed closed-event convention even at atoms.
\end{proof}

For $Y=(X_1+\cdots+X_4)/2$ and $A=6$, independence gives
$\E e^{6itY}=f(3t)^4$. Thus the lemma supplies exactly the
factor six, positive imaginary kernel, and boundary value used in the
four-summand comparison below.

\subsection{From characteristic functions to the CDF}
\label{j:fourier:prawitz}
We use Prawitz's smoothing inequality in the form recorded in
\cite[Lemma 4.1, p.~263]{Shevtsova2012}. For $n\ge2$, $c>0$, $0<\tau\le1$, put
\[
 K_2(t)=(1-t)+i\{(1-t)\cot(\pi t)+1/\pi\},\qquad 0<t<1.
\]
With $T=c\sqrt n$, that inequality gives
\begin{align*}
 \Delta_n\le{}&\int_0^\tau |K_2(t)|
   |f(ct)^n-e^{-nc^2t^2/2}|\,dt
 +\int_\tau^1|K_2(t)|\,|f(ct)|^n\,dt\\
 &+\int_0^\tau|K_2(t)-i/(\pi t)|e^{-nc^2t^2/2}\,dt
 +\frac1\pi\int_\tau^\infty e^{-nc^2t^2/2}\frac{dt}{t}.
\end{align*}
Our standardized iid laws satisfy its independence, centering and positive
total-variance hypotheses. The cited convention uses a strict CDF.
Continuity of $\Phi$ makes the supremum identical for closed CDFs: take
one-sided limits at each evaluation point. Thus atoms require no smoothing
of the original law.

For $g(u)=e^{-u^2/2}$ and envelopes $m\ge|f|$, $E_0\ge|f-g|$,
the first integrand uses
$|f^n-g^n|\le\min\{E_0\sum_{j=0}^{n-1}m^jg^{n-1-j},m^n+g^n\}$.
Every literal integer $n$ is evaluated separately. The original $20$--$24$ integral verifier
covers $[\epsilon,\tau]$ and $[\tau,1-\epsilon]$ by closed cells,
where $\epsilon=1/1024$; its numerical parameter domain is
$0<c<\pi$ and $\epsilon<\tau\le3/10<1-\epsilon$.
In particular $c\tau<1$ licenses the stable low-frequency Taylor bounds.
On a cell of width $w$ and midpoint $t_0$,
it uses either $w\sup F$ or
$w\{F(t_0)+w^2\max(0,\sup F'')/24\}$ for a fixed smooth majorant
$F$. The latter follows by Taylor expansion and integration; a failed
derivative-domain guard triggers the former bound on the entire cell.
Switching branches at a sampled midpoint would not justify this rule.

For the omitted endpoint pieces, the respective upper bounds are
\begin{align*}
 L_\epsilon&=n\left[
 \frac{rc^3}{6}\left(\frac{\epsilon^4}{4}
       +\frac{\epsilon^3}{3\pi}\right)
 +\frac{c^4}{8}\left(\frac{\epsilon^5}{5}
       +\frac{\epsilon^4}{4\pi}\right)\right],\\
 N_\epsilon&=\epsilon+\pi\epsilon^2/2,\qquad
 H_\epsilon=\epsilon^2/2+\pi\epsilon^3/3.
\end{align*}
They follow from $|K_2(t)|\le1+1/(\pi t)$ near zero, cubic Taylor
error for $f$, $e^{-u^2/2}-1+u^2/2\le u^4/8$, and
\[
 0\le1/x-\cot x\le\frac{x}{3(1-x^2/6)}\le x\quad(0<x\le1).
\]
For this last estimate use
$\sin x-x\cos x=\int_0^x t\sin t\,dt\le x^3/3$ and
$\sin x\ge x-x^3/6$. Apply the same identity with $1-t$ near one.
Finally the infinite Gaussian term is at most
$e^{-nc^2\tau^2/2}/(\pi nc^2\tau^2)$, by
$\int_a^\infty e^{-v}v^{-1}\,dv\le e^{-a}/a$.
For a strip $\rho\in[\ell,r]$, multiply the full four-term bound by
$\sqrt n/\ell$, with outward enclosures. A complete closed strip cover
and a strict rational comparison then establish the claimed CDF bound.


\subsection{Sharper scalar bounds for the same integrals}
\label{j:fourier:quadrature}
For $a>0$ write $E_1(a)=\int_a^\infty e^{-v}v^{-1}\,dv$.
Integration by parts, followed by Jensen's inequality for the convex
function $t\mapsto(a+t)^{-2}$ and an exponential random variable of mean
one, gives
\[
 E_1(a)=\frac{e^{-a}}a-e^{-a}\int_0^\infty
                \frac{e^{-t}}{(a+t)^2}\,dt
 \le e^{-a}\left(\frac1a-\frac1{(a+1)^2}\right).
\]
The last Gaussian term in Prawitz's inequality is $E_1(a)/(2\pi)$,
where $a=nc^2\tau^2/2$. This improvement has no additional moment or
support hypothesis.

Some small-index certificates require $\tau>3/10$. Their low-frequency
range is still controlled explicitly. Put $U=6/5$. On $0\le x\le U$,
the alternating sine bounds and integration give
\[
 \sin x\ge x\left(1-x^2/6+x^4/120-x^6/5040\right),\qquad
 0\le\sin x-x\cos x\le x^3/3-x^5/30+x^7/840.
\]
The denominator below is at least $1-U^2/6-U^6/5040>0$, and the
numerator factor is at least $1/3-U^2/30>0$. Consequently
\[
 0\le1/x-\cot x\le L(x):=
 \frac{x(1/3-x^2/30+x^4/840)}
      {1-x^2/6+x^4/120-x^6/5040},
 \qquad L(0)=0.
\]
It follows that
\[
 |K_2(t)-i/(\pi t)|
 \le(1-t)\sqrt{1+L(\pi t)^2}
 \quad(0<t\le19/50).
\]
At the high endpoint put $s=1-t$. The identity
$K_2(t)=s+i\,s(1/(\pi s)-\cot(\pi s))$ gives
$|K_2(t)|\le s\sqrt{1+L(\pi s)^2}$ for $t\ge3/4$.
The rational upper bound $\pi<22/7$ implies
$\pi(19/50)<6/5$ and $\pi/4<6/5$.
The numerical quadrature must therefore declare
$\epsilon<\tau\le19/50<1-\epsilon$, $c<\pi$, and $c\tau\le6/5$.
It differentiates the displayed smooth rational majorant on whole cells;
it does not differentiate an interval-selected minimum.

The cancellation series on this larger range also have explicit tails.
Truncate the even series for cosine and $e^{-u^2/2}$ at degree twenty.
For their difference, the error in each $t$ derivative of order
$j=0,1,2$, with $u=ct\in[0,U]$, is at most
\[
 \frac{2(22)_j c^j u_{\max}^{22-j}}{2^{11}11!}.
\]
Indeed, for terms of Gaussian index $k\ge11$, the ratios of consecutive
absolute derivative terms are bounded respectively by
\[
 \frac{U^2}{2(k+1)},\qquad \frac{U^2}{2k},\qquad
 \frac{U^2(2k+1)}{2k(2k-1)}.
\]
All decrease with $k$ and are less than $1/8$ at $k=11$.
The cosine ratios are
$U^2/\{(2k+2-j)(2k+1-j)\}<1/8$.
The first cosine tail is $2^{11}11!/22!<1/4$ times the first Gaussian
tail, for every $j$. Summing both geometric tails gives a factor less
than $(1+1/4)/(1-1/8)=10/7<2$.
Uniform absolute convergence justifies these differentiations.
For the removable function $2(\cos u-e^{-u^2/2})/u^3$, divide the series
symbolically, so its value tail is bounded by
$4u_{\max}^{19}/(2^{11}11!)$ and its value at zero is obtained directly
from the series. No numerical division by a nearly zero frequency is
needed.

\subsection{Uniform low-frequency remainder}

The next estimate requires actual fourth through sixth moments, acquired
from support or separately bounded. A third-moment ceiling alone is insufficient.

Assume $n\ge N>2$ and common bounds
\[
 \mu_4:=\E X^4\le M\le111/10,\quad |\gamma|\le g\le37/20,
 \quad\mu_6:=\E X^6\le M_6,
\]
and put
\[
 \mathfrak a=(3-\mu_4)/24,\quad \mathfrak z=(10\gamma-\E X^5)/120,
 \quad |\mathfrak a|\le A,\quad |\mathfrak z|\le Z.
\]
Here $A,Z,M_6$ are nonnegative common ceilings, not quantities inferred
from bounded third moments. Define
\[
 E=M_6/720+g^2/72+(1/2+g^2/144)^3/5,\qquad
 B=A+Z/2+E/4,
\]
\[
 a_0=\frac{N-1}{2N},\quad b_0=\frac{45-M}{96}(1-2/N)>0,
 \quad\lambda_N=\sqrt{\pi/2}\left(\frac N{N-1}\right)^{3/2}.
\]
\begin{proposition}\label{j:fourier:low}
For all the laws and thresholds above,
\[
 I_0(H)\le\frac{L_0+L_H|H|}{\sqrt n}+\frac{L_2}{n},
\]
where, for any $c_0\ge\lambda_N$ and $c_1\ge\lambda_N/2$,
\begin{align*}
 L_0&=c_0Z+c_1Ag+251g^3/17280,\\
 L_H&=c_0A+251g^2/4800,\\
 L_2&=E/(2a_0^2)+B^2/(2b_0^3).
\end{align*}
\end{proposition}
\begin{proof}
Let $h(t)=\exp(-t^2/2-i\gamma t^3/6)$. Taylor's integral remainder
for $e^{itX}$, followed by expectation, gives its expansion through
degree five with error at most $\mu_6t^6/720$. For
$w=-t^2/2-i\gamma t^3/6$ and $0\le t\le1/2$,
\[
 |w|\le t^2(1/2+g^2/144)<1/6,
 \quad |e^w-1-w-w^2/2|\le |w|^3e^{|w|}/6\le |w|^3/5.
\]
The first bound follows from $\sqrt{1/4+y}\le1/2+y$; for the last
use $e^{1/6}\le(1-1/6)^{-1}=6/5$. Comparing the two polynomials gives
\begin{equation}\label{j:fourier:taylor}
 f-h=-\mathfrak a t^4-i\mathfrak z t^5+r_6(t),\qquad |r_6(t)|\le Et^6,
 \qquad |f-h|\le Bt^4.
\end{equation}
The degree-six term in $w^2/2$ costs $g^2t^6/72$; it has not been
discarded.

Four integrations of $1-\cos u\ge0$ show
$\cos u\le1-u^2/2+u^4/24$ for every real $u$. Independence gives
$\E(X-X')^4=2\mu_4+6$, and hence
\[
 |f(t)|^2\le1-t^2+(M+3)t^4/12
 \le1-(45-M)t^2/48.
\]
Thus $|f|,|h|\le e^{-(45-M)t^2/96}$ throughout the low interval.
Every point on their complex joining segment obeys the same bound.
Taylor's integral formula for the polynomial $v\mapsto v^n$ proves
\[
 f^n-h^n=n(f-h)h^{n-1}+R_2,\quad
 |R_2|\le\tfrac12n(n-1)B^2t^8e^{-(n-2)(45-M)t^2/96}.
\]
The leading remainder in \eqref{j:fourier:taylor} and this power
remainder cost at most $L_2/n$ after division by $t^3$ and integration:
use $n-1\ge2a_0n$, $(n-2)(45-M)/96\ge b_0n$, $n(n-1)\le n^2$,
and $\int_0^\infty t^3e^{-At^2}dt=1/(2A^2)$,
$\int_0^\infty t^5e^{-At^2}dt=A^{-3}$.

For the remaining leading terms let $b=n-1$ and
$v=b\gamma t^3/6$. Their imaginary part after rotation by $Ht$ is
bounded by
\[
 ne^{-bt^2/2}\{(Z+A|H|)t^5+bAg t^7/6\}.
\]
The Gaussian moments of orders two and four, after dividing by $t^3$,
give $n\sqrt{\pi/2}b^{-3/2}(Z+A|H|+Ag/2)$.
This is at most $\lambda_N(Z+A|H|+Ag/2)/\sqrt n$.
Finally put $u=n\gamma t^3/6$. Direct sine and cosine remainders give
\[
 |\Im[e^{-itH}(h^n-e_n)]|
 \le e^{-nt^2/2}(|u|^3/6+|H|t u^2/2).
\]
The Gaussian moments of orders six and four respectively give
$\sqrt{\pi/2}g^3/(432\sqrt n)$ times $5$, and
$\sqrt{\pi/2}g^2|H|/(24\sqrt n)$.
Using $\sqrt{\pi/2}<251/200$ gives precisely the two displayed
rational coefficients. Every extended integrand was nonnegative.
The only index comparisons were the displayed rational inequalities;
in particular the law and $H$ can vary arbitrarily with $n$.
\end{proof}

In the selected enclosure one has actual support $|X|\le R$
with $1\le R\le6$ and $\E|X|^5\le E_5$. Then $M_6=RE_5$ and
$A=\max(2,M-3)/24$ are valid. Its sharper $Z$ is also licensed:
for $Q=X^2-\gamma X-1$ direct multiplication gives
\[
 \E Q^2=\mu_4-1-\gamma^2,\quad
 \E X^5=\gamma^3+2\gamma+\E[Q^2(X+2\gamma)].
\]
Reflect this auxiliary calculation if necessary so that
$t=|\gamma|\ge0$. For any $r\ge\max(R,2g)$ the two signs of
$\E X^5-10\gamma$ are at most
\[
 P_r(t)=r(M-1)+(10-2M)t-rt^2+t^3,
 \quad Q_r(t)=r(M-1)+(2M-10)t-rt^2-t^3.
\]
The coefficients of $\mu_4$ before replacing it by $M$ are
$r-2t,r+2t\ge0$. If $M\le5$, $P_r\ge Q_r$ and
$P_r\le r(M-1)+(10-2M)t-(r-g)t^2$, whose maximum on $[0,g]$
is at the clipped vertex used in the code. If $M>5$,
$P_r\le r(M-1)$ and
$Q_r\le r(M-1)+(2M-10)t-rt^2$, again with the stated clipped vertex.
For $r=6$, $P_6''=-12+6t<0$ on $[0,g]$. If $M\le5$ and
$P'_6(g)\ge0$, its maximum is exactly $P_6(g)$; otherwise the
preceding quadratic gives $6(M-1)+(10-2M)^2/[4(6-g)]$.
For $M>5$ the unrestrained vertex bound is
$6(M-1)+(2M-10)^2/24$. Taking the minimum of these proved ceilings
and $E_5+10g$, then dividing by $120$, gives a valid coefficient ceiling $Z$.
The coefficient ceiling therefore uses the known support radius $R$.

\subsection{The complete high-frequency half-line}

A nearly lattice law can retain characteristic-function peaks away from
zero. The high-frequency argument therefore uses an event involving the
actual difference $X-X'$ and partitions the oscillations around their
zeros. The six contributions below cover intermediate frequencies,
nonzero-cell cores, nonzero-cell wings, the small-angle part of the zero
cell, the rest of the zero cell, and the signed Gaussian comparator.
Each optional improvement has its own stated positivity and monotonicity
conditions.

For independent copies of the same marginal suppose an actual event
$\mathcal E$ satisfies
\[
 \Pp(\mathcal E)\ge\alpha>0,\quad |X-X'|\ge d>0\text{ on }\mathcal E,
 \quad\E[(X-X')^2\mid\mathcal E]\le M_E.
\]
Fix $b=9/10$, $0<r\le1/4$, and
\[
 h=333/106-(22/7)^3r^2/6>0,\quad
 \kappa=\alpha(23/48)^2d^2,\quad L(x)=\frac{x+2}{x^2+4x+2},
\]
\[
 S_3=\sum_{j=1}^{20}j^{-3}+1/800,\quad
 S_5=\sum_{j=1}^{20}(j-r)^{-5}+\frac1{4(20-r)^4}.
\]
Let $c(u)=u/2-(M+3)u^2/24$ and set
\[
 c_* =\min\{c(1/4),c(81/100)\},\quad
 c_q=1/8-(M+3)/384,\quad s=[1-(M+3)53/600]/2.
\]
\begin{proposition}\label{j:fourier:high}
Suppose $c_*>0$. For every $n\ge N$, $\sqrt n I_\infty$ is at most
the sum of the following six terms:
\begin{align*}
 T_1(n)&=2\sqrt n e^{-nc_*},\\
 T_2(n)&=\frac{M_E}{4\pi^{3/2}h\sqrt\alpha}
       \{S_3+3S_5/(n\alpha h^2)\},\\
 T_3(n)&=\frac{M_E\sqrt n}{2\pi^2}e^{-n\alpha h^2r^2},\\
 T_4(n)&=\frac{e^{-n\kappa b^2}}{2b^4\kappa\sqrt n},\\
 T_5(n)&=\frac{\sqrt n}{2b^2}e^{-n\alpha(23/48)^2},\\
 T_6(n)&=\sqrt n e^{-n/8}\{2L(n/8)+\sqrt{M-1}/3\}.
\end{align*}
If $s>0$, $T_1$ may be replaced by its minimum with
$2\sqrt n e^{-nc_q}L(ns/4)$. The fourth term may be replaced by
its minimum with $\sqrt n e^{-n\kappa b^2}L(n\kappa b^2)/(2b^2)$.
The common-floor sum at $N$ is valid if
\[
 2Nc'>1\quad\text{for }c'\in\{c_*,\alpha h^2r^2,
                      \alpha(23/48)^2,1/8\}.
\]
For the optional first term also require $2Nc_q>1$; for the optional
fourth term require $N\kappa b^2\ge1/4$. If an optional gate fails,
the simpler term remains available.
\end{proposition}
\begin{proof}
The triangle inequality separates $|f^n|$ and the entire comparator.
On $[1/2,b]$, symmetrization gives $|f|^n\le e^{-nc(t^2)}$.
The function $c$ is concave, so $c(t^2)\ge c_*$ and the integral
of $t^{-3}$ is less than two. This gives $T_1$.
Also
\[
 c(u)-c_q-s(u-1/4)=(M+3)(u-1/4)(81/100-u)/24\ge0.
\]
When $s>0$ this chord may be integrated to infinity.

For completeness, the tail calculation used here is, for $x=na b^2>0$,
\begin{equation}\label{j:fourier:tail}
 \sqrt n\int_b^\infty e^{-na t^2}t^{-3}dt
 \le\frac{\sqrt n e^{-x}}{2b^2}L(x).
\end{equation}
To see where the rational function $L$ comes from, let $U$ have density
$e^{-u}$ on $u\ge0$. Substitution followed by positive Fubini gives
\[
 \sqrt n\int_b^\infty e^{-na t^2}t^{-3}dt
 =\frac{\sqrt n e^{-x}}{2b^2}\E\frac{U}{x+U}.
\]
For example, after putting $v=na(t^2-b^2)$ the remaining factor is
$\int_0^\infty xe^{-v}(x+v)^{-2}dv$; integration by parts makes it
$\E[U/(x+U)]$. For $x>0$ put
\[
 J_x=\E\frac1{x+U},\qquad v_0=\begin{pmatrix}1\\1\end{pmatrix},\qquad
 G=\E\left[(x+U)\begin{pmatrix}1&U\\U&U^2\end{pmatrix}\right]
 =\begin{pmatrix}x+1&x+2\\x+2&2x+6\end{pmatrix}.
\]
The entries use $\E U=1$, $\E U^2=2$, and $\E U^3=6$.
For $q(U)=c_1+c_2U$, weighted Cauchy--Schwarz in its squared-residual
form gives
\[
 0\le\E\frac{[1-(x+U)q(U)]^2}{x+U}
 =J_x-2v_0^Tc+c^TGc.
\]
The matrix $G$ is positive definite: its leading entry is positive and
$\det G=x^2+4x+2>0$. Choosing $c=G^{-1}v_0$ yields
\[
 J_x\ge v_0^TG^{-1}v_0=\frac{x+3}{x^2+4x+2},\qquad
 \E\frac{U}{x+U}=1-xJ_x\le\frac{x+2}{x^2+4x+2}=L(x).
\]
This proves \eqref{j:fourier:tail} for every $a,b,n>0$ with the same
function $L$. Applying it at endpoint
$1/2$ and rate $s$ proves the optional first term.

For $t\ge b$, symmetrization, the actual-event probability bound,
and convex Jensen give
\[
 |f(t)|^n\le e^{-n\E\sin^2(t(X-X')/2)}
 \le\E[e^{-n\alpha\sin^2(tD/2)}\mid\mathcal E],
 \quad D=|X-X'|.
\]
Positive Tonelli allows a separate partition for every $D>0$.
With $\ell=2\pi/D$, use the zero cell $[0,\ell/2]$ and nonzero
cells $\ell[j-1/2,j+1/2]$, $j\ge1$. In a nonzero cell write
$t=\ell(j+v)$. On its core $|v|\le r$,
$|\sin\pi v|\ge h|v|$; on its wings $r\le|v|\le1/2$ it is
at least $hr$. These follow from $\sin x\ge x-x^3/6$ and
monotonicity of $\sin$ on $[0,\pi/2]$, with directed bounds on $\pi$.
For $|v|\le r$ Taylor's integral remainder gives
\[
 \{(j-v)^{-3}+(j+v)^{-3}\}/2
 \le j^{-3}+6(j-r)^{-5}v^2.
\]
Extend the nonnegative Gaussian majorant over the whole real line.
Its zeroth and second moments, the Jacobian
$t^{-3}dt=D^2(4\pi^2)^{-1}(j+v)^{-3}dv$, and
$\E[D^2\mid\mathcal E]\le M_E$ give $T_2$.
The two series tails are bounded by their decreasing integrals after
20. The total wing weight is at most
$\int_{\ell/2}^\infty t^{-3}dt=D^2/(2\pi^2)$, giving $T_3$.
Extending these sets below $b$ only increases their nonnegative bounds.

In the zero cell put $u=tD/2\le\pi/2$. Where $u\le1/2$,
$\sin u\ge23u/24$, and $D\ge d$ yields rate $\kappa$.
Extending the set to $[b,\infty)$ gives $T_4$ from
$t^{-3}\le t/b^4$; \eqref{j:fourier:tail} gives its optional improvement.
Where $u\ge1/2$, monotonicity gives $\sin u\ge23/48$ and the total
weight is at most $1/(2b^2)$, proving $T_5$. Endpoints have zero
Lebesgue measure even if the conditional distance distribution has atoms.

Finally $|e_n(t)|\le e^{-nt^2/2}(1+n|\gamma|t^3/6)$.
The first part on the entire $[1/2,\infty)$ uses
\eqref{j:fourier:tail}; the second uses
$\int_{1/2}^\infty e^{-nt^2/2}dt\le2e^{-n/8}/n$.
Pearson's identity $\mu_4-1-\gamma^2=\E(X^2-\gamma X-1)^2\ge0$
gives $T_6$. No comparator tail is dropped.

For floor evaluation, $\sqrt n e^{-nc'}$ decreases for $2nc'>1$,
whereas $n^{-1}$ and $n^{-1/2}e^{-nc'}$ decrease for every $c'>0$.
Also $L'(x)=-(x^2+4x+6)/(x^2+4x+2)^2<0$.
Thus the optional first term decreases under its own gate, and the
comparator decreases for $n\ge5$. The sign of the derivative of
$\sqrt{x}e^{-x}L(x)$ is the sign of
\[
 4-10x-22x^2-13x^3-2x^4,
\]
which is negative at $1/4$ and strictly decreasing on $(0,\infty)$.
This proves the optional zero-cell gate. Every comparison is of a common
scalar majorant, not of characteristic functions for different laws.
\end{proof}

\subsection{Uniform remainder conclusion}

Under the preceding moment, actual-event and scalar gates, let $K_H$ be
any complete HIGH sum evaluated at $N$, including its permitted optional
improvements. Then
\begin{equation}\label{j:fourier:conclusion}
 |E_n|\le\frac2\pi\{L_0+K_H+L_2/\sqrt N+L_H|H|\}
 \le\frac{16}{25}\{L_0+K_H+L_2/\sqrt N+L_H|H|\}.
\end{equation}
The second inequality uses $\pi>25/8$ and nonnegative coefficients.
This supplies a constant Fourier budget and a raw-threshold slope.
The scalar gates are imposed on the analytic majorants, not on the rounded
implementation. Complete HIGH bounds for different valid windows may be
minimized. Each LOW triple remains intact: a constant from one enclosure
cannot be combined with another enclosure's threshold slope without a
separate justification. These statements are uniform over every law and
threshold satisfying the displayed hypotheses; no fixed-law limit is used.

% Finite chapter. Numerical data remain in the pinned reproduction package.
\section{The finite induction through forty-two}\label{j:finite:chapter}
We prove the finite base by excluding a first violating index. Every
calculation then concerns one attained positive closed-event ratio maximizer
with score $C=C_n>\CE$, and its immediate predecessor has score at most
$\CE$. The computational propositions below assert inequalities on complete
parameter regions; the intervening arguments prove that this same marginal
lies in those regions. No bound on its number of atoms is assumed.

\subsection{Quantitative starting cases}\label{j:finite:starting}

The small cases provide quantitative starting bounds for the later
induction.  The second and third cases use the same idea: a hypothetical
violating extremizer has bounded support, and a polynomial majorizes its
closed sum event on the whole support product.  Independence then evaluates
the polynomial using moments of the same marginal.  No restriction on the
number of atoms is imposed.

\begin{theorem}\label{j:finite:n123}
For every standardized law with finite third absolute moment,
\[
 C_1\le\frac25,\qquad C_2\le\frac{2\sqrt2}{7},\qquad
 C_3\le\frac{59\sqrt3}{250}.
\]
All three upper bounds are strictly smaller than $\CE$.
\end{theorem}

\paragraph{One summand.}
The first numerical assertion is the following checked family of 150 consecutive rational cells $[a,b]$
covering $[-2,1]$, with two rational coefficients $u,v$ in each cell.
It proves
\begin{equation}\label{j:finite:n1-dual}
 L_a-v+uy+vy^2+\frac25|y|^3-\ind_{\{y\le b\}}\ge0
 \quad(y\in\R),\qquad L_a\le\Phi(a).
\end{equation}
Taking expectations cancels the linear and quadratic moment terms.
For $t\in[a,b]$ this gives
\[
 F_P(t)-\Phi(t)\le F_P(b)-L_a\le\frac25\rho.
\]
For $t\le-2$, Cantelli gives $F_P(t)\le1/(1+t^2)\le1/5$.
For $t\ge1$, the discrepancy is at most $1-\Phi(1)<2/5$.
Since $\rho\ge1$, these bounds cover both exterior rays.
Reflection gives the negative branch, including atoms.

\paragraph{A common support argument for two and three summands.}
Let $K_2=2\sqrt2/7$, $K_3=59\sqrt3/250$ and suppose $C_n>K_n$,
where $n=2$ or $3$.  Theorem~\ref{j:core:attainment} supplies a
positive closed-CDF maximizing pair.  Write
\[
 z=H/\sqrt n,\quad F=\Pp(S_n\le H),\quad D_0=F-\Phi(z),
 \quad \lambda=D_0/\rho,\quad \beta=\E[X|X|].
\]
Here $D_0$ denotes the discrepancy, not the support-gap defect used in
the main argument.  The contact condition is equivalently
\begin{align}
 n\Pp(S_{n-1}\le H-y)&\le Q_n(y),\notag\\
 Q_n(y)&=\lambda|y|^3-\frac{z\phi(z)+3D_0}{2}y^2
 -(\sqrt n\phi(z)+3\lambda\beta)y
 +nF+\frac{z\phi(z)+D_0}{2},\label{j:finite:n23-contact}
\end{align}
with equality on the entire closed support.
The published affine estimate
\begin{equation}\label{j:finite:n23-affine}
 \frac{\sqrt n\,\Delta_n(P)}{\rho(P)}\le\frac{1661}{5000}
       \left(1+\frac{429}{1000\rho(P)}\right)
\end{equation}
from \citet{Shevtsova2013} rules out $\rho\ge2$ at $n=2$ and
$\rho\ge15/8$ at $n=3$.  Together with $|\beta|\le1$,
$\phi(z)<2/5$ and $|z\phi(z)|<1/4$, it gives the following lower
bounds $Q_n(y)\ge P_n(|y|)$:
\begingroup\renewcommand{\arraystretch}{1.8}
\[
\begin{array}{c|c|c}
 n&\rho\text{ ceiling}&P_n(r)\\\hline
 2&2&\displaystyle\frac27r^3-\frac{67}{56}r^2-\frac{23}{14}r-\frac18\\[2pt]
 3&15/8&\displaystyle\frac8{35}r^3-\frac45r^2-\frac85r-\frac18
\end{array}
\]
\endgroup
Indeed, for $n=2$, $\lambda>2/7$, $D_0<5/7$ and
$\lambda<5/14$ suffice: $\sqrt2<10/7$ gives
$\sqrt2\phi(z)+3\lambda<4/7+15/14=23/14$.
For $n=3$, the corresponding bounds are
$\lambda>8/35$, $D_0<9/20$ and $\lambda<2/7$.
In the latter case \eqref{j:finite:n23-affine} gives $C_3\rho<77/100$,
and $\sqrt3>12/7$ then gives $D_0=C_3\rho/\sqrt3<9/20$.
The shifted differences are
\begin{align*}
 P_2(6+s)-2&=\frac27s^3+\frac{221}{56}s^2+
                  \frac{104}{7}s+\frac{373}{56},\\
 P_3(6+s)-3&=\frac8{35}s^3+\frac{116}{35}s^2+
                  \frac{472}{35}s+\frac{2197}{280}.
\end{align*}
Every coefficient is positive.  Equality in
\eqref{j:finite:n23-contact} therefore excludes every support point with
$|y|\ge6$, because its left side is at most $n$.

\paragraph{The two-summand majorants.}
Put $\gamma=\E X^3$ and
\[
 v=(1,y,y^2,y^3,|y|^3),\qquad \mu=\E v(X)=(1,0,1,\gamma,\rho).
\]
For a rational symmetric coefficient array indexed by $i\le j$, let
\[
 M_c(x,y)=\sum_i c_{ii}v_i(x)v_i(y)+
 \sum_{i<j}c_{ij}\{v_i(x)v_j(y)+v_j(x)v_i(y)\},
\]
and $q_c(\gamma,\rho)=\sum_i c_{ii}\mu_i^2+
2\sum_{i<j}c_{ij}\mu_i\mu_j$.
Independence gives $\E M_c(X_1,X_2)=q_c(\gamma,\rho)$.
Lemma~\ref{j:core:moments} puts the actual moments in
\[
 \mathcal D_2=\{(\gamma,\rho)\in[-2,2]\times[1,2]:
 \rho^2(\gamma^2+4)-(\gamma^2+2)^2\ge0\}.
\]
The 70 certificate rows use
$a_j=-5/2+j/20$, $b_j=a_j+1/20$, $0\le j<70$.
The checked numerical assertions, in the notation of
Lemma~\ref{j:cert:tensor}, are
\begin{align*}
 M_{c^{(j)}}(x,y)&\ge\ind_{\{x+y\le b_j\}}
                    &&((x,y)\in[-6,6]^2),\\
 q_{c^{(j)}}(\gamma,\rho)&\le\Phi(a_j/\sqrt2)+\frac27\rho
                    &&((\gamma,\rho)\in\mathcal D_2).
\end{align*}
Each sign quadrant is cut by the event line and triangulated into
closed triangles.  Degree-six Bernstein coefficients, with exact
subdivision, establish the first inequality.  The second inequality is
checked on moment rectangles; a rectangle is discarded only if
$\rho_+^2(g_*+4)-(g_*+2)^2<0$, where $g_*$ is the minimum of
$\gamma^2$ on that rectangle.  This is valid because the expression
increases with $\rho$ and decreases with $\gamma^2$ when $\rho\le2$.

For a threshold in $[a_j,b_j]$, expectation of the first inequality and
the second give discrepancy at most $2\rho/7$.
For $H\le-5/2$, Cantelli gives $2/(2+H^2)\le8/33<2/7$;
for $H\ge1$, $1-\Phi(H/\sqrt2)<2/7$.
These bounds contradict the chosen violating maximizer.  This proves
the second assertion of Theorem~\ref{j:finite:n123}.

\paragraph{The three-summand majorants.}
Use the features
\[
 (v_0,\ldots,v_6)(y)=(1,y,y^2-1,y^3,|y|^3,y^4,y^5).
\]
Let $\mathcal I$ be the 59 nondecreasing triples from $\{0,\ldots,6\}$
whose largest entry is less than 5 or which contain an entry 1 or 2.
For $\alpha\in\mathcal I$, average the associated product over its
distinct coordinate permutations, obtaining $V_\alpha$.
Set $M_c=\sum_\alpha c_\alpha V_\alpha$.
The bounded-support argument above makes every feature integrable.
Every term involving a fourth or fifth power contains a zero-mean
factor $X_i$ or $X_i^2-1$ in a different coordinate.  Independence
therefore gives the exact identity
\[
 \E M_c(X_1,X_2,X_3)=q_c(\gamma,\rho)
 =\sum_{\substack{\alpha\in\mathcal I\\\alpha\text{ uses only }0,3,4}}
     c_\alpha\gamma^{\#3(\alpha)}\rho^{\#4(\alpha)}.
\]
There are ten surviving monomials and 49 zero-expectation controls.
The controls improve the pointwise majorant without adding moment
parameters.

The second checked numerical assertion uses 111 rational rows forming a consecutive closed cover of
$[-7/2,3/2]$. They satisfy
\begin{align*}
 M_{c^{(j)}}(x_1,x_2,x_3)&\ge\ind_{\{x_1+x_2+x_3\le b_j\}}
                    &&((x_1,x_2,x_3)\in[-7,7]^3),\\
 q_{c^{(j)}}(\gamma,\rho)&\le\Phi(a_j/\sqrt3)+\frac{59}{250}\rho
                    &&((\gamma,\rho)\in\mathcal D_3),
\end{align*}
where $\mathcal D_3$ replaces $2$ by $15/8$ in the rectangular
bounds defining $\mathcal D_2$.
Symmetry reduces the sign patterns to four ordered-coordinate
orthants.  A box is discarded only when it has no ordered point.
On every other box, degree-five tensor Bernstein bounds establish
the pointwise inequalities.  On a box crossing the event plane,
the included side is checked by $M-1\ge s(x_1+x_2+x_3-b_j)$ with
$s\le0$, and the other side by the analogous inequality with $s\ge0$.
Existence of a slope is checked by intersecting exact rational intervals;
zero plane coefficients are checked directly.  This retains the closed
event boundary.

As in the two-summand case, expectation bounds the discrepancy on each
threshold interval.  The exterior bounds are
$3/(3+H^2)\le12/61<59/250$ for $H\le-7/2$ and
$1-\Phi(H/\sqrt3)<59/250$ for $H\ge3/2$.
They contradict the violating maximizer and prove the third assertion.
The first assertion was proved above.  Rational square-root and Machin
bounds give $K_2<0.405<0.409<\CE$ and $K_3<0.409<\CE$.

\subsubsection{Localization at four summands}
\label{j:finite:n4-localization-section}

The third bound is needed quantitatively.  Merely knowing $C_3\le\CE$
would not give the following support reduction.
Write
\[
 K_3=\frac{59\sqrt3}{250},\quad c_0=\frac{256}{625},\quad
 T=\frac{4097}{10000};\qquad K_3<c_0<T<\CE.
\]
For any standardized law set
$\beta=\E[X|X|]$, $u=\E X_+=\E X_->0$,
$v_\pm=(1\pm\beta)/2$, and
\begin{equation}\label{j:finite:n4-reference}
 a=\frac{v_-}{u},\quad b=\frac{v_+}{u},\quad
 w=\frac{4u^2}{1-\beta^2},\quad
 Q=\frac ua\delta_{-a}+(1-w)\delta_0+\frac ub\delta_b.
\end{equation}
Cauchy's inequality gives $0<w\le1$ and makes $Q$ a standardized
probability law.  Put $r=\rho(Q)=(1+\beta^2)/\sqrt{w(1-\beta^2)}$.
\begin{lemma}\label{j:finite:n4-excess}
The nonnegative quantity
\[
 W(X)=X_+(X_+-b)^2+X_-(X_--a)^2
\]
has expectation $\eta=\rho(P)-r$.  Thus $r\le\rho(P)$, with
equality exactly when $P=Q$.  Moreover, writing $c(X)=b$ for $X>0$
and $c(X)=a$ for $X<0$,
\begin{align}
 \E X^3&=b-a+\E[\operatorname{sgn}(X)W(X)],\notag\\
 \E X^4&=a^2-ab+b^2+\E[(|X|+2c(X))W(X)]
 \label{j:finite:n4-physical-moments}
\end{align}
whenever the fourth moment is finite.
\end{lemma}
\begin{proof}
On either half-line put $y=|X|$ and $c=a$ or $b$.  The identities
$y^3=y(y-c)^2+2cy^2-c^2y$ and
$y^4=(y+2c)y(y-c)^2+3c^2y^2-2c^3y$ integrate using
$\E[y;\text{half-line}]=u$ and
$\E[y^2;\text{half-line}]=uc$.
If $\E W=0$, the support is contained in $\{-a,0,b\}$ and the
first two moments fix the masses.  Conversely $P=Q$ has $W=0$.
\end{proof}
The three-point law is a reference;
the original law remains arbitrary.

\begin{proposition}\label{j:finite:n4-localization}
If $C_4>c_0$, a positive closed-CDF maximizing pair satisfies
\[
 \supp P\subset(-101/25,101/25),\quad\rho(P)<46/25,
 \quad 625/2116<w\le1,\quad|\beta|<321/500.
\]
In particular it belongs to the closed numerical class with
$|X|\le17/4$ and $|\beta|\le13/20$.
\end{proposition}
\begin{proof}
Use Theorem~\ref{j:core:attainment}.  The affine inequality
\eqref{j:external:affine} first gives
$\rho<979013/532500<46/25$, without a support assumption.
Since $r\le\rho$ and $w\le1$, the inequality
\[
 \frac{(1+(321/500)^2)^2}{1-(321/500)^2}>(46/25)^2
\]
gives $|\beta|<321/500$.
Put $C=C_4$, $z=H/2$, $s=\sqrt{4/3}$.  The contact expression,
using the quantitative bound on the three-summand predecessor, obeys
\begin{align*}
 I_{\rm new}(y)\le{}&\mathcal H_4(z,y)+4\rho(K_3-C)
                    +4\rho(s-1)K_3\\
 &+C\{-|y|^3+\tfrac32\rho y^2+3\beta y-\rho/2\},
\mathcal H_4(z,y)={}&8\!\left\{\Phi\left(\frac{z-y/2}{\sqrt{3/4}}\right)-\Phi(z)\right\}
                 \\
 &+4y\phi(z)+z\phi(z)(y^2-1).
\end{align*}
Taylor expansion in the translation and the variance gives, uniformly
in $z,y$,
\[
 |\mathcal H_4(z,y)|<\frac{\phi(0)}6|y|^3+\frac{|y|}{5}+\frac1{15}.
\]
Here $\sup|\phi''|=\sup|(x\phi(x))'|=\phi(0)<2/5$ and
$\sup|x(3-x^2)\phi(x)|<3/5$ suffice.  For the last bound, on
$t=x^2\le3$ compare $e^t$ with its cubic Taylor polynomial and use
\[
 \tfrac94(1+t+t^2/2+t^3/6)-t(3-t)^2
 =\tfrac34\{7(t-9/14)^2+3/28\}+\tfrac58t^2(3-t)>0.
\]
For $t\ge3$, use $|x|^3\phi(x)<1/2$.
The second derivative of
$\Phi((z-y/2)/\sqrt{1-\lambda})$ has absolute value below $4/15$
on $0\le\lambda\le1/4$, giving the remainder $1/15$ after
multiplication by eight.

Discard $4\rho(K_3-C)<0$ and use $4(s-1)<5/8$.
For $R=|y|\ge4$, the remaining coefficient of $C$ is at most
$B(R)=-R^3+69R^2/25+963R/500+23/100$.
All coefficients of $B(4+t)$ are negative, so substituting $C>c_0$
is now valid.  Consequently
\[
 I_{\rm new}(y)<-\frac{643}{1875}R^3+
 \frac{17664}{15625}R^2+\frac{77257}{78125}R+\frac{7541}{46875}.
\]
The coefficients after translation by $R=101/25+t$ are all negative.
This excludes $|y|\ge101/25$ from the closed support by contact.
Finally $r^2\ge1/w$ gives $w>625/2116$.
\end{proof}



\subsection{The physical comparison at four summands}\label{j:finite:n4physical}
Write $F_{4,P}(H)=\Pp_P(S_4\le H)$, $T=4097/10000$, and
$R_*=17/4$. The preceding localization reduces a contrary extremizer
to $|X|\le R_*$, $\rho<46/25$, $625/2116\le w\le1$,
$|\beta|\le13/20$. The reference $Q$ and excess $\eta=\rho-r$
retain the original marginal; they do not replace it by a three-point law.

\paragraph{One measure carries all frequencies.}
For $\eta>0$ set $\nu=W\,dP/\eta$, tagged by
$(\sigma,y)=(\operatorname{sgn}X,|X|)$, and write $c_+=b,c_-=a$,
$p_0=1-w$, $g(q)=\varphi_Q(q)$. For
$r_q(\sigma,y)=\exp(i\sigma q\,\cdot)[0,c_\sigma,c_\sigma,y]$,
Hermite interpolation at $0,c,c$ gives
\begin{equation}\label{j:finite:n4measure}
 \varphi_P(q)-g(q)=\eta\int r_q\,d\nu,\qquad |r_q|\le |q|^3/6.
\end{equation}
The linear and quadratic terms agree on each sign class; the constant
terms agree after both classes and the zero atom are included.
Define $k_0=(c_\sigma^2y)^{-1}$, $k_2=(y-c_\sigma)^{-2}$ and
$k_1=(y-c_\sigma)^{-1}$ off the contacts, which have zero $\nu$ mass.
Multiplication by $W=y(y-c)^2$ and integration give
\[
 \eta\int k_0\,d\nu=p_0-P\{0\}\le p_0,\quad
 \eta\int_\sigma k_2\,d\nu=u-c_\sigma P\{\sigma c_\sigma\}\le u,\quad
 \int_\sigma k_1\,d\nu=0.
\]
Cauchy and the $k_2$ bound prove absolute integrability of $k_1$.
At $\eta=0$, $P=Q$ and no measure division is used. Also $a,b\le R_*$,
since $\E[y^2;\sigma]\le R_*\E[y;\sigma]$.
Thus for fixed nonnegative $\lambda_0,\lambda_\pm$ and signed
$\theta_\pm$, any whole-domain bound
\[
 S_H\ge\sup_{\sigma,\,0<y\le R_*}
 \{A_H-\lambda_0k_0-\lambda_\sigma k_2-\theta_\sigma k_1\}
\]
with $A_H=\Re\sum_j4\kappa_j(H)g(q_j)^3r_{q_j}$ gives
\begin{equation}\label{j:finite:n4priced}
 \eta\int A_H\,d\nu\le
 \Gamma_H(\eta):=p_0\lambda_0+u(\lambda_++\lambda_-)+S_H\eta.
\end{equation}
The signed frequency terms are summed before taking this supremum.
Neither $S_H$ nor $\Gamma_H$ is required to be nonnegative.

\paragraph{The nonlinear response.}
For $\psi=g+\delta$, the exact identity is
\[
 \psi^4-g^4=4g^3\delta+\delta^2(\psi^2+2\psi g+3g^2).
\]
If $|g|\le A\le1$, $|\psi|\le1$, $|\delta|\le E$, valid
remainder bounds are
\[
 E^2(p^2+2pA+3A^2),\quad p=\min(1,A+E),\qquad
 v(4A^2+2-v),\quad v=\min(E^2,2A^2+1).
\]
The second follows from
$12\delta^2\int_0^1(1-t)(g+t\delta)^2dt$ and
$|g+t\delta|^2=(1-t)|g|^2+t|\psi|^2-t(1-t)|\delta|^2$.
The sharper disk bound used by some cells is also finite algebra.
After rotating $g=a\ge0$, put $x=\Re\delta$, $s=|\delta|^2$;
the squared remainder equals
\[
 s^2\{36a^4+4a^2s+s^2+8ax(6a^2+s)+24a^2x^2\}.
\]
Convexity in $x$ leaves collinear endpoints or $|\psi|=1$.
With $r_+=\min(E,1)$, $m=\min(A,1-r_+)$, $r_-=\min(E,1+A)$,
the collinear modulus bounds are
\[
 r_+^2(6m^2+4mr_++r_+^2),\quad r_-^2(6A^2-4Ar_-+r_-^2),\quad
 \begin{cases}r_-^4,&r_-\le1,\\r_-^2(3r_-^2-8r_-+6),&r_-\ge1.
 \end{cases}
\]
On the unit circle maximize
$Q_A(s)=s^2[18A^4+12A^2+6-(12A^2+8)s+3s^2]$
on $[(1-A)^2,\min(E^2,(1+A)^2)]$: only its endpoints and
$s_-=(3A^2+2-A\sqrt{4-3A^2})/2$, if inside, are needed.
Convexity in $a^2$ reduces the remaining disk radii to $a=A$ or a
collinear endpoint. The largest collinear value and square root of
this circle maximum therefore bound the modulus; the minimum with
either preceding complete bound is valid.
On an excess cell $[L,U]$, an affine radius $e_j(\eta)\ge0$ is
used only when proved throughout the cell. Freezing
$p_j=\min(1,a_j+\max(e_j(L),e_j(U)))$ gives a quadratic remainder
with nonnegative leading coefficient. Freezing the entire radius
instead gives a valid constant disk bound. Bounds are combined as
complete alternatives, never coefficientwise.

\paragraph{From a finite frequency sum to every threshold.}
The closed-CDF majorant of Lemma~\ref{j:fourier:closed-majorant}, with
$\widehat K(t)=(1-t)+iR(t)$ and
$R(t)=(1-t)\cot(\pi t)+1/\pi$, gives
\[
 F_{4,P}(H)\le J_P(H/2),\quad
 J_P(x)=\tfrac12+\int_0^1
 \Re\{\widehat K(t)e^{-6itx}\varphi_P(3t)^4\}\,dt.
\]
This is a real compensated integral, including the closed boundary.
Put $Z=S_4/2$, $D_x=x-Z$, $S(t)=tR(t)$. Its integrand is
$(1-t)\E\cos(6tD_x)+S(t)\E[6D_x\operatorname{sinc}(6tD_x)]$.
The cotangent partial fractions give $0\le S\le1/3$,
$|S'|<2$, $|S''|<5$, and therefore
\[
 |F_x''(t)|\le42\E|D_x|+108\E D_x^2+24\E|D_x|^3.
\]
For the derivative constants, write
$S=\pi^{-1}-2\pi^{-1}(A+MG)$ with
$A=t^2/(1+t)$, $M=t^2(1-t)$ and
$G=\sum_{j\ge2}(j^2-t^2)^{-1}$.
On $[0,1]$, $(|A'|,|A''|,|M|,|M'|,|M''|)\le(3/4,2,1/6,1,4)$
and $(G,|G'|,|G''|)\le(3/4,10/27,106/135)$, by telescoping and
$j^2-t^2\ge3j^2/4$. For $t>1/2$, put $v=1-t$ and use $v\cot(\pi v)\le1/\pi$
to obtain $R(t)\ge0$; below $1/2$ positivity is immediate.
These estimates imply the stated bounds; the sinc
integral representation justifies differentiation by the third moment.

From \eqref{j:finite:n4-physical-moments} choose whole-cell bounds
\[
 \kappa_*\ge\min(\rho,|b-a|+\eta),\quad
 M_{4,*}\ge\min(R_*\rho,a^2-ab+b^2+(R_*+2\max(a,b))\eta).
\]
Then $D_2=x^2+1$, $D_4=x^4+6x^2+2|x|\kappa_*+(M_{4,*}+9)/4$
bound $\E D_x^2,\E D_x^4$, and Cauchy bounds the odd absolute moments.
For $N=256$, $t_j=(j+1/2)/N$, $q_j=3t_j$,
$\kappa_j(H)=\widehat K(t_j)e^{-iq_jH}/N$, set
\begin{gather*}
 C_H(Q)=1+2\Re\sum_j\kappa_j(H)g(q_j)^4-2\Phi(H/2),\\
 Q_H=\min\left\{\frac{687}{8N^2},
 \frac{42\sqrt{D_2}+108D_2+24\sqrt{D_2D_4}}{12N^2}\right\}.
\end{gather*}
This midpoint reserve bounds the original physical integral once.
The constant alternative uses $|x|\le2$, $D_2\le5$, $D_4<54$.
To bridge adjacent raw thresholds $H_\ell=-4+\ell/8$,
$0\le\ell\le57$, use
$P_j=\min(1,a_j+e_j+3/(2N))$, where $a_j\ge|g(q_j)|$ and
$e_j\ge|\varphi_P(q_j)-g(q_j)|$ on the excess cell.
The increment follows from $\E|X|\le1$.
If $d_j$ is the smaller of the exact integrals over the frequency cell of
$t^2(1-t)+t/3$ and $t^2(1-t)+t(1-t)/2$, then
$|J_P''|\le B=36\sum_jd_jP_j^4$.
Indeed both polynomials bound $t^2|\widehat K(t)|$; for the second,
retaining the first cotangent pole gives $S(t)<(1-t)/2$.
The single threshold-interpolation reserve is
$E_H=(B+1/4)/(16\cdot8^2)$, since $\sup|\Phi''|<1/4$.
The radii $\eta q_j^3/6$, $\max(a,b)^2p_0q_j^2$, and $1+a_j$
are all valid; the second uses
$y f[0,c,c,y]=f[c,c,y]-f[0,c,c]$, so $|r_q|\le q^2/y$, and
then the $k_0$ budget. Thus no contact division occurs.

\paragraph{The finite inequalities.}
Let $M_H$ bound $C_H(Q)-Tr$ on a geometry rectangle, and let
$R_2(\eta)$ bound the weighted nonlinear response on $[L,U]$.
The assembled bound is
\begin{equation}\label{j:finite:n4low}
 2\{F_{4,P}(H)-\Phi(H/2)\}-T\rho
 \le P_H(\eta):=M_H+Q_H+E_H+2\Gamma_H(\eta)+2R_2(\eta)-T\eta.
\end{equation}
It is convex quadratic in $\eta$, so its two nonpositive endpoints
suffice. The optional score bound checks both denominator faces
$46/25$ and $r_L+\eta$. For the latter the derivative numerator has
derivative $P_H''(\eta)(r_L+\eta)\ge0$, hence no interior maximum.
All functions defining $M_H,S_H,R_2$ are explicit finite expressions.
Their whole-cell enclosures have the following licenses.
For the reference, Taylor's theorem uses
$\partial_i C_H=8\Re\sum_j\kappa_jg^3\partial_i g$ and
$|\partial_{ik}C_H|\le8\sum_j|\kappa_j|(3D_iD_k+D_{ik})$.
A supporting plane bounds $r$ below: for
$f(\beta)=(1+\beta^2)/\sqrt{1-\beta^2}$,
$3ff''-(f')^2=(9+9\beta^2+15\beta^4-\beta^6)/(1-\beta^2)^3>0$
proves convexity of $r=w^{-1/2}f(\beta)$.
The response is evaluated without singular divided quotients by
\[
 r_q=\sum_{m=0}^{60}(i\sigma)^{m+3}b_mh_m+E_q,\quad
 |E_q|\le(q^3/6)(qR_*)^{61}/61!,
\]
where $b_0=q^3/6$, $b_m=b_{m-1}qR_* /(m+3)$,
$h_0=1$, $h_m=(y/R_*)h_{m-1}+(m+1)(c/R_*)^m$.
Hermite--Genocchi proves this expansion and
$|\partial_y r_q|\le q^4/24$, $|\partial_c r_q|\le q^4/12$.
The 273 nearest-node cells $y_k=k/64$, clipped to $[0,R_*]$,
therefore need radius $q^4/(48\cdot64)$.
The geometry correction uses
$|g^3-g_0^3|\le|g-g_0|(|g|^2+|g||g_0|+|g_0|^2)$.
The nonnegative prices are rational data. Each fixed choice is bounded
on every tagged cell; $k_0\ge1/(c_U^2y_U)$, while the other cost is
$-\lambda_\sigma z^2-\theta_\sigma z$, $z=1/(y-c)$.
Its endpoints and admissible vertex, on both reciprocal rays if the
cell meets a contact, bound it. Require $\theta_\sigma=0$ when
$\lambda_\sigma=0$. These rules prove the priced supremum over the
whole tagged domain. Factoring the difference of fourth powers also gives
$|\psi^4-g^4|\le E(p^3+p^2A+pA^2+A^3)$, $p=\min(1,A+E)$.
This absolute bound, or the full linear-plus-remainder expression, may be
used as a complete alternative.

For large excess, condition on $j$ nonnegative summands, assigning
zero to that class. With $p=P(X\ge0)$,
\[
 \mu_j=\frac{u(j-4p)}{2p(1-p)},\quad
 V_j=\frac14\left[j\left(\frac{v_+}{p}-\frac{u^2}{p^2}\right)
 +(4-j)\left(\frac{v_-}{1-p}-\frac{u^2}{(1-p)^2}\right)\right].
\]
Cantelli bounds each conditional CDF by $1$ if $x\ge\mu_j$ and by
$V_j/(V_j+(\mu_j-x)^2)$ otherwise. The closed sign-mass interval is
$w(1-\beta)/2\le p\le1-w(1+\beta)/2$.
On each threshold cell use the upper endpoint in Cantelli and the
lower endpoint in $\Phi$; degree-four Bernstein bounds retain the
joint binomial mixture. A sign-mass cell supplies numerator $d_i\ge0$
and correlated lower $r_i\le r$, giving score $d_i/(r_i+\eta)$.
With $r_B\le r$ on the parent, set
\begin{gather*}
 c_C^0=\max\{0,\max_i(d_i/T-r_i),2/(5T)-r_B\},\qquad
 c_C=10^{-4}\lceil10^4c_C^0\rceil+10^{-4},\\
 c_S=\max(0,979013/533000-r_B),\qquad c_B=\min(c_C,c_S).
\end{gather*}
Cantelli gives the first cutoff; the affine bound gives the second.
The two threshold rays $H\le-4$ and $H\ge25/8$ have score at most
$2/5$ by Cantelli and $\Phi(25/16)>4/5$. Thus $[c_B,\infty)$ is
covered, and it meets the low-excess interval at its included endpoint.

\begin{proposition}[Four-summand numerical assertion]\label{j:finite:n4numerical}
Divide $[625/2116,1]\times[0,13/20]$ into $20\times13$ closed
rectangles. A rectangle with $\min r^2\ge(46/25)^2$ is excluded.
Every other parent, or all four of its coordinate-bisected children,
passes the preceding large-excess bound and, on sixteen equal closed
subintervals of $[0,c_B]$, the low-excess inequalities at all 58
thresholds, their joins, both tails and both reflected orientations.
All 260 dispositions and all five parent/child request slots per parent
are verified by exact acceptance. Consequently $C_4\le T<\CE$.
\end{proposition}
\begin{proof}[Derivation from the numerical assertion]
Every localized contrary law belongs to a closed parent. Its reference
and common defect measure satisfy the derived inequalities on that parent
or one of its four children. The midpoint and threshold reserves cover
all central thresholds, and the two tails cover the rest.
Every $\eta\ge0$ belongs to the low or large-excess range, including
$\eta=0$ and the common endpoint. Hence its score is at most $T$,
contradicting $C_4>T$. Reflection accounts for the other discrepancy sign.
\end{proof}

\subsection{One support criterion for the remaining finite indices}
\label{j:finite:support}

The support calculations throughout the finite range use one elementary
criterion. Its inputs are a lower bound for the attained score, upper
bounds for moments of the same marginal, and a predecessor allowance.
The Gaussian calculation keeps the actual integer $n$ throughout.

\begin{lemma}[Direct Gaussian support criterion]
\label{j:finite:direct-gaussian-support-criterion}
Let $n\ge2$, and let $(P,H)$ be an attained positive closed-event
ratio maximizer with score $C$. Write
$\rho=\E|X|^3$, $\beta=\E[X|X|]$ and
$A=\sqrt{n-1}\Delta_{n-1}(P)/\rho$.
Suppose known numbers $c,r,b_-,b_+,S,e,q$ satisfy
\[
 0<c\le C,\quad 1\le\rho\le r,\quad
 -1<b_-\le\beta\le b_+<1,\quad q\ge0,
\]
\[
 \sqrt{n/(n-1)}\le S,\qquad
 n(\sqrt{n/(n-1)}-1)-1/2\le e,\qquad
 n(A-C)_+\le q.
\]
Fix $R\ge1$ and a support sign $\eta\in\{-1,1\}$, and put
$B_+=b_+$, $B_-=-b_-$. Define
\[
 Q_\eta(v)=-v^3+\tfrac32rv^2+3B_\eta v+re,
 \qquad
 L_\eta(v)=\tfrac25(v^2/2+1/2+e)+cQ_\eta'(v).
\]
Suppose that
\begin{equation}\label{j:finite:direct-support-ray-gates}
 Q_\eta(R),Q_\eta'(R),Q_\eta''(R),
 L_\eta(R),L_\eta'(R)<0,
 \qquad 1/5-3c<0.
\end{equation}
Define the actual-index Gaussian expression
\[
 G_n(H,t)=n\sqrt n\{\Phi((H-t)/\sqrt{n-1})-\Phi(H/\sqrt n)\}
 +nt\phi(H/\sqrt n)+\tfrac H2(t^2-1)\phi(H/\sqrt n),
\]
and suppose $U_\eta$ satisfies
\[
 \sup_{H\in\R}\eta G_n(H,R)\le U_\eta,
 \qquad U_\eta+cQ_\eta(R)+rSq<0.
\]
Then $\supp P$ contains no point in the ray
$\{\eta v:v\ge R\}$. The two signs can be checked separately,
with different signed-moment bounds.
\end{lemma}

\begin{proof}
Contact gives
\[
 0\le G_n(H,t)+
 C\{-|t|^3+\tfrac32\rho t^2+3\beta t+\rho e_n\}
 +\rho s_n n(A-C),
\]
where $s_n=\sqrt{n/(n-1)}$ and $e_n=n(s_n-1)-1/2$.
The last term is at most $rSq$. At $t=\eta v$ with $v\ge R$,
the bracket is at most $Q_\eta(v)$.
Since $Q_\eta'''=-6$, the first three gates make this polynomial
strictly negative on the full ray. Thus multiplication by $C\ge c$
gives an upper bound $cQ_\eta(v)$.

The derivative is
$n\phi(z)+Ht\phi(z)-ns_n\phi(s_ny)$, where
$z=H/\sqrt n$ and $y=(H-t)/\sqrt n$.
Taylor's theorem and $\|\phi''\|_\infty=\phi(0)$ bound
$|n\phi(y)-n\phi(z)-Ht\phi(z)|$ by $\phi(0)t^2/2$.
The identity $\partial_s[s\phi(sy)]=-\phi''(sy)$ bounds the
remaining scaling error by $n\phi(0)(s_n-1)$.
This proves, for every $n\ge2$,
$|\partial_tG_n(H,t)|\le\phi(0)(t^2/2+1/2+e)$.
Therefore the derivative in $v$ of
$G_n(H,\eta v)+cQ_\eta(v)+rSq$ is at most $L_\eta(v)$.
Its second derivative as a polynomial bound is
$L_\eta''=2/5-6c<0$; the last two endpoint gates imply
$L_\eta(v)<0$ on $[R,\infty)$.
Finally $G_n(H,-R)=-G_n(-H,R)$, so the assumed $U_\eta$
bounds $G_n(H,\eta R)$ uniformly in $H$.
The strict endpoint inequality and monotonicity contradict contact
on every point of the specified support ray.
\end{proof}

\paragraph{A finite certificate for the whole threshold line.}
Let $z=H/\sqrt n$ and $a=(H-R)/\sqrt{n-1}$. Direct
differentiation gives
\[
 \partial_HG_n(H,R)=ns_n\phi(a)+
 \left[-n-RH+\tfrac12(R^2-1)(1-H^2/n)\right]\phi(z).
\]
Choose a positive rational $H_0$ such that
$H_0\ge5\sqrt n$ and $H_0-R\ge5\sqrt{n-1}$.
An outward enclosure of the square roots makes this choice exact.
On the whole real line,
\[
 |\partial_H^2G_n(H,R)|\le M_{n,R}
 :=\frac{\sqrt n}{4}\left(\frac n{n-1}+1\right)
   +\frac25R+\frac{2|R^2-1|}{5\sqrt n}.
\]
On a closed threshold cell with midpoint $h$ and half-width $w$,
let $[g_-,g_+]$ and $[d_-,d_+]$ enclose $G_n(h,R)$ and
$\partial_HG_n(h,R)$. Taylor's theorem then gives the two upper
bounds
\[
 g_++\max(|d_-|,|d_+|)w+M_{n,R}w^2/2,
 \qquad
 -g_-+\max(|d_-|,|d_+|)w+M_{n,R}w^2/2.
\]
They bound $G_n(H,R)$ and $-G_n(H,R)$ respectively throughout
the cell. The certificate includes every consecutive cell from
$-H_0$ to $H_0$, including their endpoints.

On the two exterior rays, both distribution-function arguments lie
in the same normal tail, so
\[
 |G_n(H,R)|\le n\sqrt n\,\tau+
 \left(nR+\tfrac12H_0|R^2-1|\right)\phi(H_0/\sqrt n),
 \quad |H|\ge H_0,
\]
where $\tau=3\cdot10^{-7}>\Phi(-5)$ is sufficient.
Take the largest positive-cell bound and the exterior bound for
$U_+$, and the largest negative-cell bound and the exterior bound
for $U_-$. This proves a statement about the entire threshold line,
not just about its grid points.



\begin{proposition}[Actual-index support rows]\label{j:finite:supportrows}
Under $C_{n-1}\le\CE<C_n$, each row below bounds the support of the
same positive closed maximizer by $(-R,R)$. Repeated index ranges mean
one check at every actual integer, not a transfer between integers.
\begin{center}\small
\begin{tabular}{c|c|c}
$n$&third-moment interval&$R$\\\hline
5&$[1,11/10],[11/10,23/20],[23/20,237/200],[237/200,6/5]$
 &$11/5,47/20,49/20,99/40$\\
5&$[6/5,5/4],[5/4,7/5],[7/5,3/2]$&$13/5,117/40,25/8$\\
5&$[3/2,8/5],[8/5,17/10],[17/10,7361/4000]$&$133/40,7/2,15/4$\\
6--14&$[1,11/10],[11/10,23/20]$&$23/10,12/5$\\
6--9&$[23/20,237/200],[23/20,6/5]$&$5/2,13/5$\\
6--9&$[6/5,7/5],[7/5,3/2],[3/2,8/5]$&$3,16/5,17/5$\\
6--9&$[8/5,17/10],[17/10,7361/4000]$&$18/5,19/5$\\
10--14&$[23/20,71/50]$&$3$
\end{tabular}
\end{center}
The last five lines comprise all 26 low-moment and 25 complementary
actual-index rows. Their endpoint tests are instances of the same
criterion, including both support rays and all real thresholds.
\end{proposition}
\begin{proof}[Mathematical meaning and finite inequalities]
Use $c=819/2000<\CE$, $q=0$, and $e\le91/1000$.
For each upper moment endpoint $r$, the signed-moment bound supplies
$|\beta|\le b$, where $b$ is the outward rational root of
$(1+b^2)^2\ge r^2(1-b^2)$ recorded in the support certificate.
The certificate verifies the five strict polynomial gates and the strict
Gaussian endpoint test. For $n=5$ use $H_0=16$; for $6\le n\le9$
use $H_0=35/2$, enlarged to 18 for the complementary rows at $n=9$;
for $10\le n\le14$ use $H_0=22$. All threshold meshes have width
$1/10$. The value at each endpoint, enlarged by $M_{n,R}/800$,
bounds both signs on the intervening cell: the interpolation error for
$|f''|\le M$ is $M h^2/8$, by comparison with
$\pm M(x-a)(b-x)/2$. The displayed exterior-ray argument applies
because $H_0^2\ge25n$ and $(H_0-R)^2\ge25(n-1)$.
All numerical comparisons are strict rational inequalities with the
enclosures of Section~\ref{j:cert:arithmetic}. Thus the preceding support
criterion applies before any of these radii enters a Fourier estimate.
Additional actual-index rows at $15$--$39$, and the separate two-endpoint
support test at $40$--$42$, use the same contact license; their complete
input domains are part of the numerical propositions below.
\end{proof}

\subsection{Finite-frequency refinements}
\label{j:finite:small-fourier}
At a small index, replacing a finite frequency integral by an integral
over the whole half-line can lose much of the available margin. The
following estimates retain the finite interval and the actual integer
power. They use the moments of the same supported marginal throughout.

\subsubsection{A finite-interval compensated estimate}
Let $n>2$ be an integer, $\E X=0$, $\E X^2=1$, $|X|\le R\le6$,
$\gamma=\E X^3$, $\mu_j=\E X^j$, and $f(t)=\E e^{itX}$.
Assume $\mu_4\le M\le111/10$ and $|\gamma|\le g\le37/20$.
The simultaneous coefficient and moment bounds in the signed moment atlas
of Section~\ref{j:local:moments} give
\[
 |a|\le A_*,\quad |z|\le Z,\quad \mu_6\le M_6,
 \qquad a=(3-\mu_4)/24,\quad z=(10\gamma-\mu_5)/120.
\]
Define
\[
 E=\frac{M_6}{720}+\frac{g^2}{72}
       +\frac{(1/2+g^2/144)^3}{5},\qquad
 B=A_*+Z/2+E/4,
\]
\[
 a_n=(n-1)/2,\qquad b_n=(n-2)(45-M)/96,\qquad
 I_j(A)=\int_0^{1/2}t^j e^{-At^2}\,dt.
\]
All these quantities are nonnegative, and the two damping rates are
positive. Put
\begin{align*}
 L_0={}&\sqrt n\left\{nZI_2(a_n)
       +\frac{n(n-1)A_*g}{6}I_4(a_n)
       +\frac{n^3g^3}{1296}I_6(n/2)\right\},\\
 L_H={}&\sqrt n\left\{nA_*I_2(a_n)
                              +\frac{n^2g^2}{72}I_4(n/2)\right\},\\
 L_2={}&n\left\{nEI_3(a_n)+\frac{n(n-1)B^2}{2}I_5(b_n)\right\}.
\end{align*}
For every real raw threshold $H$, the complete low-frequency estimate is
\begin{equation}\label{j:finite:finite-low}
 \int_0^{1/2}\frac{\left|\Im\left[e^{-itH}
 \{f(t)^n-e^{-nt^2/2}(1-in\gamma t^3/6)\}\right]\right|}{t^3}\,dt
 \le\frac{L_0+L_H|H|}{\sqrt n}+\frac{L_2}{n}.
\end{equation}

To prove this finite-interval version, use the Taylor decomposition
\eqref{j:fourier:taylor} with $h(t)=e^{-t^2/2-i\gamma t^3/6}$.
The same polynomial-power remainder gives
$|f^n-h^n-n(f-h)h^{n-1}|
\le n(n-1)B^2t^8e^{-b_nt^2}/2$.
Retain the integration interval $[0,1/2]$ in that proof.
The leading rotated imaginary terms are bounded by
$ne^{-a_nt^2}[(Z+A_*|H|)t^5+(n-1)A_*gt^7/6]$;
the cubic Gaussian phase remainder contributes
$e^{-nt^2/2}(|n\gamma t^3/6|^3/6+
|H|t(n\gamma t^3/6)^2/2)$.
Division by $t^3$ and integration give exactly $L_0,L_H,L_2$ above.
The damping and all Taylor hypotheses are those of the same proof;
no half-line extension or asymptotic replacement is made.

For the small indices the quantities $I_j(A)$ are enclosed directly.
The degree-20 and degree-21 Taylor polynomials of $e^{-x}$ are,
respectively, upper and lower bounds on all $x\ge0$, by the sign of
the integral remainder. Their integrated values are the rational sums
\[
 \sum_{k=0}^{m}\frac{(-A)^k}{k!(j+2k+1)2^{j+2k+1}},
 \qquad m=20,21.
\]
The selected routines require $0\le A\le8$ and $0\le j\le16$;
the rates above meet those gates for $5\le n\le9$.
Every coefficient multiplying an integral is nonnegative. When two
complete triples $(L_0,L_H,L_2)$ are available, either may be selected
by a rule independent of $H$. Selecting different coefficients from
different triples is not justified by this argument.

\subsubsection{Containing rectangles and exact finite powers}
The supported fourth-moment bound also improves the real coordinate.
If $\rho\in[\rho_l,\rho_u]$, $|X|\le R$, and $\mu_4\le M_4$, then
\[
 \Re f(t)\le1-t^2/2+M_4t^4/24.
\]
If $R^2t^2\le30$, a second lower endpoint is
\[
 \Re f(t)\ge1-t^2/2
              +\rho_l^2t^4(1-R^2t^2/30)/24.
\]
These follow from the global fourth-degree upper and sixth-degree
lower Taylor polynomials of cosine, $\mu_6\le R^2\mu_4$, and
$\mu_4\ge\rho^2$ by Cauchy's inequality. The nonnegative coefficient
gate is needed before substituting $\rho_l^2$. Intersect these real
intervals with the anchored bounds of Lemma~\ref{j:fourier:anchor}.

\begin{lemma}[Two corners of a complex-power rectangle]
\label{j:finite:two-corners}
Let $n\ge2$ be an integer, $a\ge0$, $0<x_0\le x_1$ and $Y\ge0$.
If
\[
 (n-2)Y\le3x_0/2,\qquad (2n-1)x_0^n\ge(n-1)a,
\]
then for $x\in[x_0,x_1]$ and $|y|\le Y$,
\[
 |(x+iy)^n-a|^2
 \le\max\{|(x_0+iY)^n-a|^2,|(x_1+iY)^n-a|^2\}.
\]
\end{lemma}
\begin{proof}
Write $z=x+iy$, $r=|z|$ and $F(x,y)=|z^n-a|^2$.
Since $|\arg z|\le Y/x_0$, the first gate and $3/2<\pi/2$ imply
$\Re z^{n-2}\ge0$. Direct differentiation gives
\begin{align*}
 F_{yy}&=2nr^{2n-4}\{r^2+2(n-1)y^2\}
                         +2an(n-1)\Re z^{n-2}\ge0,\\
 F_{xx}&=2nr^{2n-4}\{r^2+2(n-1)x^2\}
                         -2an(n-1)\Re z^{n-2}.
\end{align*}
Divide the latter by $2nr^{n-2}>0$. Its first term is
$r^{n-2}((2n-1)x^2+y^2)\ge(2n-1)x_0^n$; the subtracted
term is at most $a(n-1)$. Thus $F$ is convex in each coordinate.
Apply convexity first in $y$ and then in $x$, and use
$F(x,Y)=F(x,-Y)$.
\end{proof}

For Fourier smoothing take $a=e^{-nt^2/2}$ and any proved containing
rectangle for $f(t)$. Both gates must hold throughout the frequency
cell. If either fails, the geometric-sum estimate remains available.
The maximum of the two pointwise corner bounds is enclosed before
integration; the maximum of their separate integrals would not suffice.
Complex powers are finite polynomial recurrences, so zero or equal
bases create no division. A fourth-Taylor endpoint is kept fixed on
each integration cell when differentiating its smooth majorant.

Cancellation can be retained when computing the real endpoints.
For example, on $0\le u\le1$,
\[
 1-u^2/2+M_4u^4/24-e^{-u^2/2}
 =\frac{M_4-3}{24}u^4
   +\sum_{k=3}^{10}\frac{(-1)^{k+1}u^{2k}}{2^kk!}+R(u).
\]
For derivative order $j=0,1,2$ with $u=ct$, the absolute remainder
is at most $2(22)_j u_{\max}^{22-j}c^j/(2^{11}11!)$; consecutive
absolute tail terms have ratio below $1/2$ on this domain.
Alternatively, at higher frequencies enclose each explicit smooth
residual $d(u)=\cos u-e^{-u^2/2}$ or
$d(u)+(\rho_u-1)u^3/6$ by its midpoint value plus the half-width
times an enclosure of its derivative on the entire cell.
Intersecting two such valid enclosures preserves their common function.
These operations differentiate scalar envelopes, not the law's CDF.


\subsubsection{A sine improvement on the acquired $26$--$30$ region}
\label{j:finite:finite-sine}
Let the positive and negative support magnitudes lie in
$[a_+,R_+]$ and $[a_-,R_-]$, respectively. Assume
\[
 0<A\le a_+^2+a_-^2,\qquad
 B\ge R_+^2+R_-^2,\qquad M\ge\max(R_+,R_-).
\]
Put $R(u)=1-Au^2(1-Bu^2/42)/20$.
For $u\ge0$, $Mu\le\pi$ and $Bu^2\le42$,
\begin{equation}\label{j:finite:finite-sine-bound}
 |\Im f(u)|^2\le(\rho^2-\zeta^2)u^6R(u)^2/36.
\end{equation}
In particular $R(u)\ge1-21A/(40B)\ge19/40$.
Outside these frequency gates the ordinary sine-circle inequality
remains available.

Here is a proof retaining the same marginal. Write
$\operatorname{sinc}t=\int_0^1\cos(vt)\,dv$. For $u\sqrt z\le\pi$,
\[
 0\le-\frac{d}{dz}\operatorname{sinc}(u\sqrt z)
 =\frac{u^2}{2}\int_0^1v^2\operatorname{sinc}(uv\sqrt z)\,dv
 \le\frac{u^2}6-\frac{u^4z}{60}+\frac{u^6z^2}{1680}.
\]
The upper bound follows by integrating the global polynomial
$\operatorname{sinc}t\le1-t^2/6+t^4/120$; nonnegativity follows
from $0\le uv\sqrt z\le\pi$.
Integrating between $a^2$ and $b^2$ gives
\[
 |\operatorname{sinc}(au)-\operatorname{sinc}(bu)|
 \le\frac{|a^2-b^2|u^2}{6}
 \left(1-\frac{(a^2+b^2)u^2}{20}
       +\frac{(a^4+a^2b^2+b^4)u^4}{840}\right).
\]
The last bracket is at most $R(u)$, since
$a^4+a^2b^2+b^4\le B(a^2+b^2)$ and
$1-Bu^2/42\ge0$. At $a=b$ use the undivided inequality.

For the centered two-point law at $a,-b$, put
\[
 e_{a,b}=\frac{2a^2b^2}{a+b},\quad
 r_{a,b}=\frac{ab(a^2+b^2)}{a+b},\quad
 h_{a,b}=\frac{abu}{a+b}
       \{\operatorname{sinc}(au)-\operatorname{sinc}(bu)\}.
\]
The identity $(2ab)^2+(a^2-b^2)^2=(a^2+b^2)^2$ gives
$\|(e_{a,b},6h_{a,b}/(u^3R(u)))\|_2\le r_{a,b}$.
Integrate against the exact centered two-point mixture from the
ordinary sine-circle proof. The three integrals are $\zeta$, $\rho$
and $\Im f(u)$, respectively. The triangle inequality proves
\eqref{j:finite:finite-sine-bound}. The case $u=0$ is immediate.

For each of the $55$ acquired rows ($26\le n\le30$,
$32\le j\le42$), retain its closed probability interval
$[p_l,p_u]$, an actual-noise bound $s\le s_u$, and an upper
bound $d>0$ on the inward sum $a_0+a_1$.
The interface requires $d<p_l$, $d\le1-2p_u$ and
\begin{equation}\label{j:finite:b1-support-map}
 \frac{p_u^2+(1-p_u-d)^2}{p_u(1-p_u)+s_u}\ge\frac{17}{10},
 \qquad R_-^2+R_+^2\le\frac{1037}{100},
 \qquad \max(R_-,R_+)\le\frac{13}{5}.
\end{equation}
Indeed the minimum of $(p-a_0)^2+(1-p-a_1)^2$ under
$a_0,a_1\ge0$, $a_0+a_1\le d\le1-2p$ is
$p^2+(1-p-d)^2$: use the full inward sum and differentiate in
$a_0$, obtaining $2(1-2p-d)+4a_0\ge0$.
This numerator decreases in $p$ on the retained interval, while
$p(1-p)+s$ increases. Thus \eqref{j:finite:b1-support-map} gives the
required standardized inner support sum. It licenses
\eqref{j:finite:finite-sine-bound} with
$A=17/10$, $B=1037/100$, $M=13/5$ on every row.

\paragraph{A direct resource-to-sine map at indices $26$--$30$.}
The three common constants in \eqref{j:finite:b1-support-map} follow
directly from the acquired floor-$25$ resource.
Write its original simultaneous intercepts as $d_0,L$, and its prices
as $b,g,h$, so that
\[
 d_{\rm floor}-b=g d_0=hL>0,\qquad b\ge0.
\]
These identities are checked exactly against the exported residual;
$L$ is the original joint intercept, even when another bound further
clips the variance. The weighted resource implication gives
\[
 \frac n{25}(b+g\delta)+h\lambda_n\le d_{\rm floor},
 \qquad
 \frac n{25}\frac{\delta}{d_0}+\frac{\lambda_n}{L}\le1.
\]
The second inequality follows by subtracting $nb/25$, dividing by
$d_{\rm floor}-b$, and using $n\ge25$. Therefore
\[
 \delta\le \delta_u:=\min\left\{\frac{25d_0}{n},
                 \frac{25d_{\rm floor}/n-b}{g}\right\}.
\]
The bounded conditional supports give $s\le\mathcal B\delta$, where
$\mathcal B=\max(O_0,O_1)\max(p_u,1-p_l)$ and $\lambda_n=ns$.
Substitution in the joint plane yields
\[
 \lambda_n\le
 \left\{L^{-1}+(25\mathcal B d_0)^{-1}\right\}^{-1}=:\lambda_u.
\]
This step uses one joint plane before optimizing. For all $55$ rows,
the choices $d=\delta_u$ and $s_u=\lambda_u/n$ satisfy the two
preceding gates and
\[
 \frac{p_u^2+(1-p_u-\delta_u)^2}
      {p_u(1-p_u)+\lambda_u/n}>\frac{17}{10}.
\]
The smallest left side is greater than $1716/1000$ and occurs
at $n=26$ on $P_{38}$. The same rows satisfy
$R_-^2+R_+^2\le1037/100$ and $\max(R_-,R_+)\le13/5$.
The actual-index support proofs and the full $100$-panel resource
cover are checked separately. Thus the finite-sine majorant applies
on every surviving panel. Its integral uses only these common
constants and the rational moment and frequency intervals below.

Here $\eta,\zeta$ denote the symmetrization quantities of
Lemma~\ref{j:fourier:surrogate}, not the four-summand excess. Their
bounds follow directly from that lemma. Since
$(13/5)^2< (11/10)^2+8$, the general bounded-support inequality
for $\eta$ gives $\eta\le b(\rho)$ throughout $[11/10,6/5]$,
where $b(r)=(\sqrt{r^2+8}-r)/2$.
Lemma~\ref{j:fourier:surrogate} supplies $\zeta\ge b(\rho)$ and the supported bound for $\eta$.
Since $b$ decreases and $r+b(r)$ increases, on a strip $[r_l,r_u]$
we may use simultaneously
\[
 \rho+\eta\le r_u+b(r_u),\qquad
 \rho^2-\zeta^2\le r_u^2-b(r_u)^2.
\]
Rational brackets $z_l,z_u$ for $b(r_u)$ are checked by
$z_l^2+r_uz_l\le2\le z_u^2+r_uz_u$ with $0<z_l\le z_u$.
The integrator uses $r_u+z_u$ and $r_u^2-z_l^2$ in their indicated
upper-bound directions. The value $z_u$ is not separately asserted
to bound $\eta$ at the lower end of the strip.

The selected $23$ closed strips cover $[11/10,6/5]$. All five
actual indices are evaluated on each strip, giving $115$ complete
four-term Fourier inequalities. The optional fourth-Taylor real
endpoint and Lemma~\ref{j:finite:two-corners} are used only when their
whole-cell gates hold. The remaining cells use a previously valid
geometric-sum or direct interval bound.

The implementation sometimes integrates the maximum of the two
corner majorants more accurately. If $f_0,f_1$ are twice continuously
differentiable on $[t_0-h,t_0+h]$, enclose their midpoint values above
by $a_i$, their midpoint derivatives by $[d_i^-,d_i^+]$, and both
second derivatives above by $C\ge0$. Taylor's theorem bounds the
integral of their pointwise maximum by
\[
 \int_0^h\max_i(a_i+d_i^+x)\,dx
 +\int_0^h\max_i(a_i-d_i^-x)\,dx+\frac{Ch^3}{3}.
\]
Each affine maximum has at most one crossing, found exactly. This
bounds the maximum before integration and assumes nothing about
how often the original smooth curves cross.

\subsection{A common count comparison}\label{j:finite:finite-counts}
The finite argument follows the number of positive summands. Conditional on
that number, the sum is a binomial lattice point plus a centered residual.
The residual can move mass across a threshold in both directions. Keeping
its outgoing and incoming contributions together is essential near a
binomial jump. Everything in this section concerns one common marginal;
no finite-atom assumption is made.

\subsubsection{Exact conditioning and the common moments}
Use the sign normalization \eqref{j:core:normalization}, writing
$h=1/\ell=pq/m$ and $v=pq$. Thus $hX=B-p+U$,
$h^2=v+s$, $\lambda=ns$ and
$w_0=\lambda(1-\alpha)/(nq)$, $w_1=\lambda\alpha/(np)$.
At $s=0$ both conditional residuals vanish and any $\alpha\in[0,1]$
may be used; zero atoms belong to $B=0$ as before.

Write $W_j$ for a sum of $n-j$ independent copies of $U\mid B=0$
and $j$ independent copies of $U\mid B=1$, all independent. For
$b_j=\binom njp^jq^{n-j}$ and $z=np+hH=k+u$,
\begin{align}
 \Pp(S_n\le H)&=\sum_{j=0}^n b_j\Pp(W_j\le k+u-j),
 \label{j:finite:count-identity}\\
 \Pp(S_n\le H)&\le F_{\mathrm{Bin}(n,p)}(k)
 -b_k\Pp(W_k>u)+\sum_{j>k}b_j\Pp(W_j\le u-(j-k)).
 \label{j:finite:count-bound}
\end{align}
The second line discards nonnegative outgoing probabilities at $j<k$.
The outgoing event is strict and the incoming event is closed. In particular,
an atom exactly on the threshold is retained on the CDF side.
Also
\[
 V_j:=\E W_j^2=(n-j)w_0+jw_1
 =\lambda\left(1+\frac{(j-np)(\alpha-p)}{npq}\right).
\]

Suppose support acquisition gives
$U_0\in[-O_0,a_0]$ and $U_1\in[-a_1,O_1]$, with nonnegative
endpoints. For a centered $T\in[-L,R]$ of variance $w$,
\begin{equation}\label{j:finite:conditional-moments}
 w\le LR,\quad w^2/L-Lw\le\E T^3\le Rw-w^2/R,
 \quad \E T^4\le(R-L)\E T^3+LRw.
\end{equation}
The first inequality integrates $(T+L)(R-T)\ge0$. The two cubic
bounds integrate $(T+L)(T-w/L)^2\ge0$ and
$(R-T)(T+w/R)^2\ge0$. Multiplying the first polynomial by $T^2$
gives the fourth-moment bound. If $L=0$ or $R=0$, centering forces
$T=0$; the divisions are replaced by this deterministic case.

The common cubic identities \eqref{j:core:conditional-moments}, in
these coordinates, read $\rho h^3=v(1-2v)+3v(w_0+w_1)-q\mu_0+p\mu_1$,
$\gamma h^3=v(1-2p)+3v(w_1-w_0)+q\mu_0+p\mu_1$, and
$\beta h^2=v(1-2p)+s(2\alpha-1)$. The same pair
$(\mu_0,\mu_1)$ is retained in every probability bound below.

\subsubsection{Outgoing and incoming probability bounds}
A useful outgoing minorant is
\[
 G_a(x)=x(x+a)^2(2a-x),\qquad a>0.
\]
It is nonpositive outside $[0,2a]$, and its maximum is
$K a^4$, where $K=9/4+3\sqrt3/2$. Consequently
\[
 \Pp(W_k>u)\ge \frac{\E G_a(W_k-u)}{K a^4}.
\]
This is a global polynomial inequality, including $W_k=u$.
Its expectation uses
\[
 \E W_k^3=(n-k)\mu_0+k\mu_1,
\]
\[
 \E W_k^4=(n-k)\E U_0^4+k\E U_1^4
 +3\{V_k^2-(n-k)w_0^2-kw_1^2\}.
\]
Thus \eqref{j:finite:conditional-moments} makes the outgoing bound affine in
the two common cubic moments. Negative lower bounds remain valid; they
are not mistaken for negative probabilities.

For an incoming count $j>k$ and phase $u\le1/2$, put $d=j-k-u>0$.
Cantelli gives
\[
 \Pp(W_j\le-d)\le\frac{V_j}{V_j+d^2}.
\]
For completeness, for a centered $Y$ of variance $V$ and $a>0$,
$\Pp(Y\ge d)\le\E(Y+a)^2/(d+a)^2=(V+a^2)/(d+a)^2$.
Choosing $a=V/d$ proves the formula; $V=0$ is deterministic.
Exponential bounds can be sharper. If $\E Y=0$, $Y\le B$, $B>0$,
and $\E Y^2\le w$, then for $t\ge0$,
\begin{equation}\label{j:finite:one-sided-mgf}
 \E e^{tY}\le
 M(w;B,t):=\frac{B^2e^{-tw/B}+we^{tB}}{B^2+w}.
\end{equation}
Indeed, the quadratic interpolating $e^{tx}$ at $-w/B$ with multiplicity
two and at $B$ majorizes it on $x\le B$: the Hermite remainder has
sign $(x+w/B)^2(x-B)\le0$. Its quadratic coefficient is nonnegative.
The centered two-point law on $\{-w/B,B\}$ with variance $w$ evaluates
its expectation exactly. This proves the formula also with an upper
variance bound. At zero variance use $Y=0$.
Differentiation shows that $M$ is increasing and concave in $w$:
with $x=t(B^2+w)/B$,
\[
 M''(w)=\frac{B^2e^{-tw/B}}{(B^2+w)^3}
          (x^2+2x+2-2e^x)\le0.
\]
Also
\[
 M'(w)=\frac{B^2}{(B^2+w)^2}
 \{e^{tB}-e^{-tw/B}(1+x)\}\ge0,
\]
since $e^x\ge1+x$. Its tangents therefore give affine upper bounds throughout the variance
half-line. Applying \eqref{j:finite:one-sided-mgf} to $-U_0,-U_1$ and
multiplying the corresponding factors bounds $\E e^{-tW_j}$.
Chernoff then gives $e^{-td}\E e^{-tW_j}$ as an incoming upper bound.

One can also retain the common cubic moments. On $[A,B]$, $A<0<B$,
the function
\[
 r_t(x)=\frac{e^{tx}-1-tx}{x^2}
       =t^2\int_0^1(1-r)e^{trx}\,dr
\]
is increasing and convex. Its chord gives
$e^{tx}\le1+tx+c_2x^2+c_3x^3$, with $c_2,c_3\ge0$.
Applied to $-U_i$, choose $C_2\ge c_2$ and
$0\le C_{3,-}\le c_3\le C_{3,+}$, and put
\[
 Q_i(\mu_i)=\max\{0,1+C_2w_{i,+}-C_{3,-}\mu_i,
                     1+C_2w_{i,+}-C_{3,+}\mu_i\}.
\]
This is a nonnegative convex upper bound for its MGF. The maximum with zero
also makes it nonnegative at infeasible vertices of a relaxed moment polygon.
For fixed positive pivots $P_0,P_1$, weighted AM--GM yields
\[
 Q_0^{n-j}Q_1^j\le P_0^{n-j}P_1^j
 \left(\frac{n-j}{n}\frac{Q_0}{P_0}
       +\frac{j}{n}\frac{Q_1}{P_1}\right)^n.
\]
The right side is convex in $(\mu_0,\mu_1)$. A fixed exponential rate
and its phase chord produce a convex common-moment bound on an entire
closed phase interval. Choose a strategy on that whole interval and
moment polygon before comparing different strategies; a pointwise
minimum of convex functions need not be convex.

\subsubsection{Normalization and a finite closed cover}
Let $\mathcal M$ be a compact polygon containing the actual common
cubic moments, obtained from proved moment and defect restrictions.
Suppose $F$ is convex on $\mathcal M$ and $L$ is affine and strictly
positive there. Then
\begin{equation}\label{j:finite:normalized-vertex}
 \frac{F(x)}{L(x)}\le
 \max_{v\text{ vertex of }\mathcal M}\frac{F(v)}{L(v)}.
\end{equation}
To see this, write $x=\sum_v\theta_vv$. Convexity bounds the numerator
by $\sum_v\theta_vF(v)$; affinity identifies the denominator with
$\sum_v\theta_vL(v)$. The ratio is therefore bounded by the largest
vertex ratio. Segments and single points are included. If the lower
bound for $\rho$ is $\max\{r_0,L\}$, split $\mathcal M$ along
$L=r_0$ and retain both closed pieces.

\begin{proposition}[Whole-cell count certificate]\label{j:finite:count-cell}
Fix a closed box of acquired parameters, a count $k$, and a closed
phase interval. Let $B_k$ bound the binomial CDF above, $b_k^-$ the own
mass below, and $b_j^+$ the incoming masses above. Let $L_k$ be a
proved outgoing lower bound, $A_j\ge0$ proved incoming upper bounds,
and $N$ a lower bound for $\Phi((k+u-np)/(\sqrt n h))$.
Suppose the combination
\[
 B_k-b_k^-L_k+\sum_{j>k}b_j^+A_j-N
\]
is bounded on the whole phase interval by convex Bernstein control
functions $F_i(\mu_0,\mu_1)$. Let $r>0$ be the affine lower bound for
$\rho$ on each closed normalization piece and let $n_+\ge\sqrt n$.
Then the positive normalized CDF error is at most
\[
 n_+\max_{i,\,v\text{ vertex}}
       \frac{\max\{0,F_i(v)\}}{r(v)}.
\]
\end{proposition}
\begin{proof}
Insert the probability bounds into \eqref{j:finite:count-bound}. For the outgoing
term, $b_k\Pp(W_k>u)\ge b_k^-\Pp(W_k>u)\ge b_k^-L_k$,
even if $L_k<0$. The nonnegativity of the incoming bounds justifies the
upper mass substitutions. The Bernstein basis is nonnegative and sums
to one. Take the positive part and apply \eqref{j:finite:normalized-vertex};
the positive part of a convex function is convex. Finally enlarge
$\sqrt n$ to $n_+$.
\end{proof}

The acquired $p,h,H$ bounds must first put every actual threshold into
one of the original roots $k=0,\ldots,n$, $u\in[-1/2,1/2]$.
If they do not, prove the remaining phase tails separately. A tied
nearest count belongs to both closed roots. Apply
Proposition~\ref{j:cert:cover} only after every leaf is either a
proved necessary-condition contradiction or satisfies the whole-cell
score bound. The finite tree then covers infinitely many parameters.
Its leaves never assert that every feasible moment tuple is an iid law;
they need only contain all actual laws.


\subsubsection{Necessary filters for the same count domain}\label{j:finite:filters}
In the following formulas write $s_i=w_i$, $x=\mu_0$, $y=\mu_1$.
The preceding cubic identities and support-moment inequalities give all
three quantities for the same law; they may not be chosen separately.
These identities restrict the single pair $(x,y)$ to a rectangle and
linear strips. For example, replacing $p$ by the midpoint $p_m$ in
$-qx+py$ incurs error $(p-p_m)(x+y)$. On a closed cubic rectangle
its absolute value is at most the probability half-width times
$\max(|x_l+y_l|,|x_u+y_u|)$. In the signed-cubic identity the
corresponding factor is $y-x$, bounded by
$\max(|y_l-x_u|,|y_u-x_l|)$. All other products use directed interval
endpoints. Thus these strips contain every actual common pair.

\subsubsection{The physical defect as a further half-plane}
Let $\sigma=\sqrt{10}$ and retain the physical gap defect
$D=(3+\sigma)\rho-\gamma-3J$.
The endpoints of the central gap are
$A=(p-a_0)/h$, $B=(q-a_1)/h$.
The outer-gap inequality gives
\[
 D\ge\sigma\rho-\gamma-6mAB
       +6B(\E X_+^2-mB)+6A(\E X_-^2-mA).
\]
Substituting the count coordinates and the two cubic identities yields
\begin{align}
 Dh^3\ge{}&B_\sigma(p)+6vF(a_0,a_1)\notag\\
 &+qs_0\{3p(\sigma+3)-6a_0\}
  +ps_1\{3q(\sigma+1)-6a_1\}\notag\\
 &-(\sigma+1)qx+(\sigma-1)py,\label{j:finite:defect-halfplane}\\
 B_\sigma(p)&=v\{(\sigma-1)+2p-(2\sigma+6)v\},\notag\\
 F(a_0,a_1)&=a_0+a_1-a_0^2-a_1^2-a_0a_1.\notag
\end{align}
The polynomial $B_\sigma$ is the Bernoulli defect square.
Replacing $\sigma$ by a lower enclosure decreases the complete right
side, because its coefficient is exactly $\rho h^3\ge0$.
After that replacement, signed interval products and the explicit
nonnegative-coefficient gates license the termwise lower bounds.
If such a gate fails, this additional half-plane is omitted.

Substitution of \eqref{j:finite:conditional-moments} into
\eqref{j:finite:defect-halfplane} gives the stronger scalar inequality
\begin{align}
 Dh^3\ge{}&B_\sigma+6vF(a_0,a_1)\notag\\
 &+q\left\{[3p(\sigma+3)-(\sigma+7)a_0]s_0
                         +(\sigma+1)s_0^2/a_0\right\}\notag\\
 &+p\left\{[3q(\sigma+1)-(\sigma+5)a_1]s_1
                         +(\sigma-1)s_1^2/a_1\right\}.
 \label{j:finite:defect-sharp}
\end{align}
The terms with $a_i=0$ are treated by $s_i=0$, not by division.
Also $s_i\le O_i a_i$, by averaging the nonnegative quadratic
$(U_i+O_i)(a_i-U_i)$, with reflection for class one.

On a rectangle clipped by $a_0+a_1\le d<1$, the function $F$ is
nonnegative and concave: its Hessian is
$\left(\begin{smallmatrix}-2&-1\\-1&-2\end{smallmatrix}\right)$.
Its minimum therefore occurs at a vertex of the clipped polygon.
There is no requirement $d<1/2$ for this vertex argument.
The older coordinate-monotonicity shortcut requires separate gates.
For example, a lower bound for the $a_0$ derivative in
\eqref{j:finite:defect-sharp} is
\[
 6v_l(1-2d)-q_u(\sigma_u+7)(s_0)_u
                       -q_u(\sigma_u+1)O_0^2.
\]
The reflected expression uses $\sigma+5$ and $\sigma-1$.
Positivity of these gates implies $d<1/2$ and licenses that shortcut
only on its smaller domain. A bisection discards an interval only when
its entire defect lower bound exceeds the already acquired upper bound.

\subsubsection{Second moments and necessary noise floors}
The exact second-moment relations are
\begin{equation}\label{j:finite:beta-noise}
 \beta=\frac{v(1-2p)+s(2\alpha-1)}{v+s},\qquad
 m=\frac{v}{\sqrt{v+s}}.
\end{equation}
The second expression increases with $v$ and decreases with $s$.
For $s>0$, the first can be solved for $\alpha$; the case $s=0$
is tested directly. If an acquired upper bound $\beta_u$ is below
$t_l=1-2p_u$, then
\[
 v_l(t_l-\beta_u)\le s(\beta_u+1-2\alpha_l).
\]
A nonpositive right coefficient excludes the box; a positive one
gives a lower bound for $s$. If an acquired lower bound $\beta_l$
exceeds $t_u=1-2p_l$, use instead
$v_l(\beta_l-t_u)\le s(2\alpha_u-1-\beta_l)$.
These tests retain the same $\alpha$ as the conditional probabilities.

There is a further floor from the absolute third moment. The elementary
bounds $-x\le O_0s_0$, $y\le O_1s_1$ in
\eqref{j:core:conditional-moments} give
\[
 \rho_l(v+s)^{3/2}\le v(1-2v)+Ms,
\]
where $M$ bounds
$(3p+O_0)(1-\alpha)+(3q+O_1)\alpha$ at all four
$p,\alpha$ corners. Convexity of $(v+s)^{3/2}$ implies
\[
 A(v):=\rho_l v^{3/2}-v+2v^2
 \le [M-\tfrac32\rho_l\sqrt v]s.
\]
If $A(v_l)>0$, then
$\rho_l\sqrt{v_l}>1-2v_l$, so
$A'(v_l)>1/2+v_l>0$; moreover $A''>0$.
Hence a positive lower bound for $A(v_l)$ and an upper bound for
the right coefficient give a necessary noise floor on the whole box.
A nonpositive upper coefficient proves infeasibility. Directed square
roots preserve both conclusions.

Finally, the common inward bound $a_0+a_1\le d$, $d>0$,
together with \eqref{j:finite:conditional-moments}, implies
\begin{equation}\label{j:finite:common-inward}
 \rho(v+s)^{3/2}\ge v(1-2v)+B_*s+H_*s^2/d,
\end{equation}
where
\[
 B_*=3p(1-\alpha)+3q\alpha-d\max(\alpha,1-\alpha),\qquad
 H_*=\left(\frac{1-\alpha}{\sqrt q}+\frac\alpha{\sqrt p}\right)^2.
\]
For the quadratic term use
\[
 \frac{(1-\alpha)^2}{qa_0}+\frac{\alpha^2}{pa_1}
 \ge\frac{((1-\alpha)/\sqrt q+\alpha/\sqrt p)^2}{a_0+a_1}.
\]
The zero-variance restrictions cover vanishing endpoints by a limit.
The minimum of $B_*$ is at the four $p,\alpha$ corners, because it is
affine in $p$ and concave in $\alpha$.
On a $p$ interval $[p_l,p_u]$ a positive lower envelope for
$\sqrt{H_*}$ is $(1-\alpha)/\sqrt{q_u}+\alpha/\sqrt{p_u}$,
where $q_u=1-p_l$. This affine function is minimized at an
$\alpha$ endpoint. The global upper Taylor bound
\[
 (v+s)^{3/2}\le v^{3/2}+\tfrac32\sqrt v\,s
                                      +\tfrac38s^2/\sqrt v
\]
turns \eqref{j:finite:common-inward}, with $\rho\le\rho_u$, into a
necessary quadratic inequality in $s$. Its exact minimum on the
closed variance interval is tested at both endpoints and any interior
vertex. Strict positivity excludes the whole box. The implementation
defers $d=0$ to its deterministic branch.

\subsubsection{An additional outgoing event and full phase coverage}
For centered $W\in[-a_1,O_1]$ and $-a_1\le b<O_1$,
\[
 (W+a_1)(W-b)\le\ind_{\{W>b\}}(O_1+a_1)(O_1-b).
\]
Thus $\Pp(W>b)$ is at least the nonnegative part of
$(s_1-a_1b)/((O_1+a_1)(O_1-b))$.
For $b<-a_1$ the strict event is certain; for $b\ge O_1$ use zero.
If any of the $k$ class-one variables exceeds
$b=u+B+(k-1)a_1$, their sum exceeds $u+B$.
Independence gives the lower probability $1-(1-r)^k$, where $r$
is the preceding bound. The remaining centered sum, of variance $V_0$
and upper support $(n-k)a_0$, is at least $-B$ with probability at
least
\[
 \max\left\{\frac{B^2}{V_0+B^2},
                         \frac{B}{B+(n-k)a_0}\right\}\quad(B>0).
\]
These follow from Cantelli and Markov, respectively. At $B=0$ use
one if $V_0=0$ and zero otherwise. Multiplication of the two lower
bounds is valid because the two conditional groups are independent.
This supplies the optional single-event outgoing witness.

The acquired lower bound $m_l>0$ also gives
\[
 \lambda=n(v^2/m^2-v)\le n(v^2/m_l^2-v),\qquad h\le v_u/m_l.
\]
The quadratic on the right is convex, so its maximum on the $v$
interval occurs at an endpoint. If $-H_-\le H\le H_+$, the gates
\[
 np_l-H_-h_u>-1/2,\qquad np_u+H_+h_u<n+1/2
\]
ensure that all nearest counts belong to $\{0,\ldots,n\}$.
The original joint budget $d/d_0+\lambda/\Lambda\le1$ retains
its own intercept $\Lambda$ when a separate sharper noise cap clips
the count domain. Replacing that intercept by the sharper cap would
assert an additional joint inequality and is not licensed here.

\subsection{The signed fifth-index corner}
\label{j:finite:gapfree5}
The last small moment rectangle can be treated directly by conditioning
on signs. This avoids contact and Fourier estimates in that rectangle.
Let $\beta=\E[X|X|]$, $p=\Pp(X>0)$, $q=1-p$,
$m=\E X_+=\E X_-$ and $a_\pm=(1\pm\beta)/2$.
Zero atoms belong to the nonpositive class. Signwise Cauchy gives
\[
 m\ge\frac{1+\beta^2}{2\rho},\qquad
 \frac{m^2}{a_+}\le p\le1-\frac{m^2}{a_-},\qquad m^2\le pq.
\]
Set $h=pq/m$, $Y=p+hX$, $B=\ind_{\{X>0\}}$, and $U=Y-B$.
Both conditional means of $U$ vanish, with variances
\[
 v_0=p^2(qa_-/m^2-1),\qquad v_1=q^2(pa_+/m^2-1).
\]
Given $\sum_{i=1}^5B_i=j$, the noises are independent with variance
$V_j=(5-j)v_0+jv_1$. Write $b_j=\binom5j p^jq^{5-j}$ and define
\[
 T(V,d)=\begin{cases}1,&d\le0,\\V/(V+d^2),&d>0.\end{cases}
\]
Conditional Cantelli and the exact normal coordinate yield
\begin{equation}\label{j:finite:gapfree-bound}
 \frac{\sqrt5}{\rho}\{\Pp(S_5\le t)-\Phi(t/\sqrt5)\}
 \le\frac{\sqrt5}{\rho}
 \left\{\sum_{j=0}^5 b_jT(V_j,j-z)
          -\Phi\left(\frac{(z-5p)m}{\sqrt5pq}\right)\right\},
 \qquad z=5p+ht.
\end{equation}
The value $T(0,0)=1$ is essential: a deterministic sum at the threshold
belongs to the closed event.

\begin{proposition}[The positive corner]
\label{j:finite:gapfree5-result}
For every standardized iid marginal satisfying
\[
 \frac{477}{400}\le\rho\le\frac65,\qquad
 \frac{17}{50}\le\beta\le\frac{35139006743001}{10^{14}},
\]
the positive closed-CDF score in \eqref{j:finite:gapfree-bound} is strictly
below $4097/10000<\CE$ at every real threshold.
\end{proposition}
\begin{proof}
Here are the enclosures defining the finite proof. On a closed
$\beta$ slab $[b_l,b_u]$ put $m_l=(1+b_l^2)/(2(6/5))$.
The moment inequalities enclose $p$ in
$[2m_l^2/(1+b_u),\,1-2m_l^2/(1-b_l)]$.
On a closed $p,z$ box $[l,u]\times[a,b]\subset(0,1/2)\times[0,5]$
use
\begin{align*}
 v_0&\le\max\{0,u^2((1-l)(1-b_l)/(2m_l^2)-1)\},\\
 v_1&\le\max\{0,(1-l)^2(u(1+b_u)/(2m_l^2)-1)\}.
\end{align*}
Each binomial mass is maximized at the projection of $j/5$ onto
$[l,u]$, and each Cantelli factor is maximized at $z=b$ and the
variance upper bound. Clip the sum of the six upper bounds at one.
If $a-5u\ge0$, the normal argument is at least
$(a-5u)m_l/(\sqrt5\,u(1-u))$.
Otherwise it is at least $(a-5u)/\sqrt{5u(1-u)}$: use
$m\le\sqrt{pq}$ on the negative part and the fact that
$(z-5p)/\sqrt{pq}$ decreases in $p$ for $z\ge0$, $p<1/2$.
Its derivative has numerator $-[5p+z(1-2p)]/2$.
Outward square roots and a rational Gaussian lower enclosure therefore
bound the complete box. Take the nonnegative part of the CDF difference
before multiplying by $\sqrt5/\rho_l$.

The twelve closed $\beta$ slabs have endpoints
\[
 17/50,\quad 341/1000,\ldots,351/1000,\quad
 35139006743001/10^{14}.
\]
Together with the five intervals $z\in[k,k+1]$, $0\le k<5$, these
give sixty roots. The supplied binary subdivisions have 1,027 leaves.
Exact evaluation on every closed leaf gives a maximum upper bound
\[
 204848338979/500000000000<4097/10000.
\]
Both omitted rays have direct bounds. Use the full rectangle's $m_l$
and enclosing $[p_l,p_u]$, put $v_u=p_u(1-p_u)$ and
$s_u=v_u^2/m_l^2-v_u$. On the feasible domain $m^2\le v$,
$v^2/m^2-v$ increases with $v$ and decreases with $m$. Thus
$V_j\le\max\{j/p_l,(5-j)/(1-p_u)\}s_u$.
For $z\le0$, discard the nonnegative Gaussian CDF and evaluate the
Cantelli upper sum at zero, using binomial upper masses
$\binom5j p_u^j(1-p_l)^{5-j}$. The resulting score is below $331/1000$.
For $z\ge5$, the CDF is at most one and normal Cantelli gives the
score upper bound
\[
 \frac{\sqrt5}{\rho_l}\frac1{1+5m_l^2/p_u^2}<187/1000.
\]
These exact rational comparisons and the sixty roots cover all real $z$.
The proposition concerns the stated positive branch; it is not silently
extended to the opposite $\beta$ rectangle by reflection.
\end{proof}

\subsection{The numerical regions and the induction}
\label{j:finite:assembly}
For clarity we now give the application domains, separating
numerical assertions from their use. The finite acquisition theorem
\ref{j:local:acquisition} preserves one actual marginal through every
moment subdivision and update. Its state consists of simultaneous moment,
support, score, raw-remainder, defect and signed-threshold bounds. Support
is proved before higher moments are used; an initial raw remainder is
proved before defect feedback. A failed analytic gate supplies no exclusion.
Every nonexcluded state must reach a count domain or a complete Fourier
bound. These rules are the common induction invariant in every row below.

A probability panel means $P_j=[j/100,(j+1)/100]$, $0\le j<100$.
All intervals in this section are closed. Exports include every original
moment cell touching a panel, so a probability seam belongs to both panels.
An exported upper bound is the maximum over these cells; a lower bound is
the minimum. The source law is thereby retained in the export.

\paragraph{The fifth index.}
Put
\[
 b_2=84729828287999/(250000000000000),\qquad
 b_3=35139006743001/10^{14}.
\]
The signed-moment inequality gives $|\beta|\le b_2$ for
$\rho\le237/200$ and $|\beta|\le b_3$ for $\rho\le6/5$.
The necessary rational comparisons are
$(1+b_i^2)^2\ge\rho_i^2(1-b_i^2)$ at the respective upper endpoints;
the left minus right side increases with $b_i\ge0$.
There are twelve count rectangles: the two full signed bands
$[1,11/10]$ and $[11/10,23/20]$; four rectangles obtained by splitting
$[23/20,237/200]$ at $467/400$ and $[-b_2,b_2]$ at $1/5$;
and six on $[237/200,6/5]$, split at $477/400$.
In the lower half the signed intervals are
$[-b_3,1/5],[1/5,3/10],[3/10,b_3]$; in the upper half replace the
last endpoint by $17/50$. The missing upper-right rectangle is precisely
the support-free positive corner proved above, not its reflection.
The own-score and compensated HIGH inputs are
\[
\begin{array}{c|cc}
\rho\text{ and signed restriction}&\text{own score}&\text{HIGH}\\\hline
{}[1,11/10]&429/1000&337/1000\\
{}[11/10,23/20]&429/1000&369/1000\\
{}[23/20,237/200]&52/125&393/1000\\
{}[237/200,6/5],\ \beta\le1/5&207/500&193/500\\
{}[237/200,6/5],\ \beta\ge1/5&103/250&193/500
\end{array}
\]
The last two rows apply only to their count rectangles. The first score
cap is obtained without support; the other Fourier inputs use previously
proved radii from Proposition~\ref{j:finite:supportrows}. Each signed
support radius in a narrower native profile is checked on that full profile.
The twelve count rectangles produce 1,200 panel instances: 905 are
excluded by necessary resource inequalities and 295 have complete
six-root forests. Thus every panel and every count $0,\ldots,5$ is covered.
The six complementary Fourier bands and their radii are the six upper
moment intervals for $n=5$ in Proposition~\ref{j:finite:supportrows}.
They are divided into 41 closed strips and have score below $T$.
The affine bound closes $\rho\ge7361/4000$, since there its value is
at most $819/2000<T$.

\paragraph{The indices six through fourteen.}
The low-band input table is
\[
\begin{array}{c|c|c|c|c}
n&\rho&R&\text{own score}&\text{HIGH}\\\hline
6\!:\!9&[1,11/10]&23/10&423/1000&69/200\\
6\!:\!9&[11/10,23/20]&12/5&209/500&3/8\\
6\!:\!9&[23/20,237/200]&5/2&83/200&83/200\\
10\!:\!14&[1,11/10]&23/10&1047/2500&69/200\\
10\!:\!14&[11/10,23/20]&12/5&519/1250&93/250
\end{array}
\]
The third row may use the wider Fourier domain
$[23/20,6/5]$ at radius $13/5$, but its resource radius remains $5/2$.
The predecessor allowances retained in the data are $579/5000$ at six,
and $99/1000,63/1000$ in the two bands at ten; zero loss implies
these weaker nonnegative allowances. Every other predecessor allowance
here is zero. The extra outer initializer at six has the same proved
inputs and is joined separately to the resource.
For six through nine, 1,200 panels consist of 910 resource exclusions,
six exclusions by Lemma~\ref{j:local:sign-mass}, and 284 count forests.
For ten through fourteen, 1,000 panels consist of 798 exclusions and
202 surviving parent panels; five parents split into four closed
subintervals, giving 217 forests.
For six through nine, Fourier closes $[237/200,6/5]$ at radius $13/5$
and the five complementary bands in Proposition~\ref{j:finite:supportrows}.
For ten through fourteen it closes $[23/20,71/50]$ at radius three,
then $[71/50,15/8]$ without support. Every Fourier target is $T$;
the same affine tail beyond $7361/4000$ overlaps these covers.

\paragraph{The indices fifteen through thirty-four.}
The low-band input table is
\[
\begin{array}{c|c|c|c|c}
n&\rho&R&\text{own score}&\text{HIGH}\\\hline
15\!:\!19&[1,11/10]&23/10&417/1000&191/500\\
15\!:\!18&[11/10,23/20]&12/5&827/2000&409/1000\\
19&[11/10,23/20]&23/10&827/2000&409/1000\\
20&[1,11/10]&23/10&831/2000&\text{cellwise}\\
20&[11/10,9/8]&23/10&103/250&\text{cellwise}\\
21\!:\!24&[1,11/10]&23/10&419/1000&\text{cellwise}\\
21\!:\!24&[11/10,9/8]&12/5&419/1000&\text{cellwise}\\
25\!:\!30&[1,11/10],[11/10,6/5]&\text{actual-index rows}&434477/10^6&\text{cellwise}\\
31\!:\!34&[1,11/10],[11/10,6/5]&47/20,13/5&432313/10^6&\text{cellwise}
\end{array}
\]
``Cellwise'' means evaluation of the whole compensated Fourier expression
on each moment cell, not an assumed scalar HIGH ceiling. At fifteen
through nineteen a radius-three bound on $\rho\le33/25$ precedes
the own-score and HIGH estimates. At twenty, radius three on
$\rho\le6/5$ precedes its two own-score estimates. The twenty-first
through twenty-fourth own-score caps are unrestricted Fourier bounds.
The contact allowances $219/2000,57/1000$ at fifteen and
$29/250,23/500$ at twenty are weaker consequences of zero loss.
At twenty through twenty-four the unseeded initializers take eight
refinements and a ninth complete export; no earlier numerical seed is
assumed. The four old coarse contracts at twenty are accepted only after
simultaneous containment in the newly acquired domain is proved.

The Fourier complements are
\[
\begin{array}{c|c|c}
n&\rho&\text{extra premise}\\\hline
15\!:\!19&[23/20,33/25]&|X|\le3\\
15\!:\!19&[33/25,19/10]&\text{none}\\
20\!:\!24&[9/8,6/5]&|X|\le3\\
20\!:\!24&[6/5,19/10]&\text{none}\\
25\!:\!39&[6/5,19/10]&\text{none}
\end{array}
\]
Their score target is $T$, and the affine tail covers $\rho\ge19/10$.
Every unrestricted row includes the full auxiliary-moment domain.
The last row removes any need for a third count band at twenty-five.
For $26\le n\le30$, the second-band survivors $P_{32},\ldots,P_{42}$
are instead closed by the 115 finite-sine Fourier inequalities proved above.

Floors $N=25$ and $N=31$ are used respectively at $26$--$30$ and
$32$--$34$, with the weighted invariant $nD\le Nd$ of
the transfer argument in Section~\ref{j:local:transfer}. The actual index remains in every binomial
probability and normal argument, and the support gates are evaluated at
that actual index. The two floor-31 refinements use $P_{35}$ and $P_{36}$
in the second band; all other panels retain the global acquisition.

\paragraph{The extra prerequisite at thirty-nine.}
Before using the sharp-39 acquisition, run the separate local comparison
with $n=39$, $K=741/4000$, predecessor allowance $q=33/100$ and
excess allowance $1/2$. Sharpness at 37 and 38 and the recurrence give
\[
39(C_{39}-\CE)<55887/190000<3/10,\qquad C_{39}<.4226,
\]
and both deleted scores are below $.4253$. Thus its assumptions precede
its conclusion. Under $C_{39}\ge\CE+K/39$, the support input is
$(N,L,K,q,R)=(39,39,741/4000,.33,3.84)$.
The positive charge $K\rho_l$ is retained through the acquisition.
Its full moment cover excludes the complement of 78 probability rows on
$[.302,.458]$. There are 624 signed profile cells and 90 whole-count rows.
The phase cover uses nearest count throughout $(0,.05]$ and $[.05,.08]$;
on $[.08,.10]$ the nearest/whole split is $|u|=.19$, on $[.10,.12]$
it is $|u|=.27$, and above $.12$ the whole-count bound applies everywhere.
Each band is intersected with its acquired noise ceiling; all seams are
closed. The allowance tests in Section~\ref{j:local:allowance} therefore
exclude $C_{39}\ge\CE+K/39$. Since $\CE<.40974$ and
$K/39=.00475$, this supplies $C_{39}<.41449$ before sharpness at 39
is invoked. Zero noise is Bernoulli and already satisfies the weaker target.

\paragraph{The actual-index acquisitions at thirty-five through thirty-nine.}
Their initial pairs (radius, own-score cap) are
\[
\begin{array}{c|ccccc}
n&35&36&37&38&39\\\hline
R&39/10&19/5&19/5&19/5&19/5\\
c&53539/125000&1699/4000&84867/200000&211449/500000&41449/100000.
\end{array}
\]
The first four caps follow from $C_n<41/100+4177/(13824n)$ in
Lemma~\ref{j:smoothing:entry25}; the fifth was just established.
The independent support test at 35 gives $96/25<39/10$.
All five INITIAL/resource computations use the actual integer.
The saved $10,16,12,15,12$ refinement schedules end in complete
100-panel exports. At 35, the four local intervals
$[.29,.34],[.34,.39],[.39,.45],[.45,.57]$, together with the lower
probability refinements, leave precisely $P_{23},\ldots,P_{48}$.
At 36 the localized calculation on $[.30,.42]$ replaces the global
source only on $P_{30},\ldots,P_{41}$, with one-sided radii $3,37/10$.
At 37--39 the global sources suffice. Every remaining probability
panel is explicitly excluded, not omitted.

The selected count domains have the following exact census:
\begin{center}\small
\begin{tabular}{c|c|c|r}
n&\text{first low band}&\text{second low band}&\text{forests}\\\hline
15&37--56&33--44&32\\16&37--56&33--43&31\\17&37--55&33--43&30\\
18&37--54&34--43&28\\19&37--53&34--42&26\\20&37--58&35--45&49\\
21&37--57&35--44&31\\22&37--56&35--43&29\\23&37--55&35--43&28\\
24&37--54&35--42&26\\25&37--53&32--42&52\\
26&37--53&\text{Fourier}&17\\27,28,29&37--52&\text{Fourier}&16\text{ each}\\
30&37--51&\text{Fourier}&15\\31--34&36--49&30--40&25\text{ each}\\
35&\multicolumn{2}{c|}{23--48}&26\\36&\multicolumn{2}{c|}{24--48}&25\\
37&\multicolumn{2}{c|}{24--47}&24\\38&\multicolumn{2}{c|}{25--47}&23\\
39&\multicolumn{2}{c|}{26--47}&22
\end{tabular}
\end{center}
Multiple forests within a parent form a checked closed subdivision.
The shortening at $26$--$30$ has a direct explanation. On panels 52,53
the acquired outward pairs $(d,\Theta)$ in $D\le25d/n$, $\theta\le\Theta$
are $(3/10,1/500)$ and $(5/16,1/2500)$. With
$s=316227766/10^8<\sqrt{10}$ and $l=j/100$, the joint budget gives
\[
 D\ge(12+4s)(1-\Theta)^2(l-(4-s)/2)^2
 >\begin{cases}251/1000,&j=52,\\38/125,&j=53.\end{cases}
\]
Here $p\ge l>p_*$, $\sqrt{p(1-p)}\le1/2$, and
$(1-\theta)^{3/2}\ge(1-\Theta)^2$ were used.
For panel 53 at $n\ge27$, $D\le125/432<38/125$; for panel52
at $n=30$, $D\le1/4<251/1000$. These are the five excluded pairs.
At twenty, the additional panels $P_{36}$ and $P_{34}$ in the respective
bands are excluded by the sign-mass inequality, including their endpoints.

\begin{proposition}[Complete numerical finite base through thirty-nine]
\label{j:finite:numerical539}
All support, Fourier, initializer, resource and count inequalities in
these tables hold on their complete domains, with the preceding
positive-39 comparison included as a separate prerequisite.
Every exported row implies its entire count contract, and the exclusions
and retained forests cover every original moment and probability panel.
\end{proposition}
The assertion includes four logically separate checks: a complete
acquisition for the same marginal, directed containment of every scalar
and signed-phase condition, native count bounds at the actual integer,
and the complete closed union. In particular, if $q_i(H)$ are the
acquired phase polynomials and $\widetilde q_j(H)$ the contract
polynomials, for every $i$ the join proves some $j$ with
$q_i(H)\le\widetilde q_j(H)$ throughout the signed interval. Endpoints
and any interior minimum of their quadratic difference suffice. Hence
$\max_iq_i\le\max_j\widetilde q_j$; replacing a maximum by a minimum
would not have this meaning. Entire original contracts, including their
actual indices and moment domains, are retained when selecting one row.

\begin{proof}[Derivation of $C_n\le\CE$ for $5\le n\le39$]
If there is a failure, choose its least index and an attained positive
closed maximizer by Theorem~\ref{j:core:attainment}. Its immediate
predecessor is sharp, giving zero contact loss; all retained larger
allowances are valid. The displayed moment intervals cover every
$\rho\ge1$. In a Fourier or signed-corner region the score is below
$T<\CE$. In a count region the support and score premises precede
acquisition, which retains the law in an exported panel. A necessary
exclusion is impossible for this law; otherwise the directed join puts
it in a count root. The common-count identity, its whole-cell upper
bound and the complete forest then give score below $\CE$.
At 39 the separate positive comparison is applied first, as above.
Each possibility contradicts the assumed score, proving the induction.
\end{proof}

\subsection{The actual indices forty, forty-one and forty-two}
\label{j:finite:4042}
A first failure here has sharp predecessors through its two deletions.
The recurrence yields
\[
0<n(C_n-\CE)<1433/4875<3/10,\qquad C_n<167/400<.4226.
\]
Both deleted bounds are below $.4253$ and contact loss is zero, so the
weaker numerical allowances $.54/n$ and $.19$ are available.
The floor-40 argument in Section~\ref{j:local:transfer} retains
$\mathcal R_n\le r$, $nD\le40d$ for this actual law; its support
proof is checked on both index endpoints 40 and 42, not inferred from
support at 40. It acquires radius $96/25$ before its higher-moment bounds.
The initial closed moment cover and fourteen refinement passes produce
all 100 probability panels. The left refinement leaves 49 rows covering
$[.282,.38]$; the central refinement leaves 48 covering $[.38,.476]$.
Their complements are excluded by the same full acquisition.
Confinement uses initial strip widths $.55,.575,.6$ on the left and
$.55$ centrally, with cubic coefficients 30 and 24; all corresponding
width and positive-slope gates belong to that acquisition.

The central noise/phase cover is: nearest count throughout $(0,.08]$;
on $[.08,.10]$ nearest count for $|u|\ge.19$ and whole count for
$|u|\le.19$; the split is $.27$ on $[.10,.12]$; above $.12$ use
whole count. The left cover uses nearest count through $.0925$;
the splits are $.16$ on $[.0925,.10]$, $.25$ on $[.10,.12]$,
and $u=-.27$ on $[.12,.135]$, with nearest count below that phase
and whole count above; after $.135$ use whole count.
All bands are clipped to their own noise ceilings and include their seams;
zero noise is handled by the Bernoulli theorem. The 70 central and
82 left signed whole-count rows, together with the nearest-count rows,
cover all remaining positive noise.

At these three finite indices one can spend a positive binomial slack.
For $[a,b]\subset(0,1/2)$ and actual $n\in\{40,41,42\}$, let
$F_j(p)=\Pr(\operatorname{Bin}(n,p)\le j)$ and
$z_j(p)=(j-np)/\sqrt{np(1-p)}$. Differentiation shows that both
functions decrease, whereas $\sqrt{np(1-p)}/(1-2p(1-p))$ increases.
Thus both sides of every jump are bounded by
\begin{equation}\label{j:finite:binomialslack}
\frac{\sqrt{nb(1-b)}}{1-2b(1-b)}
\max\{F_j(a)-\Phi(z_j(b)),\Phi(z_j(a))-F_{j-1}(b),0\},
\qquad F_{-1}=0.
\end{equation}
Between jumps the CDF is constant and the normal CDF increases, so
$j=0,\ldots,n$ suffices. Exact endpoint probabilities and outward
Gaussian bounds check every one of the 100 panels on $[.28,.48]$,
each split into 128 closed subcells, and give a positive lower
$\epsilon_p$ for $\CE-C^{\rm bin}_{n,p}$ uniformly on its parent.

The signed finite comparisons use $I/n\le\max(I/40,I/42)$ and,
for $b_0\le\sqrt n\pi_k\le1$, $\sqrt n\pi_k J\le\max(b_0J,J)$.
These retain both signs. If $S$ bounds $\rho_B(C-C_B)/\lambda$
on $[a,b]$, $a>0$, it suffices that
$S<\rho_{B,l}\epsilon_p/b$: for $S<0$ the score is below $C_B$;
otherwise multiply by $\lambda\le b$ and divide by
$\rho_B\ge\rho_{B,l}$. The actual scalar formulas are convex in
positive noise (affine terms and nonnegative multiples of reciprocals
of positive affine terms), so their endpoint tests cover the band.
These finite-slack tests do not assert a theorem for every $n\ge40$.

\begin{proposition}[Numerical floor-40 assertion]\label{j:finite:numerical4042}
The preceding acquisition, two-endpoint support, actual-index binomial
slack and complete signed local covers satisfy their strict inequalities
at each $n=40,41,42$. Every local row consumes its acquired simultaneous
profile. Consequently $C_{40},C_{41},C_{42}\le\CE$.
\end{proposition}
\begin{proof}[Derivation from the numerical assertion]
The just-proved base through 39 and the least-failure argument give
all entry premises. Acquisition and confinement preserve the witness.
The local comparison theorem applies on its row with the signed finite
substitutions above, and excludes every positive-noise case; zero noise
is Bernoulli. This contradicts any first failure among the three indices.
\end{proof}

\begin{theorem}[The finite base]\label{j:finite:base}
For every integer $1\le n\le42$ and every standardized iid marginal
with finite third absolute moment, $\sqrt n\,\Delta_n(P)\le\CE\rho(P)$.
\end{theorem}
\begin{proof}
The quantitative starting bounds and the physical four-summand comparison
establish indices one through four. Proposition~\ref{j:finite:numerical539}
and its same-law induction establish five through thirty-nine.
Proposition~\ref{j:finite:numerical4042} supplies the last three indices.
\end{proof}

\providecommand{\phio}{\phi(0)}
\providecommand{\cR}{\mathcal R}
\providecommand{\Prb}{\mathbb P}
\providecommand{\nm}{\mathrm}
\section{Acquisition and comparison near a two-point law}
\label{j:local:section}
The defect identity has two uses. Contact and a complete Fourier bound
first confine one extremizing marginal to a neighborhood of a two-point
law. We then compare the actual sum with that law's own binomial count.
The second step needs no extremality beyond the simultaneous hypotheses
already obtained in the first. All bounds below include zero excess
$e=C_n-\CE=0$; positive conditional noise and zero noise are treated
separately at the end.

Write $\sigma=\sqrt{10}$, $m=\E X_+=\E X_-$,
$\rho=\E|X|^3$, $\gamma=\E X^3$, $\beta=\E[X|X|]$,
$\mu_j=\E X^j$ and $M_j=\E|X|^j$ (also $\eta_5=M_5$).
In the central gap we use $(-A,B)$ for the endpoints, equivalently
$A=A_-$ and $B=A_+$ in Section~\ref{j:core:section}.
The quantities $J,J_1,J_2,D,L,K_1,K_q$ and $\mathcal R$ are those of
Theorem~\ref{j:core:conservation}; $U=J_1+(\mu_4-3)/6$ below is a
coefficient, distinct from the conditional noise introduced later.
The letter $B$ denotes the positive endpoint until Bernoulli coordinates
are introduced; there its meaning is explicitly reset to a label.


\subsection{Support from the contact inequality}\label{j:local:support:theorem}
\paragraph{The scale correction.}
Write $\phi_0=\phi(0)$.
For $n>1$ put $s_n=\sqrt{n/(n-1)}$ and
$e_n=n(s_n-1)-1/2$.
\begin{lemma}\label{j:local:support:scale}
For every integer $n\ge36$, one has $0<e_n<.384/n$.
\end{lemma}
\begin{proof}
Put $x=1/n$ and $c=.384=48/125$. Direct expansion gives
\[
 (1-x)(1+x/2+cx^2)^2-1
 =x^2\{(2c-3/4)-(c+1/4)x+(c^2-c)x^2-c^2x^3\}.
\]
The bracket is decreasing for $x\ge0$, since $0<c<1$.
At $x=1/36$ its exact rational value is $823/4050000>0$.
Squaring positive quantities gives
$s_n<1+x/2+cx^2$ for $0<x\le1/36$. Similarly
$(1-x)(1+x/2)^2-1=-3x^2/4-x^3/4<0$, which gives the lower bound.
\end{proof}


\paragraph{A signed Gaussian envelope.}
Define, for every real $H,t$ and $n>1$,
\[
 \mathcal H_n(H,t)=n^{3/2}\{\Phi((H-t)/\sqrt{n-1})-
             \Phi(H/\sqrt n)\}
       +nt\phi(H/\sqrt n)+\tfrac H2(t^2-1)\phi(H/\sqrt n).
\]
Write $v=n-1$, $s=\sqrt{n/v}$, $u=|t|$, and
\[
 A=\frac{u^3-3u}{6},\quad B=\frac{u^4-6u^2+3}{24\sqrt v},
 \quad D=\frac{u^5-15u}{120v}.
\]
\begin{lemma}[Fifth-order identity with a uniform finite remainder]
If $u>\sqrt3$, then, with $z=H/\sqrt v$,
\[
 \frac{\mathcal H_n(H,\pm u)}{s^3}
 =\pm A(1-z^2)\phi(z)+B\phi^{(3)}(z)
                  \mp D\phi^{(4)}(z)+R,
\]
where the sign of the remainder need not be fixed and
\[
 |R|\le E^{(5)}_v(u):=
 \frac{51u^6}{14400v^{3/2}}
 +\frac{51(3u^2-2)}{960v^{3/2}}+\frac{u}{8v^2}.
\]
This statement is uniform over every real $H$.
\end{lemma}
\begin{proof}
Let $F_w(x)=\Phi(x/\sqrt w)$, with primes denoting $x$ derivatives.
The heat identity $\partial_w F_w=F_w''/2$ and ordinary Taylor's theorem
give, with unlabelled derivatives at $(v,H)$,
\begin{align*}
 F_{v+1}&=F_v+F_v''/2+F_v^{(4)}/8+R_0,
 &|R_0|&\le\sup_{w\in[v,v+1]}\|F_w^{(6)}\|_\infty/48,\\
 F_{v+1}'&=F_v'+F_v'''/2+F_v^{(5)}/8+R_1,
 &|R_1|&\le\sup_{w\in[v,v+1]}\|F_w^{(7)}\|_\infty/48,\\
 F_{v+1}''&=F_v''+F_v^{(4)}/2+R_2,
 &|R_2|&\le\sup_{w\in[v,v+1]}\|F_w^{(6)}\|_\infty/8.
\end{align*}
Expand $F_v(H-t)$ through degree five, with remainder at most
$u^6\|F_v^{(6)}\|_\infty/720$. In
\[
 \frac{\mathcal H_n}{n^{3/2}}=
 F_v(H-t)-F_{v+1}(H)+tF_{v+1}'(H)
             -\frac{t^2-1}{2}F_{v+1}''(H),
\]
the first three derivative coefficients cancel. The next three are
\[
 \left(\frac t2-\frac{t^3}6\right)F_v'''
 +\frac{t^4-6t^2+3}{24}F_v^{(4)}
 +\left(\frac t8-\frac{t^5}{120}\right)F_v^{(5)}.
\]
The remaining absolute value is bounded by
\[
 \left\{\frac{u^6}{720}+\frac1{48}+\frac{u^2-1}{16}\right\}
       \sup_{w\in[v,v+1]}\|F_w^{(6)}\|_\infty
 +\frac u{48}\sup_{w\in[v,v+1]}\|F_w^{(7)}\|_\infty.
\]
Fourier inversion of the Gaussian, differentiated under its absolutely
integrable Fourier transform, gives
\[
 \|\phi^{(k)}\|_\infty
 \le\frac1\pi\int_0^\infty y^k e^{-y^2/2}\,dy.
\]
Integration by parts yields
$\|\phi^{(5)}\|_\infty\le8/\pi<2.55$ and
$\|\phi^{(6)}\|_\infty\le15\phi_0<6$.
Since $F_w^{(j)}=w^{-j/2}\phi^{(j-1)}$, the displayed remainder
becomes $s^3E^{(5)}_v(u)$ after multiplication by $n^{3/2}$.
The three retained derivative coefficients give the identity.
\end{proof}

\begin{lemma}[A lower tangent bound for the fourth Gaussian derivative]
For every real $z$,
\[
 \phi^{(4)}(z)\ge\phi_0(3-15z^2/2).
\]
\end{lemma}
\begin{proof}
The same Gaussian Fourier integral gives
$\phi^{(4)}(z)=\pi^{-1}\int_0^\infty
y^4e^{-y^2/2}\cos(yz)\,dy$.
Use the globally valid inequality $\cos w\ge1-w^2/2$ inside
this integral. Its fourth and sixth Gaussian moments are $3\phi_0$
and $15\phi_0$, respectively.
\end{proof}

\begin{lemma}[Uniform signed fifth-order envelope]\label{j:local:support:envelope}
Suppose $u>\sqrt3$ and $A,B,D\ge0$. If
$W=A-15D/2>0$, put
\[
 T=A-3D+\frac{9B^2}{4W}.
\]
If $T\ge A/2+2B+3D$, then, for every real $H$ and either sign,
\[
 \mathcal H_n(H,\pm u)\le s^3\{\phi_0 T+E^{(5)}_v(u)\}.
\]
\end{lemma}
\begin{proof}
For positive $t=u$ and $|z|\le1$, divide the retained terms by
$\phi_0$ and use the preceding tangent bound to obtain
\[
 A-3D-(A-15D/2)z^2+3B|z|
 \le A-3D+9B^2/(4W).
\]
Here $(1-z^2)\phi(z)\le(1-z^2)\phi_0$ and
$|\phi^{(3)}(z)|\le3|z|\phi_0$ on this interval.
For $|z|\ge1$ the first term is nonpositive, and the Fourier bounds
$\|\phi^{(3)}\|_\infty<2\phi_0$ and
$\|\phi^{(4)}\|_\infty\le3\phi_0$ give $2B+3D$.
For negative $t=-u$, the first term has global upper bound
$2e^{-3/2}A\phi_0<A\phi_0/2$, and the other terms give
$(2B+3D)\phi_0$. The hypothesis on $T$ covers both exterior cases.
The inequality $2e^{-3/2}<1/2$ follows from four positive terms of
the exponential series at $3/2$. Finally add the uniform remainder.
\end{proof}

\begin{lemma}[A global Gaussian remainder bound]
\label{j:local:ext:global-gaussian-ray}
For every integer $n\ge2$ and all real $H,t$,
\[
 |\mathcal H_n(H,t)|
 <\frac{|t|^3}{15}+\frac{|t|}{5}+\frac14
 <\frac{|t|^3}{15}+\frac{|t|}{5}+\frac12.
\]
\end{lemma}
\begin{proof}
Write $F_w(x)=\Phi(x/\sqrt w)$, $v=n-1$ and $y=H-t$.
The identity defining $\mathcal H_n$ is
\[
 n^{-3/2}\mathcal H_n(H,t)
 =F_v(y)-F_n(H)+tF_n'(H)-\frac{t^2-1}{2}F_n''(H).
\]
Taylor expansion of $F_n(y)$ at $H$ through degree two has remainder
at most $\phi(0)|t|^3/(6n^{3/2})$, since
$\|\phi''\|_\infty=\phi(0)$.
After this expansion the remaining correction is
\[
 \frac12\{F_n''(H)-F_n''(y)\}
 +\frac12\int_v^n\{F_n''(y)-F_w''(y)\}\,dw.
\]
Here we used $\partial_wF_w=F_w''/2$. The first term is at most
$\phi(0)|t|/(2n^{3/2})$ in absolute value. Applying the heat identity
once more bounds the second term by
\[
 \frac14\int_v^n\int_w^n
       |F_r^{(4)}(y)|\,dr\,dw
 \le\frac{\|\phi'''\|_\infty}{8v^2}
 \le\frac1{4\pi v^2}.
\]
The last inequality follows by differentiating the absolutely
convergent Gaussian Fourier integral, which gives
$\|\phi'''\|_\infty\le2/\pi$. Multiply by $n^{3/2}$ and use
$\phi(0)<2/5$. The function $x^{3/2}/(x-1)^2$ decreases for $x>1$,
because its logarithmic derivative is
$3/(2x)-2/(x-1)=-(x+3)/(2x(x-1))<0$. Hence, for $n\ge2$,
\[
 \frac{n^{3/2}}{4\pi(n-1)^2}
 \le\frac{\sqrt2}{2\pi}<\frac14,
\]
using $\sqrt2<3/2$ and $\pi>3$. This proves the stronger bound and
therefore the exact ray majorant used by the original checkers.
\end{proof}

\paragraph{Finite bands and the exact support polynomial.}\label{j:local:support:polynomial}
For integers $36\le N\le L$, use the simultaneous endpoint bounds
\[
 B_+=\frac{u^4-6u^2+3}{24\sqrt{N-1}},\quad
 D_+=\frac{u^5-15u}{120(N-1)},\quad
 D_-=\frac{u^5-15u}{120(L-1)}.
\]
For an infinite index range, set $D_-=0$ instead. Put
\[
 W_*=A-15D_+/2,\qquad
 T_*=A-3D_-+\frac{9B_+^2}{4W_*}.
\]
The proof of Lemma~\ref{j:local:support:envelope} applies uniformly on the whole band
if $W_*>0$ and $T_*\ge A/2+2B_++3D_+$. In particular the favorable
term is frozen at the upper index endpoint, and the adverse terms at
the lower one. No monotonicity of a signed coefficient is assumed.

Let $(P,H)$ be an attained positive closed maximizing pair, standardized
to mean zero and variance one. Write $C=C_n\ge C_E$,
$\rho=\mathbb E|X|^3$, $\beta=\mathbb E[X|X|]$, and
$d=n(C_{n-1}-C_n)_+$. The finite contact/predecessor bound is
\[
 \rho I(t)\le\mathcal H_n(H,t)+
 C\{-|t|^3+3\rho t^2/2+3\beta t+\rho e_n\}+\rho s_nd,
\]
with $I=0$ on the entire closed support. Suppose also, for a finite
band, that $C_n-C_E\ge K/n$, where $K>0$. The finite moment bounds of~\cite[p.~125]{Shevtsova2013}, namely
$C\rho\le .3031(\rho+.646)$ and
$C\rho\le .3322(\rho+.429)$, give
\[
 c=.4095+K/L<C_n,\qquad
 \rho<r=\min\left\{\frac{.3031(.646)}{c-.3031},
                        \frac{.3322(.429)}{c-.3322}\right\}.
\]
For the unrestricted nonnegative case put $K=0$, $c=.4095$.
The sign-moment inequality
$\rho\ge(1+\beta^2)/\sqrt{1-\beta^2}$ supplies $|\beta|<b$ whenever
$(1+b^2)^2>r^2(1-b^2)$ and $0<b<1$. All these statements concern the
same current law. The lower endpoint cannot replace $L$ in $K/L$.

When a preceding resource enclosure supplies sharper actual bounds
$\rho\le r$ and $|\beta|\le b$, these may instead be used in the
same polynomials. Their validity is the only extra premise.

Use rational $\ell^2\le N-1$ and $S^2>N/(N-1)$. For $u=|t|$ define
\[
 Q(u)=-u^3+3ru^2/2+3bu+.384r/N.
\]
First check $Q<0$, before replacing $C$ by its lower bound $c$.
Replace $1/\sqrt{N-1}$ by $1/\ell$ in the adverse remainder and in
the term $2B_+$ of the domination gate. Write $E_*$ for that rational
remainder. Strict negative contact for every $n$ in the band follows
from the polynomial inequality
\[
 \frac25S^3\left\{(A-3D_-)W_*+
              \frac{(u^4-6u^2+3)^2}{256(N-1)}\right\}
 +S^3E_*W_*+\{cQ+rSq\}W_*<0.
\]
The second polynomial gate is $W_*>0$. The third is
\[
 \{A/2-3D_--3D_+-2B_+\}W_*+
         \frac{(u^4-6u^2+3)^2}{256(N-1)}>0.
\]
This is the cleared-denominator version of the domination hypothesis.
The bound on $T_*$ is positive, so the replacement $\phi_0<.4$ and
$s_n<S$ both have their correct direction.

The exact kernel tests these polynomials on all adjacent closed intervals
of length $.05$ from $R$ to $4.5$. Positive Bernstein coefficients prove
their signs on every point of each cell, including endpoints. On the
rays $u\ge4.5$ use Lemma~\ref{j:local:ext:global-gaussian-ray}, which gives the same bound
$|\mathcal H_n(H,t)|\le u^3/15+u/5+1/2$.
The resulting cubic and its first three derivatives are strictly negative
at $4.5$, as is the coefficient $Q$ before replacing $C$ by $c$.
This excludes the whole rays. Thus, if $d\le q$, actual support is in
$(-R,R)$, including exclusion of arbitrarily small exterior atoms and
exterior accumulation points.


\paragraph{The support rows used in this proof.}
For the all-index local argument, take $K=0$, $R=3.84$, $q=.19$ on
\[
[45,49],\ [50,59],\ [60,89],\ [90,134],\ [135,199],\ [200,\infty).
\]
Only indices at least 47 are used. The last row uses $D_-=0$;
there is no upper index restriction. The exact checker reconstructs
all closed $u$ cells and the two exterior rays. Missing an index band
or the infinite row is an acceptance failure.

\paragraph{One bounded-threshold series.}\label{j:local:support:series}
Put $c_k=(-1)^k/[2^kk!(2k+1)]$, and let $(x)_j$ denote the rising
factorial. Define the following rational polynomials:
\begin{align*}
 P_j(H,t)={}&\sum_{k=0}^{j+1}
 \frac{c_k(k+1/2)_{j+1-k}}{(j+1-k)!}(H-t)^{2k+1}
       -c_{j+1}H^{2j+3}\\
 &+\frac{(-1)^{j+1}tH^{2j+2}}{2^{j+1}(j+1)!}
       +\frac{(-1)^jH^{2j+1}(t^2-1)}{2^{j+1}j!}.
\end{align*}
In particular,
\[
 P_0=(t^3-3t)/6,\qquad
 P_1=\frac{3t-t^3}{4}H^2+\frac{t^4-6t^2+3}{8}H
                  +\frac{-t^5+10t^3-15t}{40}.
\]
Fix $h>0$, $E(x)=(193/71)^{\lceil x\rceil}$, and
\[
 B_r=\frac{7(h+r)}6E((h+r)^2/6)
       +\left(h+r+\frac{h(r^2+1)}8\right)E(h^2/8).
\]
\begin{lemma}\label{j:local:support:finite-series}
For every $n\ge4$, $|H|\le h$ and $|t|\le r$,
\[
 \left|\mathcal H_n/\phi(0)-\sum_{j=0}^5P_j/n^j\right|
       \le4^7B_r/n^6.
\]
\end{lemma}
\begin{proof}
Put $z=1/n$. Integration of the Gaussian exponential series gives
$\mathcal H_n/\phi(0)=S(z)/z$, where
\begin{align*}
 S(z)={}&\sum_{k\ge0}c_k\{(H-t)^{2k+1}(1-z)^{-k-1/2}
                                  -H^{2k+1}\}z^k\\
 &+te^{-H^2z/2}+\tfrac12zH(t^2-1)e^{-H^2z/2}.
\end{align*}
Its constant coefficient vanishes, and expansion of the binomial factor
gives $[z^{j+1}]S=P_j$. At $z_0=1/4$, the sum obtained by replacing
all coefficients by their absolute-value majorants is at most $B_r$:
use $(1-z_0)^{-1/2}<7/6$, $1/(2k+1)\le1$, and $e<193/71$.
The latter bound follows, for example, by summing the exponential series
through degree six and bounding the rest by a geometric series with
ratio $1/8$. All products and rearrangements are therefore absolutely
convergent. The discarded terms have degrees at least seven in $S$,
so their sum at $0\le z\le z_0$ is at most $(z/z_0)^7B_r$.
Divide by $z$ to get the result. This proves a finite uniform bound,
including every threshold in the stated interval and every $n\ge4$.
\end{proof}

\paragraph{A common contact test for radius and strip refinement.}
\label{j:local:support:panels}
Let $N\ge36$, $E,q\ge0$ and $C_0>0$. Suppose $n\ge N$,
$0\le C-C_E\le E/n$, $\delta_n:=n(C_{n-1}-C_n)_+\le q$, $C\le C_0$ and
$|X|\le R_0$. A previously acquired row gives
\[
 \rho\le r,\quad 0\le b_l\le\beta\le b_u<1,\quad |H|\le h.
\]
Define $Q(t)=-|t|^3+3\rho t^2/2+3\beta t$ and
$L_P(t)=C_EQ(t)+\phi(0)(t^3-3t)/6$.
For $t>0$, upper polynomials for $L_P(t),Q(t)$ are
\[
 (1/15-.4095)t^3+1.5(.40974)rt^2+(3(.40974)b_u-.3988/2)t,
 \quad -t^3+1.5rt^2+3b_ut.
\]
For $t<0$, write $u=|t|$ and use
\[
 (-.3988/6-.4095)u^3+1.5(.40974)ru^2+(.2-3(.4095)b_l)u,
 \quad -u^3+1.5ru^2-3b_lu.
\]
These coefficient choices are valid because $b_l\ge0$ is checked first.
On any closed cell of one sign, let $L_J$ be an upper bound for the
first polynomial and let $Q_J\ge0$ bound the positive part of the second.
For $j=1,2,3$, take $m_j$ to be the positive part of the largest tensor
Bernstein coefficient of $P_j$ on $[-h,h]\times J$. For $j=4,5$, the
absolute monomial bound
$\sum_{a,b}|[H^at^b]P_j|h^aR_J^b$ suffices, where
$R_J=\max_{t\in J}|t|$. The support contact is then bounded above by
\begin{equation}\label{j:local:support:paneltest}
 L_J+\frac{E Q_J}{N}
 +.4\left(\sum_{j=1}^5\frac{m_j}{N^j}
                +\frac{4^7B_{R_J}}{N^6}\right)
 +r\left(qS+\frac{.384C_0}{N}\right),
 \qquad S>\sqrt{N/(N-1)}.
\end{equation}
Indeed substitute the finite series in contact, split $C=C_E+(C-C_E)$,
and use the scale estimate $e_n<.384/n$. Every error retained in
\eqref{j:local:support:paneltest} is nonnegative and decreases with $n$.
A strictly negative rational result therefore excludes the entire cell
for every integer $n\ge N$, including atoms and accumulation points.
The uniform-tail specialization is $(N,E,q,C_0)=(47,.7,.19,.425)$;
its score cap follows from $C<.40974+.7/47<.425$.
The finite bands use their own acquired $E,q,C_0$ at the corresponding
floor, with exactly the same formula.

For a coordinate strip, conditional centering gives
\[
 X=\ell(Y-p),\quad
 \ell=\sqrt{(1-\theta)/(p(1-p))},\quad
 \beta=(1-\theta)(1-2p)+\kappa,\quad |\kappa|\le\theta.
\]
If $p\in[p_l,p_u]\subset(0,1/2)$ and $\theta\le\Theta$, put
$\ell_l\le\sqrt{(1-\Theta)/(p_u(1-p_u))}$.
The forbidden strip $Y\notin(-.55,1.55)$ maps into
\[
 X\le-\ell_l(p_l+.55),\qquad X\ge\ell_l(1.55-p_u).
\]
Apply \eqref{j:local:support:paneltest} between these thresholds and $\pm R_0$.
An interval beyond the known radius is empty; an equal endpoint is a
singleton that is still checked. If only an upper $\beta$ bound is given,
use $b_l=(1-\Theta)(1-2p_u)-\Theta$, after checking its nonnegativity.

The selected consequences used in the resource schedule are:
\[
\begin{array}{l|c|c|r}
\text{input rows}&\text{prior radius}&\text{conclusion}&\text{closed cells}\\\hline
10\text{ central preparation rows}&3.84&|X|<2.5&1080\\
45\text{ central rows}&2.5&-.55<Y<1.55&1431\\
38\text{ left rows}&2.7&-.55<Y<1.55&3073
\end{array}
\]
The respective largest contact bounds are below $-.084$, $-.097$ and
$-.01658$. These implications acquire no row by themselves: each is
used only after its input row is proved by the preceding resource step.
The left radii are applications of the all-threshold support polynomial,
with $(N,L,K,q)=(47,\infty,0,.19)$ and the actual moment caps
\[
 (r,b,R)=(1.390,.430,3),\qquad (1.2563,.3926,2.7).
\]
Thus the same heat proof supplies the initial radius, both left
contractions, and the five finite-base rows. The bounded-threshold
series supplies the central contraction and both coordinate strips.




\subsection{Acquiring simultaneous resources for one marginal}\label{j:local:sign}
For a ratio maximizer $C=C_n\ge\CE$, put $e=C-\CE$ and $w=H^2/(2n)$.
The core identity is
\[
 ne\rho+n\CE\rho(1-e^{-w})+n\phi(0)e^{-w}D/6+n^{3/2}L=\mathcal R.
\]
The first complete raw enclosure must give $\mathcal R<r<.4N$, for the
actual $n\ge N$. Since $C\rho\ge\CE>.4$, this licenses
Corollary~\ref{j:core:license}, including
\[
 ne\rho+n\phi(0)D/6+n^{3/2}L\le\mathcal R,\quad
 e^{-w}\ge1-5r/(2n),\quad nD\le15000r/997.
\]
The proof uses $e\ge0$ and never divides by $e$.
If a previous pass gives $nD\le Nd$, subtracting the favorable phase
term gives the sharper shared budget
\begin{equation}\label{j:local:budget}
 ne\rho+\kappa nD+n^{3/2}L\le\mathcal B(H),\qquad\kappa=997/15000,
\end{equation}
by increasing each raw quadratic curvature by
\begin{equation}\label{j:local:re:extra}
 q_{\rm extra}(\rho)=\tfrac12(1-5r/(2N))(.4095\rho-d/15)>0.
\end{equation}
Indeed $1-e^{-w}\ge we^{-w}$, $D\le d$, $\phi(0)<.4$ and
$\CE>.4095$. The positivity gate is checked before this use.
The field called \texttt{Qbudget} already contains this addition.
For a cell with $D\ge d_l\ge0$ and $\rho\ge\rho_l>0$, the same identity
also gives the score update
\[
 C\le .40974+\{r/N-.3988(1-5r/(2N))d_l/6\}/\rho_l.
\]
A negative numerator excludes the cell by $e\ge0$. Otherwise the numerator
increases with $1/n$, proving this bound at every $n\ge N$.


\paragraph{One positive-measure rule.}
Suppose zero lies in the gap $(-A,B)$, $A,B>0$. Let
\[\begin{gathered}
 p=\Pp(X>0),\quad q=1-p,\quad m=\E X_+=\E X_-,\quad V=\E[X\mid\mathbf1_{X>0}],\\
 W=X-V,\quad\theta_\pm=\E[W^2;\pm X>0],\quad\theta=\theta_++\theta_-.
\end{gathered}\]
Then $m^2=(1-\theta)pq$, and the conditional means have magnitudes
$\bar A=m/q$, $\bar B=m/p$. Set
$\Delta_A=\bar A-A\ge0$, $\Delta_B=\bar B-B\ge0$.
We also use the unnormalised endpoint displacements
\[
 M_+=\E[(X-B)\mathbf1_{X>0}]=p\Delta_B,\qquad
 M_-=\E[(-X-A)\mathbf1_{X<0}]=q\Delta_A.
\]

\begin{lemma}[Positive-measure pricing]\label{j:local:re:pricing}
For $x=|X|$, $a=|V|$, and a function $h$ constant on each sign class,
\begin{equation}\label{j:local:re:conditional}
 \E[h(|X|^j-|V|^j)]
 =\E\left[hW^2\sum_{k=0}^{j-2}(k+1)a^kx^{j-2-k}\right],\quad j\ge2.
\end{equation}
Also, for $d\mu_0=(|X|-1)^2(|X|+1/2)\,dP$,
\begin{align*}
 \mu_0(\R)&=\rho-1,\\
 M_4-1&=\int\frac{(x+1)^2}{x+1/2}\,d\mu_0,\\
 M_5-1&=\int\frac{x^3+2x^2+3x+3/2}{x+1/2}\,d\mu_0,\\
 K_1+3\rho\beta&=\int
  \frac{\operatorname{sgn}(X)x(x+2-3\rho/2)}{x+1/2}\,d\mu_0.
\end{align*}
For any nonnegative resources $z_i$ satisfying $\sum a_i z_i\le r$,
$a_i>0$, a signed linear cost obeys
\begin{equation}\label{j:local:re:simplex}
 \sum b_i z_i\le r\max\{0,b_1/a_1,\ldots,b_k/a_k\}.
\end{equation}
\end{lemma}
\begin{proof}
Divide $x^j-a^j-ja^{j-1}(x-a)$ by $(x-a)^2$; the linear term has
zero conditional mean. The radial identities follow by multiplying
out and using $\E X=0$, $\E X^2=1$. Finally each $b_i$ is at most
$a_i$ times the displayed maximum, and all $z_i$ are nonnegative.
\end{proof}

This one rule accounts for the moment, signed Taylor, endpoint movement,
outer-gap, and confinement-error estimates. In particular, no two costs
may each spend the full residual and then be added.



The gap identity and simultaneous defect budget are
\eqref{j:core:gap-defect} and \eqref{j:core:simultaneous}, with
$(A_-,A_+)=(A,B)$ and $(\Delta_-,\Delta_+)=(\Delta_A,\Delta_B)$.
The outer-gap estimate \eqref{j:core:outer-gaps} bounds the same
$T_\pm,R_\pm,J_o$. One resource pays for all their charges together.

\paragraph{A state, its update, and its exported resource.}
A state for a closed probability interval $I=[l,u]$ records
\[
 |X|\le R,\quad C\in[c_0,c_1],\quad\mathcal R\le r,
 \quad nD\le Nd,\quad\theta\le\Theta<1,
 \quad A\ge A_l>0,\quad B\ge B_l>0,
\]
together with simultaneous signed-moment and coefficient intervals.
The distinction between $nD\le Nd$ and $D\le d$ is essential.
Let $s_+=\max_Ip(1-p)$, $s_-=\min_Ip(1-p)$ and
\begin{gather*}
 b_I=\frac{(6+2\sigma)(1-\Theta)^{3/2}}{\sqrt{s_+}}
                 \operatorname{dist}(p_*,I)^2,
 \qquad r_I=d-b_I,\\
 K_+=3(\sigma+1)B_l,\quad K_-=3(\sigma+3)A_l,
 \quad K=\min(K_+,K_-),\quad m_l=\sqrt{(1-\Theta)s_-}.
\end{gather*}
If $r_I<0$, exclude the entire closed cell. Otherwise \eqref{j:core:simultaneous}
gives, simultaneously,
\begin{equation}\label{j:local:re:residual}
 6m(A_l+B_l)(\Delta_A+\Delta_B)+K_+\theta_++K_-\theta_-\le r_I.
\end{equation}
Thus the actual law satisfies the intersected bounds
\begin{gather}
 \theta\le\min\{\Theta,r_I/K\},\notag\\
 A\ge\max\left\{A_l,\sqrt{(1-\Theta)l/(1-l)}-
                  \frac{r_I}{6m_l(A_l+B_l)}\right\},\notag\\
 B\ge\max\left\{B_l,\sqrt{(1-\Theta)(1-u)/u}-
                  \frac{r_I}{6m_l(A_l+B_l)}\right\}.\label{j:local:re:geometry}
\end{gather}
Each right side is evaluated from established bounds; sequential
improvements may then be used by the next expression. A nonpositive
endpoint estimate fails this update, rather than excluding a law.
Initialization uses the same formulas with the valid lower sum
$A+B\ge\sqrt{2-d}$.

In Bernoulli coordinates set $Y=B_0+U$, $\E[U\mid B_0]=0$,
$B_0\sim\mathrm{Bernoulli}(p)$, $X=\ell(Y-p)$ and
$d_Y=(\Delta_A+\Delta_B)/(\bar A+\bar B)$.
Then $\ell^2=(1-\theta)/(pq)$ and $\lambda=n\E U^2=npq\theta/(1-\theta)$.
The \emph{weighted} versions of \eqref{j:local:re:residual} prove
\begin{equation}\label{j:local:re:exports}
 n\theta\le\frac{Nr_I}{K},\quad
 nd_Y\le\frac{Nr_I}{6(1-\Theta)(A_l+B_l)},\quad
 \lambda\le\frac{Ns_+r_I}{(1-\Theta)K}.
\end{equation}
Indeed $nD\le Nd$ whereas $nb_I\ge Nb_I$.
These exports are derived from the weighted resource; multiplying an
unweighted variance ceiling by $N$ would not prove them.

For $j=3,4,5$ set
\[
 L_3(v)=R+2v,\quad L_4(v)=R^2+2Rv+3v^2,
 \quad L_5(v)=R^3+2R^2v+3Rv^2+4v^3.
\]
Put $v_+=\sqrt{(1-l)/l}$, $v_-=\sqrt{u/(1-u)}$ and
$k_j=[1-(1-\Theta)^{j/2}]/\Theta$, with $k_j=j/2$ at $\Theta=0$.
For $u<1/2$, valid Bernoulli moment ceilings are, with $s=l(1-l)$,
\[
 B_3^*=(1-2s)/\sqrt s,\quad B_4^*=1/s-3,
 \quad B_5^*=(1-4s+2s^2)/s^{3/2}.
\]
They decrease towards $p=1/2$. Convexity places $(1-\theta)^{j/2}$
below its endpoint chord on $[0,\Theta]$; positive-measure pricing gives
\begin{equation}\label{j:local:re:moments}
 M_j\le B_j^*+r_I\max\left\{0,
 \frac{L_j(v_+)-k_jB_j^*}{K_+},
 \frac{L_j(v_-)-k_jB_j^*}{K_-}\right\}.
\end{equation}
The same rule, using
$T_+\le2mv_+\Delta_B+L_3(v_+)\theta_+$ and its reflection, bounds
\begin{equation}\label{j:local:re:outerresource}
 W_o:=(\sigma-1)T_++(\sigma+1)T_-
 \le r_I\max\left\{\begin{gathered}
 \frac{(\sigma-1)v_+}{3(A_l+B_l)},
 \frac{(\sigma+1)v_-}{3(A_l+B_l)},\\
 \frac{(\sigma-1)L_3(v_+)}{K_+},
 \frac{(\sigma+1)L_3(v_-)}{K_-}\end{gathered}\right\}.
\end{equation}
Intersect every new bound with the previous one. The versions which
merge $\theta_+$ and $\theta_-$ are obtained by replacing both costs by
their maximum and both prices by their minimum; they need no new proof.

\paragraph{The enclosure theorem.}\label{j:local:re:operator}
On each cell, the formula atlas bounds the four coefficients in
\eqref{j:core:remainder}. The result, for each actual sign $s=\operatorname{sgn}H$,
is a finite family of quadratics
\begin{equation}\label{j:local:re:quadratics}
 \mathcal R\le\max_j(A_j+F_{j,s}|H|-Q_j^{\rm raw}H^2),
 \qquad Q_j^{\rm raw}>0.
\end{equation}
The maxima retain all score--moment corners. In the initial signed
cover, intervals contain both signs of $\beta,\gamma$; no orientation is
discarded. In the probability cover, the central integral bound is a
maximum of affine functions of $\rho$, hence convex. The other
$\rho$-dependence is affine or the maximum of convex quadratics. For
fixed $H$, the score dependence is convex. Endpoints are therefore taken
\emph{before} completing the square.

\begin{theorem}[Finite resource enclosure]\label{j:local:acquisition}
Consider a finite ordered schedule of the following operations:
\begin{enumerate}
\item Split a cell into an explicitly covered family of closed cells.
\item Intersect its enclosures with an inequality proved in the atlas.
\item Exclude a cell by an empty necessary interval or a strictly
negative nonnegative resource.
\item Apply a separately proved support or confinement implication,
checking all its hypotheses against this cell's already acquired data.
\item Export a raw, weighted, or signed phase bound by
\eqref{j:local:re:quadratics}, \eqref{j:local:budget}, and \eqref{j:local:re:exports}.
\end{enumerate}
If exact outward arithmetic checks every domain, corner and cover,
every output contains all attaining laws satisfying the input. The
conclusion holds for all integers $n\ge N$ wherever the formula's
monotonicity gates are checked. It includes zero conditional variance,
singleton intersections and arbitrarily small atoms.
\end{theorem}
\begin{proof}
Induct over the schedule. A split preserves the union, intersections use
necessary inequalities, and an excluded cell contains no feasible law.
Each callback consumes only earlier hypotheses for that same law.
For a scalar raw output, use
$fh-qh^2\le(f_+)^2/(4q)$ for $h\ge0$, $q>0$; using $f^2$ is a valid
weaker bound. For a resource output first add \eqref{j:local:re:extra} to each
curvature. A negative maximum excludes the cell; a nonnegative one may
be divided by $\kappa N$ to give its next defect parameter.
The phase bound follows from the positive root of
$A+fh-qh^2\ge0$, replacing $A$ by $A_+$ if necessary.
A strict output is exported only when the certificate retains a strict
upper gap, such as its strict Fourier or target enclosure; closed-domain
coverage does not change strictness. All constants are frozen before evaluating
corners. Thus no iteration
step relies on an unproved fixed point. The stated index uniformity is
that of the displayed finite formulas, not a fixed-law limit.
\end{proof}

For the Fourier cancellation estimate only, reflect $(X,H)$ if needed:
$I_0(-X,-H)=I_0(X,H)$. Both odd coefficients change sign, and
$\int e^{-x^2/2}|zx^2-cx^4|\,dx$ is unchanged by $(z,c)\mapsto(-z,-c)$.
Consequently a lower endpoint used there bounds $|\gamma|$, not necessarily
the original signed $\gamma$. The actual contact, gap and phase data keep
their original orientation. A useful additional signed lower bound is
\begin{equation}\label{j:local:re:gamma}
 \gamma\ge(1-\Theta)^{3/2}\frac{1-2u}{\sqrt{u(1-u)}}
       -\left(R+2\sqrt{\frac{u}{1-u}}\right)\Theta.
\end{equation}
To prove it expand $\E X^3=\E V^3+\E[W^2(X+2V)]$; the positive
residual is nonnegative and the negative residual has the displayed
lower bound. The binary term decreases in both $p<1/2$ and $\theta$.
Thus $\max(0,\text{any signed lower bound})$ is a valid lower bound for
$|\gamma|$ inside this auxiliary Fourier estimate.



\paragraph{Signed moments with a common first moment.}\label{j:local:atlas}\label{j:local:moments}
Put $a_\pm=(1\pm\beta)/2$, $b_\pm=(\rho\pm\gamma)/2$.
Cauchy gives $m^2\le a_+a_-$ and $b_\pm m\ge a_\pm^2$.
Consequently
\[
 \frac{(1+\beta)^{3/2}}{\sqrt{1-\beta}}-\rho
 \le\gamma\le
 \rho-\frac{(1-\beta)^{3/2}}{\sqrt{1+\beta}},\qquad
 m\ge\max\left\{\frac{1+\beta^2}{2\rho},
                  \frac{a_+^2}{b_+},\frac{a_-^2}{b_-}\right\}.
\]
These statements also hold with a zero atom. Projection of $X|X|$ onto
the span of $X$ and $|X|$ gives
\[
 \mu_4\ge\rho^2+\frac{(\gamma-\rho\beta)^2}{1-\beta^2}.
\]
For completeness, this follows by minimizing the nonnegative expectation
$E(X|X|-uX-v|X|)^2$ over $u,v$; its Gram matrix has diagonal entries
one and off-diagonal entry $\beta$, and is positive definite.

The signwise moment matrices below, valid also for different radii,
provide the bounded-support fourth- and fifth-moment contractions.
A second useful upper bound for $m$ follows by integrating
$(x-a)^2(R-x)\ge0$ for $x=|X|\in[0,R]$. For every fixed $a>0$,
\[
 2m\le\frac{R-\rho+2a+Ra^2}{a(2R+a)}.
\]
The cover uses only the eleven rational choices $a=2/10,\ldots,12/10$.

Write $\delta=\gamma-\rho\beta$. Cauchy and a square completion on
each sign give the signed constant bound
\begin{align*}
 \frac\rho2\mu_4-\frac14\eta_5
 &\le\frac{\rho^3}{4}
 -\frac{(\rho+\gamma)(\delta_+)^2}{8(1+\beta)^2}
 -\frac{(\rho-\gamma)((-\delta)_+)^2}{8(1-\beta)^2}.
\end{align*}
Indeed on either sign, write $a=\E X_s^2$, $b=\E X_s^3$,
$c=\E X_s^4$, $f=\E X_s^5$. Cauchy gives $c\ge b^2/a$,
$f\ge c^2/b$, and
$\rho c/2-f/4\le\rho^2b/4-(c-\rho b)^2/(4b)$.
Distance to the half-line $c\ge b^2/a$ proves the stated penalty
on summing both signs. Independently,
$\mu_4^2\le\rho\eta_5$ and the projection bound above give
\[
 \frac\rho2\mu_4-\frac14\eta_5
 \le\frac{\rho^3}{4}
 -\frac{(\gamma-\rho\beta)^4}{4\rho(1-\beta^2)^2}.
\]
The larger of the two nonnegative deductions is retained.

\paragraph{An actual central gap without assuming \texorpdfstring{$D<2$}{D below 2}.}
If zero belongs to the support, the signwise gap calculation gives $J\le\rho$.
Thus $D\ge\sqrt{10}\rho-\gamma\ge(\sqrt{10}-1)\rho$.
Whenever $D<(\sqrt{10}-1)\rho$, there is therefore a central gap
$(-A,B)$, $A,B>0$, with no zero atom. Put $p=P(X>0)$, $q=1-p$,
$T_+=b_+-mB^2$ and $T_-=b_--mA^2$. The exact gap decomposition is
\[
 D=mQ+(\sigma-1)T_++(\sigma+1)T_-+3(T-J_o),\qquad
 Q=(\sqrt{\sigma+1}A-\sqrt{\sigma-1}B)^2,
 \quad\sigma=\sqrt{10}.
\]
All terms are nonnegative. Hence $T\le D/(\sigma-1)$ and
\[
 (A+B)^2\ge A^2+B^2\ge\frac{\rho-D/(\sigma-1)}m.
\]
Moreover $A\le a_-/m$, $B\le a_+/m$ and $J\le\rho+2m$.
The last bound gives the compulsory defect
\[
 D\ge \sigma\rho-\gamma-6m.
\]
The inequalities $T-J_o\ge2BR_++2AR_-$ give the useful additional
endpoint constraint
\[
 D\ge(\sigma-1)b_++6a_+B-(\sigma+5)mB^2.
\]
If $D<(\sigma-1)b_+$, this downward quadratic forces $B$ to be at
least its positive root. The reflected constraint has $a_-,b_-,A$
and the coefficients $\sigma+1,\sigma+7$. In a closed moment box,
positive coefficients are rounded down, negative coefficients up, and
the lower square-root endpoint is used in the resulting root bound.

The actual opposite-sign event has probability $2pq$.
The first two moment constraints give lower bounds for both $p$ and $q$.
For independent copies one has exactly
\[
 E[(X-X')^2\mid XX'<0]
 =\frac{1+\beta(1-2p)+4m^2}{2pq}
 =\frac{1+4m^2+\beta}{2p}
  +\frac{1+4m^2-\beta}{2q}.
\]
For an upper bound, use the upper $m$ endpoint and both endpoints of
$p$ and $\beta$. The expression is convex in $p$ and affine in $\beta$.
This improves the former $2/P(XX'<0)$ bound inside the complete HIGH
estimate. Both copies still have the same actual marginal.

\paragraph{Preserving central and outer integral costs.}
Let
\[
 P_1(A,B)=B^3/3+AB^2/2-A^3/6,\quad
 P_2(A,B)=B^4/4+AB^3/3-A^4/12,\quad G=\rho P_1-P_2.
\]
Then $J_{1,c}=mP_1$, $J_{2,c}=mP_2$. Since
$\partial_BG=B(A+B)(\rho-B)$,
\[
 G(A,B;\rho)\le G(A,\rho;\rho)
 =\rho^4/12+\rho^3A/6-\rho A^3/6+A^4/12.
\]
The last expression increases with $A$ because its derivative is
$(\rho-A)^2(\rho+2A)/6$, and is convex in $\rho$ because its second
$\rho$ derivative is $\rho^2+\rho A$. Its closed-rectangle maximum is
therefore bounded using the upper $A$ endpoint and both $\rho$ endpoints.

Positive outer gaps satisfy
$(t-a)(\rho-t)\le(\rho-B)_+^2/4$. On negative outer gaps the positive
part is bounded by $(R-\rho)^2/4$ after extracting $|t|$.
Integration gives
\[
 3(\rho J_1-J_2)_o
 \le\frac38\{(\rho-B)_+^2T_++(R-\rho)^2T_-\}.
\]
The available weighted outer budget is $D-mQ$, not the entire $D$.
The cover intersects this budget with the individual bounds
$T_+\le b_{+,u}-m_lB_l^2$ and $T_-\le b_{-,u}-m_lA_l^2$.

For the signed phase term $U=J_1+(\mu_4-3)/6$, retain
\[
 -RT_-/2\le J_{1,o}\le RT_+/2.
\]
The central minimum of $P_1$ occurs at $B_l$ and an $A$ endpoint;
its maximum occurs at $B_u$ and $A$ clipped to the interval at $A=B_u$.
A nonnegative polynomial bound is multiplied by $m_u$ for an upper
bound and $m_l$ for a lower bound; a negative one uses the opposite
endpoint. For example, the binary endpoints $A=3/2$, $B=2/3$ make this
maximum negative; the sign switch is necessary even when $D<2$.
The fourth-moment sign intervals also sharpen
$K_1=E[X|X|^3]-3\beta-3\rho\gamma/2$.
The complete $K_1,U,\phi$ rectangles are used in the signed expression
$|CK_1+\phi U|$, without replacing it by two separate absolute values.


\paragraph{A complete coarse bound before a central gap is licensed.}
Always $-\mu_4/6\le J_1\le\mu_4/2$ and
$-3J_2-\eta_5/4\le0$. On a gap of length $h=b-a$, these follow from
\[
 \int_a^b(t-a)t\,dt+\tfrac12\int_a^bt^2dt=hb^2/2,\quad
 \tfrac32\int_a^bt^2dt-\int_a^b(t-a)t\,dt=h(2a+b)^2/6,
\]
and, on a negative gap $(-v,-u)$, from
$\tfrac13\int_u^v x^3dx-\int_u^v(v-x)x^2dx=u^3(v-u)/3$.
For $1\le r\le\rho$ and actual support radius $R\le6$, set
\[
 \Gamma(R,r)\ge\max_{r\le v\le R}
 \left\{-\frac{v^2}4+\frac{(R+r)v}3-\frac{Rr}2+
                      \frac{r^3(2R-r)}{12v^2}\right\}_+.
\]
Then
\[
 3(\rho J_1-J_2)\le\frac32\Gamma(R,r)(\rho-\gamma)
                              +\frac38\rho^2J.
\]
On negative gaps extract $x=|t|$ and discard the nonpositive part
$x<\rho$. Replacing $\rho$ by $r$ and the outer endpoint by $R$
gives $\int_r^v(R-x)x(x-r)\,dx$. Dividing this integral by $v^2$
gives exactly the displayed expression. Signwise Tonelli now charges
$\Gamma\E X_-^3=\Gamma(\rho-\gamma)/2$.
On positive gaps $(t-a)(\rho-t)\le\rho^2/4$, giving $\rho^2J/8$.
The implementation uses 64 closed $v$-intervals: on each, maximize the
concave quadratic at its clipped vertex $2(R+r)/3$ and the positive
reciprocal term at the left endpoint. This proves its interval bound.
The alternative bound $\Gamma\le5/2$ follows at $r=1,R=6$ by
$3v^4-28v^3+66v^2-11>0$ on $[1,6]$: the endpoint values are $30,205$,
and the interior critical values are $(85\mp35\sqrt5)/2>0$.
These monotone enclosing parameters cover all $r\ge1,R\le6$.
\paragraph{Moment and gap coefficients for a whole cell.}
\label{j:local:mom:conditional-boxes}

The following coefficient rules are collected here to make the finite
calculation transparent. Every rule is an inequality for an actual law;
there is no extremal-law assumption inside a cell.
Write $v=1-\theta$. For $p\in[l,u]\subset(0,1)$ set
\begin{gather*}
 s_0=\min\{l(1-l),u(1-u)\},\quad
 v_* =\max\{\sqrt{u/(1-u)},\sqrt{(1-l)/l}\},\\
 r_B=(1-2s_0)/\sqrt{s_0},\quad
 g_B=\max_{t=l,u}|1-2t|/\sqrt{t(1-t)},\quad b_B=\max_{t=l,u}|1-2t|,\\
 L_3=R+2v_*,\quad L_4=R^2+2Rv_*+3v_*^2,\\
 L_5=R^3+2R^2v_*+3Rv_*^2+4v_*^3,\quad
 e_B=r_B\frac{1-4s_0+2s_0^2}{s_0(1-2s_0)}.
\end{gather*}
The polynomials $T_3=x+2a$, $T_4=x^2+2ax+3a^2$ and
$T_5=x^3+2ax^2+3a^2x+4a^3$ are nonnegative. Their maxima at
$x=R$, $a=v_*$, in Lemma~\ref{j:local:re:pricing}, give, if $\theta\le\Theta<1$,
\begin{align}
 \rho&\le r_B+(L_3-r_B)\Theta,&
 |\gamma|&\le g_B+(L_3-g_B)\Theta,\notag\\
 |\beta|&\le b_B+(1-b_B)\Theta,&
 \mu_4&\le(1-\Theta)^2(1/s_0-3)+L_4\Theta,\notag\\
 \eta_5&\le(1-\Theta)^2e_B+L_5\Theta.&&\label{j:local:mom:boxes}
\end{align}
Replacing $\theta$ by $\Theta$ uses the explicitly checked inequalities
$L_3>r_B,g_B$, $L_4>2(1/s_0-3)$ and $L_5>2e_B$.
The exact binary factors are $v^{3/2},v,v^2,v^{5/2}$; the upper bounds
use $v^{3/2}\le v$ and $v^{5/2}\le v^2$. Thus cells crossing $1/2$ and
zero conditional variances are included. The same cubic expansion, in the
later domain $p\in[.37,.47]$, $A\ge.75$, $B\ge1.05$, gives
\begin{equation}\label{j:local:mom:rho-lower}
 \rho\ge v^{3/2}\rho_B(p)+3\min(A,B)\theta
       \ge\rho_B(p)+.45\theta,\qquad
 \rho_B(p)=\frac{1-2pq}{\sqrt{pq}}.
\end{equation}
Indeed $\rho_B<1.11$ there and $v^{3/2}\ge1-3\theta/2$.
Also $\beta\ge1-2p-2(1-p)\theta$.

For the signed coefficient
$K_1=\E[X|X|^3]-3\beta-3\rho\gamma/2$, define
\[
 K(p,v)=(1-2p)\{(3v^3/2-v^2)/(pq)-3v^3+v^2+3v\}.
\]
Direct substitution gives $K_1(V)=-K(p,v)$, with the coefficient $3$
unchanged. For $v\ge2/3$, $K$ decreases in $p$: both coefficients in
braces are nonnegative and $(1-2p)/(pq)$ decreases. Its derivative in $v$
has the sign of $1-2p$, since the derivative of the braces is at least
$9v^2-6v+3>0$. Hence both variance endpoints are retained. The error is
\begin{equation}\label{j:local:mom:K-error}
 |K_1(X)-K_1(V)|\le
 \Theta\{L_4+3+\tfrac32L_3(R_\rho+g_B)\}
\end{equation}
when $\rho\le R_\rho$. This follows from the conditional expansions and
$\rho\gamma-\rho_V\gamma_V=
\rho(\gamma-\gamma_V)+\gamma_V(\rho-\rho_V)$.

Here are the remaining signed charges. They also explain the two
alternative bounds retained in the initial, coarse moment cover:
\begin{align}
 -\mu_4/6&\le J_1\le\mu_4/2,&
 -3J_2-\eta_5/4&\le0,                                  \label{j:local:mom:gap-signs}\\
 \rho J_1-J_2&\le\tfrac54(\rho-\gamma)+\tfrac18\rho^2(\rho+1),
                                                              \label{j:local:mom:oldgap}\\
 -75D/104&\le J_{1,o}\le3D/2.                           \label{j:local:mom:outer}
\end{align}
The first two lines follow from the gap integral identities above,
using $\Gamma\le5/2$ and $J\le\rho+1$. For the outer $J_1$ bounds,
$J_{1,o}\ge-3T_-$, $J_{1,o}\le3J_o\le3(T_++T_-)$ and
\eqref{j:core:gap-defect} give the stated constants.

For the previously defined $P_1,P_2$, $J_{j,c}=mP_j(A,B)$. Both polynomials increase in $B$; their $A$ derivative changes from positive to negative
at $A=B$. A rectangle minimum is therefore at $B_l$ and an $A$ endpoint.
Multiply a negative minimum by an upper bound for $m$, and a nonnegative
minimum by a lower bound. For the upper central charge, maximizing at
$(\bar a,\bar b)$ if $p\le1/2$ gives
$J_{1,c}\le v^2(2-3p)/(6pq)$, a decreasing function of $p$ on that interval.
If $p\ge1/2$, maximize first at $A=B$ to get
$J_{1,c}\le(2/3)v^2q^2/p\le1/3$. Thus every cell may use
\begin{equation}\label{j:local:mom:central-upper}
 J_{1,c}\le\max\{1/3,(2-3l)/(6l(1-l))\}.
\end{equation}

For $|X|\le R$, $R\ge1$, define
\[
 k_4(R)=\frac{(R+1)^2}{R+1/2},\qquad
 k_5(R)=\frac{R^3+2R^2+3R+3/2}{R+1/2}.
\]
Then
\begin{equation}\label{j:local:mom:universal-moments}
 \mu_4\le\min\{R\rho,1+k_4(R)(\rho-1)\},\quad
 \eta_5\le\min\{R\mu_4,1+k_5(R)(\rho-1)\},\quad
 \mu_6\le R\eta_5.
\end{equation}
The first two bounds are the supported moment estimates proved in
Section~\ref{j:fourier:compensated}; the last is pointwise.
The following coarse cancellation estimates retain the radius-six majorants.
The last identity in~\eqref{j:local:re:pricing} gives
\begin{equation}\label{j:local:mom:K-cancellation}
 |K_1+3\rho\beta|\le(96-18\rho)(\rho-1)/13.
\end{equation}
Indeed $u(u+2-3r/2)/(u+1/2)$ is convex on $[0,6]$, its maximum is
$(96-18r)/13$, and its negative part is at most $31/40$, smaller than that
maximum for $1\le r\le37/20$.
The following bounds are used only for $|X|\le6$ and
$1\le\rho\le37/20<2$, with the signed moment hypotheses below.
The other bounds needed on this domain are
\begin{align}
 |K_1|&\le6\rho-3\rho^2/2+3|\beta|,\notag\\
 K_1&\le M+3|\beta|-3\rho g_-/2,&
 -K_1&\le3\rho+3|\beta|+(3\rho/2-3)g_-,\notag\\
 -K_q&\le(\rho/2-3/8)M-\rho/2+1/8+3|\beta|g/2,&
 -K_q&\le8\rho^3/27+(3|\beta|/2-1/2)\rho,               \label{j:local:mom:K-bounds}
\end{align}
where $\mu_4\le M$, $|\gamma|\le g$, $\gamma\ge g_-$, and
$K_q=\eta_5/4-\rho\mu_4/2-3\beta\gamma/2+\rho/2$.
For the first two lines use
$|\E[X|X|^3]-3\gamma|\le3\rho$, and note the negative coefficients
of $\gamma$. For the last line use
$\eta_5\ge\mu_4^{3/2}\ge3\mu_4/2-1/2$, or maximize
$\rho y^2/2-y^3/4$ over $y\ge0$.


\paragraph{Compatible weaker resource constants.}
The implementation also permits the weaker binary charge
$b_I=(616/25)(1-\Theta)^2\operatorname{dist}(p_*,I)^2$.
It follows from $2(6+2\sigma)>616/25$ and
$(1-\theta)^{3/2}\ge(1-\theta)^2$.
With $K=\min(K_+,K_-)$ and $v_* =\max(v_+,v_-)$, the merged charge
\eqref{j:local:re:residual} gives
\[
 T_++T_-\le T_*=(d-b_I)\max\{v_*/[3(A_l+B_l)],(R+2v_*)/K\},
 \qquad d_o=\min\{d,(\sigma+1)T_*\}.
\]
Thus $W_o\le d_o$ and, for $R\le6$,
\[
 -75d_o/104\le J_{1,o}\le3d_o/2,\qquad
 3(\rho J_{1,o}-J_{2,o})\le\frac38
 \max\{\rho^2/(\sigma-1),(R-\rho)^2/(\sigma+1)\}d_o.
\]
These are precisely the merged-cost versions of the sign-specific
resource rule; they do not introduce another available budget.

\paragraph{Conditional lower moments improve two signed charges.}

Let $s_u=\max_{p\in[l,u]}pq$ and $\theta\le\Theta$.
The unsigned cases of~\eqref{j:local:re:conditional} give
\begin{align}
 \mu_4&\ge M_l:=\max\{1,(1-\Theta)^2(1/s_u-3)\},\label{j:local:entry201:m4lower}\\
 \eta_5&\ge E_l:=
 (1-\Theta)^{5/2}\frac{1-4s_u+2s_u^2}{s_u^{3/2}}.\label{j:local:entry201:e5lower}
\end{align}
Both binary moment functions decrease in $s\in(0,1/4]$.
For the fifth this follows because the derivative has the sign of
$s^2+2s-3/2<0$; the fourth is immediate.
The fourth-moment lower bound sharpens the lower endpoint of
$J_1+(\mu_4-3)/6$; its explicit use appears in the cell map below.

The constant signed charge is
\[
 -K_q=\frac{\rho\mu_4}{2}-\frac{\eta_5}{4}
       +\frac32\beta\gamma-\frac{\rho}{2}.
\]
The old coherent estimate used $\eta_5\ge3\mu_4/2-1/2$ and
$\mu_4\le M$. A second estimate uses \eqref{j:local:entry201:e5lower}.
Their difference is independent of $\rho$. Consequently the improved
bound is exactly the old one minus
\begin{equation}\label{j:local:entry201:eta-improvement}
 \frac14\max\{0,E_l-3M/2+1/2\}.
\end{equation}
This is a constant shift on the cell and does not spoil convexity.
It does not assert the generally invalid implication
$\eta_5\ge3M/2-1/2$ from an upper bound on $\mu_4$.

\paragraph{Signed conditional Taylor cancellation.}

Let $K_1=\E[X|X|^3]-3\beta-\tfrac32\rho\gamma$.
On a sign class put $x=|X|$, $a=|V|$, $W=X-V$ and
$\varepsilon=1$ on the positive class, $-1$ on the negative class.
In the actual central-gap domain there is no atom at zero.
With $w_3=x+2a$, $w_4=x^2+2ax+3a^2$ and
$\gamma_V=\E V^3$, exact conditional expansion gives
\begin{equation}\label{j:local:entry201:Kexact}
 K_1(X)-K_1(V)
 =\E\!\left[W^2
 \left\{\varepsilon(w_4-3-\tfrac32\rho w_3)
                   -\tfrac32\gamma_Vw_3\right\}\right].
\end{equation}
Apply~\eqref{j:local:re:conditional} to $j=2,3,4$, with $h=1$
or $h=\varepsilon$, and use
$\rho\gamma-\rho_V\gamma_V
=\rho(\gamma-\gamma_V)+\gamma_V(\rho-\rho_V)$.

The positive coefficient in \eqref{j:local:entry201:Kexact} is
$F(x,a;t)=x^2+2ax+3a^2-3-\tfrac32t(x+2a)$
with $t=\rho+\gamma_V$; the negative coefficient is $-F$
with $t=\rho-\gamma_V$. Their ranges are bounded on complete
rectangles of actual conditional means and support magnitudes.
The positive-support magnitude is at most $R$.
The same formula permits separate actual bounds on the two signed support magnitudes.
On the late domain $p<1/2$, bounds for $\gamma_V$ are
\[
 \sqrt{(1-\Theta)^3(1-2u)^2/[u(1-u)]}
 \le\gamma_V\le(1-2l)/\sqrt{l(1-l)}.
\]
The positive mean lies in
$[\sqrt{(1-\Theta)(1-u)/u},\sqrt{(1-l)/l}]$, and the negative
magnitude mean lies in
$[\sqrt{(1-\Theta)l/(1-l)},\sqrt{u/(1-u)}]$.
Use the lower bound $\rho\ge\rho_B(u)$ only after checking
$3\min(A_l,B_l)\ge(3/2)\sup_I\rho_B$. Indeed conditional expansion
and $(1-\theta)^{3/2}\ge1-3\theta/2$ then imply
$\rho\ge(1-\theta)^{3/2}\rho_B(p)+3\min(A,B)\theta\ge\rho_B(p)$.
Otherwise retain the universal lower bound one.

Since $x+2a>0$, $F$ decreases with $t$.
For fixed $t$ its Hessian is
$\left(\begin{smallmatrix}2&2\\2&6\end{smallmatrix}\right)>0$.
Its rectangle maximum is at a corner. Its minimum is among the four
corners, the minimizer on each edge, and the stationary point
$(x,a)=(3t/8,3t/8)$ if that lies in the rectangle.
On fixed-$x$ edges use $a=t/2-x/3$, clipped to the edge; on
fixed-$a$ edges use $x=3t/4-a$, also clipped.
This is a complete finite extremum list.
Multiplication by $0\le\theta\le\Theta$ uses
$\Theta\min(0,L)$ and $\Theta\max(0,U)$ for the lower and upper
coefficient bounds. It intersects this new interval with the retained
absolute-error interval; neither need be discarded.

\paragraph{The complete central rectangle.}
The central part of $\rho J_1-J_2$ is exactly
$mG(A,B;\rho)$, where $G=\rho P_1-P_2$.
Let a proved probability cell give
$A\in[A_l,A_u]$, $B\in[B_l,B_u]$, $m\in[m_l,m_u]$.
Here
\[
 A_u=\sqrt{u/(1-u)},\quad B_u=\sqrt{(1-l)/l},\quad
 m_l=\sqrt{(1-\Theta)s_0},\quad
 m_u=\sqrt{\max_{p\in[l,u]}pq}.
\]

\begin{lemma}[Complete central rectangle]\label{j:local:entry201:central}
Write $G(A_l+(A_u-A_l)x,B_l+(B_u-B_l)y;\rho)
=\sum t_{kh}x^ky^h$. Define the 25 rational coefficients
\[
 b_{ij}(\rho)=\sum_{\substack{k\le i\\h\le j}}
 t_{kh}\frac{\binom{i}{k}}{\binom4k}
          \frac{\binom{j}{h}}{\binom4h},\qquad 0\le i,j\le4,
\]
and
\[
 \mathcal G(\rho)=\max_{i,j}
       \{m_l b_{ij}(\rho),\,m_u b_{ij}(\rho)\}.
\]
Then $\rho J_{1,c}-J_{2,c}\le\mathcal G(\rho)$.
Moreover $\mathcal G$ is convex in $\rho$.
\end{lemma}
\begin{proof}
Apply Lemma~\ref{j:cert:bernstein} separately to the
two coordinates after the affine substitution. The Bernstein basis is nonnegative and sums to one on the closed square,
so $G$ is a convex combination of its 25 coefficients.
For each coefficient retain both $m$ endpoints, or equivalently use
$m_u$ when it is nonnegative and $m_l$ otherwise. This gives a uniform
bound on the entire three-dimensional rectangle, including boundaries.
Each endpoint-multiplied coefficient is affine in $\rho$; their maximum
is convex. No interior stationary point is omitted.
\end{proof}



\paragraph{One interval-feedback theorem.}
\label{j:local:entry:feedback}\label{j:local:fourier:uniform}

The signed identity from the preceding section is
\begin{equation}\label{j:local:entry:remainder}
 \cR=C\{3\beta J/2+3\rho J_1-3J_2-K_q\}
 +H(CK_1+\varphi U)-H^2(C\rho+\gamma\varphi/2)+E_n,
\end{equation}
where $U=J_1+(\mu_4-3)/6$, $\varphi=\phi(H/\sqrt n)$.
For any one coherent low-frequency triple and complete high-frequency value $K_H$ at floor $N$,
\begin{equation}\label{j:local:entry:fourier-budget}
 |E_n|\le K_F+s_F|H|,\quad
 K_F=c_\pi(L_0+K_H+L_2/\sqrt N),\quad s_F=c_\pi L_H,\quad c_\pi\ge2/\pi.
\end{equation}
The implementation uses $16/25$ before the signed phase and $212/333$
afterward; both are valid since $\pi>333/106$.

\begin{lemma}[Continuous cells and feedback]\label{j:local:entry:cell-theorem}\label{j:local:entry:corners}
Suppose bounds for one actual cell give
\[
 r_0\le\rho\le r_1,\quad0<c_0\le C\le c_1,\quad
 |\beta|\le b,\quad |\gamma|\le g,\quad \mu_4\le M,\quad D\le d,
\]
$K_1\in[k_l,k_u]$, $U\in[u_l,u_u]$, and
$0\le\varphi_-\le\varphi\le2/5$. Freeze the shared outer charge $d_o$,
lower fifth moment $E_l$, upper $s_u=\max_{p\in[l,u]}p(1-p)$, and
central envelope $\mathcal G$ from the preceding bounds.
Define
\begin{align*}
 A(r)&=\tfrac32 b(r+2\sqrt{s_u})+3\mathcal G(r)
 +\tfrac38\max\{r^2/(\sigma-1),(R-r)^2/(\sigma+1)\}d_o\\
 &\quad +(r/2-3/8)M-r/2+1/8+\tfrac32bg
 -\tfrac14\max\{0,E_l-3M/2+1/2\},\\
 B(c)&=\max_{k\in\{k_l,k_u\},\,u\in\{u_l,u_u\},\,
                      \psi\in\{\varphi_-,2/5\}}|ck+\psi u|,\\
 Q(c,r)&=cr+\tfrac12\varphi_-(3.16r-3-d)
       -\tfrac12(2/5-\varphi_-)(d-.16)_+.
\end{align*}
If all four corner values of $Q$ are positive, then, for every real $H$,
\begin{align}
 \cR&\le\max_{(c,r)\in\{c_0,c_1\}\times\{r_0,r_1\}}
       \{K_F+cA(r)+(B(c)+s_F)|H|-Q(c,r)H^2\}\notag\\
 &\le\max_{c,r\text{ corners}}
       \{K_F+cA(r)+(B(c)+s_F)^2/(4Q(c,r))\}.           \label{j:local:entry:four-corners}
\end{align}
For the following feedback, assume in addition that the cell contains a
ratio extremizer with $e=C_n-\CE\ge0$, $n\ge N$, and that
$0<T<.4N$ is a strict upper bound for the final maximum in
\eqref{j:local:entry:four-corners}. Then $T<nC_n\rho$, since
$C_n\rho\ge\CE>.4$, and the weight-removal step in
Corollary~\ref{j:core:license} gives
\begin{equation}\label{j:local:entry:feedback-values}
 D<\frac{15000T}{997N},\quad
 \varphi>\frac{997}{2500}\left(1-\frac{5T}{2N}\right),\quad
 C<\CE+V/N\quad\text{if also }\cR/\rho<V.
\end{equation}
Only positive density floors are used. These statements remain valid at
$H=0$, on closed cell edges, and for zero conditional variance.
\end{lemma}
\begin{proof}
The constant coefficient follows from the shared outer budget,
Lemma~\ref{j:local:entry201:central}, and~\eqref{j:local:mom:K-bounds}, using
\[
 J\le\rho+2ABm\le\rho+2(1-\theta)^{3/2}\sqrt{pq}
       \le\rho+2\sqrt{s_u}.
\]
This follows from $A\le\bar a$, $B\le\bar b$, $m^2=(1-\theta)pq$;
it is used only in the charge $3|\beta|J/2$. Since $\gamma\ge3.16\rho-3-d$, its negative
part is at most $(d-.16)_+$. Using the upper density only for that negative
part gives $C\rho+\gamma\varphi/2\ge Q$. Absolute value is separately
convex in each coordinate of $ck+\psi u$, proving the formula for $B$.
For fixed $r$ and $|H|$, the resulting upper bound is convex in $c$.
Take those endpoints \emph{before} completing the square. At either
endpoint $c>0$, it is convex in $r$, since $\mathcal G$ and the maximum of the two outer quadratics are
convex, every other term of $A$ is affine, and $Q$ is affine in $r$. This proves the first line of~\eqref{j:local:entry:four-corners};
$fh-qh^2\le f^2/(4q)$ proves the second. Bilinearity of $Q$ proves
its positivity throughout the rectangle from its corners.

\end{proof}

\paragraph{Signed phase curvature in the resource schedule.}
\label{j:local:re:signed-phase-repair}
The auxiliary Fourier orientation does not change the signed third moment
in the actual transfer identity. The paired phase estimate must therefore
retain its sign. On a fixed resource cell put
\[
 L=\frac{997}{2500}\left(1-\frac{5r}{2N}\right),\qquad
 g_0(\rho)=3.16\rho-3-d,\qquad U=.4,
\]
where $r,d$ are the preceding raw and defect upper bounds, and the
checker requires $0\le L\le U$. The actual density lies in $[L,U]$
and the actual signed moment satisfies $\gamma\ge g_0(\rho)$. Consequently
\begin{equation}\label{j:local:re:sign-aware-price}
 C\rho+\tfrac12\gamma\phi
 \ge q_{\rm s}(C,\rho):=
 C\rho+\tfrac12\min\{L g_0(\rho),U g_0(\rho)\}.
\end{equation}
This follows by first multiplying the lower bound on $\gamma$ by the
nonnegative actual density, and then minimizing over its two endpoints.
The simpler expression $C\rho+Lg_0(\rho)/2$ is licensed if $g_0\ge0$,
or if nonnegativity of the actual $\gamma$ has been separately proved.
It is not licensed merely by reflecting the auxiliary Fourier integral.

The corrected price preserves the four-corner reduction. For fixed $C$,
$q_{\rm s}$ is concave in $\rho$, being affine plus the minimum of two
affine functions. For fixed $\rho$ it is affine in $C$. In the phase
envelope $A+f|H|-q_{\rm s}H^2$, the remaining terms are separately
convex in $C,\rho$, by the moment and coefficient bounds already proved.
Thus the entire envelope is separately convex, and two successive
one-dimensional endpoint bounds place its maximum at the four corners.
The same argument applies after adding the affine conserved-resource
curvature $e_\rho$ and subtracting a barrier price $K\rho$.
If the four corrected divisors are positive, concavity makes them
positive throughout the rectangle, and square completion gives the
valid raw and resource bounds
\[
 \max_{C,\rho\text{ corners}}\left(A+\frac{f^2}{4q_{\rm s}}\right),
 \qquad
 \max_{C,\rho\text{ corners}}
       \left(A+\frac{f^2}{4(q_{\rm s}+e_\rho)}-K\rho\right).
\]
Any threshold-root export likewise uses the corrected corner curvature.

For completeness, nonnegativity in the later acquired domains can be checked
without replaying their iterations. Write $B=\mathbf1_{\{X>0\}}$,
$V=\mathbb E(X\mid B)$,
$W=X-V$, $\theta=\mathbb E W^2$, and $p=\mathbb P(X>0)$. Conditional
centering and $W\ge-R-V$ imply
\[
 \gamma=\mathbb EV^3+3\mathbb E(VW^2)+\mathbb EW^3
 \ge\mathbb EV^3+2\mathbb E(VW^2)-R\theta.
\]
For $p\in[l,u]\subset(0,1/2)$ and $\theta\le\Theta<1$ this gives
\begin{equation}\label{j:local:re:actual-gamma-license}
 \gamma\ge
 \frac{(1-\Theta)^{3/2}(1-2u)}{\sqrt{u(1-u)}}
 -\left(R+2\sqrt{\frac{u}{1-u}}\right)\Theta.
\end{equation}
Here the binary cubic decreases in $p<1/2$, and the right side decreases
as its variance bound increases. These are bounds on the same actual law
throughout every later call. The fixed acquired input domains license
\eqref{j:local:re:actual-gamma-license} before those calls are made.
The exact domain check covers 302 input rows: 45 central-input rows,
45 left-input rows, 45 first-confinement feedback rows, 45 Sperner feedback
rows, 84 signed-resource rows, and 38 final left-radius rows. With the
common already acquired radius $R=3.84$, each right side is greater than
$.096$. Thus all later paired calls have an actual nonnegative signed
third moment, independently of any auxiliary Fourier orientation.

There is also a convenient cell-constant correction. Since $\rho\ge1$,
$(-g_0(\rho))_+\le(d-.16)_+$, so the weaker price
\[
 C\rho+\tfrac12L g_0(\rho)
       -\tfrac12(U-L)(d-.16)_+
\]
is valid. This is the correction selected for the additional index--46
calculation. Both forms keep the raw remainder separate from the
barrier-subtracted resource budget.

The broad preliminary and complementary resource passes are replayed with
\eqref{j:local:re:sign-aware-price}; their subsequent consumers retain the same
fixed target tables only when the complete corrected replay verifies every
exported resource. A diagnostic equality of two corner maxima alone does
not replace this check, since threshold bounds also consume the phase
coefficients. Later uses of the simpler formula require their own explicit
nonnegative signed-moment or small-defect license.




\paragraph{Why the floor calculations retain their force at later indices.}
\label{j:local:transfer}
The reuse of a finite resource calculation needs a statement about each
analytic bound, rather than monotonicity of the optimal constants $C_n$.
Let $N$ be its fixed evaluation index and $n\ge N$ the actual index.
Fix one contrary maximizing marginal and all of its acquired moments,
endpoints and signed phase constraints. Throughout a finite refinement
sequence retain
\begin{equation}\label{j:local:ext:floor-invariant}
 \mathcal R_n\le r,\qquad nD\le Nd.
\end{equation}
Here $\mathcal R_n$ is the complete raw Fourier remainder, not one
selected summand. The smaller inequality $D\le d$ is available for
index-free moment and gap estimates, but does not replace the second
inequality in \eqref{j:local:ext:floor-invariant}.

The index-free operations are intersections of necessary moment
intervals, signwise Cauchy bounds for $\mu=\E X_+$, central-gap
bounds, and physical endpoint prices. For example,
\[
 \mu^2\le p_u(1+\beta_u)/2,\qquad
 \mu^2\le(1-p_l)(1-\beta_l)/2,\qquad
 \mu^2\le\max_{p\in[p_l,p_u]}p(1-p)
\]
do not depend on $n$. If a proved gap inequality gives $6\mu AB\ge P>0$,
then $(A+B)^2\ge2P/(3\mu_u)$, by $4AB\le(A+B)^2$.
Such an improvement requires the already acquired gap and positive mean.
It cannot be used to manufacture a missing initial gap.

For the six-piece high-frequency window bound of Proposition~\ref{j:fourier:high}, keep its admissible
event parameters and window fixed. Its nonconstant index factors are
positive sums and products of
\[
 \sqrt n e^{-cn}\quad(2Nc>1),\qquad n^{-1},\qquad
 e^{-knb^2}/\sqrt n,
\]
and of the decreasing function
\[
 r(x)=\frac{x+2}{x^2+4x+2},\qquad
 r'(x)=-\frac{x^2+4x+6}{(x^2+4x+2)^2}<0.
\]
Thus the ordinary first piece, peak-core correction, peak tails and
zero-peak tail decrease. The normal comparator has the form
$\sqrt n e^{-n/8}\{2r(n/8)+g_c/3\}$, and decreases for $N\ge5$.
The improved first piece has an additional factor $r(n\ell/4)$
with $\ell>0$, and uses its own positive damping and $2Nc>1$ gates.
For the optional improved zero piece, the derivative sign of
$f(x)=\sqrt{x}e^{-x}r(x)$ is the sign of
\[
 4-10x-22x^2-13x^3-2x^4.
\]
This is $-11/128$ at $x=1/4$ and is strictly decreasing for $x>0$.
Hence the optional bound decreases whenever the native gate
$Nkb^2\ge1/4$ holds. All these statements retain the original
window conditions $0<a\le1$, $m_e>0$, $0<d_0\le12$,
$1\le M\le111/10$, $0<w\le1/4$, and positive sine slope.
A minimum of complete valid window bounds stays valid; one need not
assume that the minimizing window remains the same at every index.
Rational numerical upper bounds are taken at the fixed floor $N$.

The low-frequency envelopes have index factors
\[
 (n/(n-1))^{3/2},\quad (n/(n-1))^{3/2}/(n-1),\quad
 a_n^{-j},\quad b_n^{-j},\quad n^{-1/2},
\]
where $a_n=(n-1)/(2n)$, $b_n=(45-M)(1-2/n)/96$ and $j>0$.
The first two decrease and $a_n,b_n$ increase positively on the
retained moment domain. Sums and products with nonnegative fixed
numerators therefore decrease. This applies to the full remainders
$e/(2a_n^2)+B^2/(2b_n^3)$ and $L_2/\sqrt n$.
Every paired-moment eligibility condition remains necessary; the
index comparison cannot repair a failed moment-rectangle hypothesis.

The phase estimate supplies
\[
 \phi(H/\sqrt n)\in[\phi_{0,l}(1-5r/(2N)),2/5]
\]
under its raw-remainder and positivity gates. The actual lower
endpoint increases as $n$ increases. Similarly, for $r,D_l\ge0$,
\[
 r/n-\phi_{0,l}(1-5r/(2n))D_l/6
\]
increases with $1/n$, so evaluation at $N$ is an upper bound.
Put $g_0=(79/25)\rho-3-d$ and $t=(d-4/25)_+$.
The general phase curvature has derivative $(g_0+t)/2\ge0$
with respect to the lower density: if $d\le4/25$, use
$g_0\ge4/25-d$; otherwise $g_0+t=(79/25)(\rho-1)$.
The signed paired curvature is
\[
 C\rho+\tfrac12\min\{\phi_lg_0,(2/5)g_0\}.
\]
It increases with $\phi_l$ for $g_0\ge0$ and is constant in
$\phi_l$ for $g_0<0$. This proves the required direction with
both signs retained. The positive additions and all native
curvature gates are checked before using the phase maximum.

For the sharp contrary score, $C>\CE$ and $\CE<40974/10^5$ give
\[
 n(C-\CE)\rho
 \ge N\max\{0,C-40974/10^5\}\rho_l.
\]
The breakpoint $C=40974/10^5$ is therefore included in the score
cover. On either side the charge is affine in $C$; the same coherent
$(C,\rho)$ corner argument bounds the remaining phase expression.
A charged phase budget $B$ with $nD\le B/\kappa$,
$\kappa=997/15000$, exports $d=B/(\kappa N)$.
A negative budget excludes the law before division. Maximum over a
complete cover, upward rounding and minimum with a previous valid
entry all preserve \eqref{j:local:ext:floor-invariant}. Induction over the
finite saved passes proves the invariant without a limiting argument.
A positive $K/n$ calculation is separate: its $K\rho$ charge must
be retained and is not covered by this sharp-target simplification.

Finally, if the physical affine resource is
\[
 D\ge b+g d_{\rm in}+(N/n)h_N\lambda_n,
 \qquad \lambda_n=ns,\quad b,g,h_N\ge0,
\]
then \eqref{j:local:ext:floor-invariant} gives
\[
 (n/N)(b+g d_{\rm in})+h_N\lambda_n\le d,
 \qquad b+g d_{\rm in}+h_N\lambda_n\le d.
\]
The original joint intercepts therefore bound the actual noise.
An equivalent directed contract records $D\le Nd/n$, noise price
$Nh_N/n$, and actual phase normalization $15000/(997n)$.
The actual index remains in every binomial mass and count probability.
The present applications are $N=25$, $26\le n\le30$, and $N=31$,
$32\le n\le34$, with a separate support calculation for every
actual index and every support state used by the refinements.

\paragraph{Extracting a phase bound from the whole charged envelope.}
The two local B1 refinements at floor $31$ also tighten the signed
phase intervals. Suppose that, for either sign and $h=|H|\ge0$,
the same acquired law satisfies
\[
 0\le D\le c_n\max_j\{A_j+F_jh-Q_jh^2\},\qquad
 c_n>0,\quad Q_j>0\quad\hbox{for every }j.
\]
The maximum is over the complete charged corner family; its members
cannot be selected independently of the law. At least one corner
polynomial is nonnegative at the actual $h$. Its discriminant
$\Delta_j=F_j^2+4Q_jA_j$ is therefore nonnegative, and solving that
quadratic inequality gives
\[
 h\le\frac{F_j+\sqrt{\Delta_j}}{2Q_j}.
\]
Choose a nonnegative rational enclosure $u_j$ with
$u_j^2\ge\Delta_j$ whenever $\Delta_j\ge0$, and set
\[
 r_j=\begin{cases}
 0,&\Delta_j<0,\\
 \max\{0,(F_j+u_j)/(2Q_j)\},&\Delta_j\ge0.
 \end{cases}
\]
A negative discriminant cannot supply the nonnegative corner, so
$h\le\max_j r_j$. Intersect this bound with the previously acquired
bound for the same sign. The construction is valid at a double root
and at $h=0$. Its checker requires every $Q_j>0$; it has no unchecked
linear or constant branch.

The fixed floor-$31$ normalization is $15000/(997\cdot31)$.
Under the weighted transfer, the actual normalization becomes
$c_n=15000/(997n)$, which stays positive. The implication from
$D\ge0$ to the nonnegativity of at least one charged corner is
therefore unchanged. Separate maxima of the roots of the minus and
plus families give the two signed phase caps used by the actual-index
count contracts.

\paragraph{Transporting the physical variance refinement.}
One resource step uses the bounded conditional supports to obtain
$s\le B\delta$, where $s$ is the rescaled conditional variance,
$\delta=d_{\rm in}$, and
\[
 B=\max(O_0,O_1)\max(p_u,1-p_l)>0.
\]
Here $O_0,O_1$ are outward support offsets in the same conditional
coordinates. If the acquired joint plane is
$\delta/d_0+\lambda_n/L\le1$, with $\lambda_n=ns$, then
\[
 \lambda_n\le \frac{1}{L^{-1}+(nBd_0)^{-1}}.
\]
The physical variance parameter is
$\theta_n=\lambda_n/(n p(1-p)+\lambda_n)$. Writing
\[
 v_l=\min_{p\in[p_l,p_u]}p(1-p)>0
\]
gives
\[
 \theta_n\le
 \frac{1}{1+n v_l/L+v_l/(Bd_0)}
 \le \frac{1}{1+N v_l/L+v_l/(Bd_0)}
 \quad(n\ge N).
\]
The last quantity is precisely the floor calculation's variance
refinement. Downward rounding of the support scale enlarges $O_0,O_1$;
upward rounding of this last bound and intersection with the previous
physical $\theta$ bound preserve its direction. This transports the
bound for $\theta$, rather than treating the floor's diagnostic
$\lambda_N$ as an actual-index bound for $\lambda_n$.
The original joint intercept $L$ remains the one used in the affine
plane. For the selected floor-$25$ first-band resource there are $193$
such saved refinements, each checked at every actual index $26$--$30$.


Here is the version of the moment matrices that permits different
positive and negative support radii. On either sign write $Y=|X|$,
$0\le Y\le r$, and let $p_s,m,a,b,c_4,c_5$ be its unnormalized
moments of orders zero through five. Positivity of the weights
$r-Y$, $Y(r-Y)$, and $Y^2(r-Y)$ gives
\[
 \begin{pmatrix}rp_s-m&rm-a\\rm-a&ra-b\end{pmatrix},\quad
 \begin{pmatrix}rm-a&ra-b\\ra-b&rb-c_4\end{pmatrix},\quad
 \begin{pmatrix}ra-b&rb-c_4\\rb-c_4&rc_4-c_5\end{pmatrix}
 \succeq0.
\]
For positive denominators their determinant inequalities imply
\begin{align*}
 p_s&\ge m/r+(rm-a)^2/[r(ra-b)],\\
 c_4&\le rb-(ra-b)^2/(rm-a),\\
 c_5&\le rc_4-(rb-c_4)^2/(ra-b).
\end{align*}
A zero diagonal entry forces the corresponding off-diagonal entry to
vanish. One may check that condition or discard the nonnegative
correction, but one cannot divide by zero. In interval form a correction
is bounded below by the square of a positive-part lower numerator
divided by an upper positive denominator. The sign radius $r$ is used
only for its own class. Summing the first inequality over the two signs
and using $p_++p_-\le1$ can improve the common upper mean.

One further bound used at small indices strengthens the signed constant
coefficient. Set $\delta=\gamma-\rho\beta$ and $Y=|X|$. Then
\[
 V:=\E\{Y(Y^2-\rho Y)^2\}
   =\E|X|^5-2\rho\mu_4+\rho^3,
 \qquad
 \delta=\E\{Y\operatorname{sgn}(X)(Y^2-\rho Y)\}.
\]
Here $m>0$ by centering and unit variance.
Cauchy's inequality for the nonnegative measure $Y\,dP$ gives
$\delta^2\le2mV$. Consequently
\begin{equation}\label{j:local:ext:weighted-cubic-loss}
 \frac\rho2\mu_4-\frac14\E|X|^5
 \le\frac{\rho^3}{4}-\frac{(\gamma-\rho\beta)^2}{8m}.
\end{equation}
This deduction can be combined with the two other deductions in
Section~\ref{j:local:atlas} by taking their maximum. Their sum would spend
the same nonnegative quantity more than once.



\paragraph{A universal lower bound for each sign probability.}
\label{j:local:ext:sign-mass}
Six probability panels at index six are excluded before the count
calculation. The exclusion is the following elementary moment inequality;
it does not require support, a central gap, or absence of a zero atom.

\begin{lemma}[Sign mass and third moment]\label{j:local:sign-mass}
If $\E X=0$, $\E X^2=1$, and $\rho=\E|X|^3<\infty$, then
\[
 \rho^2\Pp(X>0),\quad \rho^2\Pp(X<0)
 \ \ge\ \frac{8(7\sqrt7-17)}{27}
 \ >\ \frac{11261}{25000}.
\]
\end{lemma}
\begin{proof}
Write $p=\Pp(X>0)$, $a=\E X_+^2$, and
$m=\E X_+=\E X_->0$. Centering and unit variance imply
$0<a<1$. Cauchy's inequality, first with factors
$X_+^{1/2}$ and $X_+^{3/2}$ and then with their negative-sign
counterparts, gives
\[
 a^2\le m\E X_+^3,\qquad (1-a)^2\le m\E X_-^3.
\]
Also $m^2\le pa$. Consequently
\[
 p\rho^2\ge f(a):=\frac{(a^2+(1-a)^2)^2}{a}.
\]
The derivative has the sign of $6a^2-2a-1$, because
\[
 f'(a)=\frac{(2a^2-2a+1)(6a^2-2a-1)}{a^2}.
\]
Thus the unique minimum on $(0,1)$ is at
$a=(1+\sqrt7)/6$, where
$f(a)=8(7\sqrt7-17)/27$. The strict rational comparison follows
from $2645751^2<7\cdot10^{12}$ and
$8(7\cdot2645751/10^6-17)/27>11261/25000$.
Apply the same argument to $-X$ for the other sign probability.
\end{proof}



\paragraph{Scalar consequences of the common inward budget.}
\label{j:local:ext:inward-scalar}
Some finite count domains use a stronger absolute-third-moment floor or
noise ceiling before their five-dimensional subdivision. The following
calculation supplies those scalar restrictions. It uses the sign
conditioning of \eqref{j:core:conditional-moments}.

Reflect only for the scalar calculation, so that
$0<p\le1/2$, $q=1-p$, $v=pq$. This preserves $\rho$, total
conditional variance $s$ and total inward displacement. The count
probability, its phase and its original orientation are not reflected.
Suppose the inward endpoints obey $a_0+a_1\le d$ with $0<d\le p$.
The conditional third-moment inequalities and Cauchy's inequality give
\begin{equation}\label{j:local:ext:inward-I}
 \rho\ge I(p,s,d):=
 \frac{v(1-2v)+(3p-d)s+s^2/(qd)}{(v+s)^{3/2}}.
\end{equation}
Indeed, writing $b=qs_0$, $c=ps_1$ and $b+c=s$, the absolute-cubic
numerator is at least
\[
 v(1-2v)+3pb+3qc-a_0b-a_1c
                       +\frac{b^2}{qa_0}+\frac{c^2}{pa_1}.
\]
Its linear part is at least $(3p-d)s$, while its quadratic part is
at least $(b+c)^2/(q(a_0+a_1))\ge s^2/(qd)$.
Vanishing inward endpoints force the corresponding centered variance
to vanish; those terms are handled before division. If $d=0$, then
$s=0$ directly.

The derivative has the exact numerator
\begin{equation}\label{j:local:ext:inward-I-derivative}
 (v+s)^{5/2}I_s
 =v(6p-3p^2-d-3/2)
   +\left\{\frac{2p}{d}-\frac{3p-d}{2}\right\}s
   +\frac{s^2}{2qd}.
\end{equation}
Consequently, on $p\in[l,u]\subset(0,1/2]$, the two gates
$d\le l$ and $f:=6l-3l^2-d-3/2\ge0$ imply
\[
 \rho\ge\rho_B(p):=\frac{1-2p(1-p)}{\sqrt{p(1-p)}}
                  \ge\rho_B(u).
\]
The last function is decreasing on this half-interval. If a previously
acquired ceiling is $\rho\le R$, a point with
\[
 R^2p(1-p)-(1-2p(1-p))^2<0
\]
is impossible. This polynomial increases on $0<p\le1/2$, because
its derivative is $(1-2p)(R^2+4-8p(1-p))\ge0$.
Exact bisection may therefore raise the closed lower probability
endpoint to a known infeasible point; retaining that endpoint is a
conservative enlargement.

When $f>0$ and $s\le S$, formula~\eqref{j:local:ext:inward-I-derivative}
also gives
\[
 I_s\ge
 \frac{l(1-l)f}{[u(1-u)+S]^{5/2}}.
\]
Integration from zero supplies a new variance ceiling. Directed square
roots and outward division make every finite update valid separately.

A useful weaker floor survives failure of the monotonicity gate.
Assume still $0<d\le l$ and set
\[
 a=(3/2-6l+3l^2+d)_+,
 \qquad B=2l/d-3u/2\ge5/4.
\]
The only possible decrease of $I$ is bounded below by
$-(va-Bs)_+/v^{5/2}$. Integrating this negative triangle gives
\begin{equation}\label{j:local:ext:inward-deficit}
 \rho\ge\rho_B(u)-\frac{a^2}{2B\sqrt{l(1-l)}}.
\end{equation}
In particular, failure of the monotonicity gate must not be replaced
by the unsupported assertion $\rho\ge\rho_B(p)$.

There is a second variance contraction which works with either sign
of $f$. The global Taylor upper bound for $(v+s)^{3/2}$, combined
with \eqref{j:local:ext:inward-I}, gives the necessary inequality
\begin{equation}\label{j:local:ext:scalar-noise-quadratic}
 Gs^2+fs\le\epsilon\,[u(1-u)+S]^{3/2},
\end{equation}
where
\[
 \epsilon=(R-\rho_B(u))_+,\qquad
 G=\frac1{(1-l)d}-\frac{3(1-2l(1-l))}{8l(1-l)}
       \ge\frac5{8l(1-l)}+\frac34>0.
\]
To see this, subtract $\rho_B(p)(v+s)^{3/2}$ from the lower
numerator in \eqref{j:local:ext:inward-I}, use the quadratic Taylor upper
bound on the subtracted term, and then use $p\ge l$ and
$\rho_B(p)\ge\rho_B(u)$. The positive part in $\epsilon$ is
necessary when $f<0$. The larger real root of
\eqref{j:local:ext:scalar-noise-quadratic}, enclosed from above, is a valid
new upper bound for $s$, even for negative $f$.

Finally suppose the original joint resource has the stronger form
\[
 d_{\rm actual}\le d_0(1-s/S_0),\qquad S_0=\Lambda/n.
\]
To test a proposed upper endpoint $a\in(0,S_0)$, assume $s\ge a$.
Then $d_{\rm actual}\le d_0(1-a/S_0)$, so the preceding quadratic
calculation can be repeated with this smaller inward bound. If its
resulting upper bound is strictly less than $a$, the whole region
$s\ge a$ is impossible. Otherwise the probe proves no lower bound
on $s$. A finite sequence of such hypothetical tests is sound
without any convergence assertion. Throughout it, the original
$S_0$ remains in the joint budget; a separately improved noise
ceiling clips that budget and does not replace its intercept.


\subsection{Paired Fourier cancellation}
\label{j:local:paired}
Use the common Taylor and power estimates of Proposition~\ref{j:fourier:low}.
Let $a=(3-\mu_4)/24$, $z=(10\gamma-\E X^5)/120$, $c=a\gamma/6$,
$|a|\le a_*$, $|z|\le Z$, and take a rectangle containing the actual
$(z,c)$. Put $I_u\ge\sqrt{\pi/2}$ and
$F_N\ge(N/(N-1))^{3/2}$. The common error remains
\[
 L_2=E_6/(2a_0^2)+B_6^2/(2b_0^3),\quad
 a_0=(N-1)/(2N),\quad0<b_0\le(45-M)(1-2/N)/96,
\]
where $E_6=E$, $B_6=B$ are the Taylor-error constants of Proposition~\ref{j:fourier:low}.
For $J(z,c)=\int_0^\infty e^{-x^2/2}|zx^2-cx^4|\,dx$ set
\[
 U(z,c)=\begin{cases}I_u(|z|+3|c|),&zc\le0,\\
 I_u\bigl||z|-3|c|\bigr|+6\sqrt{|zc|}e^{-|z|/(2|c|)},&zc>0,
 \end{cases}\qquad J_u=\max_{\rm corners}U(z,c).
\]
A complete admissible low-frequency triple is
\begin{align}\label{j:local:fourier:paired}
 L_0&=F_NJ_u+\frac{F_NI_u}{N-1}
 (35a_*g^3/48+35Zg^2/24)+251g^3/17280,\\
 L_H&=F_NI_u\{a_*+5Zg/[2(N-1)]\}+251g^2/4800.\nonumber
\end{align}
To verify it, write $b=n-1$, $v=b\gamma t^3/6$. After division by $t^3$
and rotation by $Ht$, the leading imaginary part is bounded by
\[
 |zt^2-atv|+a_*t|\sin v-v|+Zt^2|\cos v-1|
 +|H|t\{a_*t+Zt^2|\sin v|\}.
\]
Multiply by $ne^{-bt^2/2}$ and integrate. Gaussian moments give the four
coefficients $35/48,35/24,1,5/2$ above; the leading integral is
$nJ(z,c)/b^{3/2}$. The exact norm $J$ is convex and invariant under
$(z,c)\mapsto(-z,-c)$, so its rectangle maximum is at a corner.
For $zc\le0$, its value is $\sqrt{\pi/2}(|z|+3|c|)$.
For $z,c>0$, splitting at $r=\sqrt{z/c}$ and integrating by parts gives
\[
 J=\sqrt{\pi/2}(z-3c)+6cr e^{-r^2/2}
       +2(3c-z)\int_r^\infty e^{-x^2/2}\,dx.
\]
Discard a nonpositive last term and otherwise bound its integral by
$\sqrt{\pi/2}$ to obtain $U$. The common remainder and comparator terms
are unchanged. All inverse-index factors decrease with $n$ on their
stated positive domains. No sign of the actual contact moment $\gamma$
is changed by this auxiliary reflection.


\begin{lemma}[Retaining the signed Gaussian tail]\label{j:local:entry189:mills}
For $r>0$, put
\[
 \mathcal M(r)=e^{r^2/2}\int_r^\infty e^{-t^2/2}\,dt,
 \qquad
 \mathcal M_k(r)=\cfrac1{r+\cfrac1{r+\cfrac2{\ddots+\cfrac{k}{r}}}}
 \quad(k\ge1),\qquad \mathcal M_0(r)=1/r.
\]
Odd truncations are lower bounds for $\mathcal M$; even truncations
are upper bounds. For $z,c>0$, set $x=z/c$ and
\[
 d_{k+1}=1,\qquad d_j=1+\frac j{xd_{j+1}}\quad(j=k,\ldots,1),
 \qquad m_k(x)=d_1^{-1}.
\]
Then the cancellation norm in Section~\ref{j:local:paired}
satisfies
\begin{equation}\label{j:local:entry189:mills-bound}
 J(z,c)\le \sqrt{\pi/2}(z-3c)+\sqrt{zc}\,e^{-x/2}
 \left(6+\frac{2(3-x)}x m_k(x)\right),
\end{equation}
where $k=11$ if $x\ge3$ and $k=12$ if $x<3$.
The factor in parentheses is positive. One may take the minimum with
the previous bound $U(z,c)$ at each corner of the coefficient rectangle.
\end{lemma}
\begin{proof}
The positive integrals
$I_j(r)=\int_0^\infty t^j e^{-rt-t^2/2}\,dt$ satisfy
$I_1=1-rI_0$ and $I_{j+1}=jI_{j-1}-rI_j$ by integration by parts.
Thus, writing $q_j=I_j/I_{j-1}>0$,
\[
 I_0=\frac1{r+q_1},\qquad q_j=\frac j{r+q_{j+1}}.
\]
Replacing the final positive $q_{k+1}$ by zero reverses the inequality
at each substitution. This proves the alternating finite bounds, since
translation gives $I_0=\mathcal M(r)$. No convergence of an infinite
continued fraction is required. The displayed rational recursion gives
$m_k(x)=r\mathcal M_k(r)$ when $r=\sqrt x$.

Splitting the norm at $r=\sqrt{z/c}$, as in the preceding proof, gives
the exact identity
\[
 J(z,c)=\sqrt{\pi/2}(z-3c)+\sqrt{zc}\,e^{-x/2}
 \left(6+\frac{2(3-x)}x r\mathcal M(r)\right).
\]
The multiplier of $\mathcal M$ is nonpositive when $x\ge3$ and
positive when $x<3$. Substitute the corresponding finite lower or
upper bound. Positivity of the remaining factor is immediate for
$x<3$; for $x\ge3$, use $0<m_{11}(x)<1$ to bound it below by
$6-2(x-3)/x>4$. Finally the exact norm $J$, rather than either
upper-bound formula, is convex, so its containing rectangle is still
exhausted by its four corners.
\end{proof}

At every positive coefficient corner use the displayed finite Mills bound.
A negative Gaussian coefficient uses a lower enclosure for
$\sqrt{\pi/2}$; a nonnegative coefficient uses an upper one. The remaining
factor is positive, so outward upper square-root and exponential enclosures
are valid. A simultaneous reflection of both odd coefficients leaves the
norm unchanged.

\paragraph{The two coefficient choices in the certificates.}
All coefficients come from the actual marginal. For $Q=X^2-\gamma X-1$,
multiplication gives
\begin{align}\label{j:local:fourier:pearson}
 \E Q^2&=\mu_4-1-\gamma^2,&
 \E X^5&=\gamma^3+2\gamma+\E[Q^2(X+2\gamma)],\notag\\
 \mu_6&=1+3\gamma^2+\gamma^4+
       \E[Q^2\{(X+\gamma)^2+2\gamma^2+2\}].
\end{align}
For any radius $R_0\ge|X|$ with $R_0\ge2g$, the two signs of
$\E X^5-10\gamma$ are bounded, at $t=|\gamma|$, by
\begin{align*}
 P_{R_0}(t)&=R_0(M-1)+(10-2M)t-R_0t^2+t^3,\\
 Q_{R_0}(t)&=R_0(M-1)+(2M-10)t-R_0t^2-t^3.
\end{align*}
The gate $R_0-2g\ge0$ permits substituting $M$ for $\mu_4$ in both
signs. The triangle choice uses the valid enclosing radius $R_0=6$, $a_* =\max(2,M-3)/24$,
$M_6=RE_5$, $b_0=(45-M)(1-2/N)/96$ and
\begin{equation}\label{j:local:fourier:z0}
\begin{gathered}
 120Z=\min\{E_5+10g,Z_0\},\\
 Z_0=\begin{cases}
 P_6(g),&M\le5, P'_6(g)\ge0,\\
 6(M-1)+(10-2M)^2/[4(6-g)],&M\le5, P'_6(g)<0,\\
 6(M-1)+(2M-10)^2/24,&M>5.
 \end{cases}
\end{gathered}
\end{equation}
If $M\le5$, $P_6\ge Q_6$ and $P_6''<0$ on $[0,g]$; use the
increasing endpoint or $t^3\le gt^2$ and complete the square.
If $M>5$, $P_6\le6(M-1)$ and the upper quadratic for $Q_6$ gives
the third row. One may also use the actual radius by setting
$r_0=\max(R,2g)>g$ and
\[
 Z_r=\begin{cases}
 r_0(M-1)+(10-2M)t_*-(r_0-g)t_*^2,
 & M\le5,\\
 & t_*=\min\{g,(10-2M)/[2(r_0-g)]\},\\
 r_0(M-1)+(2M-10)t_*-r_0t_*^2,
 & M>5,\\
 & t_*=\min\{g,(2M-10)/(2r_0)\}.
 \end{cases}
\]
Indeed for $M\le5$, $P_{r_0}\ge Q_{r_0}$ and $t^3\le gt^2$;
for $M>5$, $P_{r_0}\le r_0(M-1)$ and dropping $-t^3$ bounds
$Q_{r_0}$. Each displayed quadratic is maximized at its clipped vertex.
Thus the selected implementation takes
$120Z=\min(E_5+10g,Z_0,Z_r)$.
It uses $c_0\ge\sqrt{\pi/(16a_0^3)}=\lambda_N$ and $c_1=c_0/2$.
The weaker choice $c_1\ge\sqrt{\pi/(256a_0^5)}
=\lambda_N N/[2(N-1)]$ is also permitted by the same transfer lemma.

The positive paired rectangle below applies on the full domain $M<3$.
Its selected $M\le2$, $g\le1$ specialization may use the weaker
$b_0=43(1-2/N)/96$. The sixth-moment and signed-coefficient bounds
are otherwise the same. Before applying it, the lower bound on $\rho$
must also remain valid when the conditional variance is positive.

\begin{proposition}[The actual absolute third moment]
Suppose $0<l\le p\le u<1/2$, $0\le\theta\le\Theta<1$, and the actual
conditional cubic identity gives
\[
 \rho=\E|X|^3\ge(1-\theta)^{3/2}\rho_B(p)+3a\theta,
 \qquad a\ge0,\quad
 \rho_B(p)=\frac{1-2p(1-p)}{\sqrt{p(1-p)}}.
\]
Here $a$ can be the minimum of the two acquired lower support endpoints.
With outward bounds $r_-\le\rho_B(u)$ and $r_+\ge\rho_B(l)$, a valid
lower bound is
\[
 r_*:=\max\{1,r_--(3r_+/2-3a)_+\Theta\}\le\rho.
\]
\end{proposition}
\begin{proof}
The function $\rho_B$ decreases on $(0,1/2)$ and convexity gives
$(1-\theta)^{3/2}\ge1-3\theta/2$. Therefore
\[
 \rho\ge\rho_B(p)-(3\rho_B(p)/2-3a)\theta
 \ge r_--(3r_+/2-3a)_+\Theta.
\]
The bound $\rho\ge1$ follows from $\E X^2=1$ and monotonicity of the
absolute moment norms. No separate requirement $3a\ge3r_+/2$ is needed.
\end{proof}

\begin{proposition}[A positive cancellation rectangle]\label{j:local:paired-rectangle}
Let $\E X=0$, $\E X^2=1$, $|X|\le R$ almost surely, and let
\[
 1\le\mu_4\le M<3,\qquad
 0\le\gamma_l\le|\gamma|\le g,\qquad g^2<2,\quad R>2g,
 \quad M-1-\gamma_l^2\ge0.
\]
The paired Fourier formulas are valid with damping constant
\[
                 b_0=\frac{45-M}{96}(1-2/N),\qquad N>2.
\]
In particular, after simultaneous reflection of both odd Fourier
coefficients, the valid rectangle for
$a_0=(3-\mu_4)/24$, $z=(10\gamma-\E X^5)/120$,
$c=a_0\gamma/6$ is
\begin{gather*}
 a_*=(2-\gamma_l^2)/24,\qquad
 c_l=(3-M)\gamma_l/144,\quad
 c_u=(2-\gamma_l^2)g/144,\\
 t_* =\min\{g,(10-2M)/(2(R-g))\},\\
 120Z=\min\{E_5+10g,R(M-1)+(10-2M)t_*-(R-g)t_*^2\},\\
 z_l=\max\{-Z,[8\gamma_l-\gamma_l^3
        -(R+2\gamma_l)(M-1-\gamma_l^2)]/120\},\quad z_u=Z,
\end{gather*}
where $E_5\ge\E|X|^5$. If $3g^2\le2$, $c_u$ may also be decreased
to $\min(c_u,(2g-g^3)/144)$.
The same sixth-moment bound remains valid:
\[
 M_6=\min\{RE_5,1+3g^2+g^4+
 (R^2+2Rg+3g^2+2)(M-1-\gamma_l^2)\}.
\]
\end{proposition}
\begin{proof}
Pearson's identity gives $\mu_4\ge1+\gamma^2$, so $0\le a_0\le a_*$.
Both stated bounds on $c$ follow by multiplying nonnegative quantities.
The optional refinement uses the monotonicity of $2t-t^3$ on
$[0,g]$ when $3g^2\le2$.

The identities \eqref{j:local:fourier:pearson} give the upper absolute
fifth-coefficient bound by maximizing
\[
 R(M-1)+(10-2M)t-(R-g)t^2,\qquad0\le t\le g.
\]
The larger of the two Pearson sign bounds is still this one: $M<3<5$,
and $R>2g$ permits replacing $\mu_4$ by its upper bound. The quadratic
has its maximum at the displayed clipped vertex. The numerator in the
lower bound for $z$ has derivative
$10-2M+2Rt+3t^2>0$, so its lower endpoint is obtained at $\gamma_l$.
The last identity in \eqref{j:local:fourier:pearson} gives $M_6$.

The signed Gaussian transfer is applicable because $M<3<11.1$ and
$g<\sqrt2<1.85$. Its norm is convex in $(z,c)$, so the rectangle is
checked at its four corners; all Gaussian and Mills estimates remain
unchanged. Its damping hypothesis is
$b_0\le(45-M)(1-2/N)/96$. The older value
$43(1-2/N)/96$ is unavailable when $M>2$ and must be replaced as stated.
\end{proof}

\paragraph{A necessary implementation seam.}
The same new $r_*$ must be used in the four Fourier corners, the signed
phase-curvature correction, all conserved budgets, and exported quadratic
resources. The latter correction remains
\[
 q(C,\rho)=C\rho+\tfrac12\min\{\phi_{\rm low}g_0,.4g_0\},
 \qquad g_0=3.16\rho-3-d.
\]
Recomputing the old $\rho_B$ lower endpoint in any of these locations would
invalidate the rectangle association. The new adapter returns $r_*$ as
part of its paired result and the corrected callback uses that value.



For the signed fourth-moment bounds, use the same $r_*$ in
\eqref{j:local:entry201:Kexact}. The monotonicity of $F$ in its third argument and its already proved rectangle extrema license this substitution without assuming $\rho\ge\rho_B$. Retain $1-\Theta>2/3$ when using
the separate binary-$K$ monotonicity bound.


\subsection{Confinement from the remaining defect budget}
Two different support implications are useful. Gaussian contact excludes
a point by a strict negative upper bound on its contact expression;
closed cells cover the finite exterior and a derivative-sign certificate
covers any remaining ray. The shared Gaussian support theorem supplies
this callback, including the bounded-threshold radius and coordinate-strip
contracts; it is a prior theorem, not a test on moments alone. Noise confinement proves $|U|\le a$ on an
already acquired row. It then implies, without a small-noise limit,
\begin{equation}\label{j:local:re:radiusfeedback}
 |X|\le R_*:=
 \frac{\max\{u+a,1-l+a\}}{\sqrt{\min_Ip(1-p)}}.
\end{equation}
The denominator uses $\E U^2\ge0$, so zero noise is included.
The Fourier/moment formulas remain valid for $R\ge1$; the paired Pearson
choice additionally checks $R>2g$. In particular there is no intrinsic
lower radius $2.3$: the radial ratios in Lemma~\ref{j:local:re:pricing} increase
on all $x\ge0$, and the other formulas use their explicit domains.

Here is the selected proof schedule. Each row names \emph{an application}
of Theorem~\ref{j:local:acquisition}, not a different analytic theorem.
\begin{center}\small
\begin{tabular}{p{.13\linewidth}p{.32\linewidth}p{.45\linewidth}}
\toprule
Node&Input and operation&Output consumed later\\\midrule
I&$N=47$, $R=3.84$, $C\le.4247$; 13 signed passes plus export&
100 closed probability cells; 35 survive; $\mathcal R<1.885$, $nD<47(.590)$\\
C&Central preparation, contact radius $2.5$, probability refinement&
45 cells $[.38,.47]$, width $.002$; canonical simultaneous input\\
L&Complete complementary cover, contact radius $3$, refinement&
39 cells $[.302,.38]$; all other complementary cells excluded\\
C1&First confinement from C; radius formula; enclosure&
First improved central table\\
C2&Antichain confinement from C; intersect C1; enclosure&
Selected central table: 45 cells, $[.38,.47]$\\
S&Separate sign-resource pricing on C2 and L&
One left cell excluded; 38 left cells remain\\
L2&Left contact radius $2.7$ from S; enclosure&
Selected left table: 38 cells, $[.304,.38]$\\\bottomrule
\end{tabular}
\end{center}
Every shared endpoint belongs to its adjacent closed cells. The
confinement producing C2 consumes C, not C2. The selected central local
comparison uses C2; it does not need the optional further central
refinement in S. The final cover is exactly $[.304,.47]$.

\begin{corollary}[Uniform acquired entry]\label{j:local:re:acquired}
Let an attained positive maximum have $n\ge47$, $C_E\le C\le.4247$,
$0\le n(C-C_E)\le.7$, decrement at most $.19$, and both deleted
score caps at most $.4253$.
Given the shared contact support and strip implications, the exact
schedule above places its law in one of the 38 selected left rows or
45 selected central rows, with every resource in that row holding
simultaneously. This assertion has no upper bound on $n$.
\end{corollary}
\begin{proof}
The shared all-threshold support theorem supplies radius $3.84$.
Apply the finite enclosure theorem in the displayed order, retaining
the weighted bounds at every edge. The initialization's stored upper
index 98 occurs only through $.4095+K/98$; here $K=0$, so it imposes
no upper-index restriction. All subsequent index-dependent inequalities
are checked for every $n\ge47$ by their stated monotonicity gates.
The accepted closed covers exclude the complementary cells and give
the stated 83 rows.
\end{proof}

\paragraph{The local-comparison interface.}
For each final row supply $p\in I$, $\theta\le\Theta$, $\lambda\le\Lambda$,
the actual support/amplitude, $d_Y\le d_0$, both signed phase limits,
and all four coherent quadratics bounding \eqref{j:local:budget}.
On a phase/noise subcell, let $T_H$ be the maximum of each concave
quadratic at its clipped vertex, then the maximum over corners and the
minimum with the scalar bound. If a proved moment inequality is
\[
 D\ge\alpha d_Y+b\theta+P,\qquad\alpha>0,
 \quad \ell\ge\ell_-,\quad\lambda\ge L_0,
\]
then the exact strict resource supplied to the comparison is
\begin{equation}\label{j:local:re:localG}
 nd_Y<G:=\frac{T_H/\kappa-NP-b\ell_-^2L_0}{\alpha}.
\end{equation}
This uses $n\theta=\ell^2\lambda$ and $nP\ge NP$. If $G\le0$,
the entire closed subcell is excluded. Formula~\eqref{j:local:re:localG} is
the same resource as \eqref{j:local:budget}, not an additional allowance.

\paragraph{One confinement lemma with two concentration choices.}\label{j:local:confinement}
A proved amplitude is used as a support update, never as a premise of
its own proof. The following derivation applies to every acquired row,
with the same formulas for the central and left probability ranges.
\paragraph{Statement and hypotheses.}
Let $n\ge N=47$, $Y=B+U$, $\Pp(B=1)=p<1/2$ and $\E(U\mid B)=0$.
Put $q=1-p$, $s=\E U^2$, $\lambda=ns$ and $X=\ell(Y-p)$, where
$\E X^2=1$. Thus $\ell^2=(1-\theta)/(pq)$ and $\theta=\ell^2s$.
The conditional label-zero support is at most $a\ge0$ and the label-one
support is at least $b\le1$, where $a,b$ are support points and
$b-a=1-d$. In particular $0\le a\le d$, $1-d\le b\le1$.
Assume the initial strip $\operatorname{supp}Y\subset(-.55,1.55)$.
Write $\rho=\E|X|^3$, $\beta=\E[X|X|]$ and $\epsilon=\E|U|$.

Select an acquired closed probability panel $[l,h]$.
Its fields give upper bounds $\Theta,d_0,\rho_0,\beta_+,T,\Lambda$,
a lower bound $\beta_-$, and separate bounds $H_+,H_-$ for the two signs of
$H$. Assume
\[
 \theta\le\Theta,\quad d\le d_0,\quad \rho\le\rho_0,
 \quad \beta_-\le\beta\le\beta_+,\quad 0\le\lambda\le\Lambda,
 \quad .4<C\le C_0=.4253.
\]
For the positive sign $0\le H\le H_+$; for the negative sign
$-H_-\le H\le0$. Let $x=np+H/\ell$ and
$\nu=\mathcal L(x-\sum_{i<n}Y_i)$. Assume a contact majorant $Q$ for
$g(y)=\nu([y,\infty))$, with equality at each support point of $Y$,
and derivative
\begin{equation}\label{j:local:conf:derivative}
 \sqrt n(-Q'(y))=\ell\phi(H/\sqrt n)
 +\frac\ell n\{3C\beta+(H\phi(H/\sqrt n)+3C\rho)t-3Ct|t|\},
 \quad t=\ell(y-p).
\end{equation}
Both actual deleted convolutions satisfy, separately,
\begin{equation}\label{j:local:conf:concentration}
 Q_j(w)\le\frac{\ell\phi(0)w+2C_0\rho}{\sqrt j},
 \qquad j=n-1,n-2,
\end{equation}
for every interval of width $w\ge0$.

Here is the precise common resource assumption. Put
$\kappa=997/15000<\phi(0)/6$, $J=\int_{[a,b)}(y-a)\nu(dy)$,
and note that the central gap contributes $m\ell J=(1-\theta)J$ to $L$.
Thus $L\ge(1-\theta)J$; retain $e\ge0$ and
\begin{equation}\label{j:local:conf:budget}
 n e\rho+n\phi(0)D/6+n^{3/2}L<\mathcal B(H),\qquad
 \mathcal B(H)=\min\{T,\max_j(A_j+F_{j,\pm}|H|-Q_jH^2)\}.
\end{equation}
The four quadratics are the table's \texttt{quadratics} rows, using
\texttt{Qbudget}, not \texttt{Qraw}. The Gaussian phase contribution
has already been included in \texttt{Qbudget}.

For each panel let $E,F_\ell,m_0$ be respectively strict lower, upper,
and lower rational bounds for
\[
 \sqrt{(1-\Theta)/(h(1-h))},\qquad
 1/\sqrt{l(1-l)},\qquad \sqrt{(1-\Theta)l(1-l)}.
\]
The bounds use the grid $10^{-6}$, with an extra unit outward. Define
\begin{align*}
 c_+&=2(5.16)E^2(1-h-d_0)^2,&
 c_-&=2(7.16)E^2(l-d_0)^2,\\
 b_*&=\min\{3(4.16)E(1-h-d_0),3(6.16)E(l-d_0)\},\\
 P&=m_0E^2\pos{3.5\operatorname{dist}([l,h],[.41886,.41887])-2.041d_0}^2.
\end{align*}
The simultaneous moment inequality has the following usable form:
\begin{equation}\label{j:local:conf:moment}
 D\ge c_+M_++c_-M_-+b_*\theta+P,\quad
 d=(M_+/p+M_-/q)/\ell,\quad
 \epsilon\le2(M_++M_-)/\ell,\quad M_\pm\ge0.
\end{equation}
To derive it, the endpoint expansion of \eqref{j:core:gap-defect} gives
$D\ge m(\sqrt{\sigma+1}A-\sqrt{\sigma-1}B)^2
+2(\sigma+2)B^2M_++2(\sigma+4)A^2M_-
+3(\sigma+1)B\theta_++3(\sigma+3)A\theta_-$.
Insert $A\ge E(l-d_0)$, $B\ge E(1-h-d_0)$ and the rational lower
coefficients displayed above. For the square, write
$A=\ell(p-a)$, $B=\ell(1-p-(1-b))$, with
$a+(1-b)=d\le d_0$. Its unshifted coefficient is
$(\sqrt{\sigma+1}+\sqrt{\sigma-1})(p-p_*)$;
$\sqrt{\sigma+1}+\sqrt{\sigma-1}>3.5$ and
$\sqrt{\sigma+1}<2.041$ prove the stated lower square $P$.
For completeness, on the positive class put $Z=X-B\ge0$.
Then $T_+=2B^2M_++3B\E[Z^2;X>0]+\E[Z^3;X>0]$ and
$R_+=BM_++\E[Z^2;X>0]$. Consequently
$(\sigma-1)T_++6BR_+\ge2(\sigma+2)B^2M_+
+3(\sigma+1)B\theta_+$, since $\E[Z^2;X>0]\ge\theta_+$.
The reflected expansion gives the negative-class bound.
The centered absolute deviation on each class is at most twice its
mean distance from the inner endpoint, proving the bound on $\epsilon$.
All quantities refer to the same marginal.

\paragraph{Concentration from an antichain.}
\begin{lemma}[Strict-gap concentration]
Suppose $Y=U_0+B(1+U_1-U_0)$, where $B$ is Bernoulli$(p)$ independent
of the pair $(U_0,U_1)$, and $1+U_1-U_0\ge c>0$ almost surely.
For independent copies $Y_1,\ldots,Y_j$ and any interval $I$ of diameter
strictly less than $c$,
\[
 \Pp\left(\sum_{i=1}^jY_i\in I\right)
 \le \max_{0\le k\le j}\binom jk p^k(1-p)^{j-k}.
\]
The same bound holds after adding any independent random variable.
\end{lemma}
\begin{proof}
Condition on all noise pairs. The sum is a fixed constant plus
$\sum B_iw_i$, with every $w_i\ge c$. The family $\mathcal F$ of subsets
whose weighted sums lie in $I$ is an antichain: distinct comparable
subsets have sums differing by at least $c$, larger than the diameter.
A uniformly random permutation of $\{1,\ldots,j\}$ defines the maximal
chain of its initial subsets. The chain contains a fixed $k$-subset with
probability $1/\binom jk$, and meets an antichain at most once. Hence
\[
 \sum_{S\in\mathcal F}\binom j{|S|}^{-1}\le1.
\]
Multiplying each summand by its binomial level mass proves the bound
conditionally. Integration removes the conditioning. Conditioning on an
additional independent variable proves the last assertion. This argument
includes all atom collisions and any open or closed interval endpoints.
\end{proof}
For the actual marginal we may couple the two conditional noises
independently and then sample an independent label. Since $U_0\le a$ and
$U_1\ge b-1$, the step is at least $b-a=1-d$.
Strict diameter is necessary: a Bernoulli law on $\{0,c\}$ assigns mass
one to $[0,c]$, exceeding its largest atom when $0<p<1$.

\paragraph{A rational cap for every remaining index.}
Put $J_0=N-2=45$ and define
\[
 r_L=l-(1-l)/J_0,\qquad r_R=h+h/J_0,\qquad
 v_* =\min\{r_L(1-r_L),r_R(1-r_R)\}.
\]
The checker verifies $0<r_L\le r_R<1$. Adjacent binomial-mass ratios
show that any mode $k$ at index $j\ge J_0$ satisfies
$(j+1)p-1\le k\le(j+1)p$, hence $r=k/j\in[r_L,r_R]$.
In particular $1\le k\le j-1$. Finite Stirling bounds give
\[
 \sqrt j\binom jk p^k(1-p)^{j-k}
 \le\frac{\exp(1/(12j)-jD(r\Vert p))}{\sqrt{2\pi r(1-r)}}
 <\frac{.4}{\sqrt{v_*}(1-1/(12J_0))}.
\]
Here $D(r\Vert p)\ge0$ follows from Jensen's inequality for the logarithm;
$\exp x\le(1-x)^{-1}$ for $0\le x<1$ follows term by term from the
power series. The finite Stirling inequalities used are
$\sqrt{2\pi m}(m/e)^m<m!<\sqrt{2\pi m}(m/e)^m\exp(1/(12m))$.
Choose the strict $10^{-6}$ grid upper square root $M$ of
$.16/[v_*(1-1/(12J_0))^2]$. Thus, whenever $w<1-d$,
\begin{equation}\label{j:local:conf:antichainconcentration}
 Q_j(w)<M/\sqrt j,\qquad j\ge n-2\ge J_0.
\end{equation}
This covers the $n-1$ full summands of $\nu$ and the $n-2$ full summands
remaining after one mark or deletion. An independent remaining conditional
noise preserves the cap. We do not replace the latter index by $n-1$.


Either the score-derived cap or the strict-gap cap may be selected on
an interval where its hypotheses hold. The first confinement uses the
score cap; later central and left schedules use the strict-gap cap.
For the left panels replace the conditional cubic coefficient 24 below
by 28; in general any $c_\rho$ satisfying
$c_\rho^2[l(1-l)]^3>[3(1-l)+.55]^2$ is valid.
\paragraph{Signed deletion, including atoms.}
Let $m=n-1$, $Z=x-\sum_{i=1}^mY_i$, $\mu=x-mp$ and let $K=\sum B_i$.
For a fixed removed label put
$T_b(I)=\Pp(x-\sum_{i<m}Y_i-b-U_b\in I)$.
Exchangeability and $mp-K=Z-\mu+\sum U_i$ give the exact identities
\[
 \nu(I)-T_1(I)=\frac{\E[(Z-\mu+\sum U_i)1_I(Z)]}{mp},\qquad
 T_0(I)-\nu(I)=\frac{\E[(Z-\mu+\sum U_i)1_I(Z)]}{mq}.
\]
There is no absolute value around the sum of noises.
For a full removed pair $(Y_i,U_i)$, condition on the pair. The remaining
sum gives a function $f(Y_i)\in[0,Q_{m-1}(w)]$. Since $\E U_i=0$,
\[
 \E[U_i f(Y_i)]\le Q_{m-1}(w)\E(U_i)_+
 =\tfrac12\epsilon Q_{m-1}(w).
\]
This also works when one conditional law is deterministic.

Let independent conditional noises satisfy $\Delta=U_1-U_0\ge-d$.
For $\tau>0$ put $I^+=I+1+[-d,\tau]$. On $\Delta\le\tau$, the
point for $T_1$ in $I$ maps to the point for $T_0$ in $I^+$. On the
complement, conditioning on both noises gives error at most
$Q_{m-1}(w)\Pp(\Delta>\tau)$. Moreover
\begin{equation}\label{j:local:conf:bad}
 \Pp(\Delta>\tau)\le
 \min\left\{\frac{s}{s+p\tau^2},\frac d{d+\tau}\right\}.
\end{equation}
Indeed $\E\Delta=0$, $\operatorname{Var}\Delta\le s/p$, and
$\Delta+d\ge0$ has mean $d$. Cantelli and Markov prove the two bounds.
The exceptional event is strict, so $\Delta=\tau$ belongs to the good event.

If $I,I^+$ are disjoint, combine the two signed noise terms before
bounding them. For the same remaining-sum law their conditional
probabilities $f_I,f_J$ satisfy
\[
 0\le f_I\le Q_I,\quad0\le f_J\le Q_J,\quad f_I+f_J\le Q_{\rm hull}.
\]
Since $p<q$,
\[
 \frac{f_I}{p}+\frac{f_J}{q}
 \le\min\left\{\frac{Q_I}p+\frac{Q_J}q,
 \frac{Q_{\rm hull}}q+\left(\frac1p-\frac1q\right)Q_I\right\}.
\]
Multiplication by $\E U_+=\epsilon/2$ proves the combined bound.
Closed, open and half-open intervals are all covered by containment in
the corresponding closed interval. No density or atom separation is used.

\paragraph{Slope bounds and the two exterior intervals.}
Split $|H|$ into 16 closed equal cells $[h_l,h_h]$ and $\lambda$ into
24 closed equal cells $[s_l,s_h]$. On each box put
\[
 \begin{gathered}
 \rho_* =\min\{\rho_0,\rho_B(l)+24s_h/N\},\\
 b_- =\max\{\beta_-,1-2h-2(1-h)\Theta\},\qquad
 b_+ =\min\{\beta_+,1-2l+\Theta\}.
 \end{gathered}
\]
Here $\rho_B(p)=(1-2p(1-p))/\sqrt{p(1-p)}$. Write $s_i=\E U_i^2$. Conditional centering and the two sign supports give
\[
 \rho=\ell^3\{pq(1-2pq)+3pq(s_0+s_1)
                  +p\E U_1^3-q\E U_0^3\}.
\]
The initial strip and $d_0<.55$ imply $|U_i|<.55$.
Thus $\rho\le\rho_B+\ell^3(3q+.55)s\le\rho_B+24s$; the last
coefficient inequality is checked as
$24^2[l(1-l)]^3>[3(1-l)+.55]^2$.
Similarly $\beta=(1-\theta)(1-2p)+\ell^2(ps_1-qs_0)$, and
$|\ell^2(ps_1-qs_0)|\le\theta$. These identities prove the stated
outward bounds on $\rho$ and $\beta$.
Set $f=.4F_\ell$, $\varphi_-=.39886(1-h_h^2/(2N))>0$.
The exact upper bound for $H\phi(H/\sqrt n)$ used below is
\[
 v_+=\begin{cases}.4h_h&H\ge0,\\-\varphi_-h_l&H\le0.\end{cases}
\]
For all $t\le0$, completing the square in \eqref{j:local:conf:derivative} gives
\[
 \mathcal M_-=F_\ell\pos{(v_++3C_0\rho_*)^2/(12(.4))-3(.4)b_-},
 \qquad L_-=E\varphi_- -\mathcal M_-/N.
\]
A positive bound $B_*>\sqrt{(1-l)/l}$ gives valid central-gap lower
and positive-half-line upper slope bounds
\begin{align*}
 \mathcal M_+&=F_\ell\pos{.4h_hB_*+3C_0B_*(B_*-1)-1.2b_-},\\
 L_g&=E\varphi_--\max(\mathcal M_-,\mathcal M_+)/N,\\
 U_+&=f+\frac{F_\ell}{N}
 \left\{3C_0b_++\frac{(.4h_h+3C_0\rho_*)^2}{12(.4)}\right\}.
\end{align*}
They follow by completing a square on the positive half-line and using
$t(t-1)\le B_*(B_*-1)$ for $1\le t\le B_*$; on $0\le t\le1$ that
term has the favorable sign. Negative values are covered by $\mathcal M_-$.
All divisions and positivity requirements are explicit gates.

\paragraph{The sharper lower-ray slope.}
Take $\tau_-=5r_0/12$, $\eta=10^{-5}$ and
$I=[a-r,a-d-\tau_--\eta]$, where $r\ge r_0$.
Since $a\le d$, the entire interval has
$t\le t_*=-E(l+\tau_-+\eta)<0$.
Write $B(t)=b_-+\rho_*t+t^2$ and
$B_*^{\min}=B(\min(t_*,-\rho_*/2))$. For the actual $C$ and $t\le t_*$,
\[
 3C(\beta+\rho t+t^2)+H\phi t
 \ge1.2B(t)+v_+t+3(C_0-.4)\min(0,B_*^{\min}).
\]
The inequality uses $t<0$ and separates the possible negative part of
$B(t)$ before varying $C$. Minimize the right side at the clipped vertex
$\min(t_*,-(v_++1.2\rho_*)/2.4)$ and call its value $q_*$. Then
\[
 L=E\varphi_--F_\ell\max(0,-q_*)/N
\]
is a lower slope on the whole interval, for every $n\ge N$.
Positive corrections are discarded. Therefore contact gives
\[
 \sqrt n\nu(I)\ge L(r-d-\tau_--\eta).
\]
This one slope replaces the previous separate bounds on the long and
short parts of the interval.

\paragraph{Lower interval geometry.}
Its translate is $I^+=[a-r+1-d,b-\eta]$, of width $r-\eta$.
The initial strip implies $I^+\subset[a+g,b)$,
$g=1-.55-2d_0>0$, hence $\nu(I^+)\le J/g$.
The intervals are disjoint throughout $r\le.55+d_0$, since
$\inf I^+-\sup I=1-r+\tau_-+\eta>0$.
At $r=r_0$, concentration and directed distance bounds are
\begin{align*}
 Q_I&=f(r_0-\tau_--\eta)+2C_0\rho_*,&
 Q_J&=f(r_0-\eta)+2C_0\rho_*,\\
 Q_{\rm hull}&=f(1+r_0-\eta)+2C_0\rho_*,\\
 D_I&=\pos{v_H-l-\tau_--\eta},&
 D_J&=\pos{v_H+1-l-\eta},
\end{align*}
where $v_H=-h_l/F_\ell$ for $H\ge0$ and $v_H=h_h/E$ for $H\le0$.
These signs follow from $\mu=p+H/\ell$.

\paragraph{Upper interval geometry.}
For a support point $y\ge1+r_0$, take $\tau_+=7r_0/15$,
$I=[a,a+\tau_+]$ and $I^+=[b,a+1+2\tau_+]$.
The strict gates $d_0+2\tau_+<r_0$ and $1-d_0-\tau_+>0$ give
$I^+\subset[b,y)$ and disjointness. Contact, including the atom at $b$,
and the leakage definition give
\[
 \sqrt n\nu(I)\ge L_g(1-d)-\sqrt nJ/\tau_+,\qquad
 \sqrt n\nu(I^+)\le U_+(.55+d).
\]
Use
\begin{align*}
 Q_I&=f\tau_++2C_0\rho_*,&Q_J&=f(2\tau_++d_0)+2C_0\rho_*,\\
 Q_{\rm hull}&=f(1+2\tau_+)+2C_0\rho_*,\\
 D_I&=\pos{d_0+\tau_++h_h/E-l},&
 D_J&=\pos{1+d_0+2\tau_++h_h/E-l}.
\end{align*}
The upper side deliberately keeps the absolute phase bound.


\paragraph{The strict-gap alternative.}
If $1-.55-2d_0>0$ and the displayed lower intervals have width below
$1-d$ throughout $r\in[r_0,.55+d_0]$, take $Q_I=Q_J=M$.
For the upper intervals the corresponding gate is
$1-2d_0-2\tau_+>0$; then again $Q_I=Q_J=M$.
Use the same directed distances. Their concentration caps are now
independent of $r$, and use
$a_e=F_dM\{1/l+1/(1-h)\}/2$ in the next formulas.
The lower-ray gate reduces to $L>0$: its contact side increases with
$r$, while the entire error is constant on the acquired finite strip.
For the score-derived cap retain the explicit affine slope gate below.
No assertion outside that strip is needed.

\paragraph{One budget, two error bounds.}
Let $F_d$ be the strict grid upper bound for
$\sqrt{\max\{(47/46)^3,47/45\}}$. It simultaneously covers the two
concentration factors in \eqref{j:local:conf:concentration} after multiplication by
$\sqrt n$, uniformly for $n\ge47$. For either pair of intervals set
\begin{align*}
 a_e&=\frac{F_d}{2}\min\left\{\frac{Q_I}l+\frac{Q_J}{1-h},
 \frac{Q_{\rm hull}}{1-h}+\left(\frac1l-\frac1{1-l}\right)Q_I\right\},\\
 A&=F_d\left\{Q_ID_I/l+Q_JD_J/(1-h)\right\},\qquad B=F_dQ_I.
\end{align*}
The probability denominators are outward because $1/p-1/(1-p)$ is
decreasing on $(0,1/2)$.
Use $a_d=L$, $a_J=1/g$ on the lower side, and
$a_d=L_g+U_+$, $a_J=1/\tau_+$ on the upper side.
For the phase cell, maximize each concave budget quadratic at its
vertex clipped to $[h_l,h_h]$, take their maximum, then the minimum
with $T$; call this upper bound $T_*$. Put
\[
 R=T_*/\kappa-b_*E^2s_l,
 \quad\alpha=\min(c_+lE,c_-(1-h)E),\quad G=(R-NP)/\alpha.
\]
If $G\le0$, the strict common budget excludes the entire closed box.
Otherwise $nd<G$ and
\[
 B_\tau=\min\{s_h/(s_h+Nl\tau^2),G/(N\tau+G)\}
\]
bounds \eqref{j:local:conf:bad}. With
\[
 z=\max\left\{\frac{a_d/l+2a_e}{Ec_+},
 \frac{a_d/(1-h)+2a_e}{Ec_-},\frac{\kappa a_J}{1-\Theta}\right\},
\]
one valid total error bound is
\begin{equation}\label{j:local:conf:oldbound}
 \mathcal E_0=(zR+A)/N-zP+B B_\tau.
\end{equation}
This follows by multiplying the single budget by $z/\kappa$.
The inverse-index coefficient $zR+A$ is checked positive.

A stronger bound spends the same resource on the pair event. Let
$R_0=R-NP>0$ and define the three pairs
\[
 (v_+,w_+)=\left(\frac{a_d/l+2a_e}{Ec_+}\frac{R_0}{N},
                    \frac{R_0}{Ec_+l}\right),
\]
\[
 (v_-,w_-)=\left(\frac{a_d/(1-h)+2a_e}{Ec_-}\frac{R_0}{N},
                    \frac{R_0}{Ec_-(1-h)}\right),\qquad
 (v_J,w_J)=\left(\frac{\kappa a_J}{1-\Theta}\frac{R_0}{N},0\right).
\]
They are upper cost/displacement coordinates at the vertices of the
simplex with resources $nc_+M_+$, $nc_-M_-$ and $n^{3/2}L/\kappa$.
For any $t\ge0$,
\[
 \frac{D}{N\tau+D}\le\frac{N\tau}{(N\tau+t)^2}D+
 \frac{t^2}{(N\tau+t)^2}\qquad(D\ge0).
\]
The difference is $N\tau(D-t)^2/[(N\tau+D)(N\tau+t)^2]$.
Consequently a further valid error bound is
\[
 \mathcal E_t=\frac AN+
 \max_{i\in\{+,-,J\}}\left(v_i+B\frac{N\tau}{(N\tau+t)^2}w_i\right)
 +B\frac{t^2}{(N\tau+t)^2}.
\]
The checker takes the minimum of \eqref{j:local:conf:oldbound} and these eleven
bounds with $t=jG/10$, $j=0,\ldots,10$. Each is individually valid;
no numerical minimization or derivative of their minimum is assumed.

The lower contradiction is the strict gate
$Lr_0-L(\tau_-+\eta)>\mathcal E$.
To extend it to every $r\ge r_0$, bound increments directly in the original
affine-in-width signed deletion bound for the same law. The budget gives
$n\epsilon<E_*=2T/(\kappa E\min(c_+,c_-))$. The sufficient slope gate is
\[
 L>\frac{F_df}{N}\left\{D_I/l+D_J/(1-h)
 +\frac{E_*}{2}\left(1/l+1/(1-h)\right)\right\}+F_dfB_\tau.
\]
Using the uncoupled increment is safe even when the hull bound was used
at $r_0$. The upper contradiction is $L_g-.55U_+>\mathcal E$.
Thus no support point lies below $a-r_0$ or above $1+r_0$.
Since $a\ge0$, $b\ge1-d_0$ and $d_0<r_0$, the label-zero and label-one
supports both satisfy $|U|<r_0$.


\paragraph{The selected schedules.}
The ordinary, central-antichain, and left-antichain certificates substitute
their recorded $r_0$ into these same gates. Noise-dependent amplitudes
are the same theorem with a positive lower noise endpoint in
$R=T_*/\kappa-b_*E^2s_l$. In the central rows the additional amplitudes
are $.35,.30,.275$ on $[.05,.08],[.08,.10],[.10,\Lambda]$;
on the left they are $.375,.325,.30$. Intersect each interval with its
own row's $\Lambda$ and each amplitude with that row's global amplitude.
The fixed acceptance checks both signs of $H$, all closed phase/noise
cells, clipped endpoints, and strict budget exclusions. No proof of
LOW, MID, or a binomial score comparison is used by this confinement lemma.




\paragraph{Finite-index substitution in confinement.}
\label{j:local:finite-confinement}
Here the numerical constants in the confinement proof of
Section~\ref{j:local:confinement} are parameters. More precisely,
replace its floor $47$ by $N\ge39$, its initial strip $.55$ by a
number $w\in(0,1)$, and its conditional cubic coefficient by
$c_\rho>0$. Retain a probability panel $[l,h]\subset(0,1/2)$ and assume
\[
 \operatorname{supp}Y\subset(-w,1+w),\quad d_0<w,
 \quad c_\rho^2[l(1-l)]^3>[3(1-l)+w]^2.
\]
The strip and the inward bounds $U_0\le d_0$, $U_1\ge-d_0$
together give $|U_i|<w$. The displayed conditional cubic identity gives
$\rho\le\rho_B+\ell^3(3q+w)s\le\rho_B+c_\rho s$.
Use $J_0=N-2$ for the binomial mode interval and the deletion factor
$\sqrt{\max\{(N/(N-1))^3,N/(N-2)\}}$. In the interval geometry
use $g=1-w-2d_0>0$ and $r\le w+d_0$; in the upper contradiction
use $L_g-wU_+>\mathcal E$. The other width, disjointness,
positive-slope and budget gates are unchanged. These substitutions
follow directly through the proof: $w$ enters only the conditional
cubic bound and the endpoints of the known strip, while the index
enters only the decreasing factors already bounded above. They
therefore prove the same confinement implication whenever all the
resulting gates pass.

The positive-allowance calculation at actual index $39$ uses $N=39$,
$J_0=37$ and deletion factor
$\sqrt{\max\{(39/38)^3,39/37\}}$. Its acquired initial widths are
$w=.55$ or $.575$; the central and left cubic coefficients are
$24$ and $30$, respectively. All the same explicit gates are checked
with these parameters before a contraction is used.
The central floor-$40$ sequence uses $N=40$, $w=.55$ initially and
$c_\rho=24$. The left rows start with their separately acquired
strip $w\in\{.55,.575,.6\}$ and use $c_\rho=30$. After a strict
contraction, the proved amplitude becomes the next strip width.
Thus every later use has an already established strip, and the
left sequence does not assume the narrower uniform-tail strip.



\subsection{The common comparison identity and the nearest count}\label{j:local:shortloc:section}

A law concentrated near two points has a natural integer count and a small
centered error. We compare it with its own Bernoulli count. The same exact
identity handles every noise regime. There are two useful estimates of its
right side: discard all but the nearest outgoing count, or retain all counts
through a random cutoff. The estimates share their normalization, moments,
and tail bounds; only the final scalar witness changes.

\paragraph{The identity and the common numerical test.}

Let $0<p\le1/2$, $q=1-p$, $v=pq$, and let
\[
 Y=B+U,\qquad \Pr(B=1)=p,\qquad \mathbb E(U\mid B)=0.
\]
The two actual conditional laws $U_0,U_1$ obey
\begin{equation}\label{j:local:shortloc:support}
 |U_j|\le A,\quad U_0\le a_0,\quad U_1\ge-a_1,\quad
 a_0,a_1\ge0,\quad d=a_0+a_1\le d_0<\min(p,q),\quad d_0\le A.
\end{equation}
Write $s_j=\mathbb E U_j^2$, $s=qs_0+ps_1$, $\lambda=ns$,
$\ell=(v+s)^{-1/2}$, $\theta=s/(v+s)$ and
$\rho=\mathbb E|\ell(Y-p)|^3$. At the actual threshold choose
\[
 x=np+H/\ell=k+u,\qquad k\in\{1,\ldots,n-1\},\quad |u|\le\tfrac12,
 \qquad \delta=k-np=H/\ell-u.
\]
Either nearest integer is allowed at a tie. Put
\[
 \pi_j=\binom njp^jq^{n-j},\quad F_B(j)=\Pr(\operatorname{Bin}(n,p)\le j),
 \qquad \rho_B=(1-2v)/\sqrt v,
\]
with $\pi_j=0$ outside $[0,n]$. Let $W_j$ be an independent sum of $j$
copies of $U_1$ and $n-j$ copies of $U_0$. The signed scores are
\[
 C=\frac{\sqrt n}{\rho}\{\Pr(\textstyle\sum Y_i\le x)-\Phi(H/\sqrt n)\},
 \qquad C_B=\frac{\sqrt n}{\rho_B}\{F_B(k)-\Phi(\delta/\sqrt{nv})\}.
\]
Throughout this section $C\ge c_*>0$; the selected application uses $c_*=.4$.
Define
\begin{align}
 f(u)&=\frac{\Phi((\delta+u)/\sqrt{n(v+s)})-
                    \Phi(\delta/\sqrt{n(v+s)})}{\pi_k},\label{j:local:shortloc:f}\\
 Q&=\Pr(\textstyle\sum Y_i\le k+u)-F_B(k),\nonumber\\
 \mathcal R&=\sqrt n\{\Phi(\delta/\sqrt{nv})-
                 \Phi(\delta/\sqrt{n(v+s)})\}+C(\rho_B-\rho).\nonumber
\end{align}
Direct subtraction gives the exact identity
\begin{equation}\label{j:local:shortloc:master}
 \rho_B(C-C_B)=\sqrt n\pi_k\{Q/\pi_k-f(u)\}+\mathcal R.
\end{equation}
In particular, the phase $u$, the sign of $H$, and both conditional laws
refer to the same marginal throughout.

\begin{samepage}
\begin{theorem}[Parameterized local comparison]\label{j:local:shortloc:theorem}
Fix $n\ge N$, $\lambda>0$ and a closed parameter cell on which
\eqref{j:local:shortloc:support} holds. Suppose $\sqrt n\pi_k\ge m>0$.
Either of the following witnesses proves $\rho_B(C-C_B)<-\varepsilon\lambda$.
\begin{enumerate}
\item \emph{Nearest-count witness.} Valid bounds $O,T,R_0,\kappa$ satisfy
\[
 \Pr(W_k>u)+f(u)\ge O,\quad
 \sum_{i\ge1}\Pr(W_{k+i}\le u-i)\le T,\quad
 \pi_{k+i}/\pi_k\le R_0,\quad \mathcal R\le\kappa\lambda/n.
\]
Only legal counts occur in the sum; adding positive upper bounds for
other counts is allowed. If $D=O-R_0T>0$ and
\begin{equation}\label{j:local:shortloc:directtest}
             mD-\kappa\lambda/n>\varepsilon\lambda,
\end{equation}
the conclusion follows.
\item \emph{Whole-count witness.} Bounds $B(\lambda)<0$ and $I(\lambda)\ge0$
satisfy
\[
 \rho_B(C-C_B)\le\sqrt n\pi_k\lambda B(\lambda)
                                      +\lambda I(\lambda)/n.
\]
It is enough that
\begin{equation}\label{j:local:shortloc:whole-test}
                  mB(\lambda)+I(\lambda)/N<-\varepsilon.
\end{equation}
\end{enumerate}
The witnesses may differ from cell to cell. A finite closed cover and a
separate treatment of $\lambda=0$ therefore suffice for a complete local
comparison.
\end{theorem}
\end{samepage}
\begin{proof}
Conditioning on the binomial count gives, with all atom conventions intact,
\begin{equation}\label{j:local:shortloc:count}
 Q=-\sum_{i=0}^{k}\pi_{k-i}\Pr(W_{k-i}>u+i)
       +\sum_{i=1}^{n-k}\pi_{k+i}\Pr(W_{k+i}\le u-i).
\end{equation}
Dropping the outgoing terms with $i>0$ gives $Q/\pi_k-f(u)\le-D$.
Use \eqref{j:local:shortloc:master} and $D>0$. For the second witness, use
$B<0$ before replacing $\sqrt n\pi_k$ by $m$, and $I\ge0$ before replacing
$n$ by $N$. These sign requirements are part of the test, not conventions.
\end{proof}

\paragraph{Normalization and the repair, once for every witness.}

Here is an explicit finite recipe for the constants required above.
Let $p\in[p_-,p_+]\subset(0,1/2)$, $n\ge N$, $\lambda\le\Lambda$,
$\theta\le\Theta$, and $|\delta|\le\Delta$, where $\Delta\ge1$ also includes
every binomial mode. Choose $0<r_-<p_- -\Delta/N$ and
$p_++\Delta/N<r_+<1/2$. Let $v_*>0$ be a lower bound for $p(1-p)$.
Choose positive $E,L,R$ bounding, on $[r_-,r_+]$,
\[
 |D'''(r\Vert p)|=\frac{1-2r}{r^2(1-r)^2},\qquad
 |(\log(r(1-r)))'|=\frac{1-2r}{r(1-r)},\qquad
 \frac1{12r(1-r)},
\]
respectively, where $D(r\Vert p)=r\log(r/p)+(1-r)\log((1-r)/q)$.
Then, with $Z(t)=Lt/2+R+Et^3/(6N)$,
\begin{equation}\label{j:local:shortloc:mass}
 \sqrt n\pi_k=\frac{\phi(0)}{\sqrt v}
       e^{-\delta^2/(2nv)+\zeta},\qquad n|\zeta|<Z(|\delta|).
\end{equation}
The finite Stirling estimate
$\log(j!)=(j+1/2)\log j-j+\log(2\pi)/2+r_j$, $0<r_j<1/(12j)$,
used in the concentration argument above gives this formula on substitution
in $n!/(k!(n-k)!)$ and Taylor expansion of relative entropy and the logarithmic
prefactor. Here is a finite proof of the estimate. For
$a_j=\log(j!)-(j+1/2)\log j+j$ and $z=(2j+1)^{-1}$,
\[
0<a_j-a_{j+1}=\sum_{i\ge1}\frac{z^{2i}}{2i+1}<\frac1{12j(j+1)}.
\]
Summation bounds $a_j-\lim a_i$. The integration-by-parts recursion for
$\int_0^{\pi/2}\sin^j t\,dt$, and the inequalities between successive
integrals, give Wallis's product and $\lim a_i=\log(2\pi)/2$.
The entropy Taylor error and factorial error give respectively
$E|\delta|^3/(6N)$ and $1/[12r(1-r)]$ in $Z$.

Consequently a valid mass lower bound is any positive
\begin{equation}\label{j:local:shortloc:m}
 m<\frac{.39886}{\sqrt{p_+(1-p_+)}}
          \left[1-\frac{\Delta^2/(2v_*)+Z(\Delta)}N\right].
\end{equation}
The same formula at a mode proves $\sqrt n\max_j\pi_j<1$ whenever
$v_*[1-Z(\Delta)/N]^2>.4^2$ and $Z(\Delta)<N$.
The rational bounds $.39886<\phi(0)<.4$ follow, for example, by the
alternating arctangent series in Machin's identity.

Differentiating \eqref{j:local:shortloc:f} gives
\begin{equation}\label{j:local:shortloc:logder}
 \log f'=-\tfrac12\log(1+s/v)+\delta^2/(2nv)
                   -(\delta+u)^2/(2n(v+s))-\zeta.
\end{equation}
Thus $n|\log f'|\le B_0\lambda+P(D,a)$ on $|\delta|\le D\le\Delta$,
$|u|\le a\le1/2$, where
\[
 B_0>\frac1{2v_*}+\frac{\Delta^2}{2Nv_*^2},\qquad
 P(D,a)=\frac{2Da+a^2}{2v_*}+Z(D).
\]
Choose $0<b<1$ above $[B_0\Lambda+P(\Delta,1/2)]/N$ and put
$F=(1-b)^{-1}$. The elementary estimates $e^z\le(1-z)^{-1}$ and
$1-e^{-z}\le z$ imply
\begin{equation}\label{j:local:shortloc:derivative}
 |f'-1|\le(B_1\lambda+P_1)/n,\qquad B_1=FB_0,\quad P_1=FP(D,a).
\end{equation}
This supplies both a derivative interval $l_f\le f'\le h_f$ and the
finer phase-error bound $|u-f(u)|\le a(B_1\lambda+P_1)/n$.

The exact count variance is
\begin{equation}\label{j:local:shortloc:variance}
 t_j=\mathbb E W_j^2=\lambda+(j-np)(s_1-s_0),\qquad
 |s_1-s_0|\le s/p\le c_v s,\quad c_v\ge1/p_-.
\end{equation}
It gives $\omega_-\lambda\le t_k\le\omega_+\lambda$ with
$\omega_\pm=1\pm c_vD/N$. All uses below require the checked condition
$\omega_->0$. Also
$t_{k+i}\le iV_*\lambda$ for $V_*=1+c_v(\Delta+1)/N$.
The binomial recursion gives, for every legal $i\ge1$,
\begin{equation}\label{j:local:shortloc:ratio}
 \frac{\pi_{k+i}}{\pi_k}\le R_0,
 \qquad R_0>1+\frac{p_++\Delta-1}
 {Nv_*[1-(\Delta-1)/(Np_-)]},
\end{equation}
provided the denominator is positive and $p_++\Delta<2$.
Indeed the first ratio has this bound and all subsequent rising factors
are at most one. No count outside $[0,n]$ is used in this recursion.

Finally \eqref{j:local:shortloc:support} fixes the signs inside the absolute cubic.
Expansion and centering give
\[
 \mathbb E|Y-p|^3=v(1-2v)+3pq(s_0+s_1)
                   +p\mathbb E U_1^3-q\mathbb E U_0^3
 \ge v(1-2v)+(3p-d)s.
\]
The inequality $(1+z)^{3/2}\le1+3z/2+3z^2/8$ for $z\ge0$ yields
\begin{equation}\label{j:local:cubic-bound}
 (v+s)^{3/2}(\rho-\rho_B)\ge s
     \left(6p-3p^2-\tfrac32-d-\frac{3(1-2v)s}{8v}\right).
\end{equation}
Let $Q_*>0$ be a lower bound for the bracket and
$m_c\le c_*/[p_+(1-p_+)+\Lambda/N]^{3/2}$.
The normal variance correction is nonpositive for $\delta\le0$;
for $\delta>0$ its derivative in $s$ is at most $.4\delta/(2v^{3/2})$.
Therefore
\begin{equation}\label{j:local:shortloc:repair}
 \mathcal R\le(P_*D_+-m_cQ_*)\lambda/n,
 \qquad P_*\ge\frac{.4}{2v_*^{3/2}},\quad D_+\ge\delta_+.
\end{equation}
One may take $D_+=D$, but retaining the positive part is often decisive.
All bounds concern $v+s$, which can exceed $1/4$.

\paragraph{Two elementary tools for probabilities.}

\begin{lemma}[A shifted quartic]\label{j:local:shortloc:quartic}
Suppose $W$ is a sum of independent centered variables bounded in absolute
value by $A$, and $t=\mathbb E W^2$. Let $K_0\ge(9+6\sqrt3)/4$ and $a>0$.
For $r\ge0$ and $\sigma\in\{-1,1\}$,
\begin{equation}\label{j:local:shortloc:quarticformula}
 \Pr(W+\sigma r>0)\ge
 \frac{(3a^2-A^2)t-3t^2-4Art-6r^2t-r^4
                      +3a^2r^2+2\sigma a^3r}{K_0a^4}.
\end{equation}
\end{lemma}
\begin{proof}
The polynomial $h_a(z)=(-z^4+3a^2z^2+2a^3z)/(K_0a^4)$ is below
$\mathbf1_{\{z>0\}}$: it is nonpositive for $z\le0$, and its maximum for
$z>0$ is at $z/a=(1+\sqrt3)/2$ and is at most one. Independence gives
$|\mathbb E W^3|\le At$ and $\mathbb E W^4\le A^2t+3t^2$.
Expand $\mathbb Eh_a(W+\sigma r)$ and apply these two inequalities.
The value at the threshold is zero, so atoms require no limiting argument.
\end{proof}

\begin{lemma}[One common tail bound]\label{j:local:shortloc:tails}
If independent centered $X_i\le A$ have total variance at most $V$, then,
for $b>0$, their sum has upper tail bounded by the minimum of one and
\begin{equation}\label{j:local:shortloc:tailmin}
 \begin{gathered}
 \frac{V}{V+b^2},\qquad
 \mathcal B(A,V,b)=e^{b/A}\left(\frac{V}{V+Ab}\right)^{b/A+V/A^2},\qquad
 \frac{\mathbb E(P_\mu-j)_+^3}{(b/A+\mu-j)^3},\\
 \mu=V/A^2,\qquad 0\le j<b/A+\mu.
 \end{gathered}
\end{equation}
where $P_\mu$ is Poisson. The same holds for a lower tail when $|X_i|\le A$.
All bounds apply to closed events, hence also to strict events.
\end{lemma}
\begin{proof}
Cantelli's bound follows by applying Markov to $(\sum X_i+V/b)^2$.
Taylor's integral remainder gives
$\mathbb E e^{zX_i}\le\exp\{\mathbb E X_i^2(e^{zA}-1-zA)/A^2\}$
for $z\ge0$. Multiply and use $z=A^{-1}\log(1+Ab/V)$ to get $\mathcal B$.
The zero-variance case is immediate.

For the third bound, smooth $f(x)=(x-t)_+^3$ by a nonnegative smooth
compactly supported density. Then $f'',f''',f''''\ge0$. For
$\mathcal Af(x)=f(x+A)-f(x)-Af'(x)$, the Taylor quotient
\[
 \frac{f(x+y)-f(x)-yf'(x)}{y^2}
       =\int_0^1(1-r)f''(x+ry)\,dr
\]
is nondecreasing in $y$. A centered $X\le A$ of variance $v$ consequently
satisfies $\mathbb E f(x+X)\le f(x)+(v/A^2)\mathcal Af(x)$.
For $\mathcal P_zf(x)=\mathbb E f(x+A(P_z-z))$,
$\partial_z\mathcal P_zf=\mathcal P_z\mathcal Af$ and
\[
 \partial_z^2\mathcal P_zf=\mathcal P_z\mathcal A^2f\ge0,\qquad
 \mathcal A^2f(x)=\int_0^A\!\int_0^A(A-r)(A-w)f''''(x+r+w)\,dr\,dw.
\]
Thus $\mathcal P_{v/A^2}f\ge f+(v/A^2)\mathcal Af$. Replace the independent
variables successively; translations and expectations preserve the derivative
signs and Poisson rates add. Increasing the total variance to $V$ is permitted
by Jensen after adding independent centered Poisson noise. Removing the
smoothing is justified by bounded actual upper supports and finite Poisson
third moments. Markov at $t=A(j-\mu)$ proves the displayed bound.
\end{proof}

These bounds have finite rational evaluations. For example,
\begin{align}\label{j:local:shortloc:poissonfinite}
 \mathbb E(P_\mu-j)_+^3={}&\mu^3+3(1-j)\mu^2+(1-3j+3j^2)\mu-j^3\\
 &+e^{-\mu}\sum_{h=0}^{\lfloor j\rfloor}\frac{\mu^h(j-h)^3}{h!}.\nonumber
\end{align}
Replacing $e^{-\mu}$ by $(\sum_{h=0}^{24}\mu^h/h!)^{-1}$ enlarges the
nonnegative correction. Alternatively sum the positive Poisson series
through 32; the remaining tail is at most its first term divided by
$1-\mu(34-j)^3/[34(33-j)^3]$, provided $j<33$ and that ratio is below one.
For Bennett, put $x=Ab/V$ and
$L=[(b/A+\mu)\log(1+x)-b/A]$. The positive series
$\log(1+x)=2\sum_{h\ge0}z^{2h+1}/(2h+1)$, $z=x/(2+x)$, gives lower
bounds for $L$; any positive lower exponential series gives an upper bound
for $e^{-L}$. Rounding a positive lower $L$ down before exponentiation is
also valid. These are choices of rational witnesses, not optimization claims.

\paragraph{The nearest-count witness, including vanishing noise.}

For $u=-r\le0$, add $-h_fr$ to \eqref{j:local:shortloc:quarticformula} with
$\sigma=1$; for $u=r\ge0$, add $l_fr$ with $\sigma=-1$.
Both bound $\Pr(W_k>u)+f(u)$ from below. For a cell
$t\in[t_-,t_+]$, $r\in[r_-,r_+]$, this expression is concave in $t$
(second derivative $-6/(K_0a^4)$), so only $t_-,t_+$ are needed.
At each endpoint, bound a monomial with positive coefficient at $r_-$
and one with negative coefficient at $r_+$. Any rational $a>0$ is a valid
witness. Other valid choices are
\begin{equation}\label{j:local:shortloc:ownother}
 \frac{r_-^2}{t_++r_-^2}-h_fr_+\quad(u\le0),\qquad l_fr_-\quad(u\ge0).
\end{equation}
They follow from Cantelli and nonnegativity, respectively. Their maximum
or one selected witness is valid without an optimality proof.

For incoming counts use $t_{k+i}\le iV_*\lambda$. On $u\in[-r_+,-r_-]$
the first distance is $1+r_-$ and every later distance is at least $i$.
On $u\in[r_-,r_+]$ they are $1-r_+$ and $i(1-r_+/2)$ for $i\ge2$.
Let $b,b_o$ be these first and later distances. On $\lambda\in[L,U]$,
let $T_1$ be any bound from Lemma~\ref{j:local:shortloc:tails} at $(A,V_*U,b)$.
Since $\mathcal B(A,iV,ib)=\mathcal B(A,V,b)^i$, all later counts total at most
$B^2/(1-B)$ for $B=\mathcal B(A,V_*U,b_o)<1$. Thus
\begin{equation}\label{j:local:shortloc:midtest}
 D=O-R_0\{T_1+B^2/(1-B)\}>0,\qquad
 mD-\kappa U/N>\varepsilon U
\end{equation}
proves the entire cell when $\kappa\ge0$. The selected MID normalizers
explicitly verify this sign. This is the common central and left MID proof.

A cell touching $\lambda=0$ requires estimates proportional to $\lambda$;
a positive lower endpoint may not be substituted. The following rescaling
provides them. Put $r=zt$, $t_*=\omega_+U$ and $d_a=K_0a^4$.
The shifted quartic plus $br$, divided by $t$, becomes
\begin{equation}\label{j:local:shortloc:scaledown}
 F_{\sigma,b}(z,t)=\frac{3a^2-A^2-3t}{d_a}
 +\left(\frac{2\sigma}{K_0a}+b-\frac{4At}{d_a}\right)z
 +\frac{3a^2t-6t^2}{d_a}z^2-\frac{t^3z^4}{d_a}.
\end{equation}
It extends to $t=0$ and is concave in $t$. On $0\le z\le z_0$, replacing
$z^4$ by $z_0^2z^2$ gives a lower quadratic at $t=t_*$; at $t=0$ it is
linear. Endpoints and a clipped quadratic vertex therefore give an exact
lower bound $J(\sigma,b,z_0)$ for the whole rectangle.

For negative phase choose $z_-,c_->0$ and verify
\begin{gather}\label{j:local:shortloc:tinynegative}
 J(1,-h_f,z_-)>c_-,\quad
 \frac{z_-^2}{1+z_-^2t_*}-h_fz_->c_-,\quad
 \frac{2z_-}{(1+z_-^2t_*)^2}>h_f,\\
 \frac1{1+4t_*}-\frac{h_f}{2}>c_-t_*,\qquad z_-t_*<1/2.\nonumber
\end{gather}
The quartic handles $r\le z_-t$. Beyond that point use
$r^2/(t+r^2)-h_fr$. Its derivative at $r=z_-t$ is positive by the third
gate; its second derivative has just one sign change, so any later critical
point is a maximum. Its minimum is at $z_-t$ or $1/2$, covered by the other
gates. Hence the own count exceeds $c_-t$.
For positive phase choose $z_+,c_+>0$ with
\begin{equation}\label{j:local:shortloc:tinypositive}
 J(-1,l_f,z_+)>c_+,\qquad l_fz_+>c_+,\qquad z_+t_*<1/2.
\end{equation}
The quartic handles $r\le z_+t$ and the normal term handles the rest.
The own count exceeds $c_+t$ and also exceeds $l_fr$.

For incoming tails proportional to $\lambda$, fix a Poisson anchor $j$.
A positive series term of index $h$ divided by
$\lambda(b/A+\mu-j)^3$ has logarithmic derivative, multiplied by $\lambda$,
\[
 h-1-\mu-\frac{3\mu}{b/A+\mu-j},\qquad \mu=V_*\lambda/A^2.
\]
If $b/A-j>0$ and this is nonnegative at $\lambda=U$ for the first positive
index, it is nonnegative for every term and all $0<\lambda\le U$.
Therefore the Poisson bound divided by $\lambda$ is largest at $U$.
Likewise $B^2/[\lambda(1-B)]$ is increasing whenever
$\lambda B'/B\ge1/2$. Differentiating the Bennett formula gives
\[
 \lambda B'/B=b/A-\mu\log(1+b/(A\mu));
\]
The function $g(\mu)=b/A-\mu\log(1+b/(A\mu))$ decreases because
$\log(1+x)>x/(1+x)$ for $x>0$. Moreover the logarithmic derivative of
$B^2/[\lambda(1-B)]$, multiplied by $\lambda$, is
$2g-1+gB/(1-B)\ge0$ when $g\ge1/2$. Thus checking the gate at $U$
proves the whole interval. A rational upper bound for the logarithm
verifies that gate. On cells using the fixed positive bound $l_fr_-$ instead,
unscaled endpoint tails suffice: after a positive bracket is checked, divide
it by $U$ to obtain a lower bound per $\lambda$. These arguments, followed by
\eqref{j:local:shortloc:directtest}, prove TINY uniformly down to arbitrary positive
noise, including a deterministic conditional component.

\subsection{Whole-count comparison}

Assume $c=1-d_0>1/2$. Couple all $2n$ conditional variables independently and
set $W_j=\sum_{i\le j}U_{1,i}+\sum_{i>j}U_{0,i}$, extending it constantly
outside $[0,n]$. Every increment of $j+W_j$ is at least $c$. Define
\[
 K=\max\{j\in\mathbb Z:j+W_j\le k+u\},\qquad R=K-k.
\]
An independent binomial selector has the correct actual count laws, so
\begin{equation}\label{j:local:shortloc:cutoff}
                  \Pr(\textstyle\sum Y_i\le k+u)=\mathbb EF_B(k+R).
\end{equation}
Let $\chi=(n+1)p-k<2$ and $m_B=\max_j\pi_j$. For every integer $r$,
\begin{equation}\label{j:local:shortloc:curvature}
 F_B(k+r)-F_B(k)-r\pi_k\le\frac{m_B}{(n+1)v}
 \left[(\chi-1)_+r_++\chi_+\binom{r_-}{2}+\binom{r_-}{3}\right].
\end{equation}
Indeed, the positive part of $\pi_{k+i}-\pi_{k+i-1}$ is at most
$m_B(\chi-i)_+/((n+1)v)$, by
$\pi_j-\pi_{j-1}=\pi_j((n+1)p-j)/(p(n-j+1))$.
Double summation gives \eqref{j:local:shortloc:curvature}: for $r>0$ only $i=1$
can contribute; for $r=-h<0$ use $(\chi+i)_+\le\chi_++i$ and sum the
multiplicities $h-1-i$. Zero extension makes the same calculation valid
outside the legal count range.

Write $t=\mathbb EW_k^2$. Crossing monotonicity and Cauchy--Schwarz give
\begin{equation}\label{j:local:shortloc:crossingmoments}
 \mathbb ER_+\le\frac{t}{(c-u)^2},\quad
 \mathbb E\binom{R_-}{2}\le\frac{t}{[c+\min(u,c/3)]^2},\quad
 \mathbb E\binom{R_-}{3}\le\frac{t\sqrt{A^2+3t}}{(2c-1/2)^3}.
\end{equation}
For the first, $R=i\ge1$ implies $-W_k\ge ci-u$ and
$(ci-u)^2\ge i(c-u)^2$. For $R=-h\le-2$, the next count is outside the
closed event, so $W_k>c(h-1)+u$. Put $j=h-1$.
The function $j(j+1)/(2(cj+u)^2)$ decreases for $j\ge1$ if $u\le c/3$;
for larger $u$ reduce to that endpoint. Also
$(j^3-j)/(6(cj-1/2)^3)$ decreases for $j\ge2$, since its derivative has
the sign of $2cj-(3j^2-1)/2<0$. These give the last two bounds using
$\mathbb E(W_k)_+^3\le\sqrt{\mathbb EW_k^2\mathbb EW_k^4}$.

\paragraph{The marked crossing.}
Let $D_i=U_{1,i}-U_{0,i}$, $\vartheta=\mathbb E(D_i)_+=\mathbb E(D_i)_-$,
$0<\eta\le1$, and $R_f=\mathbb E\{u+\eta-W_k\}$, where
$\{x\}=x-\lfloor x\rfloor$ is the fractional part. The two deleted laws are
\[
 A_0\overset{d}=kU_1+(n-k-1)U_0,\qquad
 A_1\overset{d}=(k-1)U_1+(n-k)U_0.
\]
Then
\begin{align}
 \mathbb ER&\le u+\eta-R_f+J,\label{j:local:shortloc:mark}\\
 J&=\frac{\vartheta}{\eta}\sum_{i\ge1}
 \{\Pr(A_0\le u+d_0-ci)+\Pr(A_1>u+\eta+i-A)\}.\nonumber
\end{align}
Here are the pathwise details. Put $z=u-W_k$ and mark
\[
 Z_f=\sum_{i=1}^{n-k}(D_{k+i})_-\mathbf1_{\{W_k\le u-ci\}},\quad
 Z_b=\sum_{i=1}^{k}(D_{k-i+1})_+\mathbf1_{\{W_k>u+\eta+i\}}.
\]
If $R>z+\eta$, the crossing activates the marks through $R$: for $R>0$
use $z\ge cR$, while for $R<0$ use $z<R-\eta$. In either case
$R\le z+Z_f+Z_b$. Hence if $h=R-\lfloor z+\eta\rfloor>0$,
$Z_f+Z_b>h-1+\eta\ge\eta h$. Thus
$R\le\lfloor z+\eta\rfloor+(Z_f+Z_b)/\eta$ on every path.
On an active forward mark $U_0>U_1\ge-d_0$; deleting its pair implies
$A_0\le u+d_0-ci$. On a backward mark delete $U_1\le A$ to obtain the
strict event for $A_1$. The remaining sum is independent of the entire
pair, so the expectations factor. Extending to all $i\ge1$ adds only
nonnegative bounds. Finally $|R|\le1+(|W_k|+|u|)/c$ proves integrability
and permits expectations. No coupling has replaced either deleted marginal.

\paragraph{A finite bound for every mark.}
Fix $\lambda\in[L,U]$, $L>0$, and $t\le\omega_+\lambda$; put
$T_L=\omega_+L$, $T_U=\omega_+U$. On a phase subcell $u\in[v_l,v_h]$,
for $i=1,2$ let $b_1=ci-d_0-v_h>0$ and $b_2=i-A+v_l>0$.
Lemma~\ref{j:local:shortloc:tails} gives a bound $\omega_+\lambda K_i^{\rm ind}$,
where
\[
 K_i^{\rm ind}=\sum_{j=1}^2\min\left\{\frac1{b_j^2+T_L},
                          \frac{\overline P(A,T_U,b_j)}{T_L}\right\}.
\]
The denominator $T_L$ is a lower bound for the chosen upper envelope
$\omega_+\lambda$, not for any actual variance.
Alternatively couple $A_0=Z+U_1$, $A_1=Z+U_0$ with
$Z\overset d=(k-1)U_1+(n-k-1)U_0$. The two events project into
$Z\le-\alpha$ and $Z>\beta$, where
$\alpha=ci-2d_0-v_h>0$, $\beta=i-A-d_0+v_l>0$.
For any $z\in(-\alpha,\beta)$ their sum is at most
$(\omega_+\lambda+z^2)/\min(z+\alpha,\beta-z)^2$.
Consequently a second kernel is
\[
 K_i^{\rm com}=\frac{1+z^2/T_L}{\min(z+\alpha,\beta-z)^2},\qquad
 z=\max\{-T_L/\beta,\min((\beta-\alpha)/2,T_L/\alpha)\}.
\]
This center is inside the interval; differentiating the two quadratic
ratios on either side of the midpoint proves its optimality, although
any interior rational center suffices. Take the smaller kernel on each
phase subcell and the maximum over their closed cover. Denote it $K_i$.
For $i\ge3$ use fourth moments and the convex midpoint integral for
$(cx-b)^{-4}$. On $u\in[u_l,u_h]$ this gives
\begin{equation}\label{j:local:shortloc:K}
 K=K_1+K_2+\frac{A^2+3T_U}{3c}
  \left\{\frac1{((5/2)c-A-u_h)^3}
                 +\frac1{((5/2)c-A+u_l)^3}\right\},\qquad
 J\le\frac{\vartheta}{\eta}\omega_+\lambda K.
\end{equation}
All denominators must be positive. Indeed, the tail fourth moment is at
most $A^2t+3t^2$, and integration from $5/2$ to infinity covers every
remaining integer, including the harmless added illegal counts.

If the same-law resource supplies $nd<G$ and
$G^2<c_vNL$, then
\begin{equation}\label{j:local:shortloc:theta}
 \vartheta\le\frac{c_vG\lambda}{c_vn\lambda+G^2}.
\end{equation}
To prove this, for centered $X\ge-b$, $\mathbb EX^2\le V$ and $b^2\le V$,
the quadratic $b(X-V/b)^2/(b+V/b)^2$ majorizes $X_-$ on $[-b,\infty)$.
Its expectation is at most $bV/(V+b^2)$. Apply this to $D_i\ge-d$,
using $\operatorname{Var}(D_i)=s_0+s_1\le c_v\lambda/n$, and monotonicity
in $b\le\sqrt V$. The case $d=0$ is immediate. With $\eta=z\lambda$,
$z>0$, $zU\le1$, \eqref{j:local:shortloc:K} becomes
\begin{equation}\label{j:local:shortloc:Jfinal}
 \frac J\lambda\le\frac{\omega_+Kc_vG}{z(c_vN\lambda+G^2)}.
\end{equation}

\paragraph{The fractional-part gain.}
If each centered summand has support span at most $b<1/2$ and the sum has
variance $t$, then, for any real $w$ and
$0<r_*\le2\sin^2(\pi b)/b^2$,
\begin{equation}\label{j:local:shortloc:harmonic}
             \mathbb E\{w-W\}\ge\frac5{23}(1-e^{-r_*t}).
\end{equation}
Indeed $1-\cos(2\pi x)\le(23/5)x$ for $0\le x\le1$.
A finite proof of this sole scalar inequality is recorded below.
For an independent copy $Z'$ of one summand, concavity of sine on
$[0,\pi b]$ gives
\[
 |\mathbb Ee^{2\pi iZ}|^2
 =1-2\mathbb E\sin^2\pi(Z-Z')
 \le e^{-4\sin^2(\pi b)\operatorname{Var}(Z)/b^2}.
\]
Multiply square roots over the independent summands and use periodicity.
This argument includes zero variances and integer collisions.
Here the conditional spans are at most $A+d_0$, so this is the $b$ to use.

For a reproducible scalar proof, $x\ge10/23$ follows from $1-\cos\le2$.
On $[0,10/23]$ use the degree-ten alternating upper polynomial for
$1-\cos(2\pi x)$, with directed coefficients:
\[
 Q(x)=\sum_{j=1}^5\frac{(-1)^{j+1}(2\pi_jx)^{2j}}{(2j)!},\qquad
 \pi_j=\begin{cases}3.141592653589794,&j\text{ odd},\\
                       3.141592653589793,&j\text{ even}.\end{cases}
\]
Put
$P(x)=23/5-Q(x)/x$. On the eight intervals with left endpoint $a_j=5j/92$
and width $w=5/92$, its degree-nine Bernstein coefficients are
\[
 b_{jk}=\sum_{i=0}^k\frac{\binom ki}{\binom9i}
           \sum_{h=i}^9[x^h]P(x)\binom hi a_j^{h-i}w^i,
 \qquad 0\le j<8,\quad0\le k\le9.
\]
The 80 rational coefficients are positive; the retained scalar acceptor
reconstructs them. Expansion in powers of the unit-interval coordinate
proves the formula; Bernstein weights are nonnegative and sum to one.
The alternating remainder and the stated rational bounds on $\pi$ follow
from the same elementary series above. This completely specifies the
finite scalar obligation. A rational rate follows by replacing
$\sin(\pi b)$ with $x-x^3/6+x^5/120-x^7/5040$, $x=3.14159b$, on the
checked alternating-series domain.

Set $r=r_*\omega_-$, require $rU<3$, and define
\begin{equation}\label{j:local:shortloc:gain}
 g=\frac5{23}\left(\frac{r^3L^2}{6}-\frac{r^4L^3}{24}
                     +\frac{r^5L^4}{120}-\frac{r^6L^5}{720}\right).
\end{equation}
The sixth-order Taylor lower bound for $1-e^{-r\lambda}$ gives
$R_f/\lambda\ge5r/23-5r^2\lambda/46+g$. The discarded higher-order
part is increasing as a function of $\lambda$ for $r\lambda<3$:
its derivative has the sign of $48-18x+(24/5)x^2-x^3>0$ on $[0,3]$.
For instance this polynomial has positive degree-three Bernstein
coefficients $48,30,132/5,51/5$ there. Thus freezing it at $L$ is valid.

\paragraph{One final expression.}
On a signed phase cell find an interval $[x_l,x_h]$ containing $H/\ell$
by its four endpoint quotients, using positive bounds for $\ell$.
Let
\[
 D=\max(|x_l-u_h|,|x_h-u_l|),\quad
 D_+=\max(0,x_h-u_l),\quad a=\max(|u_l|,|u_h|),\quad b_\chi=p_+-x_l.
\]
The same formula allows coarser upper bounds $D,D_+,b_\chi$, as used by
the central certificate. Check $D\le\Delta$ and $b_\chi+u_h<2$.
With $\omega_\pm=1\pm c_vD/N$, the curvature coefficient is
\begin{equation}\label{j:local:shortloc:V}
 V_0=\frac{\omega_+}{v_*}\left[
 \sup_{u_l\le u\le u_h}\left\{
 \frac{(b_\chi+u-1)_+}{(c-u)^2}
 +\frac{(b_\chi+u)_+}{[c+\min(u,c/3)]^2}\right\}
 +\frac{\sqrt{A^2+3T_U}}{(2c-1/2)^3}\right].
\end{equation}
On a rational phase subcell, both numerators are bounded at the right
endpoint; the first denominator is smallest there, and the second is
smallest at the left endpoint. This proves the supremum bound, including
both kinks. Equations \eqref{j:local:shortloc:cutoff}--\eqref{j:local:shortloc:crossingmoments},
$\sqrt n m_B\le1$, and \eqref{j:local:shortloc:master} now give the whole-count
witness with
\begin{align}\label{j:local:shortloc:BI}
 B(\lambda)&=z+\frac{\omega_+Kc_vG}{z(c_vN\lambda+G^2)}
                      -\frac{5r}{23}+\frac{5r^2\lambda}{46}-g,\\
 I(\lambda)&=V_0+aB_1+\frac{aP_1}{\lambda}+P_*D_+-m_cQ_*.
 \nonumber
\end{align}
We used $|u-f(u)|\le a(B_1\lambda+P_1)/n$ and $\sqrt n\pi_k\le1$.
Every constant is fixed before $\lambda$ varies. Both $B$ and
$mB+I/N$ are convex on $\lambda>0$; $I$ is decreasing. Thus their two
endpoint gates $B<0$, $I>0$, $mB+I/N<-.001$ prove the entire closed noise
cell. This proves all three LOW variants by exactly the same argument.

\subsection{Signed refinements and the complete local covers}

The retained tables contain 45 central probability intervals covering
$[.38,.47]$ and 38 left intervals covering $[.304,.38]$. Each row supplies
its own $\Theta$, signed $H$ bounds, $d_0$, $\Lambda$, and global amplitude;
the confinement profiles may improve the amplitude with $\lambda$.
All these are hypotheses of the present section, acquired earlier without
using its conclusion. The shared resource has the form
\[
 nD_{\rm mom}<\frac{15000}{997}T_H,\qquad
 D_{\rm mom}\ge\alpha d+\beta\theta+P.
\]
Consequently a noise cell $[L,U]$ has
\[
 nd<G=\{(15000/997)T_H-NP-\beta\ell_-^2L\}/\alpha.
\]
Here $T_H$ is the maximum of the acquired concave quadratic budget on the
closed signed $H$ cell, further bounded by its scalar ceiling. It already
contains the phase charge. If $G\le0$, the entire cell is impossible and
receives no score test; the exclusion is strict even at $G=0$.

For positive noise the complete cover is the following. Each interval is
intersected with the row ceiling $\Lambda$; shared endpoints are included.
\begin{center}\small
\begin{tabular}{@{}llll@{}}\toprule
Noise & Central phase & Left phase & Witness\\\midrule
$(0,.05]$ & all & all & nearest count (scaled)\\
$[.05,.06]$ & $|u|\ge.09$ & $|u|\ge.09$ & nearest count\\
$[.06,.08]$ & $|u|\ge.13$ & $|u|\ge.14$ & nearest count\\
$[.08,.10]$ & $|u|\ge.18$ & $|u|\ge.19$ & nearest count\\
$[.10,.12]$ & $|u|\ge.29$ & $u\le-.29$ & nearest count\\
$[.10,.12]$ & --- & $u\in[.29,.5]$ & whole count\\
$[.05,\Lambda]$ & remaining phase & remaining phase & whole count\\\bottomrule
\end{tabular}
\end{center}
Beyond $.12$ the last row covers the full phase interval. The three LOW
objects are one central cover, one left cover with phase slices, and one
extra positive left cover. Their distinction is only the closed domain
and the chosen outward constants. Signed $H$ cells meet at zero; nearest
count ties meet at $u=\pm1/2$. The fixed coverage checker also checks every
row-ceiling singleton and every entirely excluded LOW row.

The common TINY parameters are
\[
 N=47,\quad p\in[.304,.47],\quad |H|\le1.393,\quad
 \Theta=.0154,\quad d_0=.0181,\quad A=.425,\quad\Lambda=.05.
\]
For its scaled own witness use $a=.4784$, $K_0=9/(4\cdot.4641)$,
$(z_-,c_-)=(1.66,.514)$ and $(z_+,c_+)=(1.2,1.036)$.
Its six phase cells are $[-.5,0],[0,.1],\ldots,[.4,.5]$.
The first two use scaled Poisson anchors $1.25,1$, respectively; the other
four use the fixed positive normal bound. The exact normalizers and gates
specified in the accompanying certificate map verify these constants.
The central MID amplitudes are $.35,.35,.30,.275$ in its four rows.
The left ones are $.375,.375,.325,.30$, each further reduced by the row's
global amplitude. No interpolation of the tables is used.

All six accepted comparisons give the same strict conclusion
\[
                        \rho_B(C-C_B)<-.001\lambda.
\]
Their proofs have now reduced to Theorem~\ref{j:local:shortloc:theorem}, the shared
normalization, and the two explicitly derived witnesses. The arithmetic
objects still have their original domains and targets; no certificate
branch has been silently deleted or relabeled.

If $\lambda=0$, conditional centering and $p,q>0$ imply $U_0=U_1=0$
almost surely. The law is exactly Bernoulli. Schulz's strict finite-binomial
bound~\cite[Theorem~1]{Schulz2016} states, for every $0<p<1$, finite $n$
and real $x$,
\[
 \sqrt n\left|\Pp(\operatorname{Bin}(n,p)\le x)
 -\Phi((x-np)/\sqrt{npq})\right|<\CE\rho_B.
\]
It applies directly to the actual threshold, including atoms. At positive
noise the same theorem closes $C<C_B<\CE$. No estimate divides by zero
or replaces a uniform small-noise bound with a fixed-law limit.

\paragraph{A triangle of allowed noise and displacement.}
Introduce Bernoulli coordinates
\[
 X=\ell(B_0+U-p),\quad B_0=\mathbf1_{\{X>0\}},\quad
 \ell=\bar A+\bar B=\frac m{pq},\quad \E[U\mid B_0]=0.
\]
Then
\[
 \ell^2=\frac{1-\theta}{pq},\quad m\ell=1-\theta,
 \qquad d=\frac{\Delta_A+\Delta_B}{\ell},\quad
 \lambda=n\E U^2=\frac{npq\theta}{1-\theta}.
\]

\begin{theorem}[Joint inward/noise constraint]\label{j:local:triangle}
Suppose $n\ge N>0$, and bounds for this same law give
\[
 nD\le N d_D,\quad \theta\le\Theta<1,\quad pq\le v_*,
 \quad A+B\ge S_*>0,
\]
and
\[
 D\ge b+6mS_*(\Delta_A+\Delta_B)+K\theta,
 \qquad b\ge0,\quad K>0.
\]
Set $r=d_D-b$. If $r>0$, define
\[
 d_* =\frac{r}{6(1-\Theta)S_*},\qquad
 \Lambda=\frac{Nv_*r}{K(1-\Theta)}.
\]
Then
\begin{equation}\label{j:local:eq:triangle}
                 \boxed{\displaystyle
                   \frac d{d_*}+\frac\lambda\Lambda\le1.}
\end{equation}
If $r<0$, the premises are impossible. If $r=0$, they force
$d=\lambda=0$.
\end{theorem}
\begin{proof}
Multiply the assumed joint lower bound for $D$ by $n$. Since all
coefficients being weakened are nonnegative,
\begin{align*}
 Nd_D\ge nD
 &\ge nb+6n(1-\theta)S_*d+K\frac{(1-\theta)\lambda}{pq}\\
 &\ge Nb+6N(1-\Theta)S_*d+\frac{K(1-\Theta)}{v_*}\lambda.
\end{align*}
Subtract $Nb$. For $r>0$, division by $Nr$ proves
\eqref{j:local:eq:triangle}. For $r\le0$, positivity of both coefficients gives
the stated alternatives without dividing by zero.
\end{proof}

Apply this with the binary charge in \eqref{j:core:simultaneous},
$S_*=A_l+B_l$ and $K=\min\{3(\sigma+1)B_l,3(\sigma+3)A_l\}$.
\begin{corollary}[Closed noise bands and outward rounding]
If $0\le L\le\lambda\le U\le\Lambda$, then
\[
                       d\le d_*(1-L/\Lambda).
\]
Both intercepts in \eqref{j:local:eq:triangle} may first be replaced by larger
positive numbers. The same conclusion then holds using those enlarged
intercepts and the enlarged noise ceiling.
\end{corollary}
\begin{proof}
Rearrange \eqref{j:local:eq:triangle} and use $\lambda\ge L$.
Increasing either denominator weakens its nonnegative summand, so it
preserves the joint inequality. This also proves the rounding assertion.
\end{proof}



\paragraph{The coefficient must be bounded in the correct direction.}
Suppose $n\ge N$, $p\in[p_-,p_+]\subset(0,1/2)$,
$0\le d\le d_0$, and $0\le\lambda\le\Lambda$.
Choose $0<v_*\le p_-(1-p_-)$ and put
\[
 V^*=p_+(1-p_+)+\Lambda/N,\qquad
 Q_*=6p_--3p_-^2-\frac32-d_0
                     -\frac{3(1-2v_*)\Lambda}{8v_*N}.
\]
Then $Q\ge Q_*$: the first term increases with $p$, and
$(1-2v)/v$ decreases with $v$. In particular, $Q_*$ may be negative.
Already with $d=s=0$ its sign changes at $p=1-1/\sqrt2$.

Let $C$ be the score multiplying the moment correction, and assume
$0<c_*\le C\le c^*$. Choose positive multipliers satisfying
\[
        0<m_-\le\frac{c_*}{(V^*)^{3/2}},\qquad
        m_+\ge\frac{c^*}{v_*^{3/2}},
\]
and define the signed gain
\begin{equation}\label{j:local:gain}
             G_*:=\min\{m_-Q_*,m_+Q_*\}.
\end{equation}

\begin{proposition}[Signed gain and nonnegative repair]\label{j:local:repair}
Let $\Phi$ denote the standard normal distribution function. For an actual
count displacement $\delta$, set
\[
 \mathcal R=\sqrt n\left\{
 \Phi\!\left(\frac{\delta}{\sqrt{nv}}\right)
 -\Phi\!\left(\frac{\delta}{\sqrt{n(v+s)}}\right)\right\}
                         +C(\rho_B-\rho).
\]
If $D_+\ge\max(\delta,0)$ and $P_*\ge .4/(2v_*^{3/2})$, then
\begin{equation}\label{j:local:repair_bound}
 \mathcal R\le(P_*D_+-G_*)\frac{\lambda}{n}
 \le \kappa\frac{\lambda}{N},\qquad
 \kappa=\max\{0,P_*D_+-G_*\}.
\end{equation}
\end{proposition}
\begin{proof}
Write $A=C/(v+s)^{3/2}$. The hypotheses give
$0<m_-\le A\le m_+$. By \eqref{j:local:cubic-bound},
$C(\rho-\rho_B)\ge sAQ\ge sAQ_*$.
If $Q_*\ge0$, $AQ_*\ge m_-Q_*$; if $Q_*<0$,
$AQ_*\ge m_+Q_*$. This proves $C(\rho-\rho_B)\ge sG_*$
in both cases. One must use the \emph{upper} multiplier in the negative
case.

The normal part is nonpositive when $\delta\le0$.
For $\delta>0$, differentiation with respect to an intermediate variance
$t\in[0,s]$ gives
\[
 \frac{d}{dt}\left[-\sqrt n\,
 \Phi\!\left(\frac{\delta}{\sqrt{n(v+t)}}\right)\right]
 =\frac{\delta}{2(v+t)^{3/2}}
   \phi\!\left(\frac{\delta}{\sqrt{n(v+t)}}\right)
 \le\frac{.4\delta}{2v_*^{3/2}}\le P_*D_+.
\]
Integrate over $[0,s]$, and combine with the moment bound.
Finally $s=\lambda/n$ and $\kappa\ge0$, so
$(P_*D_+-G_*)\lambda/n\le\kappa\lambda/n\le\kappa\lambda/N$.
The clamp at zero is needed before replacing $n$ by its lower bound.
\end{proof}



\paragraph{A binomial mass bound on the actual phase cell.}
Suppose the finite Stirling argument has supplied positive constants
$E,L_0,R$ on a fixed segment containing every $k/n$ under consideration, with
\[
 \sqrt n\,\pi_k=\frac{\phi(0)}{\sqrt v}
       \exp\{-\delta^2/(2nv)+\zeta\},\qquad
 n|\zeta|\le Z_N(|\delta|),\qquad
 Z_N(t)=\frac{L_0t}{2}+R+\frac{Et^3}{6N},
 \quad n\ge N.
\]
The constants may have been computed on a much wider global phase range.
Let $v_-\le v\le v_+$ and let $q_+\ge\sqrt{v_+}$.
\begin{lemma}[Mass localization]\label{j:local:lem:mass}
On any smaller closed cell where $|\delta|\le D$,
\[
 \sqrt n\,\pi_k\ge
 \frac{.39886}{q_+}
 \left[1-\frac{D^2/(2v_-)+Z_N(D)}N\right].
\]
The maximum of this expression and any previously proved mass lower bound
is also a valid mass lower bound. There is no requirement $D\ge1$.
\end{lemma}
\begin{proof}
Set $A_*=D^2/(2v_-)+Z_N(D)\ge0$. The logarithmic formula gives
$\sqrt n\pi_k\ge(.39886/q_+)e^{-A_*/n}$.
Now $e^{-z}\ge1-z$ and $n\ge N$. If the bracket is negative the displayed
inequality still holds, and the previously positive bound can be retained.
The global Stirling segment, rather than the smaller value $D$, continues
to justify $E,L_0,R$.
\end{proof}

\paragraph{A sharper elementary harmonic constant.}
\begin{lemma}\label{j:local:lem:cos}
For every $x\ge0$,
\[
               1-\cos(2\pi x)\le\frac{4553}{1000}x.
\]
\end{lemma}
\begin{proof}
Put $K=4553/1000$. For $x\ge2/K$ the assertion follows from
$1-\cos\le2$. On $[0,2/K]$ take the degree-fourteen alternating upper
Taylor polynomial $Q(x)$ for $1-\cos(2\pi x)$. Its remainder is nonpositive:
the first omitted term is negative, and the remaining term magnitudes
decrease, since $(4\pi/K)^2<16\cdot15$.

Use the directed endpoints
\[
 \pi_-=3.141592653589793<\pi<3.141592653589794=\pi_+.
\]
In each positive coefficient of $Q$ replace $\pi$ by $\pi_+$; in each
negative coefficient replace it by $\pi_-$. This only increases $Q$.
The exact checker expands $P(x)=K-Q(x)/x$, a polynomial of degree thirteen,
on each of the 64 intervals $[2j/(64K),2(j+1)/(64K)]$. If
$P(x)=\sum_{h=0}^{13}p_hx^h$, its Bernstein coefficients on $[a,a+w]$ are
\[
 b_k=\sum_{i=0}^k\frac{\binom ki}{\binom{13}i}
        \sum_{h=i}^{13}p_h\binom hi a^{h-i}w^i,
               \qquad k=0,\ldots,13.
\]
All 896 rational coefficients are positive; the smallest is greater than
$.00012312$. Thus $P>0$ on the full interval. The value at zero follows by
continuity (or directly).

For completeness the stated $\pi$ endpoints are checked using alternating
series for $\arctan(1/5)$ (32 terms) and $\arctan(1/239)$ (12 terms).
The tangent addition formula gives
$\tan(4\arctan(1/5)-\arctan(1/239))=1$.
The angle lies between zero and $4/5<\pi/2$, so it equals $\pi/4$.
The positivity follows already from
$4(1/5-1/(3\cdot5^3))-1/239>0$; the upper bound follows from
$\arctan z<z$. This proves Machin's identity on the correct branch.
The full rational recipe is \texttt{check\_harmonic4553.py}; no sampled
positivity is used.
\end{proof}

If the original characteristic-function argument gives
$|\E e^{2\pi iW}|\le e^{-rt}$, Lemma~\ref{j:local:lem:cos} upgrades its fractional
part estimate from $(5/23)(1-e^{-rt})$ to
\[
                 \frac{1000}{4553}(1-e^{-rt}).
\]
Consequently every occurrence of $5/23$ and $5/46$ in that harmonic gain is
replaced by $1000/4553$ and $500/4553$, respectively. The existing degree-six
Taylor gain and its monotonicity argument remain unchanged.



\paragraph{The signed Gaussian correction in nearest-count comparison.}

Let $n\ge N$, $v=p(1-p)>0$, $s\ge0$, $\lambda=ns$, and let the
actual nearest count and phase satisfy
\[
 \delta=k-np=h-u,\qquad
 f(t)=\frac{\Phi((\delta+t)/\sqrt{n(v+s)})-
                   \Phi(\delta/\sqrt{n(v+s)})}{\pi_k}.
\]
The actual $k$ and $\delta$ are held fixed when $t$ varies. Assume the
proved finite Stirling formula on this count range is
\[
 \sqrt n\,\pi_k=\frac{\phi(0)}{\sqrt v}
       \exp\{-\delta^2/(2nv)+\zeta\},\qquad
 n|\zeta|\le Z_N(D),\quad |\delta|\le D.
\]
Write $v\ge v_->0$ and $Z=Z_N(D)\ge0$.

\begin{lemma}[Negative phase]
Suppose $u=-r$, $0\le a\le r\le b$, and $h\le h_*$. Define
\[
 t_*=\max(0,\min(b,h_*+b)),\qquad G_-=2(h_*+b)t_*-t_*^2,\qquad
 \eta_-=G_-/(2v_-)+Z.
\]
If $D^2\le Nv_-$ and $\eta_-<N$, then
\[
 f'(-t)\le\frac1{1-\eta_-/N}\quad(0\le t\le r),\qquad
 u-f(u)\le\frac{b\eta_-}{n(1-\eta_-/N)}.
\]
\end{lemma}
\begin{proof}
Direct differentiation and the Stirling formula give
\[
 \log f'(-t)=-\tfrac12\log(1+s/v)
 +\frac{\delta^2s}{2nv(v+s)}
 +\frac{2\delta t-t^2}{2n(v+s)}-\zeta.
\]
Since $\log(1+x)\ge x/(1+x)$ for $x\ge0$, the first two terms are at most
\[
 \frac{s}{2(v+s)}\left(\frac{\delta^2}{nv}-1\right)\le0.
\]
This cancellation removes the variance term from the upper derivative
bound. Because $\delta=h+r$, $r\le b$ and $0\le t\le r\le b$,
\[
 2\delta t-t^2\le 2(h_*+b)t-t^2\le G_-.
\]
The last maximum is the concave quadratic maximum on $[0,b]$, achieved
at the displayed projection $t_*$. It is nonnegative because $t=0$ is
available. Therefore $\log f'(-t)\le\eta_-/n$.
Using $e^x\le(1-x)^{-1}$ for $0\le x<1$ gives
\[
 f'(-t)-1\le\frac{\eta_-/n}{1-\eta_-/n}
             \le\frac{\eta_-}{n(1-\eta_-/N)}.
\]
Integrate $u-f(u)=\int_0^r(f'(-t)-1)\,dt$.
\end{proof}

Equivalently, $G_-=b^2+2h_*b$ for $h_*\ge0$, and
$G_-=\max(0,h_*+b)^2$ for $h_*<0$. This sharper vertex is used in the
selected finite left LOW implementation.

\begin{lemma}[Positive phase]
Suppose $u\in[a,b]\subset[0,1/2]$, $h\le h_*$, and define
\[
 G_+=\max\left(0,\max_{z\in[a,b]}(2h_*z-z^2)\right).
\]
Then
\[
 u-f(u)\le\frac b n
       \left(\frac\lambda{2v_-}+\frac{G_+}{2v_-}+Z\right).
\]
The inner maximum is at the projection of $h_*$ onto $[a,b]$.
\end{lemma}
\begin{proof}
The exact derivative identity and $\log(1+s/v)\le s/v$ give
\[
 \log f'(t)\ge-\frac{\lambda+2\delta t+t^2}{2nv}-\zeta.
\]
For fixed actual $u$, the function $2\delta t+t^2$ is convex on $[0,u]$.
Its maximum is at an endpoint, and $\delta=h-u$ gives
\[
 2\delta t+t^2\le\max(0,2hu-u^2)\le G_+.
\]
Use $e^{-x}\ge1-x$ and integrate from zero to $u$.
\end{proof}

\paragraph{Whole-count insertion and interval coverage.}
The whole-count identity is split as
$Q/\pi_k-f(u)=(Q/\pi_k-u)+(u-f(u))$.
Thus one needs an upper bound on the signed correction $u-f(u)$;
an absolute-value estimate is unnecessary. After division by $\lambda$,
the new coefficients replacing the old Gaussian correction are
\[
 \begin{array}{ll}
 u\ge0:&\displaystyle
 \frac b{2v_-}+\frac{b\{G_+/(2v_-)+Z\}}\lambda,\\[6pt]
 u\le0:&\displaystyle
 \frac{b\eta_-}{(1-\eta_-/N)\lambda}.
 \end{array}
\]
On each fixed phase cell these have the form $A+B/\lambda$ with
$A,B\ge0$. They therefore preserve the original convex endpoint check on
each closed noise interval. Every phase cell is confined to one half-line;
zero is a shared closed endpoint. The count displacement bound $D$ comes
from the same actual $h-u$ cell, and the Stirling constants continue to
come from the already proved global count segment.

The other all-index steps are unchanged: the event bracket is checked
negative before substituting a binomial mass lower bound, and the entire
inverse-$n$ correction is checked positive before replacing $n$ by $N$.
The inverse correction is $I(\lambda)=K+A+B/\lambda$, hence decreases with $\lambda$; positivity at both endpoints (in particular the upper endpoint) licenses the all-index step throughout the interval. The final score is convex in $\lambda$: its event fraction is a positive constant divided by an affine positive function, its $B/\lambda$ coefficient is nonnegative, and its remaining terms are affine or constant. These are independent gates in the exact kernel. The formulas apply to
both signs of $h$ without reflecting the law or changing the count during
integration. They also supply an improved upper derivative for the
negative-phase MID inequalities.
\paragraph{Full-phase MID below the whole-count threshold.}
For $\lambda\in[.05,.08]$, the selected shifted-quartic MID comparisons
are checked on every closed phase interval of length $.001$ in
$[0,1/2]$, for both phase signs. The shifted-quartic polynomial and
nearest-count identity are defined at phase zero, so the two closed covers
include $u=0$ directly. This permits the whole-count LOW comparison to
start at $\lambda=.08$. The four MID bands use phase starts
$0,0,.19,.27$ at noise bands $[.05,.06]$, $[.06,.08]$,
$[.08,.10]$, and $[.10,.12]$, respectively.




\paragraph{Positive allowances and finite inverse-index corrections.}
\label{j:local:allowance}
For actual index $39$ put $K=741/4000$. The target is
$C< C_B+K/39$, so the scalar tests differ from
the zero-allowance tests in Section~\ref{j:local:shortloc:section}.
All of them follow from the same identity
\eqref{j:local:shortloc:master}. Write $b=\sqrt{39}\pi_k$ and retain
the verified two-sided bound $m\le b\le1$. A nearest-count
bracket $D$ may now have either sign; consequently
\[
 bD\ge\min\{mD,D\}.
\]
If the repair is at most $r\lambda/39$, with $r\ge0$, a noise
band $[a,b_\lambda]$ is excluded by the intermediate-noise test
\[
 \min\{mD,D\}-r b_\lambda/39+K/39>0.
\]
Indeed the identity gives $\rho_B(C-C_B)<K/39$.
The elementary inequality $\rho_B\ge1$ follows from
$(1-2v)^2-v=(1-v)(1-4v)\ge0$, $0<v\le1/4$.
Thus $C<C_B+K/39\le\CE+K/39$.

At arbitrarily small positive noise the rescaled quartic and tail
estimates give the same inequality per unit $\lambda$. The tiny-noise
test adds the allowance
\[
             \frac{K}{39\max\{.05,\Lambda\}}
\]
to that per-noise reserve. Multiplying by
$0<\lambda\le\min\{.05,\Lambda\}$ again gives a bound
strictly below $K/39$. No positive lower noise cutoff is inserted.

For the whole-count test, write
\[
 S(\lambda)=\max\{mB(\lambda),B(\lambda)\}
                         +I(\lambda)/39.
\]
Both the bracket and the inverse correction may have either sign;
the first maximum uses $m\le b\le1$, while the division by $39$
uses the actual integer. Hence
$\rho_B(C-C_B)\le\lambda S(\lambda)$.
The saved gate on $[a,b_\lambda]$ is
\[
        K\rho_{B,l}/39-b_\lambda S(\lambda)>0,
\]
where $\rho_{B,l}$ is a positive lower bound for $\rho_B$.
If $S<0$ the conclusion is immediate; otherwise it follows by
$\lambda\le b_\lambda$ and $\rho_B\ge\rho_{B,l}$.
The formula is convex in positive $\lambda$, since its only
nonaffine terms are positive multiples of
$(\alpha\lambda+\beta)^{-1}$ and $1/\lambda$, and maxima of
convex functions preserve convexity. Both endpoint gates therefore
cover the whole noise cell.



At zero noise, the strict finite-binomial theorem gives $C<\CE$,
so every positive allowance is excluded there as well.
For a bounded integer band $N\le n\le L$, a signed inverse correction
has the valid bound $I/n\le\max(I/N,I/L)$. If
$b=\sqrt n\pi_k\in[b_0,1]$, a signed bracket satisfies
$bJ\le\max(b_0J,J)$ and $bJ\ge\min(b_0J,J)$.
These two-sided rules replace the sign restrictions needed by the
infinite-index test. The actual index remains in every binomial slack
calculation. With a uniform binomial reserve $\epsilon_p>0$ and an
upper bound $S$ on $\rho_B(C-C_B)/\lambda$, it is enough to verify
$\rho_{B,l}\epsilon_p>U\max(S,0)$ on a noise cell $\lambda\le U$.
No negative inverse correction is enlarged by replacing its denominator
with the lower index.

\paragraph{The local exclusion, including equality.}\label{j:local:comparison}
\begin{proposition}[Finite local exclusion]\label{j:local:finite42}
For an integer $42\le n\le48$, no attained positive ratio maximum satisfies
\begin{gather*}
 0<n(C_n-\CE)\le54/100,\qquad n(C_{n-1}-C_n)_+\le19/100,\\
 C_n<4226/10000,\qquad C_{n-1},C_{n-2}<4253/10000.
\end{gather*}
\end{proposition}
\begin{proof}
Apply the support polynomial with $(N,L,K,q,R)=(42,48,0,.19,3.84)$.
In particular its favorable fifth-order term is evaluated at $48$; the
infinite-range support row cannot be substituted for this finite band.
The subsequent bounded-threshold tests use excess allowance .54, loss .19
and own-score cap .4226. The resource invariant is
$\mathcal R_n\le r$, $nD\le42d$. The transfer argument of
Section~\ref{j:local:transfer} applies at floor $42$, including the signed
curvature, retained constant penalties and actual-noise triangle.
Confinement uses the already acquired strips and the substitutions in
Section~\ref{j:local:finite-confinement}, with deletion floor $40$.
The left radius first contracts to $3.025$; contact then proves the
coordinate strip $(-.575,1.575)$ before confinement is invoked.
On the four noise bands $[.05,.08]$, $[.08,.10]$, $[.10,.12]$ and
$[.12,\Lambda]$, the finite central amplitude ceilings are
$.375,.325,.275,.25$, and the left ceilings are $.40,.35,.30,.275$.
Each ceiling is intersected with its row's independently acquired
global amplitude. These finite profiles use the same proved
confinement gates with their own parameters; the uniform-tail profiles
are not substituted for them.

The accepted acquisition cover has 184 closed cells: 92 exclusions,
45 left rows covering $[.29,.38]$ and 47 central rows covering
$[.38,.474]$. Thus the moment and probability complements are excluded
before the comparison. Every row retains its own resource and both
threshold signs. The central positive-noise cover uses TINY on $(0,.05]$;
above .05 the whole-count phase radii are .09,.14,.19,.29,.5 on bands
separated by .06,.08,.10,.12, with the complementary nearest-count
phases through .12. The left cover uses TINY on $(0,.05]$, full-phase
MID on $[.05,.08]$, then nearest count outside and whole count inside
phase radii .19 and .27 on $[.08,.10]$ and $[.10,.12]$; above .12 the
whole-count test covers every phase. Each band is intersected with its
actual acquired noise ceiling. The 78 central and 74 left signed
whole-count rows, and the nearest-count covers, include every seam.

The signed cubic multipliers, noise-dependent inward bound, localized
binomial mass, $4553/1000$ harmonic inequality and signed Gaussian
correction proved above give the strict comparison on these entire
closed cells. The two complete LOW families retain their inverse-index
sign gates before any floor substitution. Zero noise is Bernoulli.
The strict local comparisons and Schulz's bound contradict $C_n>\CE$.
This proposition has the stated upper index $48$; it is separate from the
floor-$40$ comparison used only at actual indices $40,41,42$.
\end{proof}

\begin{proposition}\label{j:local:uniform}
There is no attained positive closed-event ratio maximizer satisfying
\begin{gather*}
 n\ge47,\qquad 0\le n(C_n-\CE)\le7/10,\qquad
 n(C_{n-1}-C_n)_+\le19/100,\\
 C_n<4247/10000,\qquad C_{n-1},C_{n-2}<4253/10000.
\end{gather*}
\end{proposition}
\begin{proof}
The support, acquisition and confinement implications proved above and
their accepted closed numerical covers place the same law in one of the
45 central or 38 left rows. Their six local tests cover every positive
noise, both signs of $H$ and every phase, and give
$\rho_B(C_n-C_B)<-.001\lambda$. Schulz's theorem gives $C_B<\CE$.
For zero noise the law is exactly Bernoulli, and the same finite-binomial
theorem gives a strict bound at every threshold. Both cases contradict
$C_n\ge\CE$.
The initial score interval is $[.4095,.4247]$, the native barrier is zero,
and the resource identity uses $e=C_n-\CE$ only as a nonnegative summand.
Strict numerical reserves are independent of $e>0$. The stored value $98$
in initialization occurs only in $.4095+K/98$ with $K=0$; the support's
last band and all later monotonicity gates have no upper index limit.
\end{proof}
\begin{corollary}\label{j:local:equality}
Once $C_j\le\CE$ has been proved at every index, $C_n<\CE$ for every
$n\ge47$, in particular for every $n\ge700$.
\end{corollary}
\begin{proof}
If $C_n=\CE$, ratio attainment supplies the positive closed-event maximum.
The predecessor loss and scaled excess are zero, and both deleted bounds
are at most $\CE$. All hypotheses of the proposition hold. Attainment is
essential here: strictness for each individual law alone would not rule
out a nonattained supremum.
\end{proof}

\section{Completion of the proof}\label{j:composition}

We now put together the finite induction and the two local exclusions.
Each exclusion concerns one attained ratio maximizer. The induction
supplies its predecessor hypotheses before any acquisition or comparison
is applied to that law.

\begin{proof}[Proof of Theorem~\ref{j:main}]
The finite induction, Theorem~\ref{j:finite:base}, gives
$C_j\le\CE$ for $1\le j\le42$. If a later strict failure exists,
choose its least index $n$. Both deleted constants are then at most
$\CE$, and Corollary~\ref{j:smoothing:entry} yields
\[
 0<n(C_n-\CE)<\frac{3003}{10250}<.3,\qquad
 C_n<\frac{73}{175}<.4226,\qquad
 n(C_{n-1}-C_n)_+=0.
\]
Also $C_{n-1},C_{n-2}\le\CE<.4253$. For $43\le n\le48$,
these are the hypotheses of the finite local exclusion,
Proposition~\ref{j:local:finite42}. For $n\ge49$ they are the
hypotheses of the uniform exclusion, Proposition~\ref{j:local:uniform}.
In each case attainment supplies the maximizer and the local comparison
rules it out. Thus $C_n\le\CE$ at every index.

Now suppose $C_n=\CE$ for some $n\ge47$. The just-proved upper bound
makes the two deleted constants at most $\CE$ and the predecessor
loss zero. The scaled excess is zero as well. The uniform local
exclusion includes this boundary case, so attainment again gives a
contradiction. This proves the strict assertion about the supremum.
Finally, Proposition~\ref{j:external:limit} supplies a fixed Bernoulli
law whose normalized discrepancies converge to $\CE$. Hence every
smaller universal coefficient fails.
\end{proof}

There is no assumption of a sharp tail in this induction. The numerical
finite base uses previously established smaller indices and its stated
loose ceilings. The finite local exclusion is proved at its actual
indices, and the uniform exclusion uses analytic estimates valid for
every index above its floor. Their certificates establish numerical
inequalities on those independently justified domains. This order is
also why equality is considered only after the weak all-index bound
has been proved.

\begingroup\small\raggedright
\urlstyle{same}
\bibliographystyle{plainnat}
\bibliography{references-classical,references-modern-a,references-modern-b,references-added}
\endgroup
\end{document}
